\documentclass[a4paper,10pt]{article}
\usepackage[top=24mm,bottom=24mm,left=26mm,right=26mm,headheight=14pt]{geometry}
\usepackage{fontspec,amsmath,amsthm,unicode-math}
\usepackage[english,french]{babel}
\usepackage{microtype,xcolor,graphicx,tikz,booktabs,array,needspace,fancyhdr}
\usepackage{longtable,tabularx,enumitem,float,import,xurl,multicol}
\usetikzlibrary{arrows.meta,calc,positioning,decorations.pathreplacing,matrix}
\usepackage[unicode,hidelinks]{hyperref}
\usepackage[nameinlink,noabbrev]{cleveref}
\crefname{theorem}{théorème}{théorèmes}\Crefname{theorem}{Théorème}{Théorèmes}
\crefname{proposition}{proposition}{propositions}\Crefname{proposition}{Proposition}{Propositions}
\crefname{lemma}{lemme}{lemmes}\Crefname{lemma}{Lemme}{Lemmes}
\crefname{section}{section}{sections}\Crefname{section}{Section}{Sections}
\crefname{equation}{équation}{équations}\Crefname{equation}{Équation}{Équations}
\setmainfont{STIXTwoText-Regular.otf}[BoldFont=STIXTwoText-Bold.otf,ItalicFont=STIXTwoText-Italic.otf,BoldItalicFont=STIXTwoText-BoldItalic.otf]
\setmathfont{STIXTwoMath-Regular.otf}
\newtheorem{theorem}{Théorème}[section]
\newtheorem{proposition}[theorem]{Proposition}
\newtheorem{lemma}[theorem]{Lemme}
\newtheorem{corollary}[theorem]{Corollaire}
\newtheorem{falsifier}[theorem]{Critère de réfutation}
\newtheorem{ablation}[theorem]{Preuve de nécessité structurelle}
\newtheorem{criterion}[theorem]{Critère}
\theoremstyle{definition}\newtheorem{definition}[theorem]{Définition}
\newtheorem{axiom}[theorem]{Axiome}
\newtheorem{principle}[theorem]{Principe}
\theoremstyle{remark}\newtheorem{remark}[theorem]{Remarque}
\newcommand{\RR}{\mathbb R}\newcommand{\QQ}{\mathbb Q}
\newcommand{\ZZ}{\mathbb Z}\newcommand{\CC}{\mathbb C}
\newcommand{\Gr}{\operatorname{Gr}}\newcommand{\Res}{\operatorname{Res}}
\newcommand{\APP}{\mathrm{APP}}\newcommand{\TPK}{\mathrm{TPK}}
\newcommand{\TRIT}{\mathrm{TRIT}}\newcommand{\id}{\operatorname{id}}
\newcommand{\diag}{\operatorname{diag}}\newcommand{\conv}{\operatorname{conv}}
\newcommand{\FF}{\mathbb F}\newcommand{\Id}{\operatorname{id}}
\newcommand{\tr}{\operatorname{tr}}\newcommand{\dd}{\mathrm d}
\newcommand{\Span}{\operatorname{span}}\newcommand{\Cyl}{\operatorname{Cyl}}
\providecommand{\llbracket}{\lBrack}\providecommand{\rrbracket}{\rBrack}
\providecommand{\square}{\mdlgwhtsquare}
\providecommand{\dashrightarrow}{\mathrel{\rightdasharrow}}
\providecommand{\index}[1]{}
\newcommand{\HMT}{\mathrm{HMT}}\newcommand{\Tr}{\operatorname{Tr}}\newcommand{\Spec}{\operatorname{Spec}}
\newcommand{\Cstar}{C^{*}}\newcommand{\R}{\mathbb R}\newcommand{\Z}{\mathbb Z}\newcommand{\I}{\mathrm I}
\setcounter{tocdepth}{2}
\setlist{nosep,topsep=4pt}
\numberwithin{equation}{section}
\setlength{\parskip}{4pt plus 1pt}\setlength{\parindent}{1em}
\setlength{\emergencystretch}{2em}
\renewcommand{\normalsize}{\fontsize{10.5}{13.5}\selectfont}\normalsize
\clubpenalty=10000\widowpenalty=10000\displaywidowpenalty=10000
\AddToHook{cmd/section/before}{\Needspace{10\baselineskip}}
\AddToHook{cmd/subsection/before}{\Needspace{6\baselineskip}}
\AddToHook{env/theorem/before}{\Needspace{6\baselineskip}}
\AddToHook{env/proposition/before}{\Needspace{6\baselineskip}}
\pagestyle{fancy}\fancyhf{}
\fancyhead[L]{\footnotesize Grassmanniennes positives et amplituèdres : formes canoniques et raffinement}
\fancyfoot[C]{\thepage}\renewcommand{\headrulewidth}{.2pt}
\hypersetup{pdftitle={Grassmanniennes positives et amplituèdres : formes canoniques et raffinement},pdfsubject={Dualité réticulaire et écart de masse de Yang–Mills en holographie modulaire triadique},pdfauthor={Oumar Haidara Fall; Rubén Ramos Balsa}}
\makeatletter
\newcommand{\compactcontents}{%
  \begingroup
  \fontsize{9.2}{11.1}\selectfont
  \setlength{\parskip}{0pt}\setlength{\columnsep}{7mm}
  \renewcommand{\l@section}{\@dottedtocline{1}{0pt}{1.8em}}
  \renewcommand{\l@subsection}{\@dottedtocline{2}{1.5em}{2.7em}}
  \renewcommand{\l@part}[2]{\addpenalty{-\@highpenalty}\addvspace{6pt}%
    {\bfseries\noindent ##1\nobreak\hfill\nobreak ##2\par}\nobreak}
  \section*{Table des matières}
  \@starttoc{toc}
  \endgroup}
\makeatother
% BEGIN NUCLEO_ADITIVO_20260921_PREAMBULO
\makeatletter
\@ifundefined{example}{\theoremstyle{definition}\newtheorem{example}[theorem]{Exemple}}{}
\@ifundefined{notatabla}{\newcommand{\notatabla}[1]{\par\smallskip{\footnotesize #1\par}}}{}
\makeatother
\providecommand{\FF}{\mathbb F}

% END NUCLEO_ADITIVO_20260921_PREAMBULO
% Apertura independiente del índice: corrección quirúrgica 20260928.
\let\HMTIndiceOriginal\tableofcontents
\renewcommand{\tableofcontents}{\clearpage\begingroup\setlength{\parskip}{0pt}\fontsize{9.2}{10.5}\selectfont\HMTIndiceOriginal\endgroup\clearpage}
\begin{document}
\thispagestyle{plain}
\begin{center}
{\small Manuscrit de compilation pour le transfert académique}\par\vspace{6mm}
{\LARGE\bfseries Grassmanniennes positives et amplituèdres :\\formes canoniques et raffinement}\par\vspace{3mm}
{\large Dualité réticulaire et saut de masse de Yang--Mills\\en holographie modulaire triadique}\par\vspace{5mm}
Oumar Haidara Fall \textperiodcentered\ Rubén Ramos Balsa\par
{\fontsize{8}{10}\selectfont Coécriture sans préséance de contribution\par
Intégration et compilation documentaires générées par ChatGPT - Astra.}
\end{center}
\begin{abstract}
L'Holographie modulaire triadique construit la structure discrète conjointe du continu au moyen d'APP, de TRIT et du transport avec mémoire du TPK. La réduction compatible des cylindres produit la réalisation arithmétique, tandis que la projection conjuguée conserve l'incidence exceptionnelle. Sur cette généalogie, on démontre une collection de réalisations positives récupérables. La charge uniforme et deux mémoires nonadiques déterminent le registre à douze coordonnées ; ses séparations ordonnées produisent une mesure planaire de frontière et une représentation grassmannienne avec inverse marqué. Les lecteurs qui se factorisent par le registre sont transportés vers cette représentation et conservés sous les changements de coordonnées du même point. Le raffinement nonadique construit des drapeaux positifs et des points amplituédraux dans tout rang et toute dimension externe finis admissibles. En rang un, on obtient une couverture simpliciale complète et l'annulation des résidus intérieurs en toute dimension finie. Les formes à valeurs dans la droite d'action satisfont une condition exacte de recollement sur les frontières communes. Les moyennes cellulaires conservent une mesure chronologique et décomposent l'incrément des moments en détails orthogonaux ; leur limite détermine un opérateur de Jacobi dont les compressions finies ont un spectre simple, des résidus positifs et des résolvantes en fraction continue. La réduction de Schur conserve l'énergie effective de la partition, avec reconstruction du détail et réalisation limite de Sobolev. Le flot du registre induit en outre une déformation réticulaire à vingt-quatre dimensions qui conserve le volume et la symétrie ternaire, et transforme l'inversion du paramètre en dualité réticulaire et en réciprocité thêta. On réunit les hamiltoniens elliptique, hyperbolique et parabolique, leurs coordinatisations symplectiques, leurs termes de connexion et leurs transformations de Legendre, ainsi que les lecteurs matériels correspondants. Le développement de Yang--Mills incorpore la mesure commune, l'échelle aire--volume, la coercivité uniforme, le témoin spectral persistant et la reconstruction de l'hamiltonien. L'articulation de ces résultats distingue et relie les positivités de Plücker, des moments, des constitutives et du spectre au moyen d'opérations explicites sur le même état enrichi. La représentation incidentielle conserve les douze coordonnées au moyen d'une isométrie pondérée avec inverse explicite. La transition hexade--octade permet de distinguer l'aire totale de la composition sectorielle et de reconstruire les occupations ; sa purification détermine une entropie d'intrication avec loi modulaire finie. La composition longitudinale du retour marqué détermine l'échelle aréale et conserve sa naturalité sous les raffinements hérités ; les formes canoniques à coefficient opératoriel d'aire transmettent l'additivité et les résidus de frontière.
\par\smallskip\noindent\textbf{Mots-clés :} Grassmanniennes positives ; réseaux planaires ; coordonnées de Plücker ; géométrie d’incidence ; holographie modulaire triadique.
\end{abstract}
\clearpage
\compactcontents
\clearpage

\section{Introduction : généalogie de la forme et conservation du raffinement}
\label{sec:genealogy}
La question directrice de ce travail est la conservation d'une structure engendrée lorsque sa réalisation mathématique change. L'Holographie modulaire triadique produit, au moyen d'opérations arithmétiques relevées, des états retenant simultanément l'orientation, l'incidence, la frontière et la mémoire. Sa structure discrète du continu réunit les prolongements compatibles de ces états. La réalisation archimédienne exprime des coordonnées par raffinement de cylindres ; la réalisation conjuguée conserve des relations exceptionnelles d'incidence. La généalogie précède donc les valeurs et les formes sous lesquelles elles sont représentées. C'est la thèse de départ que le développement formel suivant expose avec ses opérations et ses dépendances.

Le problème géométrique propre à l’article commence lorsqu’une sortie
positive de cette construction devient une séparation entre nœuds,
un poids de réseau, une coordonnée de Plücker, un moment d’une mesure, un coefficient
d’énergie ou une direction de déformation réticulaire. Pour justifier ce passage,
il faut déterminer quelle information retrouve chaque représentation, comment
elle se compare aux autres et ce que conserve son raffinement. L’intérêt
d’une réalisation positive augmente de manière décisive lorsqu’elle admet un
inverse marqué : la géométrie cesse alors de n’offrir qu’une évaluation
de l’état et conserve désormais l’information utilisée par ses lecteurs.
La marque comprend l’ordre des positions, l’échelle totale et les données
enrichies qu’exige le lecteur particulier.

L’apport central s’organise autour de trois opérations. La première
est la reconstruction : la charge et la mémoire récupèrent les coordonnées
du registre, et les quotients de mineurs récupèrent ses séparations.
La seconde est le raffinement compatible : une partition se subdivise
en conservant ses extrémités, ses marques et une archive exacte du détail.
La troisième est le transport des structures : une forme, un opérateur ou un
lecteur est transporté à une autre représentation au moyen d’une application explicite. Ces
opérations permettent de démontrer la conservation des résidus, de l’énergie effective,
des moments et des spectres dans les domaines concrets où chaque identité a
un sens. Le qualificatif positif acquiert ainsi plusieurs réalisations
mathématiques précises, articulées par la reconstruction de la même donnée.

La géométrie positive possède ses propres antécédents, qui permettent de situer précisément le niveau de cette contribution. Les mesures de frontière des réseaux orientés paramètrent des cellules de la grassmannienne totalement non négative dans le travail de Postnikov~\cite{postnikov} ; Scott décrit la structure d'algèbre amassée de l'anneau de coordonnées grassmannien~\cite{scott}. Arkani-Hamed et Trnka introduisent l'amplituèdre dans la description géométrique des amplitudes de la théorie supersymétrique \(\mathcal N=4\) de Yang--Mills à la limite planaire~\cite{amplituhedron}. La formulation des géométries positives caractérise ensuite leurs formes canoniques par des pôles logarithmiques et des résidus récursifs de frontière~\cite{positive}. Le présent développement compose ces langages de réalisation avec des données préalablement engendrées par HMT. Son résultat distinctif est une collection marquée, récupérable et raffinable, accompagnée des applications reliant sa forme positive et ses lectures matérielles.

Le raffinement étend en outre la portée dimensionnelle. Le nombre de nœuds croît par subdivisions nonadiques et permet de construire des matrices positives de tout rang et de toute dimension externe finis compatibles avec ce nombre. La preuve repose sur les déterminants de Vandermonde et sur des formes de moments de rang plein. En rang un, on obtient la géométrie convexe complète, avec couverture par simplexes et résidus en toute dimension finie. Pour les rangs supérieurs, on construit une collection explicite d'images amplituèdrales et l'on préserve leurs drapeaux marqués. La couverture globale d'un amplituèdre de rang supérieur constitue un énoncé de portée différente du théorème d'existence et de compatibilité de cette collection ; les preuves ultérieures conservent cette distinction dans leur formulation.

La dimension physique est introduite par les opérateurs correspondants. Le calendrier tritique distingue rotation elliptique, ouverture hyperbolique et transport parabolique de mémoire. Leurs réalisations hamiltoniennes conservent la forme symplectique et permettent de calculer la transformation de Legendre ainsi que les termes de connexion issus d'une coordinatisation variable. Dans le secteur de jauge, l'échelle aire--volume, la mesure commune et les formes de réflexion précèdent la reconstruction de l'hamiltonien de Yang--Mills. L'argument spectral est présenté avec la coercivité uniforme, le témoin qui atteint le premier niveau positif et les opérateurs d'entrelacement qui le préservent. La relation avec les formes canoniques est établie au moyen de leurs lecteurs et de leurs résidus, tandis que le saut appartient à l'hamiltonien et à son espace physique. Cette hiérarchie situe chaque résultat à l'endroit où sa démonstration produit effectivement la conclusion.

L'exposé incorpore d'abord la construction formelle nécessaire pour obtenir l'état et son registre. Il développe ensuite les réalisations positives, l'extension dimensionnelle, l'énergie et la dualité réticulaire ; il réunit enfin les réalisations hamiltoniennes et le développement de jauge. Les programmes de vérification accompagnent des identités déterminées et des contrôles négatifs explicites. La démonstration symbolique conserve la portée générale ; l'exécution exacte vérifie les spécialisations déclarées et leur correspondance avec les formules imprimées.

\subsection{Généalogie arithmétique et système de réalisations géométriques}
La construction commence par les deux feuillets arithmétiques d’APP. Si
\(p,q\in\{1,\ldots,9\}\), chaque opération conserve une division complète,
\[
 p+q=R_\Sigma+9q_\Sigma,\qquad
 pq=R_\Pi+9q_\Pi.
\]
TRIT ajoute le régime local et l'orientation. Le sélecteur tritique de phase opératoire choisit l'opération compatible avec cet état, tandis que les deux feuillets demeurent dans le support. Sur l'espace des états enrichis, le Topological Phase Kernel compose sélection, transport et mise à jour :
\[
 U_t=\operatorname{Upd}_t\circ\operatorname{Tra}_t\circ\operatorname{Sel}_t.
\]
Le registre conserve les résidus, quotients, retenues, feuillets, routes, frontières, mémoire et survie utilisés par chaque prolongement. L'holonomie nonadique \(\Gamma_9=U_{t+8}\circ\cdots\circ U_t\) appartient à cette dynamique enrichie. Sa représentation réduite de mémoire est effectuée ensuite. Le retour de phase conserve l'avancement du registre et distingue donc des états situés à des profondeurs différentes.

Le prolongement
\[
 w_6\longrightarrow w_{12}\longrightarrow w_{18}\longrightarrow
 w_{24}\longrightarrow w_{30}\longrightarrow R_{36}\longrightarrow G_9
\]
organise les raffinements compatibles du même état. Les cinq constructions du continu émergent conjointement, avec leurs aspects spectral, solénoïdal, de jauge, cohomologique et déterminantal. Les lecteurs de cet état expriment corrélativement \(\pi,\varphi,e,\alpha\). Leur origine commune admet un ordre interne d'évaluation : les coordonnées régionales déjà engendrées agissent avec le registre dodécaphasique dans les lecteurs de clôture de \(\alpha\). Dans cet article, les grassmanniennes, polygones et formes différentielles sont des réalisations ultérieures de données engendrées par cette chaîne. Leurs paramètres sont fixés avant toute comparaison géométrique ou métrologique.

Le problème concret consiste à décrire ce que conserve une réalisation
positive et ce qui peut varier lorsque sa forme change. Une même liste positive
admet plusieurs usages mathématiques : poids d’arêtes, longueurs d’intervalles,
coefficients d’une forme quadratique ou coordonnées d’un lecteur matériel.
Chaque usage possède un domaine, une opération et des invariants propres.
L’expression \emph{multimorphologie} désigne ici le système de ces réalisations
avec leurs applications explicites de compatibilité. Une identité de leurs nombres
exige une vérification moindre qu’une identité de leurs opérateurs ; le
développement suivant formule et démontre chaque niveau séparément.

La base de provenance est l'article sur la génération coinductive des constantes et la structure algébrique de l'électron, avec le développement de l'extrême et moyenne raison holographique. Leurs constructions antérieures fixent le registre utilisé ci-dessous. Les résultats géométriques de ce texte commencent par cette sortie finie et prouvent ses conséquences au moyen de données explicites. Les propositions de transfert matériel déclarent en outre les lecteurs qu'elles emploient. Cette organisation préserve la portée spécifique de chaque preuve.

\clearpage
\part{Construction arithmétique, état enrichi et continu conjoint}
\begingroup
\sloppy
\renewcommand{\TRIT}{\{-1,0,+1\}}
\import{foundation/}{sections/nucleo.tex}
% BEGIN NUCLEO_ADITIVO_20260921_CENSO
\clearpage
\section{Noyau topologique de phase et génération du catalogue de germes}
\label{act21:tpk:capitulo}

Les deux évaluations de l'APP acquièrent une efficacité générative lorsque l'on précise
quelle évaluation agit, en quelle position elle s'effectue, comment cette position se déplace
et quelle information chaque transition conserve. Le \emph{Topological Phase Kernel},
abrégé en TPK et dénommé en espagnol \emph{núcleo topológico de fase}, réunit
ces opérations dans une dynamique sur des états enrichis. Sa définition
comprend la sélection, le transport, la mise à jour arithmétique et la conservation du
registre. La construction finie développée dans ce chapitre fournit
les germes sur lesquels agiront les relèvements régionaux et le prolongement
coinductif.

Le résultat principal de cette étape est un catalogue généré et classé :
les conditions initiales des deux curseurs produisent des émissions entières ;
leur lecture ternaire détermine des mots orientés ; l'action diédrale organise
ces mots en orbites. Chaque étape conserve ses préimages et les relations
qui permettent de revenir à la description précédente. La chaîne de cardinaux
\[
 104\,976\longrightarrow468\longrightarrow243\subset729,
 \qquad 243\longrightarrow43,
\]
désigne sous une forme abrégée des objets différents. Chaque objet est ensuite construit et chaque dénombrement
est justifié, avant que l'un d'eux soit utilisé comme donnée de
la sélection régionale.

\subsection{Sélection de phase, transport et état enrichi}
\label{act21:tpk:operaciones}

\subsubsection{Sélecteur tritique de phase opérationnelle}

La phase nonadique prend ses valeurs en \(D_9=\{1,\ldots,9\}\). Le sélecteur  
tritique est la fonction
\begin{equation}
 M_{\mathrm{ph}}:D_9\longrightarrow\{-1,0,+1\},\qquad
 M_{\mathrm{ph}}(g)=
 \begin{cases}
 +1,&g\in\{1,4,7\},\\
 -1,&g\in\{2,5,8\},\\
 0,&g\in\{3,6,9\}.
 \end{cases}
 \label{act21:tpk:selector}
\end{equation}
Dans la réalisation émettrice, la valeur \(+1\) active l'évaluation additive,
\(-1\) active l'évaluation multiplicative et \(0\) enregistre l'événement
neutre. L'orientation cardinale du curseur est une autre coordonnée de l'état.
Cette séparation permet de fixer simultanément quelle opération est lue et dans quelle
direction est transporté son point d'évaluation.

Pour les transitions numérotées par \(t\geq1\), la phase précédant la lecture est
\begin{equation}
 g_t=1+[(t-1)]_9,\qquad \epsilon_t=M_{\mathrm{ph}}(g_t),
 \label{act21:tpk:fase}
\end{equation}
où \([n]_b\) désigne le représentant de \(n\) modulo \(b\) dans
\(\{0,\ldots,b-1\}\). Une fenêtre de neuf transitions contient la suite
\((+1,-1,0,+1,-1,0,+1,-1,0)\).

\subsubsection{Curseurs orientés et transport cardinal}

Des indices \(I_9=\{0,\ldots,8\}\) sont utilisés pour les positions et des étiquettes positives \((a,b)=(i+1,j+1)\) pour les évaluations APP. Un curseur fini est
\[
 c=(i,j;d)\in\mathcal S:=I_9^2\times\mathcal D,
 \qquad \mathcal D=\{N,E,S,O\}.
\]
Les quatre directions agissent par
\[
 v_N=(-1,0),\quad v_E=(0,1),\quad
 v_S=(1,0),\quad v_O=(0,-1).
\]
L'application cardinale correspondant à \(d\) est
\[
 T_d(i,j)=\bigl([i+(v_d)_1]_9,[j+(v_d)_2]_9\bigr).
\]
Ainsi, \(M_{\mathrm{ph}}\) et \(T_d\) ont des domaines, des codomaines et des actions différents. L'activation tritique sélectionne l'un des deux curseurs ; le transport cardinal modifie sa position dans la direction que ce curseur conserve déjà.

Le relèvement entier du curseur est
\(\widetilde c=(x,y;d)\in\mathbb Z^2\times\mathcal D\), avec position finie
\(([x]_9,[y]_9)\). Ses nombres de tours sont récupérés en
\begin{equation}
 x=9\lfloor x/9\rfloor+[x]_9,\qquad
 y=9\lfloor y/9\rfloor+[y]_9.
 \label{act21:tpk:vueltas}
\end{equation}
Ces deux divisions concernent le transport spatial. La division d'une
évaluation additive ou multiplicative constitue un registre arithmétique
distinct, qui sera conservé conjointement avec elles.

\subsubsection{Composition de la transition}

L'état de la réalisation finie comprend les deux curseurs relevés,
leurs directions, la phase, les accumulateurs de la fenêtre active, la liste
ordonnée des blocs déjà émis et le registre des transitions. Chaque registre
local contient les positions antérieures, les évaluations entières des deux
feuillets, leurs restes et quotients, le régime actif et les orientations.
Les champs de préfixe et de frontière qui apparaîtront dans un prolongement seront
incorporés au même état avec leurs règles de compatibilité.

La transition s'ordonne comme
\begin{equation}
 \mathcal U_t=\mathsf{Mem}_t\circ\mathsf{Tra}_t\circ\mathsf{Sel}_t.
 \label{act21:tpk:composicion}
\end{equation}
La sélection détermine le régime et conserve l'échantillon antérieur à tout déplacement. Le transport applique le mouvement actif et l'inversion cardinale prescrite par le calendrier. L'actualisation incorpore cet échantillon aux accumulateurs et au registre ; à la fin d'une fenêtre, elle émet son bloc. La composition utilise donc un échantillon antérieur au transport bien que son archivage définitif s'effectue à la fin de la transition.

\subsection{Conditions initiales et énumération réversible}
\label{act21:tpk:semillas}

Les feuillets additif et multiplicatif disposent de curseurs indépendants. Le
domaine des conditions initiales est le produit ordonné
\[
 \Omega_{\mathrm{seed}}=\mathcal S_+\times\mathcal S_\times,
 \qquad |\mathcal S|=9^2\cdot4=324.
\]
Par conséquent,
\begin{equation}
 |\Omega_{\mathrm{seed}}|=324^2=104\,976.
 \label{act21:tpk:censo-inicial}
\end{equation}
Le nombre \(81\) compte des positions ; \(324\) compte des positions orientées d'un curseur ; \(104\,976\) compte des paires ordonnées de curseurs. Le choix d'une position en APP conserve uniquement une partie d'une condition initiale de l'émetteur.

On fixe la numérotation
\[
 \operatorname{ind}(N)=0,\quad\operatorname{ind}(E)=1,\quad
 \operatorname{ind}(S)=2,\quad\operatorname{ind}(O)=3,
\]
et l'on définit
\begin{equation}
 \nu(i,j;d)=4(9i+j)+\operatorname{ind}(d).
 \label{act21:tpk:indice-semilla}
\end{equation}

\begin{proposition}[Énumération exacte des conditions initiales]
\label{act21:tpk:enumeracion}
L'application \(\nu\) est une bijection de \(\mathcal S\) sur
\(\{0,\ldots,323\}\). Le couple d'indices se code de manière réversible
par \(324\nu_++\nu_\times\in\{0,\ldots,104\,975\}\).
\end{proposition}
\begin{proof}
La division par quatre restaure
\[
 \operatorname{ind}(d)=[\nu]_4,\qquad
 j=[\lfloor\nu/4\rfloor]_9,\qquad i=\lfloor\nu/36\rfloor.
\]
Les bornes du domaine garantissent que ces trois quantités appartiennent à
leurs ensembles d'indices et reconstruisent \(\nu\) au moyen de
\eqref{act21:tpk:indice-semilla}. Pour le couple, la division euclidienne par
\(324\) restitue le quotient \(\nu_+\) et le reste \(\nu_\times\).
\end{proof}

L'énumération parcourt le domaine complet avant tout regroupement
régional. Les indices servent à localiser les préimages d'une
émission, en conservant explicitement la condition qui a produit chaque trajectoire.

\subsection{Évaluation des feuillets et registre arithmétique}
\label{act21:tpk:lecturas}

En une position positive \((a,b)\), les évaluations entières sont
\(A_+(a,b)=a+b\) et \(A_\times(a,b)=ab\). Pour \(n>0\), on conserve
\begin{equation}
 \rho_9(n)=1+[n-1]_9,\qquad
 q_9(n)=\lfloor(n-1)/9\rfloor,\qquad
 n=\rho_9(n)+9q_9(n).
 \label{act21:tpk:division-evaluacion}
\end{equation}
L'activation additive incorpore \(\rho_9(A_+)\) à l'accumulateur correspondant. L'activation multiplicative utilise un observable du résidu \(\rho_9(A_\times)\), défini ci-dessous.

\subsubsection{Exposant discret dans les unités modulo neuf}

Les puissances de \(2\) modulo \(9\) parcourent
\[
 2^0,2^1,\ldots,2^5\equiv1,2,4,8,7,5\pmod9,
 \qquad 2^6\equiv1\pmod9.
\]
Le logarithme discret dans ce groupe de six unités se représente par
\begin{equation}
 \ell_9(1)=0,\quad\ell_9(2)=1,\quad\ell_9(4)=2,
 \quad\ell_9(8)=3,\quad\ell_9(7)=4,\quad\ell_9(5)=5.
 \label{act21:tpk:logaritmo}
\end{equation}
La loi émettrice étend l'observable par le biais de \(\ell_9(3)=\ell_9(6)=\ell_9(9)=0\). Pour conserver son origine, l'échantillon enregistré inclut le résidu original et l'indicateur d'appartenance au groupe d'unités. Ainsi, une valeur observée nulle peut provenir de l'unité \(1\) ou de l'un des trois résidus non inversibles, et son origine demeure récupérable.

\subsubsection{Ordre de lecture et mouvement}

Les évaluations s'effectuent aux positions antérieures à la transition.
Si \(\epsilon_t=+1\), on accumule la lecture additive et on déplace son curseur.
Si \(\epsilon_t=-1\), on accumule l'exposant discret multiplicatif et on
déplace le second curseur. Si \(\epsilon_t=0\), on incrémente le compteur
neutre et on conserve les deux positions.

À la fin des étapes \(27\) et \(54\), les deux directions sont remplacées
par leurs opposées. Les fenêtres s'enchaînent en conservant leurs positions et
orientations. Au début de chaque fenêtre, seuls ses trois accumulateurs locaux sont remis à zéro ; le registre précédent reste disponible.

\begin{proposition}[Inversion du transport relevé]
\label{act21:tpk:inversion}
Soient \(a,b\in\{0,1\}\) les indicateurs d'activation et d'inversion d'un
curseur. Si \(\operatorname{opp}\) échange \(N\leftrightarrow S\) et
\(E\leftrightarrow O\), la transformation
\[
 F_{a,b}(z,d)=\bigl(z+a v_d,\operatorname{opp}^{\,b}(d)\bigr)
\]
admet pour inverse
\[
 F_{a,b}^{-1}(z',d')=
 \bigl(z'-a v_{\operatorname{opp}^{\,b}(d')},
       \operatorname{opp}^{\,b}(d')\bigr).
\]
La réduction spatiale modulo neuf entrelace cette transformation avec le
transport fini.
\end{proposition}
\begin{proof}
Appliquer deux fois l'opposition restaure la direction initiale. Une fois
cette direction récupérée, la soustraction de \(a v_d\) annule le déplacement.
Pour la projection spatiale, il suffit de l'égalité
\([z+a v_d]_9=[[z]_9+a v_d]_9\), composante par composante. L'inversion de
direction conserve cette égalité.
\end{proof}

La phase détermine les indicateurs utilisés par l'inverse. Le registre
spatial et l'arithmétique peuvent ainsi être conservés conjointement, même en
traversant le bord du représentant \(I_9^2\).

\subsection{Six fenêtres et émission décimale par blocs}
\label{act21:tpk:emision}

\subsubsection{Accumulation dans une fenêtre nonadique}

La fenêtre \(m\), avec \(1\leq m\leq6\), comprend les étapes
\(9m-8,\ldots,9m\). Soient \(S_m\) la somme de ses résidus positifs
additifs actifs, \(E_m\) la somme de ses exposants discrets actifs et
\(Z_m\) son nombre d'événements neutres.

\begin{lemma}[Multiplicités d'activation]
\label{act21:tpk:tres-eventos}
Chaque fenêtre contient trois activations additives, trois multiplicatives et
trois neutres. En particulier, \(Z_m=3\).
\end{lemma}
\begin{proof}
Les phases d'une fenêtre parcourent exactement \(1,\ldots,9\). Chacun
des trois ensembles de \eqref{act21:tpk:selector} a pour cardinal trois.
\end{proof}

Les six fenêtres forment \(54\) transitions et produisent six blocs.  
Cette longueur appartient à l'émetteur initial défini ici. La continuation  
en profondeur agit ensuite sur les mots et leurs registres compatibles.

\subsubsection{Règle de codification décimale}

La loi émettrice spécifie les chiffres
\begin{align}
 d_{m,1}&=[S_m]_{10},\notag\\
 d_{m,2}&=[S_m+E_m]_{10},\notag\\
 d_{m,3}&=[3S_m+5E_m+7Z_m]_{10},\notag\\
 U_m&=100d_{m,1}+10d_{m,2}+d_{m,3}.
 \label{act21:tpk:bloque-decimal}
\end{align}
Les poids \(100,10,1\) sont les positions des centaines, des dizaines et des unités.
L'indice \(10\) indique le reste décimal, et chaque bloc satisfait
\(0\leq U_m\leq999\). Le mot émis est la suite ordonnée
\(\mathbf U=(U_1,\ldots,U_6)\), de largeur fixe de trois chiffres par bloc.

Les coefficients \(3,5,7\) constituent les poids explicites de cette loi
émettrice. Une fois la règle fixée, chaque chiffre se calcule à partir des
accumulateurs, et les résultats de ce chapitre se déduisent de cette règle.
Le nombre \(1\) qui apparaît sous sa forme réduite a une origine
additionnelle concrète : \(7Z_m=21\equiv1\pmod{10}\).

Si \(s_m=[S_m]_{10}\) et \(e_m=[E_m]_{10}\), on obtient
\begin{equation}
 G(s_m,e_m)=100s_m+10[s_m+e_m]_{10}
                   +[3s_m+5e_m+1]_{10}=U_m.
 \label{act21:tpk:acoplamiento-firmas}
\end{equation}
La lettre \(e_m\) désigne ici une composante de la signature d'exposants discrets. Sa définition appartient à l'émetteur fini.

\begin{proposition}[Récupération des deux signatures]
\label{act21:tpk:recuperacion-firmas}
L'application \(G:\{0,\ldots,9\}^2\to\{0,\ldots,999\}\) est injective.
Pour un bloc \(U=100a+10b+c\) de son image, l'inverse est
\[
 s=a,\qquad e=[b-a]_{10},\qquad
 c=[3s+5e+1]_{10}.
\]
En conséquence, le mot \(\mathbf U\) détermine les deux signatures de six composantes \(\mathbf s\) et \(\mathbf e\).
\end{proposition}
\begin{proof}
Les deux derniers termes de \eqref{act21:tpk:acoplamiento-firmas} appartiennent à
\(\{0,\ldots,99\}\), de sorte que la centaine récupère \(s\). La dizaine
est \([s+e]_{10}\) ; en soustrayant \(s\) et en réduisant modulo dix, on récupère
\(e\). L'unité vérifie la troisième congruence. La reconstruction s'applique
indépendamment aux six positions du mot.
\end{proof}

Les totaux entiers \(S_m,E_m\) se retrouvent à partir des lectures enregistrées ;
les centaines et les dizaines retrouvent leurs restes. Par exemple, un total
\(S_m=15\) produit \(s_m=5\). Cette différence identifie exactement
quelle information chaque niveau de codage conserve.

\subsubsection{Codification positionnelle du mot de six blocs}

Les six blocs admettent une codification entière à largeur fixe,
\begin{equation}
 \mathscr N(\mathbf U)=\sum_{m=1}^{6}U_m1000^{6-m},
 \qquad0\leq\mathscr N(\mathbf U)<1000^6.
 \label{act21:tpk:entero-base1000}
\end{equation}
Le premier bloc occupe la position de plus haut poids et le sixième, celle des unités de base mille. Un bloc \(058\) occupe une position complète avec une valeur entière \(58\) ; sa largeur de trois chiffres fait partie de la présentation décimale de cette position.

\begin{proposition}[Récupération positionnelle des émissions]
\label{act21:tpk:inversa-base1000}
La codification \eqref{act21:tpk:entero-base1000} détermine le mot complet
par le biais de
\[
 U_m=\left[\left\lfloor
 \frac{\mathscr N(\mathbf U)}{1000^{6-m}}\right\rfloor\right]_{1000},
 \qquad1\leq m\leq6.
\]
\end{proposition}
\begin{proof}
En divisant par \(1000^{6-m}\), les blocs précédents produisent des multiples entiers de \(1000\), le bloc \(m\) produit \(U_m\), et la contribution des suivants appartient à \([0,1)\). La partie entière élimine cette dernière contribution et le résidu modulo \(1000\) élimine les premières.
\end{proof}

Cette codification conserve les six émissions. Son domaine possède des positions ordonnées de base mille, tandis que le domaine APP possède des positions spatiales modulo neuf. La composition maintient les deux indices : \(m\) identifie une fenêtre d'émission et \((i,j)\) identifie la cellule visitée lors d'une transition de cette fenêtre.

\subsection{Calcul complet d'une émission qui commence par 501}
\label{act21:tpk:ejemplo501}

On fixe les conditions initiales
\[
 c_{+,0}=(0,4;N),\qquad c_{\times,0}=(2,3;N),
 \qquad \nu_+=16,\quad\nu_\times=84.
\]
Les positions positives initiales sont \((1,5)\) et \((3,4)\). L'indice
chronologique précédent est \(t=0\), la première phase lue est \(g_1=1\) et le
registre est vide. Ces coordonnées sont les entrées
du calcul ; les blocs s'obtiennent en exécutant la récurrence.

\begin{table}[htbp]
\caption{Première fenêtre : lectures précédentes et accumulation active.}
\label{act21:tpk:tabla-501}
\small
\begin{tabular}{@{}rcccrrr@{}}
\toprule
\(t\)&\(\epsilon_t\)&position \(+\)&position \(\times\)&
\(\Delta S\)&\(\Delta E\)&\((S,E,Z)\)\\
\midrule
1&\(+1\)&\((1,5)\)&\((3,4)\)&6&0&\((6,0,0)\)\\
2&\(-1\)&\((9,5)\)&\((3,4)\)&0&0&\((6,0,0)\)\\
3&0&\((9,5)\)&\((2,4)\)&0&0&\((6,0,1)\)\\
4&\(+1\)&\((9,5)\)&\((2,4)\)&5&0&\((11,0,1)\)\\
5&\(-1\)&\((8,5)\)&\((2,4)\)&0&3&\((11,3,1)\)\\
6&0&\((8,5)\)&\((1,4)\)&0&0&\((11,3,2)\)\\
7&\(+1\)&\((8,5)\)&\((1,4)\)&4&0&\((15,3,2)\)\\
8&\(-1\)&\((7,5)\)&\((1,4)\)&0&2&\((15,5,2)\)\\
9&0&\((7,5)\)&\((9,4)\)&0&0&\((15,5,3)\)\\
\bottomrule
\end{tabular}
\notatabla{Les positions sont des étiquettes positives des cellules antérieures au mouvement. Les événements neutres augmentent uniquement \(Z\).}
\end{table}

Les trois évaluations additives actives sont
\[
 1+5=6,\qquad9+5=14=5+9,\qquad8+5=13=4+9,
\]
de telle sorte que \(S_1=6+5+4=15\). Les trois évaluations multiplicatives actives
sont \(3\cdot4=12\), \(2\cdot4=8\) et \(1\cdot4=4\). Leurs résidus
positifs sont \(3,8,4\), avec des observables \(0,3,2\) ; par conséquent,
\(E_1=5\). Enfin \(Z_1=3\). La règle décimale donne
\[
 d_{1,1}=[15]_{10}=5,\quad
 d_{1,2}=[15+5]_{10}=0,\quad
 d_{1,3}=[45+25+21]_{10}=1,
\]
et on obtient \(U_1=501\).

La seconde lecture additive précédente a pour quotient arithmétique \(1\),
tandis que son curseur relevé se trouve en \((-1,4;N)\), avec un nombre de tours
verticaux \(-1\). Ce sont deux données différentes du même registre. Toutes deux
interviennent dans la récupération intégrale de l'évaluation et de sa position.

\begin{table}[htbp]
\caption{Les six fenêtres de la condition initiale \((16,84)\).}
\label{act21:tpk:tabla-seis501}
\begin{tabular}{@{}rrrrrrrr@{}}
\toprule
\(m\)&pas&\(S_m\)&\(E_m\)&\(Z_m\)&\(s_m\)&\(U_m\)&\([-U_m]_3\)\\
\midrule
1&1--9&15&5&3&5&501&0\\
2&10--18&6&5&3&6&614&1\\
3&19--27&24&5&3&4&498&0\\
4&28--36&21&5&3&1&169&2\\
5&37--45&12&5&3&2&272&1\\
6&46--54&12&5&3&2&272&1\\
\bottomrule
\end{tabular}
\end{table}

Le mot résultant et ses signatures sont
\begin{align*}
 \mathbf U&=(501,614,498,169,272,272),\\
 \mathbf s&=(5,6,4,1,2,2),&
 \mathbf e&=(5,5,5,5,5,5).
\end{align*}
L'inversion cardinale après l'étape \(27\) modifie les parcours des
trois fenêtres finales. Le bloc répété \(272\) correspond à deux
fenêtres successives possédant leurs propres positions et leur propre registre chronologique.

Comme deuxième cas, la condition initiale minimale \((\nu_+,\nu_\times)=(0,0)\)
produit
\[
 \mathbf U=(252,109,252,903,850,850),\quad
 \mathbf s=(2,1,2,9,8,8),\quad
 \mathbf e=(3,9,3,1,7,7).
\]
Dans ce cas, les deux curseurs retournent à leur position et orientation initiales
après \(54\) transitions. La phase nonadique effectue six tours, le
compteur chronologique avance de \(54\) unités et le registre contient
\(54\) échantillons. La lecture positionnelle du retour et
l'histoire enregistrée décrivent des propriétés distinctes de l'exécution.

\subsection{Factorisation du recensement fini}
\label{act21:tpk:censo}

Les signatures \(f_+(\nu)=\mathbf s\) et \(f_\times(\nu)=\mathbf e\) se
calculent en exécutant le curseur correspondant pendant les six fenêtres.
L'indépendance des positions et des directions entre les deux curseurs donne
\begin{equation}
 T_{\mathrm{em}}(\nu_+,\nu_\times)
   =G^{(6)}\bigl(f_+(\nu_+),f_\times(\nu_\times)\bigr),
 \label{act21:tpk:factorizacion}
\end{equation}
où \(G^{(6)}\) applique \eqref{act21:tpk:acoplamiento-firmas} à chaque position.
La factorisation permet d'énumérer d'abord \(324+324\) trajectoires et
de combiner ensuite les signatures distinctes, en conservant la multiplicité de
leurs préimages.

\par\noindent\begin{minipage}{\linewidth}
\subsubsection{Images des deux signatures}

Les tableaux suivants comprennent toutes les signatures. Leur multiplicité est le
nombre de conditions initiales qui les produisent ; l'indice minimal localise
l'une de ces conditions. La préimage complète est définie par
\[
 f_\sigma^{-1}(a)=\{\nu\in\{0,\ldots,323\}:f_\sigma(\nu)=a\}.
\]
La récurrence des sections précédentes détermine chaque élément de cet
ensemble par un calcul entier fini.

\begin{table}[H]
\centering\fontsize{9.5}{11.4}\selectfont
\setlength{\abovecaptionskip}{8pt}\setlength{\belowcaptionskip}{6pt}
\renewcommand{\arraystretch}{1.12}
\caption{Les dix-huit signatures additives.}
\label{act21:tpk:firmas-aditivas}
\begin{tabular}{@{}lrr@{\qquad}lrr@{}}
\toprule
signature&multiplicité&\(\nu_{\min}\)&signature&multiplicité&\(\nu_{\min}\)\\
\midrule
\texttt{122564}&18&17&\texttt{564122}&18&16\\
\texttt{122898}&18&24&\texttt{645212}&18&4\\
\texttt{212645}&18&5&\texttt{654889}&18&33\\
\texttt{212988}&18&0&\texttt{889221}&18&13\\
\texttt{221456}&18&29&\texttt{889654}&18&32\\
\texttt{221889}&18&12&\texttt{898122}&18&25\\
\texttt{456221}&18&28&\texttt{898465}&18&20\\
\texttt{465898}&18&21&\texttt{988212}&18&1\\
\texttt{546988}&18&9&\texttt{988546}&18&8\\
\bottomrule
\end{tabular}

\par\vspace{12pt}

\caption{Les vingt-six signatures multiplicatives.}
\label{act21:tpk:firmas-multiplicativas}
\begin{tabular}{@{}lrr@{\qquad}lrr@{}}
\toprule
signature&multiplicité&\(\nu_{\min}\)&signature&multiplicité&\(\nu_{\min}\)\\
\midrule
\texttt{000000}&108&8&\texttt{555555}&24&84\\
\texttt{177393}&8&1&\texttt{555717}&8&14\\
\texttt{177555}&8&54&\texttt{555771}&8&18\\
\texttt{177771}&8&11&\texttt{555933}&8&24\\
\texttt{339555}&8&6&\texttt{717555}&8&12\\
\texttt{339771}&8&15&\texttt{717717}&8&21\\
\texttt{339933}&8&59&\texttt{717933}&8&19\\
\texttt{393177}&8&0&\texttt{771177}&8&9\\
\texttt{393393}&8&45&\texttt{771339}&8&13\\
\texttt{393555}&8&40&\texttt{771555}&8&16\\
\texttt{555177}&8&52&\texttt{933339}&8&57\\
\texttt{555339}&8&4&\texttt{933555}&8&26\\
\texttt{555393}&8&41&\texttt{933717}&8&17\\
\bottomrule
\end{tabular}
\end{table}
\end{minipage}\par


\begin{proposition}[Image et multiplicités de l'émetteur]
\label{act21:tpk:468}
L'image \(\mathcal U_6=\operatorname{im}T_{\mathrm{em}}\) contient
\(468\) mots distincts. Ses préimages se répartissent en \(432\)
fibres de cardinal \(144\), \(18\) de cardinal \(432\) et \(18\) de
cardinal \(1944\).
\end{proposition}
\begin{proof}
L'évaluation exhaustive de l'algorithme de la section ~\ref{act21:tpk:algoritmo} sur les indices \(0,\ldots,323\) donne les deux tableaux ci-dessus. La première partition a \(18\) classes de cardinal \(18\), et la seconde \(24\) classes de cardinal \(8\), une de cardinal \(24\) et une de cardinal \(108\). Les sommes \(18\cdot18=324\) et \(24\cdot8+24+108=324\) vérifient la masse totale des deux partitions. La proposition ~\ref{act21:tpk:recuperacion-firmas} rend injectif le couplage de ses images ; par conséquent, il y a \(18\cdot26=468\) émissions. Par \eqref{act21:tpk:factorizacion}, la préimage d'un couple de signatures est le produit de ses deux préimages. On obtient \(18\cdot24=432\) produits de cardinal \(18\cdot8=144\), \(18\) de cardinal \(18\cdot24=432\) et \(18\) de cardinal \(18\cdot108=1944\). Enfin,
\[
 432\cdot144+18\cdot432+18\cdot1944=104\,976.
\]
La partie énumérative est une vérification finie exacte de la récurrence ;
la factorisation explique pourquoi ce décompte couvre toutes les paires.
\end{proof}

L'émission de l'exemple \ref{act21:tpk:ejemplo501} a une fibre de cardinal
\(18\cdot24=432\). L'exemple d'indices \((0,0)\) a une fibre de
cardinal \(18\cdot8=144\). Dans les deux cas, le mot émis récupère
les signatures, tandis que l'identification d'une condition initiale particulière
utilise son indice ou son registre de trajectoire.

\subsection{Lecture ternaire orientée et conservation des fibres}
\label{act21:tpk:reduccion-ternaria}

La réduction orientée agit séparément sur chaque bloc :
\begin{equation}
 \Psi:\mathcal U_6\longrightarrow\mathbb F_3^6,\qquad
 \Psi(U_1,\ldots,U_6)=([-U_1]_3,\ldots,[-U_6]_3).
 \label{act21:tpk:psi}
\end{equation}
Son signe fixe une orientation du catalogue. L'identification équilibrée
\(0\mapsto0\), \(1\mapsto+1\), \(2\mapsto-1\) s'applique ensuite,
composante par composante.

Le module neuf organise les évaluations et positions d'APP ; les restes modulo dix forment chaque chiffre ; la base mille regroupe trois positions décimales ; le module trois détermine la signature ternaire de chaque bloc. Ces quatre usages ont des objets et des fonctions définis. En particulier, \(\Psi\) lit six blocs et produit six trits ; une conversion positionnelle de l'entier formé par dix-huit chiffres serait une autre opération.

Pour les deux conditions calculées, on obtient
\begin{align*}
 \Psi(501,614,498,169,272,272)&=(0,1,0,2,1,1),\\
 \Psi(252,109,252,903,850,850)&=(0,2,0,0,2,2).
\end{align*}
L'émetteur décimal est défini avant cette réduction. La fonction de
\(\Psi\) est de produire l'objet ternaire sur lequel agissent la symétrie
orbitale et les sélecteurs régionaux. Cette fonction explique sa position dans la
généalogie : elle relie une émission déjà calculée à une classification structurale.

L'image \(\mathcal W_6:=\Psi(\mathcal U_6)\) contient exactement
\(243\) mots au sein de l'ensemble ambiant de \(3^6=729\) mots.
L'énumération de l'algorithme final donne l'histogramme
\begin{equation}
\begin{array}{c|rrrrrr}
 |\Psi^{-1}(w)|&1&2&3&4&5&6\\\hline
 \text{nombre de mots }w&132&66&18&3&6&18.
\end{array}
\label{act21:tpk:fibras-ternarias}
\end{equation}
Les sommes des deux lectures sont
\[
132+66+18+3+6+18=243,
\qquad132+2\cdot66+3\cdot18+4\cdot3+5\cdot6+6\cdot18=468.
\]
La première compte les mots ternaires ; la seconde récupère le nombre d'émissions réparties entre ses fibres.

\begin{example}[Deux émissions avec la même signature ternaire]
\label{act21:tpk:colision}
Les émissions
\[
 (498,501,614,252,252,109),\qquad
 (870,810,923,252,252,109)
\]
sont distinctes et toutes deux ont pour image \((0,0,1,0,0,2)\) par \(\Psi\).
Chacune a \(144\) conditions initiales. Le mot ternaire conserve
une classe de lecture ; les émissions et leurs préimages distinguent les données
qui convergent vers cette classe.
\end{example}

La conservation des fibres s'exprime concrètement par
\[
 \mathcal F_w=\{\mathbf U\in\mathcal U_6:\Psi(\mathbf U)=w\},
 \qquad
 \Omega_{\mathbf U}=T_{\mathrm{em}}^{-1}(\mathbf U).
\]
Ainsi, la préimage complète de \(w\) dans le domaine initial est l'union disjointe des ensembles \(\Omega_{\mathbf U}\), avec \(\mathbf U\in\mathcal F_w\). Le registre enrichi peut conserver l'émission, ses signatures et la condition initiale ainsi que sa lecture ternaire. L'inverse exacte se réfère alors à la donnée enregistrée avec son type ; le mot de six trits pris isolément conserve la portée précise de \eqref{act21:tpk:psi}.

\subsection{Action diédrale et les quarante-trois orbites}
\label{act21:tpk:orbitas}

Sur six coordonnées se définissent les permutations
\begin{align}
 r(u_1,u_2,u_3,u_4,u_5,u_6)
   &=(u_3,u_1,u_2,u_5,u_6,u_4),\\
 s(u_1,u_2,u_3,u_4,u_5,u_6)
   &=(u_4,u_5,u_6,u_1,u_2,u_3).
 \label{act21:tpk:generadores-diedrales}
\end{align}
La première fait tourner les deux groupes de trois coordonnées dans des sens conjugués ;
la seconde échange les deux groupes. Leur composition satisfait
\(r^3=s^2=1\) et \(srs=r^{-1}\). Elles engendrent ainsi une action du groupe diédral
\(D_3\), à six éléments.

\begin{proposition}[Équivariance et stabilité du catalogue]
\label{act21:tpk:equivarianza}
La réduction \(\Psi\) commute avec l'action de \(D_3\). Les ensembles
\(\mathcal U_6\) et \(\mathcal W_6\) sont stables sous leurs générateurs.
\end{proposition}
\begin{proof}
L'égalité \(\Psi(g\mathbf U)=g\Psi(\mathbf U)\) se vérifie en chaque
coordonnée parce que \(g\) permute les positions et \(\Psi\) applique la même
réduction à chaque position. La stabilité du catalogue constitue une autre
vérification : en générant ses \(468\) éléments par
\eqref{act21:tpk:factorizacion}, on vérifie que \(r\mathbf U\) et
\(s\mathbf U\) appartiennent à nouveau à cet ensemble. L'algorithme de la
section~\ref{act21:tpk:algoritmo} effectue ces vérifications d'appartenance.
L'équivariance transmet ensuite la stabilité à son image ternaire.
\end{proof}

\begin{proposition}[Décomposition orbitale complète]
\label{act21:tpk:43}
L'ensemble \(\mathcal W_6\) se décompose en \(43\) orbites : \(38\)
de cardinal six et \(5\) de cardinal trois.
\end{proposition}
\begin{proof}
Pour chaque \(w\in\mathcal W_6\) on forme
\[
 \mathcal O(w)=\{w,rw,r^2w,sw,srw,sr^2w\}.
\]
On choisit comme représentant le mot lexicographiquement minimal de cet ensemble. Deux mots ont le même représentant exactement lorsqu'ils appartiennent à la même orbite. L'algorithme génère la liste des représentants du tableau~\ref{act21:tpk:representantes} ; la somme de leurs cardinaux vaut \(38\cdot6+5\cdot3=243\). Comme chaque mot du catalogue participe à cette opération et que le catalogue est stable, la liste couvre toute l'image.
\end{proof}

\begin{table}[htbp]
\caption{Représentants minimaux et cardinaux des quarante-trois orbites.}
\label{act21:tpk:representantes}
\small
\begin{tabular}{@{}lr@{\qquad}lr@{\qquad}lr@{}}
\toprule
représentant&cardinal&représentant&cardinal&représentant&cardinal\\
\midrule
\texttt{001002}&6&\texttt{002112}&6&\texttt{012222}&6\\
\texttt{001011}&6&\texttt{002211}&6&\texttt{021022}&6\\
\texttt{001012}&6&\texttt{002220}&6&\texttt{021121}&6\\
\texttt{001021}&6&\texttt{011021}&6&\texttt{021202}&6\\
\texttt{001022}&6&\texttt{011112}&6&\texttt{022022}&3\\
\texttt{001100}&3&\texttt{011120}&6&\texttt{022112}&6\\
\texttt{001101}&6&\texttt{011201}&6&\texttt{022121}&6\\
\texttt{001110}&6&\texttt{011211}&6&\texttt{022202}&3\\
\texttt{001112}&6&\texttt{011220}&6&\texttt{022211}&6\\
\texttt{001202}&6&\texttt{011222}&6&\texttt{022222}&6\\
\texttt{001210}&6&\texttt{012112}&6&\texttt{112112}&3\\
\texttt{001211}&6&\texttt{012121}&6&\texttt{112211}&3\\
\texttt{001220}&6&\texttt{012202}&6&\texttt{112222}&6\\
\texttt{001222}&6&\texttt{012210}&6&&\\
\texttt{002012}&6&\texttt{012220}&6&&\\
\bottomrule
\end{tabular}
\end{table}

Les orbites héritent des invariants entiers de leurs fibres. Pour chaque mot,
\(n(w)=|\mathcal F_w|\) compte les émissions, tandis que
\[
 M(w)=\sum_{\mathbf U\in\mathcal F_w}|\Omega_{\mathbf U}|
\]
compte les conditions initiales. Le recensement des symboles
\(c(w)=(n_0,n_1,n_2)\) enregistre combien de fois chaque classe ternaire apparaît.
Ces données, conjointement avec l'orientation et le support des parcours,
fournissent les prédicats de la sélection régionale ultérieure. La
classification conserve ainsi plus d'information que la forme d'un mot
représentant.

\subsection{Algorithme autonome d'énumération et de contrôle}
\label{act21:tpk:algoritmo}

L'énumération utilise les entiers, les listes finies et les opérations déjà définies.  
Le procédé suivant précise son exécution ; \(\sigma\) prend les valeurs  
\(+\) et \(\times\).

\begin{enumerate}
\item Pour chaque \(\nu=0,\ldots,323\), récupérer \((i,j;d)\) par
l'inverse de \eqref{act21:tpk:indice-semilla}. Initialiser une liste vide de signatures.
\item Pour chaque fenêtre \(m=1,\ldots,6\), initialiser l'accumulateur \(A=0\).
Parcourir \(t=9m-8,\ldots,9m\), calculant \(\epsilon_t\) par
\eqref{act21:tpk:fase}.
\item Si \(\sigma=+\) et \(\epsilon_t=+1\), ajouter
\(\rho_9(i+j+2)\) à \(A\) et mettre à jour
\((i,j)\leftarrow T_d(i,j)\). Si \(\sigma=\times\) et
\(\epsilon_t=-1\), ajouter
\(\ell_9(\rho_9((i+1)(j+1)))\) et effectuer le même transport.
\item À la fin de \(t=27\) ou \(t=54\), remplacer \(d\) par son opposée.  
À la conclusion de la fenêtre, ajouter \([A]_{10}\) à la liste des signatures.  
Conserver la position et la direction pour la fenêtre suivante.
\item Grouper les \(324\) indices par égalité de leurs listes de six
composants. Ces deux groupements donnent \(f_+\), \(f_\times\) et toutes
leurs préimages.
\item Pour chaque paire de signatures distinctes, appliquer \(G^{(6)}\).
Attribuer à l'émission la multiplicité produit des deux préimages.
Vérifier la récupération des signatures dans chaque bloc et la somme des
multiplicités \(104\,976\).
\item Appliquer \(\Psi\) à chaque émission et regrouper par égalité de mots.
Compter les \(243\) valeurs et les multiplicités de
\eqref{act21:tpk:fibras-ternarias}.
\item Appliquer \(r,s\) à chaque émission et vérifier son appartenance au catalogue.
Pour chaque mot ternaire, former ses six images diédriques, éliminer les
répétitions et enregistrer son minimum lexicographique. Compter les \(43\)
orbites et les deux tailles du tableau~\ref{act21:tpk:representantes}.
\end{enumerate}

L'exécution se termine : le premier tronçon évalue \(648\) trajectoires de \(54\) transitions ; le deuxième combine \(468\) couples de signatures ; le troisième regroupe un ensemble de \(243\) mots. Une vérification directe supplémentaire peut exécuter les \(104\,976\) couples et vérifier la factorisation \eqref{act21:tpk:factorizacion} ; sa portée coïncide avec le même domaine fini.

\subsection{Conservation du registre et fonction des germes}
\label{act21:tpk:conclusion}

Soit \(\mathsf{Sample}(q_t)\) l'échantillon antérieur à la transition et \(\mathcal M_t\) le registre accumulé. La mise à jour de la mémoire est
\begin{equation}
 \mathcal M_{t+1}=\mathcal M_t\mathbin{\Vert}\mathsf{Sample}(q_t),
 \label{act21:tpk:archivo}
\end{equation}
où \(\Vert\) désigne la concaténation ordonnée. La récurrence donne
\[
 \mathcal M_n=\mathcal M_0\mathbin{\Vert}
 (\mathsf{Sample}(q_0),\ldots,\mathsf{Sample}(q_{n-1})).
\]
Les accumulateurs se reconstruisent en additionnant les contributions actives de
chaque segment de neuf enregistrements. Leurs quotients, les tours des
curseurs et les orientations restent associés aux mêmes transitions.
L'exécution détermine ce registre au moyen de la récurrence ; le fichier
préserve la provenance concrète des lectures.

À la fin de cette étape sont construits le domaine complet des conditions
initiales, l'émission décimale, ses deux signatures récupérables, l'image ternaire,
ses fibres et sa classification orbitale. Les recensements \(104\,976\), \(468\),
\(243\), \(729\) et \(43\) ont désormais un domaine et un mécanisme d'obtention.
Les sélections régionales recevront ces structures avec leurs invariants,
et les élévations agiront sur des mots orientés dont la provenance est
conservée. Cette continuité fait du catalogue fini un antécédent
constructif du développement coinductif.

% END NUCLEO_ADITIVO_20260921_CENSO
\import{foundation/}{sections/extension.tex}
\import{foundation/}{sections/continuo_conjunto.tex}
\import{foundation/}{sections/generacion.tex}
% BEGIN NUCLEO_ADITIVO_20260921_SELECCION
\clearpage
% Ampliación demostrativa DF003--DF007. No modifica la fuente integral.
% Propietario común de los archivos Lean citados en estos comentarios:
% output/PAQUETE_ARTICULO_I_INTEGRACION_ESTADO_CAMPOS_20260919/.
% Residencia causal de la primera sección: después de los lectores regionales
% coinductivos y antes de emplear sus registros como entrada de la selección K.
% Residencia causal de la segunda sección: después de la construcción marcada
% Paley--Witt--Golay--Leech y antes de la reconstrucción de la estructura VOA.
% Los dos cortes expositivos conservan el mismo origen APP--TRIT--TPK.

\section[Génération régionale et raffinement]{Compatibilité intégrale de la génération régionale, du raffinement cylindrique et de l'incidence}
\label{act21:sec:df-lectores-cilindros}

La prolongation régionale produit des suites ordonnées de blocs, ainsi que
leurs profondeurs de lecture. La projection archimédienne et l'incidence
ternaire se calculent sur ces mêmes suites. Le lien entre ces deux
réalisations peut s'exprimer entièrement par des opérations entières :
sélection par intervalles rationnels, division euclidienne, élévation de six
trits, lecture de retenue et transformation de Witt. Cette formulation permet
de suivre chaque coefficient depuis son bloc générateur jusqu'à son emploi ultérieur.

La construction utilise les trois lecteurs régionaux déjà obtenus de l'état APP--TRIT--TPK. Nous les désignerons par les indices \(c\in\{\mathrm P,\mathrm E,\mathrm F\}\), correspondant respectivement aux régions de clôture, de propagation et d'auto-échelle. Leur réalisation réelle s'écrira \(v_c\). Les identifications ultérieures \(v_{\mathrm P}=\pi_{\mathrm{HMT}}\), \(v_{\mathrm E}=e_{\mathrm{HMT}}\) et \(v_{\mathrm F}=\varphi_{\mathrm{HMT}}\) conservent cet ordre causal. L'opération qui émet un bloc inspecte les intervalles rationnels du lecteur régional ; la valeur réelle intervient dans la démonstration de correction de cette opération.

\subsection{Émission positionnelle par séparation rationnelle de cellules}
\label{act21:subsec:df-emision-racional}

% Fuente: lean/biblioteca/RegionalPublicationComposition.lean,
% produced_brackets, eventual_publication, joint_eventual_publication,
% search_terminates, stopped_cell_correct, publications_compatible.

Pour chaque région, on dispose de deux suites rationnelles
\(\ell_c(j),u_c(j)\), construites par son lecteur, qui satisfont
\[
 \ell_c(j)\leq v_c\leq u_c(j).
\]
La propriété de séparation des cellules démontrée pour ces lecteurs affirme
que, pour chaque entier \(s>0\), il existe \(J\) tel que, pour tout \(j\geq J\),
\begin{equation}
 \lfloor s\ell_c(j)\rfloor
 =\lceil su_c(j)\rceil-1
 =\lfloor sv_c\rfloor.
 \label{act21:eq:df-separacion-celdas}
\end{equation}
L'évitement des frontières rationnelles par les valeurs régionales est l'hypothèse des théorèmes précédents qui garantit cette séparation. Ici, on emploie sa conclusion opérationnelle, avec les lecteurs et les intervalles qui la produisent.

\begin{definition}[Indice de séparation et préfixe émis]
Pour \(s>0\), soit
\[
 j_c(s)=\min\{j\in\mathbb N:
 \lfloor s\ell_c(j)\rfloor=\lceil su_c(j)\rceil-1\}.
\]
Pour une base entière \(b\geq2\) et une profondeur \(n\geq0\), on définit
\begin{equation}
 P_c(b,n)=\lfloor b^n\ell_c(j_c(b^n))\rfloor.
 \label{act21:eq:df-publicacion-prefijo}
\end{equation}
Le préfixe contient la partie entière et les premières \(n\) positions en base \(b\) ; par conséquent, \(P_c(b,0)\) est la partie entière régionale.
\end{definition}

\begin{theorem}[Terminaison et correction de l'émission]
\label{act21:thm:df-emision-correcta}
L'indice \(j_c(s)\) existe pour tout \(s>0\). Les préfixes émis satisfont
\begin{equation}
 P_c(b,n)=\lfloor b^n v_c\rfloor,
 \qquad
 P_c(b,n)\leq b^n v_c<P_c(b,n)+1.
 \label{act21:eq:df-prefijo-correcto}
\end{equation}
Pour chaque échelle \(s\) il est possible en outre de choisir un unique indice de suffisance \(J(s)\) valide simultanément pour les trois régions.
\end{theorem}

\begin{proof}
L'égalité à partir d'un certain rang de \eqref{act21:eq:df-separacion-celdas} rend non vide
l'ensemble qui définit \(j_c(s)\). Son minimum existe par le bon ordre de
\(\mathbb N\). À tout indice où l'égalité se produit, la borne
inférieure rationnelle fournit
\(\lfloor s\ell_c(j)\rfloor\leq\lfloor sv_c\rfloor\).
La borne supérieure, jointe à l'évitement d'une frontière de cellule, donne
\(\lfloor sv_c\rfloor\leq\lceil su_c(j)\rceil-1\).
Les deux extrémités entières coïncidant, l'entier intermédiaire coïncide avec
elles. Cela démontre \eqref{act21:eq:df-prefijo-correcto}. Pour l'affirmation
conjointe, il suffit de prendre le maximum des trois indices de suffisance obtenus
pour la même échelle. L'indice conjoint contrôle simultanément les trois
émissions, tandis que chaque région conserve ses propres intervalles.
\end{proof}

\begin{proposition}[Compatibilité exacte des préfixes]
\label{act21:prop:df-prefijos-compatibles}
Pour \(n,m\geq0\),
\begin{equation}
 \left\lfloor\frac{P_c(b,n+m)}{b^m}\right\rfloor=P_c(b,n).
 \label{act21:eq:df-truncamiento}
\end{equation}
En particulier, le bloc suivant
\(d_c(b,n)=P_c(b,n+1)\bmod b\) satisfait
\begin{equation}
 P_c(b,n+1)=bP_c(b,n)+d_c(b,n),\qquad 0\leq d_c(b,n)<b.
 \label{act21:eq:df-paso-posicional}
\end{equation}
\end{proposition}

\begin{proof}
Écrivons \(b^n v_c=N+\theta\), avec
\(N=P_c(b,n)\) et \(0\leq\theta<1\). Alors
\(P_c(b,n+m)=b^mN+\lfloor b^m\theta\rfloor\), et le second terme
appartient à \(\{0,\ldots,b^m-1\}\). La division euclidienne par \(b^m\)
récupère \(N\). La spécialisation \(m=1\) produit
\eqref{act21:eq:df-paso-posicional} et sa borne résiduelle.
\end{proof}

\subsection{Blocs régionaux de six trits à profondeur arbitraire}
\label{act21:subsec:df-bloques-arbitrarios}

% Fuente: deltas/integracion/JointRegionalFrontier.lean,
% publish_base729_eq_base3, generated_block_from_longer_prefix,
% generated_block_eq_finite, generated_signature_eq_finite.

L'égalité entière \(729=3^6\) détermine le regroupement de six positions ternaires en un bloc. Pour \(t\geq0\), nous définissons
\begin{equation}
 u_c(t)=P_c(729,t+1)\bmod729,
 \qquad
 B_{c,j}(t)=
 \left\lfloor\frac{u_c(t)}{3^{5-j}}\right\rfloor\bmod3,
 \quad 0\leq j<6.
 \label{act21:eq:df-bloque-y-banda}
\end{equation}
Ainsi, \(B_c(t)\) est un mot ordonné de six trits et
\begin{equation}
 u_c(t)=\sum_{j=0}^{5}B_{c,j}(t)3^{5-j}.
 \label{act21:eq:df-elevacion-banda}
\end{equation}
Le mot et son élévation entière sont deux représentations mutuellement récupérables du même bloc émis.

\begin{theorem}[Stabilité d'extraction à toute profondeur]
\label{act21:thm:df-bloque-estable}
On vérifie, pour tout \(n\),
\[
 P_c(729,n)=P_c(3,6n).
\]
Si \(t<n\), alors
\begin{equation}
 u_c(t)=
 \left\lfloor\frac{P_c(729,n)}{729^{n-t-1}}\right\rfloor\bmod729.
 \label{act21:eq:df-bloque-prefijo-largo}
\end{equation}
En conséquence, toute extension du préfixe préserve tous les
blocs précédemment émis et toutes les signatures calculées dessus.
\end{theorem}

\begin{proof}
La première égalité résulte de ce que les deux préfixes se séparent à la même
échelle, \(729^n=3^{6n}\), et du
théorème~\ref{act21:thm:df-emision-correcta}. Pour la seconde, nous appliquons
\eqref{act21:eq:df-truncamiento} aux profondeurs \(n\) et \(t+1\), puis
nous prenons le résidu modulo \(729\). L'extraction ternaire de
\eqref{act21:eq:df-bloque-y-banda} est déterministe. Toute fonction explicite
de ces dix-huit coordonnées conserve, par conséquent, le même résultat en
étendant le préfixe.
\end{proof}

La construction est quantifiée sur \(t\in\mathbb N\). Une table qui
matérialise cent blocs constitue une évaluation finie de cette
construction ; la formule \eqref{act21:eq:df-bloque-prefijo-largo} détermine
exactement sa compatibilité avec toute extension ultérieure.

\begin{example}[Premier bloc conjoint]
L'évaluation des lecteurs régionaux sélectionnés donne les mots
\[
 B_{\mathrm P}(0)=010211,\qquad
 B_{\mathrm E}(0)=201101,\qquad
 B_{\mathrm F}(0)=121200.
\]
Leur élévation par \eqref{act21:eq:df-elevacion-banda} produit
\[
 u_{\mathrm P}(0)=103,\qquad
 u_{\mathrm E}(0)=523,\qquad
 u_{\mathrm F}(0)=450.
\]
Par exemple,
\(201101_3=2\cdot243+27+9+1=523\).
Les six positions, y compris les zéros initiaux, sont conservées dans la
représentation en bande. Le zéro initial de \(010211\) possède une position
définie, bien que l'écriture décimale de \(103\) en soit dépourvue.
\end{example}

\subsection{Registres de transition et reconstruction des élévations linéaires}
\label{act21:subsec:df-registros-transicion}

% Fuentes: deltas/integracion/RegionalTransitionRecords.lean y
% GeneratedTransitionRecords.lean; lean/biblioteca/TPKLifts.lean.
% Los enlaces reconstruidos son R(0)->R(1) y R(2)->R(3).

Sur \(\mathbb F_3\), nous réunissons deux émissions consécutives de chaque région dans la matrice
\begin{equation}
 R(t)=
 \begin{pmatrix}
 B_{\mathrm P}(t)\\ B_{\mathrm P}(t+1)\\
 B_{\mathrm E}(t)\\ B_{\mathrm E}(t+1)\\
 B_{\mathrm F}(t)\\ B_{\mathrm F}(t+1)
 \end{pmatrix}.
 \label{act21:eq:df-registro-matricial}
\end{equation}
Cette définition fixe à la fois l'ordre régional et l'ordre temporel. Les
quatre premiers registres sont
\[
 X_0=R(0),\quad Y_0=R(1),\quad X_1=R(2),\quad Y_1=R(3).
\]
Ses valeurs, obtenues par extraction des préfixes régionaux, sont
\begin{align*}
 X_0&=\begin{pmatrix}
0&1&0&2&1&1\\0&1&2&2&2&2\\2&0&1&1&0&1\\
1&2&1&2&2&1\\1&2&1&2&0&0\\1&1&2&2&0&2
\end{pmatrix},&
 Y_0&=\begin{pmatrix}
0&1&2&2&2&2\\0&1&0&2&1&1\\1&2&1&2&2&1\\
1&0&2&0&1&1\\1&1&2&2&0&2\\1&2&1&0&2&0
\end{pmatrix},\\[1ex]
 X_1&=\begin{pmatrix}
0&1&0&2&1&1\\0&0&2&1&1&1\\1&0&2&0&1&1\\
0&1&2&2&2&2\\1&2&1&0&2&0\\0&1&0&2&1&0
\end{pmatrix},&
 Y_1&=\begin{pmatrix}
0&0&2&1&1&1\\1&1&0&2&2&1\\0&1&2&2&2&2\\
1&0&2&0&1&1\\0&1&0&2&1&0\\0&1&0&2&0&0
\end{pmatrix}.
\end{align*}

\begin{theorem}[Reconstruction unique des deux liens initiaux]
\label{act21:thm:df-elevaciones-reconstruidas}
Les équations \(X_0L_0=Y_0\) et \(X_1L_1=Y_1\) ont, chacune, une
unique solution dans \(M_6(\mathbb F_3)\). Ces solutions sont
\begin{align*}
 L_0&=\begin{pmatrix}
2&2&2&1&2&1\\2&1&2&2&1&1\\1&0&0&1&0&2\\
0&0&1&1&1&0\\2&2&2&0&2&1\\2&1&2&1&0&0
\end{pmatrix},&
 L_1&=\begin{pmatrix}
2&1&2&0&2&2\\1&2&1&1&2&0\\1&2&1&1&1&2\\
0&1&0&0&2&1\\2&0&2&1&0&1\\0&2&2&2&1&1
\end{pmatrix}.
\end{align*}
Les deux matrices sont inversibles.
\end{theorem}

\begin{proof}
Les matrices
\begin{align*}
 A_0&=\begin{pmatrix}
1&1&2&0&1&2\\0&0&1&0&0&1\\2&1&2&1&0&1\\
0&2&0&1&0&2\\2&1&1&0&2&2\\2&1&1&1&1&2
\end{pmatrix},&
 A_1&=\begin{pmatrix}
1&0&1&2&0&0\\0&0&1&1&2&2\\0&1&2&0&1&1\\
1&1&0&2&0&0\\1&1&2&1&1&2\\1&0&0&0&0&2
\end{pmatrix}
\end{align*}
satisfont \(A_iX_i=X_iA_i=I\), par multiplication dans \(\mathbb F_3\).
Par conséquent, les solutions sont imposées par \(L_i=A_iY_i\) ; la
multiplication produit les deux matrices montrées. Leurs inverses sont
\begin{align*}
 L_0^{-1}&=\begin{pmatrix}
1&2&0&2&0&2\\2&0&2&2&0&0\\2&1&2&0&2&0\\
1&0&0&0&2&0\\0&2&1&1&2&0\\2&2&2&2&2&2
\end{pmatrix},&
 L_1^{-1}&=\begin{pmatrix}
2&0&2&1&2&1\\2&1&0&0&2&0\\2&0&2&2&0&2\\
2&0&0&1&1&0\\1&1&1&1&1&0\\2&0&1&2&2&0
\end{pmatrix}.
\end{align*}
La vérification bilatérale \(L_iL_i^{-1}=L_i^{-1}L_i=I\) démontre la
récupération exacte de chaque mot transformé. Tous les produits
utilisent des résidus \(0,1,2\), avec leur interprétation en \(\mathbb F_3\).
\end{proof}

La lecture causale de ce résultat est précise : les émissions régionales produisent \(R(t)\) ; les registres \(R(0),\ldots,R(3)\) déterminent les deux élévations précédentes ; leurs inverses récupèrent les mots transformés.  
La continuité pour tout \(t\) appartient aux lecteurs coinductifs et au théorème ~\ref{act21:thm:df-bloque-estable}. La reconstruction matricielle démontrée ici a pour domaine les deux liens spécifiés dans son énoncé.

\subsection{Raffinement simultané dans les bases ternaire et décimale}
\label{act21:subsec:df-cilindros-mixtos}

% Fuente: deltas/integracion/RegionalCylinderTransition.lean.

Une émission ternaire de six trits et une émission décimale de trois chiffres
constituent deux lectures de bases différentes. Pour conserver leur relation, on
enregistre l'état
\[
 s=(n,k,N,D),
\]
où \(N\) est un préfixe de profondeur \(n\) en base \(729\), et \(D\) est
un préfixe de profondeur \(k\) en base \(1000\). Leurs cylindres
archimédiens sont
\begin{equation}
 I_{729}(s)=\left[\frac{N}{729^n},\frac{N+1}{729^n}\right),\qquad
 I_{1000}(s)=\left[\frac{D}{1000^k},\frac{D+1}{1000^k}\right).
 \label{act21:eq:df-cilindros}
\end{equation}
L'extrémité supérieure semi-ouverte attribue chaque valeur à une unique cellule d'une base et d'une profondeur données.

\begin{definition}[Défauts entiers de compatibilité]
On définit
\begin{align}
 A(s)&=N1000^k-D729^n,\label{act21:eq:df-defecto-inferior}\\
 B(s)&=(N+1)1000^k-D729^n=A(s)+1000^k.
 \label{act21:eq:df-defecto-superior}
\end{align}
L'état est compatible lorsque
\begin{equation}
 A(s)<729^n,\qquad B(s)>0.
 \label{act21:eq:df-compatibilidad}
\end{equation}
\end{definition}

\begin{lemma}[Critère entier d'intersection]
Les cylindres de \eqref{act21:eq:df-cilindros} ont une intersection non vide si et
seulement si \eqref{act21:eq:df-compatibilidad} est vérifiée.
\end{lemma}

\begin{proof}
Deux intervalles semi-ouverts \([a,b)\) et \([c,d)\), avec \(a<b\) et
\(c<d\), se coupent exactement lorsque \(a<d\) et \(c<b\). En appliquant cette
condition à \eqref{act21:eq:df-cilindros} et en multipliant par les dénominateurs
positifs, nous obtenons
\(N1000^k<(D+1)729^n\) et
\(D729^n<(N+1)1000^k\). Ce sont les deux inégalités de
\eqref{act21:eq:df-compatibilidad}.
\end{proof}

\begin{definition}[Mise à jour des profondeurs et des préfixes]
Pour \(0\leq u<729\) et \(0\leq d<1000\), nous définissons
\begin{align*}
 T_u(n,k,N,D)&=(n+1,k,729N+u,D),\\
 T_{u,d}(n,k,N,D)&=(n+1,k+1,729N+u,1000D+d).
\end{align*}
La première opération raffine uniquement le cylindre ternaire. La seconde  
raffine les deux lectures et conserve les deux profondeurs dans l'état.
\end{definition}

\begin{proposition}[Récurrences exactes de frontière]
\label{act21:prop:df-rec-frontera}
Les mises à jour précédentes satisfont
\begin{align}
 A(T_u s)&=729A(s)+u1000^k,\label{act21:eq:df-rec-ternaria}\\
 A(T_{u,d}s)&=729000A(s)+u1000^{k+1}-d729^{n+1}.
 \label{act21:eq:df-rec-conjunta}
\end{align}
\end{proposition}

\begin{proof}
Pour la première identité,
\[
 (729N+u)1000^k-D729^{n+1}
 =729(N1000^k-D729^n)+u1000^k.
\]
Pour la seconde,
\begin{align*}
 &(729N+u)1000^{k+1}-(1000D+d)729^{n+1}\\
 &\qquad=729000(N1000^k-D729^n)
      +u1000^{k+1}-d729^{n+1}.
\end{align*}
Les deux sont des identités entières antérieures à toute approximation réelle.
\end{proof}

\begin{theorem}[Raffinement des cylindres régionaux générés]
\label{act21:thm:df-refinamiento-generado}
Soit
\[
 s_c(n,k)=(n,k,P_c(729,n),P_c(1000,k)).
\]
Alors \(s_c(n,k)\) est compatible quels que soient \(n,k\). De plus,
\begin{align*}
 T_{u_c(n)}s_c(n,k)&=s_c(n+1,k),\\
 T_{u_c(n),d_c(1000,k)}s_c(n,k)&=s_c(n+1,k+1).
\end{align*}
La division euclidienne des préfixes raffinés récupère exactement les
préfixes précédents. Pour tout \(a,b\geq0\),
\[
 \frac{P_c(729,n+a)-P_c(729,n+a)\bmod729^a}{729^a}
 =P_c(729,n),
\]
et l'identité analogue est vérifiée en base \(1000\).
\end{theorem}

\begin{proof}
Le théorème~\ref{act21:thm:df-emision-correcta} place \(v_c\) simultanément dans
les deux intervalles \eqref{act21:eq:df-cilindros} ; le critère d'intersection donne
la compatibilité. Les identités de mise à jour sont
\eqref{act21:eq:df-paso-posicional} dans chaque base. La récupération s'obtient à partir de
\eqref{act21:eq:df-truncamiento}. Par conséquent, la récurrence du défaut et la
récupération du préfixe sont des conséquences d'une même mise à jour.
\end{proof}

\begin{example}[Premier couple de cylindres régionaux]
Pour \(n=k=1\), l'évaluation des trois régions donne
\[
\begin{array}{c|rr|rr}
 c&N&D&A&B\\\hline
 \mathrm P&2290&3141&211&1211\\
 \mathrm E&1981&2718&-422&578\\
 \mathrm F&1179&1618&-522&478
\end{array}
\]
Dans les trois lignes \(A<729\) et \(B>0\). La première colonne de préfixes
se décompose comme \(2290=3\cdot729+103\),
\(1981=2\cdot729+523\) et \(1179=729+450\) : réapparaissent exactement les
blocs de six trits déjà émis. Les signes différents de \(A\) expriment
la position relative des deux extrémités inférieures, sans altérer la
appartenance commune de la valeur régionale.
\end{example}

\subsection{Sélection du bloc suivant par le lecteur régional}

La compatibilité de deux intervalles exprime une relation entre lectures. La détermination du bloc suivant utilise en outre l'intervalle rationnel produit par la région. Cette opération peut se formuler sans éliminer l'histoire de profondeur.

\begin{definition}[Admissibilité de la prolongation régionale]
\(c,n\) étant fixés, soient \(s=729^{n+1}\) et \(j=j_c(s)\). Un résidu
\(u\in\{0,\ldots,728\}\) est admissible lorsque
\begin{equation}
 \lfloor s\ell_c(j)\rfloor
 \leq729P_c(729,n)+u
 \leq\lceil su_c(j)\rceil-1.
 \label{act21:eq:df-admisibilidad}
\end{equation}
\end{definition}

\begin{theorem}[Unicité de la prolongation compatible avec le lecteur]
\label{act21:thm:df-prolongacion-unica}
La condition \eqref{act21:eq:df-admisibilidad} est équivalente à \(u=u_c(n)\).
Par conséquent, il existe exactement une prolongation admissible de chaque région
à chaque profondeur.
\end{theorem}

\begin{proof}
À l'indice de séparation, les deux extrémités entières de
\eqref{act21:eq:df-admisibilidad} coïncident avec \(P_c(729,n+1)\).
La condition est ainsi équivalente à
\(729P_c(729,n)+u=P_c(729,n+1)\).
Par \eqref{act21:eq:df-paso-posicional}, le côté droit est
\(729P_c(729,n)+u_c(n)\). L'annulation du terme commun donne l'égalité des résidus. Réciproquement, cette égalité satisfait la condition.
\end{proof}

La dépendance effective de cette sélection reste complètement déterminée :
région sélectionnée, lecteur rationnel, profondeur, indice de séparation,
préfixe antérieur et résidu suivant. Le défaut \(A\) résume une relation
de frontière ; l'état \((n,k,N,D)\) et les intervalles régionaux conservent
les données employées par la mise à jour.

\subsection{Signature ternaire, retenue, incidence de Witt et horloge de lecture}
\label{act21:subsec:df-firma-conjunta}

% Fuentes: lean/regional/N69RegionalSignature.lean,
% deltas/integracion/JointRegionalFrontier.lean y JointCylinderSignature.lean.

Une profondeur \(t\) étant fixée, écrivons \(B_{c,j}=B_{c,j}(t)\). Pour chaque colonne \(j\), on calcule les entiers et les résidus suivants :
\begin{align}
 S_j&=B_{\mathrm P,j}+B_{\mathrm E,j}+B_{\mathrm F,j},
 &q_j&=S_j\bmod3,
 &c_j&=\left\lfloor S_j/3\right\rfloor,
 \label{act21:eq:df-firma-qc}\\
 a_j&=(B_{\mathrm P,j}+B_{\mathrm E,j}-B_{\mathrm F,j})\bmod3,
 &w_j&=\sum_c\mathbf1_{B_{c,j}\ne0},
 &\bar w_j&=w_j\bmod3.
 \label{act21:eq:df-firma-axial}
\end{align}
Le résidu de l'expression axiale est pris dans les entiers. Une implémentation
avec des naturels le calcule comme
\((B_{\mathrm P,j}+B_{\mathrm E,j}+3-B_{\mathrm F,j})\bmod3\), qui le conserve
exactement.

La matrice d'incidence ternaire utilisée par le lecteur est
\begin{equation}
 W=\begin{pmatrix}
0&1&1&1&1&1\\
1&0&1&1&2&2\\
1&1&0&2&1&2\\
1&1&2&0&2&1\\
1&2&1&2&0&1\\
1&2&2&1&1&0
\end{pmatrix}.
\label{act21:eq:df-matriz-witt}
\end{equation}
Pour chaque région, deux charges résiduelles sont conservées :
\begin{equation}
 v_c^{\mathrm{res}}=\sum_j B_{c,j}\bmod3,
 \qquad
 v_c^{\mathrm W}=\sum_j(B_cW)_j\bmod3.
 \label{act21:eq:df-cargas-residuales}
\end{equation}
La signature complète du lecteur est
\(\mathcal S(B)=(q,a,c,w,\bar w,v^{\mathrm{res}},v^{\mathrm W})\).
Sa structure conserve séparément le reste de colonne, la retenue, le
contraste axial, le poids de support et les deux lectures de charge.

\begin{theorem}[Récupération exacte et bornes de la signature]
\label{act21:thm:df-recuperacion-firma}
Pour chaque colonne et chaque région, on a
\begin{align}
 S_j&=q_j+3c_j,&0\leq q_j,a_j&<3,&0\leq c_j&\leq2,
 \label{act21:eq:df-recuperacion-total}\\
 B_{\mathrm F,j}&=2(q_j+3-a_j)\bmod3,&0\leq w_j&\leq3,
 &0\leq\bar w_j&<3.
 \label{act21:eq:df-recuperacion-eje}
\end{align}
De plus,
\begin{equation}
 v_c^{\mathrm W}
 =\left(\sum_j\bigl((B_cW)_j\bmod3\bigr)\right)\bmod3.
 \label{act21:eq:df-witt-reduccion}
\end{equation}
Les récupérations énoncées déterminent le total de chaque colonne et sa
coordonnée d'auto-échelle ; le registre régional conserve en outre
les trois bandes ordonnées.
\end{theorem}

\begin{proof}
L'identité \eqref{act21:eq:df-recuperacion-total} est la division euclidienne de
\(S_j\) par \(3\). Chaque terme de \(S_j\) est entre \(0\) et \(2\),
de telle sorte que \(0\leq S_j\leq6\) et \(c_j\leq2\). En \(\mathbb F_3\),
la soustraction des deux congruences qui définissent \(q_j\) et \(a_j\) donne
\(q_j-a_j=2B_{\mathrm F,j}\). Comme l'inverse de \(2\) est \(2\), on
obtient \eqref{act21:eq:df-recuperacion-eje} ; le terme \(3\) permet
d'écrire celle-ci avec des nombres naturels. Les bornes de \(w_j\) résultent de sommer trois
indicateurs. Enfin, réduire une somme entière modulo \(3\) équivaut
à réduire chaque terme et à réduire à nouveau la somme, ce qui prouve
\eqref{act21:eq:df-witt-reduccion}.
\end{proof}

\begin{example}[Signature du premier bloc conjoint]
Les trois bandes initiales produisent
\begin{align*}
 q&=(0,0,2,2,1,2),&a&=(1,2,0,1,1,2),\\
 c&=(1,1,0,1,0,0),&w&=(2,2,2,3,1,2),\\
 v^{\mathrm{res}}&=(2,2,0),&v^{\mathrm W}&=(2,1,1).
\end{align*}
La quatrième colonne a un total \(2+1+2=5\) ; sa signature enregistre \(q_3=2\) et \(c_3=1\), par conséquent \(q_3+3c_3=5\).
La récupération axiale donne \(2(2+3-1)\bmod3=2\), précisément la coordonnée d'auto-échelle.
Les images de Witt, réduites composante par composante, sont
\[
 (2,0,2,1,1,2),\qquad
 (0,0,0,2,0,2),\qquad
 (2,1,1,2,1,0).
\]
Sa somme résiduelle par ligne produit les trois charges \((2,1,1)\).
\end{example}

\begin{definition}[Horloge entière de lecture décimale]
L'indice de transition \(t\) et la profondeur décimale de représentation sont
reliés par
\begin{equation}
 \kappa(t)=\max\{k\in\mathbb N:1000^k\leq3^{6(t+1)}\}.
 \label{act21:eq:df-reloj}
\end{equation}
\end{definition}

\begin{proposition}[Caractérisation et monotonie de l'horloge]
\label{act21:prop:df-reloj}
L'entier \(\kappa(t)\) est déterminé de manière univoque par
\begin{equation}
 1000^{\kappa(t)}\leq729^{t+1}<1000^{\kappa(t)+1}.
 \label{act21:eq:df-reloj-cotas}
\end{equation}
La fonction \(\kappa\) est monotone et satisfait
\(\kappa(20)=\kappa(21)=20\).
\end{proposition}

\begin{proof}
Les puissances de \(1000\) sont strictement croissantes et non bornées ;
il existe donc un unique intervalle de puissances consécutives qui
contient \(729^{t+1}\). La monotonie de cette dernière expression démontre
celle de \(\kappa\). Les deux évaluations finales sont les comparaisons
entières
\[
 1000^{20}\leq729^{21}<1000^{21},\qquad
 1000^{20}\leq729^{22}<1000^{21}.
\]
Ces comparaisons s'effectuent avec des puissances entières exactes.
\end{proof}

L'égalité de profondeur décimale en deux transitions distinctes explique la nécessité de conserver les deux indices : l'horloge de lecture peut rester constante tandis que la bande ternaire avance et que sa signature est mise à jour. L'information chronologique se trouve dans le registre ordonné des transitions.

\subsection{Théorème de composition du bloc, du cylindre et de la signature}

\begin{theorem}[Compatibilité conjointe des lectures régionales]
\label{act21:thm:df-composicion-conjunta}
Pour chaque profondeur \(n\), il existe un unique triplet de blocs
\((u_{\mathrm P},u_{\mathrm E},u_{\mathrm F})\) qui satisfait
l'admissibilité régionale \eqref{act21:eq:df-admisibilidad} dans les trois régions.
Ce triplet est \((u_{\mathrm P}(n),u_{\mathrm E}(n),u_{\mathrm F}(n))\).
Simultanément :
\begin{enumerate}
\item chaque bloc met à jour son cylindre par les récurrences
\eqref{act21:eq:df-rec-ternaria} et \eqref{act21:eq:df-rec-conjunta} ;
\item ses six trits produisent la signature
\(\mathcal S(B_{\mathrm P}(n),B_{\mathrm E}(n),B_{\mathrm F}(n))\);
\item les totaux de colonne se retrouvent comme \(q+3c\), et les charges de Witt sont obtenues à partir des mêmes bandes par
\eqref{act21:eq:df-witt-reduccion} ;
\item tout préfixe de plus grande profondeur renvoie exactement ces
blocs, cylindres tronqués et signatures.
\end{enumerate}
\end{theorem}

\begin{proof}
Nous appliquons le théorème ~\ref{act21:thm:df-prolongacion-unica} à chacun des trois
indices régionaux. L'unicité coordonnée par coordonnée donne l'unicité
conjointe. Nous substituons ces mêmes résidus dans les opérations
\(T_u,T_{u,d}\), ce qui permet d'appliquer le
théorème ~\ref{act21:thm:df-refinamiento-generado}. L'application déterministe de
\eqref{act21:eq:df-bloque-y-banda} produit les bandes utilisées par
\eqref{act21:eq:df-firma-qc}--\eqref{act21:eq:df-cargas-residuales}. Les identités
de récupération sont celles du théorème ~\ref{act21:thm:df-recuperacion-firma}.
Finalement, le théorème ~\ref{act21:thm:df-bloque-estable} et le troncage de
chaque base démontrent la stabilité à plus grande profondeur. À chaque étape, on a conservé l'identité du bloc émis ; aucune des deux
lectures ne requiert de sélectionner un autre bloc.
\end{proof}

La conclusion est une compatibilité opérationnelle entre la contraction
cylindrique et la conservation de l'incidence. Les deux utilisent le même
registre généré, maintiennent les profondeurs et admettent la récupération
entière. Ce résultat est la forme explicite du lien
\[
 \text{bloc régional émis}
 \longrightarrow
 \begin{cases}
 \text{cylindre et frontière entière},\\
 \text{bande, signature, retenue et charge de Witt}.
 \end{cases}
\]
Le développement ultérieur de l'incidence et des champs emploie la deuxième lecture, conservant l'origine commune avec la première.

\begingroup\fontsize{10}{11.8}\selectfont\setlength{\abovedisplayskip}{5pt}\setlength{\belowdisplayskip}{5pt}
\subsection{Recensement des calendriers et compatibilité simultanée des biographies régionales}
Les lignes précédentes se regroupent, sans appliquer encore aucun lecteur décimal,
dans les trois biographies de blocs produites par les transitions :
\begin{align}
 B^{\rm cl}
  &=010211|012222|010211|002111|110221,\notag\\
 B^{\rm pr}
  &=201101|121221|102011|012222|102011,
 \label{act21:eq:biografias-tpk-tres-caracteres}\\
 B^{\rm au}
  &=121200|112202|121020|010210|010200.\notag
\end{align}
En effet, les paires consécutives des deux premiers tronçons sont les lignes de \(X_0,Y_0\), et celles des deux derniers sont les lignes de \(X_1,Y_1\). Ainsi, les biographies sont une autre écriture du registre de douze transitions, non une collection de préfixes conventionnels.

Pour vérifier le mot opératoire du calendrier, soit
\(c=c_1c_2c_3c_4\in\{0,1\}^4\) et, pour chaque régime
\(\chi\in\{{\rm cl},{\rm pr},{\rm au}\}\), définissons
\[
 u_0^\chi=b_0^\chi,\qquad
 u_j^\chi=u_{j-1}^\chi L_{c_j}\quad(1\le j\le4).
\]
Pour condenser le recensement, notons par
\[
 N_{\mathrm{ar}}(c):=
 \#\{(j,\chi):u_j^\chi=b_j^\chi,\ 1\le j\le4\},\qquad
 N_{\mathrm{bio}}(c):=
 \#\{\chi:(u_0^\chi,\ldots,u_4^\chi)=B^\chi\}.
\]
La multiplication exacte dans \(\FF_3\) des seize calendriers donne :
\[
\begin{array}{c|c|c}
 c&N_{\mathrm{ar}}(c)&N_{\mathrm{bio}}(c)\\ \hline
0000,0001&6&0\\
0010&9&0\\
0011&12&3\\
0100,0101,0110,0111&3&0\\
1000,1001,1010,1011&0&0\\
1100,1101,1110,1111&0&0
\end{array}
\tag{D}
\label{act21:eq:censo-dieciseis-calendarios}
\]
Le tableau épuise les \(2^4=16\) possibilités et chacune de ses entrées s'obtient
par quatre multiplications d'une ligne par l'une des deux matrices reproduites.
Par conséquent, \(0011\) est l'unique calendrier qui reconstruit
simultanément les trois biographies complètes. En notation opératorielle,
\begin{equation}
 (L_0,L_0,L_1,L_1)
 \label{act21:eq:calendario-0011-interno}
\end{equation}
ce n'est pas une prescription indépendante : c'est la seule lecture binaire compatible
avec les douze arêtes déjà produites par \(U_t\).

Le calendrier \eqref{act21:eq:calendario-0011-interno} produit quatre nouveaux
blocs de six composants. Leurs concaténations forment
\[
 w_6\longrightarrow w_{12}\longrightarrow w_{18}
 \longrightarrow w_{24}\longrightarrow w_{30}.
\]
Les sous-indices de \(w\) indiquent la longueur accumulée. La frontière \(R_{36}\) conserve la structure terminale et ouvre la relation de prolongement ; elle n'est pas renommée comme un dernier bloc périodique qui pourrait se répéter indéfiniment.

L'identité \(L=X^{-1}Y\) démontre l'unicité une fois les
transitions fixées ; elle ne remplace pas leur production par \(U_t\). Réciproquement,
\eqref{act21:eq:censo-dieciseis-calendarios} prouve dans l'article l'unicité
de l'ordre des quatre transports. La chaîne causale est ainsi close :
le recensement APP produit les régions initiales, le TRIT fixe le régime et
l'orientation, le TPK produit le registre des transitions, et ce registre
détermine \(L_0,L_1\) et \(0011\) avant toute réalisation archimédienne.


\par\endgroup

\clearpage
% Desarrollo reunido K31-07--K31-15. Inserción anterior a registro_k y alpha.
% Mantiene explícita la interfaz del repertorio terminal de ocho bloques.
% Propietarios y localizadores: metadata/K_ALPHA_PROPIETARIOS.md.
\section{Détermination régionale du descripteur dodécaphasique et sélection du registre de couplage}
\label{act21:chap:despliegue-selector-k}

La détermination du registre de couplage articule deux constructions
qu'il convient de déployer dans leur ordre effectif. La première produit les
antécédents du sélecteur à partir des représentations régionales du TPK :
bandes de six trits, signatures de colonne, calendrier de conversion,
descripteur orienté et tableau fonctionnel. La seconde opère sur ces
antécédents et sur le répertoire terminal de huit blocs, énumère leurs
ordonnancements admissibles et sélectionne un unique registre au moyen de deux
conditions simultanées : incidence exceptionnelle et hétérotypie de lecture.
La valeur du registre apparaît au terme de cette composition.

Le présent développement présuppose les représentations régionales obtenues
dans les chapitres précédents à partir d'APP--TRIT--TPK. Son propos est de rendre
explicite chaque opération intermédiaire qui conduit de ces représentations
au sélecteur fini. Sont ensuite incorporées la réalisation du même
registre par incidences signées, l'inversion de Hadamard et la
représentation de la constante de structure fine. Les trois constructions
conservent leurs domaines respectifs : sélection du registre, récupération
à partir de ses canaux et évaluation arithmétique de ses représentations.

\subsection{Bandes régionales, codification et signatures entières}

\subsubsection{Domaines et codification des bandes régionales}

On fixe l'ordre des canaux
\[
 \chi\in\{\pi,e,\varphi\},
 \qquad B_\chi(t)\in\mathbb F_3^6,
 \qquad t\in\mathbb N.
\]
Les symboles de canal identifient les régions déjà engendrées. Chaque
$B_\chi(t)$ conserve la place qu'il occupe dans sa suite de représentations ;
l'indice $t$ compte les blocs ternaires. L'évaluation naturelle d'une
bande $b=(b_0,\ldots,b_5)$ est
\begin{equation}
 \operatorname{enc}_6(b)=
 243b_0+81b_1+27b_2+9b_3+3b_4+b_5,
 \qquad 0\leq\operatorname{enc}_6(b)<729.
 \label{act21:eq:kdes-enc6}
\end{equation}
Son inverse, avec une largeur fixée, est
\[
 \operatorname{dig}_j(n)=
 \left\lfloor\frac{n}{3^{5-j}}\right\rfloor\bmod3,
 \qquad 0\leq j<6.
\]
Ainsi sont retenus les zéros initiaux : les mots $002001$ et $2001$
pourraient partager une écriture entière abrégée, mais le type du
premier contient six positions. Toutes les concaténations qui suivent
utilisent cette largeur.

Dans la fenêtre de cent blocs, la représentation de six cents trits du
canal $\chi$ détermine un entier $P_\chi^{(600)}$. L'extraction s'effectue
par
\begin{equation}
 b_\chi(t)=
 \left\lfloor\frac{P_\chi^{(600)}}{729^{99-t}}\right\rfloor\bmod729,
 \qquad 0\leq t<100,
 \qquad B_\chi(t)=\operatorname{enc}_6^{-1}(b_\chi(t)).
 \label{act21:eq:kdes-extraccion}
\end{equation}
L'entier de six cents trits est une représentation du producteur régional ;
la table de cent lignes est l'évaluation de
\eqref{act21:eq:kdes-extraccion}. Cette séparation permet de remplacer une table
historique par son producteur sans modifier le sélecteur qui l'utilise.

\subsubsection{Résidus, quotients, occupation et orientation}

Pour chaque colonne $j$ et temps $t$, on introduit
\begin{align}
 T_j(t)&=B_\pi(t)_j+B_e(t)_j+B_\varphi(t)_j,\\
 q_j(t)&=T_j(t)\bmod3,
 &c_j(t)&=\left\lfloor T_j(t)/3\right\rfloor,\\
 a_j(t)&=\bigl(B_\pi(t)_j+B_e(t)_j-B_\varphi(t)_j\bigr)\bmod3,\\
 w_j(t)&=\mathbf1_{B_\pi(t)_j\ne0}
       +\mathbf1_{B_e(t)_j\ne0}
       +\mathbf1_{B_\varphi(t)_j\ne0},
 &\bar w_j(t)&=w_j(t)\bmod3.
 \label{act21:eq:kdes-firmas}
\end{align}
La soustraction de la troisième expression s'effectue sur les entiers et se réduit
ensuite. Dans une implémentation sur les naturels, l'expression équivalente
est $(B_\pi+B_e+3-B_\varphi)\bmod3$ ; la soustraction tronquée altérerait le lecteur.

\begin{proposition}[Reconstruction de la somme de colonne]
Pour toutes les bandes admissibles,
\[
 T_j(t)=q_j(t)+3c_j(t),\qquad
 0\leq q_j(t)<3,\quad 0\leq c_j(t)\leq2,
 \quad0\leq w_j(t)\leq3.
\]
\end{proposition}
\begin{proof}
Chaque colonne additionne trois représentants de $\{0,1,2\}$, de sorte que
$0\leq T_j(t)\leq6$. La division euclidienne par trois fournit le reste
et le quotient indiqués. L'occupation compte trois conditions binaires ;
sa borne se déduit directement de cette définition.
\end{proof}

Cette identité conserve le quotient de la somme locale. La reconstruction
des trois bandes individuelles requiert les bandes elles-mêmes ou un lecteur
supplémentaire qui les distingue : une somme de colonne a moins de coordonnées
que le triplet qui la produit. Par conséquent, la ligne temporelle utilisée plus
tard conserve conjointement les trois bandes et les quatre champs
$q,a,c,\bar w$.

\subsection{Horloge de conversion et premier avancement nul}

\subsubsection{Capacité ternaire et profondeur décimale}

Une représentation de $t+1$ bandes contient $6(t+1)$ trits. Le calendrier
de conversion associé est défini par l'entier
\begin{equation}
 d_t=\max\{k\in\mathbb N:1000^k\leq3^{6(t+1)}\}
     =\max\{k\in\mathbb N:1000^k\leq729^{t+1}\}.
 \label{act21:eq:kdes-reloj}
\end{equation}
La définition compare des capacités entières. Le compteur $t$ et la profondeur $d_t$ enregistrent des grandeurs distinctes : une transition ajoute une bande de six trits ; l'horloge compte les puissances complètes de la base décimale mille contenues dans la capacité ternaire accumulée.

\begin{lemma}[Caractérisation entière de l'horloge]
Pour $t,k\in\mathbb N$ on a
\[
 d_t=k\quad\Longleftrightarrow\quad
 1000^k\leq729^{t+1}<1000^{k+1}.
\]
De plus, $d_t\leq d_{t+1}$.
\end{lemma}
\begin{proof}
Les puissances de mille forment une suite strictement croissante et non bornée. La puissance $729^{t+1}$ se situe entre deux termes consécutifs ; l'exposant inférieur est exactement le maximum de \eqref{act21:eq:kdes-reloj}. L'inégalité $729^{t+1}<729^{t+2}$ implique la monotonie du maximum.
\end{proof}

\begin{proposition}[Premier avancement nul]
Le plus petit temps qui satisfait $d_{t+1}=d_t$ est
\[
 t_*=20.
\]
En particulier,
\[
 d_0=0,\quad d_{19}=19,\quad d_{20}=20,
 \quad d_{21}=20,\quad d_{22}=21.
\]
\end{proposition}
\begin{proof}
Les comparaisons entières exactes
\[
 1000^{20}\leq729^{21},\qquad
 729^{22}<1000^{21},\qquad
 1000^{21}\leq729^{23}<1000^{22}
\]
donnent les valeurs aux temps vingt, vingt-et-un et vingt-deux. Pour
$0\leq t\leq20$, l'expression
$729^{t+1}/1000^t=729(729/1000)^t$ est décroissante et sa valeur en vingt
est supérieure ou égale à un. Par conséquent,
$1000^t\leq729^{t+1}<1000^{t+1}$ et $d_t=t$. Aucun temps inférieur à
vingt ne peut satisfaire $d_{t+1}=d_t$ ; en vingt, c'est le cas.
Toutes les comparaisons utilisées sont de puissances entières.
\end{proof}

Le nombre vingt est donc une sortie du critère minimal. La règle qui sélectionne l'instant est ``première avance nulle de l'horloge'' ; l'algorithme et sa preuve peuvent utiliser cette règle sans recevoir vingt comme paramètre externe.

\subsubsection{Retour nonadique et antécédent temporel}

La longueur du retour de phase détermine
\begin{equation}
 r=t_*-9=11.
 \label{act21:eq:kdes-retorno}
\end{equation}
Les requêtes du descripteur ont lieu en $t_*+1$, $r+1$ et $r$.
Cette distribution conserve la relation entre l'instant de pause, la
translation nonadique et la lecture de phase immédiatement associée. Les
temps onze, douze et vingt-et-un désignent des positions du calendrier
régional ; leur différence temporelle reste enregistrée même lorsque deux
temps produisent la même profondeur décimale.

\subsection{Construction orientée du descripteur de vingt-quatre trits}

\subsubsection{Défauts opposés de semi-bande}

On considère les défauts orientés de la coordinatisation régionale
\[
 \delta_+=(0,0,0,1,1,1),\qquad
 \delta_-=(0,0,0,2,2,2)\in\mathbb F_3^6.
\]
Son opposition se vérifie composante par composante :
\[
 \delta_++\delta_-=0\quad\text{en }\mathbb F_3^6.
\]
Le représentant $2$ réalise $-1$ dans la coordinatisation ternaire. La concaténation
des queues de trois positions, dans l'ordre négatif--positif, produit
le mot $222111$. La somme des défauts et la concaténation de leurs
queues sont des opérations distinctes : la première a pour codomaine
$\mathbb F_3^6$, tandis que la seconde conserve l'ordre de deux
segments de longueur trois.

\subsubsection{Charge duale et sélection de phase}

Dans la coordinatisation régionale, la transformation de six coordonnées utilise
\[
 A_{\rm reg}=
 \begin{pmatrix}
 0&1&1&1&1&1\\
 1&0&1&1&2&2\\
 1&1&0&2&1&2\\
 1&1&2&0&2&1\\
 1&2&1&2&0&1\\
 1&2&2&1&1&0
 \end{pmatrix}.
\]
Avec des vecteurs ligne, la charge duale de $b\in\mathbb F_3^6$ est la somme
des coordonnées de $bA_{\rm reg}$. Les sommes entières des lignes de
$A_{\rm reg}$ sont $(5,7,7,7,7,7)$ ; d'où résulte l'expression locale
\begin{equation}
 \eta^\vee(b)=2b_0+b_1+b_2+b_3+b_4+b_5\pmod3.
 \label{act21:eq:kdes-cargadual}
\end{equation}
La phase de lecture à l'instant $r$ est
$\theta_r=(r-1)\bmod3$. Pour chaque canal, on définit
\[
 \varepsilon_\chi(r)=
 \begin{cases}
 1,&\eta^\vee(B_\chi(r))=\theta_r,\\
 2,&\eta^\vee(B_\chi(r))\ne\theta_r.
 \end{cases}
\]
La représentation phase--charge est
\begin{equation}
 D(r)=\varepsilon_\pi(r)B_\pi(r)
      +\varepsilon_e(r)B_e(r)
      +\varepsilon_\varphi(r)B_\varphi(r)
      \quad\text{en }\mathbb F_3^6.
 \label{act21:eq:kdes-fasecarga}
\end{equation}
Le coefficient se détermine par une égalité de charges, avant d'évaluer le mot qui résultera de la combinaison.

\subsubsection{Composition du descripteur}

On désigne par $\Vert$ la concaténation de mots de largeur fixée.
La construction complète est
\begin{equation}
 \begin{split}
 \Omega_{24}={}&
 (B_e(t_*+1)+\delta_-)
 \Vert (\delta_-|_{3,4,5}\Vert\delta_+|_{3,4,5})\\
 &\Vert(B_\pi(r+1)+B_e(r+1))\Vert D(r).
 \end{split}
 \label{act21:eq:kdes-omega}
\end{equation}
Chaque terme de longueur six se calcule en $\mathbb F_3^6$ ; la
concaténation finale a une longueur $6+6+6+6=24$. Cette construction
sépare le descripteur $\Omega_{24}$ des étapes régionales
$w_{24}^{\chi}$ : les deux objets partagent la même longueur, mais leurs producteurs
et leurs places dans la composition sont différents.

\begin{proposition}[Évaluation régionale du descripteur]
Les représentations régionales et le critère de la première avancée nulle
déterminent
\[
 \Omega_{24}=020022\,222111\,211201\,021101,
 \qquad [\Omega_{24}]_3=66241910521.
\]
\end{proposition}
\begin{proof}
Les bandes nécessaires, conservant l'ordre $(\pi,e,\varphi)$, sont
\[
 \begin{array}{c|ccc}
 t&B_\pi(t)&B_e(t)&B_\varphi(t)\\\hline
 11&022212&110001&100021\\
 12&012012&202222&000120\\
 21&221022&020100&211021
 \end{array}
\]
et proviennent de l'extraction \eqref{act21:eq:kdes-extraccion}. Au temps
vingt-et-un,
$020100+000222=020022$ en $\mathbb F_3^6$. Les queues des défauts
apportent $222111$. Au temps douze,
$012012+202222=211201$.

Pour le temps onze, les charges duales sont respectivement $0,1,2$ ; la phase vaut $(11-1)\bmod3=1$. Les coefficients sont $(2,1,2)$ et
\[
 2(022212)+(110001)+2(100021)=021101
 \quad\text{en }\mathbb F_3^6.
\]
La concaténation de ces quatre résultats donne le mot annoncé.
Son évaluation entière s'obtient par Horner en base trois, ou par trois
concaténations successives en base $729$.
\end{proof}

Le tableau précédent joue un rôle vérificateur : chacune de ses entrées provient du producteur régional. La règle de construction \eqref{act21:eq:kdes-omega} est celle qui permet de reproduire le descripteur sans introduire en tant que donnée la chaîne finale de vingt-quatre trits.

\subsection{Panneau fonctionnel et répertoire terminal}

\subsubsection{Neuf positions régionales}

Le panneau fonctionnel conserve les trois premières bandes de chacun des trois canaux, dans l'ordre régional déclaré :
\begin{equation}
 \mathcal P_9=
 \bigl(b_\pi(0),b_\pi(1),b_\pi(2),
       b_e(0),b_e(1),b_e(2),
       b_\varphi(0),b_\varphi(1),b_\varphi(2)\bigr).
 \label{act21:eq:kdes-panel}
\end{equation}
Son évaluation est
\[
 \mathcal P_9=(103,161,103,523,457,301,450,398,438).
\]
Les neuf positions conservent l'origine régionale, bien que la valeur
103 apparaisse deux fois. Le panneau comporte neuf incidences ordonnées et
huit valeurs distinctes. L'énumération ultérieure utilise les incidences ;
la déduplication des candidats s'effectue à une étape postérieure.

\subsubsection{Répertoire non ordonné de huit blocs}

Le répertoire terminal qui intervient dans cette construction est
\begin{equation}
 \mathcal S_{\rm term}=
 \{002001,020111,200110,121012,
   010122,001001,221110,222110\}\subset\mathbb F_3^6.
 \label{act21:eq:kdes-sterm}
\end{equation}
Ses huit éléments sont distincts. La codification entière, ordonnée par
valeur uniquement pour fixer une représentation informatique, est
\[
 \{28,55,98,175,437,498,687,714\}.
\]
La provenance documentaire de ce répertoire est la segmentation terminale du mot de soixante-douze trits et la récupération par positions au moyen des incidences des lecteurs phase--charge et affine. Le sélecteur développé ci-après reçoit le répertoire sans son ordonnancement historique et met à l'épreuve ses $8!$ ordonnancements.

La formulation exacte de l'implémentation réunie maintient ici une
interface explicite : les producteurs régionaux remplacent matériellement
le descripteur, le panneau et les lignes temporelles ; le répertoire
\eqref{act21:eq:kdes-sterm} conserve son producteur documentaire antérieur. Cette
distinction identifie où intervient chaque antécédent. L'appartenance de
huit mots à un répertoire temporel de lecteurs, la sélection de ces
huit mots et la détermination de leur ordre terminal sont trois
énoncés différents. La preuve de sélection terminale démontre le
troisième pour le répertoire déclaré.

La récupération documentaire par positions peut s'exprimer de manière
exacte. Soit $\mathcal I$ le registre des incidences phase--charge et soit
$\mathcal J$ le registre de l'extension affine. Chaque incidence conserve
la position $\ell$ du mot de soixante-douze trits et le mot
émis $u$. Pour $5\leq\ell\leq12$, on forme
\[
 \mathcal S_\ell=
 \{u:(\ell,u)\in\mathcal I\cup\mathcal J\}.
\]
Le récupérateur vérifie que les huit positions apparaissent et que chaque
$\mathcal S_\ell$ est un singleton. Si $s_\ell$ est son élément unique,
la sortie livrée au sélecteur est
\[
 \mathcal S_{\rm term}=\{s_5,s_6,\ldots,s_{12}\}.
\]
La preuve de singleton par position préserve la cohérence des témoins historiques ; l'exportation ultérieure préserve les huit mots et supprime uniquement la mise en ordre historique. Le sélecteur reconstruit l'ordre par l'énumération complète et les conditions terminales, au lieu de le recevoir du registre d'origine.

% RESULTADO_RECUPERADO: K31-10/P10 del inventario REV05.
% Fuente de las incidencias:
% 16_CIERRE_GLOBAL_HMT_MD_2026-07-22/02_LECTURA_DODECAFASICA/continuacion_fases/
% verificar_obstruccion_y_dato_minimo_w72.py, lineas 45--60 y 141--233.
% RESULTADO_OBSTRUCCION_Y_DATO_MINIMO_W72.json, phase_charge_reader.
% Fuente de las bandas: N69_signatures_t0_t99.csv, tiempos impresos abajo.
% Consumidor: PAQUETE_CONTINUIDAD_K_UNIDAD_20260918_BASE_COMPARTIDA/python/
% selector_orbital_original.py, verify_historical_inputs, lineas 159--193.
% No se transfiere la restriccion del lector comparativo al generador completo.
\subsubsection{Incidences temporelles du répertoire terminal et récupération par position}
\label{act21:sec:repertorio-incidencias-completas}

La représentation du répertoire terminal utilise les mêmes blocs régionaux,
les mêmes retenues et le même calendrier qui produisent le descripteur initial
de vingt-quatre trits. Il convient d'expliciter séparément deux opérations de
cette composition. La première récupère huit mots à partir d'incidences
temporelles enregistrées. La seconde présente ces huit mots au sélecteur
dodécaphasique comme un ensemble non ordonné. Le sélecteur rétablit
l'ordre par ses conditions d'incidence, de cylindre et de clôture.
Cette distinction permet de conserver les témoins de provenance sans transformer
l'ordre d'une segmentation historique en donnée d'entrée du sélecteur.

Soit $B_t=(b_{t,0},b_{t,1},b_{t,2})\in(\mathbb F_3^6)^3$ la frontière
régionale déjà engendrée au temps $t$, avec $0\leq t<100$. Dans la coordinatisation
utilisée par le registre historique, la matrice d'incidence est
\[
 A_{\mathrm{hist}}=
 \begin{pmatrix}
 0&1&1&1&1&1\\
 1&0&1&1&2&2\\
 1&1&0&2&1&2\\
 1&1&2&0&2&1\\
 1&2&1&2&0&1\\
 1&2&2&1&1&0
 \end{pmatrix}.
\]
Les mots sont considérés comme des vecteurs ligne. Leurs deux charges et la phase
chronologique sont
\[
 r_{t,i}=\sum_{j=1}^{6}b_{t,i,j},\qquad
 d_{t,i}=\sum_{j=1}^{6}(b_{t,i}A_{\mathrm{hist}})_j,
 \qquad \vartheta_t=t-1\pmod 3.
\]
Toutes les opérations de ces trois expressions ont lieu dans
$\mathbb F_3$. L'indice historique identifie une coordinatisation précise de
l'incidence de Witt ; il maintient explicite la transformation de coordonnées
lorsqu'une autre présentation de cette même incidence est utilisée.

Pour $c\in\{r_t,d_t\}$, $\delta\in\mathbb F_3$ et
$\epsilon\in\{1,-1\}$, nous définissons
\[
 \eta_i(c,\delta,\epsilon)=
 \begin{cases}
  \epsilon,&c_i=\vartheta_t+\delta,\\
  -\epsilon,&c_i\ne\vartheta_t+\delta,
 \end{cases}
 \qquad
 L_{t,c,\delta,\epsilon}
   =\sum_{i=0}^{2}\eta_i(c,\delta,\epsilon)b_{t,i}.
\]
L'orientation $-1$ s'écrit $2$ lors de l'impression des coefficients
ternaires. Outre ces douze lecteurs comparatifs par frontière, le
registre dispose du lecteur affine
\[
 L^{\mathrm{af}}_t
    =\sum_{i=0}^{2}(d_{t,i}+\vartheta_t)b_{t,i}.
\]
L'addition de la phase dans le lecteur affine et la comparaison des charges dans
le lecteur précédent sont des opérations distinctes, toutes deux sur des données déjà
produites par la frontière TPK.

\begin{table}[htbp]
\centering\small
\begin{tabular}{rccc cc}
\toprule
$t$&$b_{t,0}$&$b_{t,1}$&$b_{t,2}$&$r_t$&$d_t$\\
\midrule
8 &200121&212212&110110&011&202\\
11&022212&110001&100021&001&012\\
30&210112&110202&211110&100&012\\
34&121022&000011&112110&220&021\\
37&120121&011202&112221&100&201\\
49&110212&200212&100122&110&201\\
58&111121&022121&122201&122&220\\
67&010012&001220&011120&122&122\\
75&220000&020111&100002&120&021\\
81&122201&020202&201122&202&001\\
90&022112&022110&121010&202&200\\
\bottomrule
\end{tabular}
\caption{Frontières temporelles qui réalisent les incidences imprimées du
répertoire terminal. Chaque entrée de six chiffres conserve l'ordre des
six coordonnées de la frontière régionale.}
\label{act21:tab:terminal-bandas-testigo}
\end{table}

Le tableau suivant déploie toutes les incidences comparatives du
registre sur les douze blocs de la segmentation considérée. Les
lettres $r$ et $d$ désignent respectivement la charge visible et la charge
duale. La colonne $j$ conserve l'indice de position du témoin historique ;
l'opération de livraison au sélecteur oublie cet indice après avoir formé
l'ensemble de mots.

\begin{table}[htbp]
\centering\small
\begin{tabular}{rrcrrc c}
\toprule
$j$&$t$&charge&$\delta$&$\epsilon$&$(\eta_0\eta_1\eta_2)$&$L_{t,c,\delta,\epsilon}$\\
\midrule
4&11&$d$&0& 1&212&021101\\
5&67&$d$&1&$-1$&211&002001\\
5&67&$d$&2& 1&211&002001\\
5&67&$r$&1&$-1$&211&002001\\
5&67&$r$&2& 1&211&002001\\
6&81&$d$&0& 1&222&020111\\
6&81&$r$&2& 1&222&020111\\
7&34&$r$&1&$-1$&111&200110\\
7&75&$d$&1&$-1$&211&200110\\
7&75&$r$&2&$-1$&211&200110\\
8&90&$r$&0& 1&121&121012\\
8&90&$r$&1&$-1$&121&121012\\
9&49&$d$&0&$-1$&121&010122\\
11&11&$d$&2&$-1$&211&221110\\
11&37&$d$&0&$-1$&121&221110\\
12&8&$d$&0&$-1$&111&222110\\
12&8&$r$&1&$-1$&111&222110\\
12&58&$d$&1&$-1$&111&222110\\
12&58&$r$&0&$-1$&111&222110\\
\bottomrule
\end{tabular}
\caption{Incidences phase--charge complètes sur la segmentation terminale  
aux cent frontières du registre. Les coïncidences de mots entre  
lignes conservent leurs temps, orientations et types de charge.}
\label{act21:tab:terminal-incidencias-completas}
\end{table}

Le dixième bloc dispose d'un témoin affine. Dans $t=30$ on a
$d_{30}=(0,1,2)$ et $\vartheta_{30}=2$, d'où
\[
 (d_{30,0}+\vartheta_{30},d_{30,1}+\vartheta_{30},
 d_{30,2}+\vartheta_{30})=(2,0,1)
\]
et, en utilisant les bandes du tableau ~\ref{act21:tab:terminal-bandas-testigo},
\[
 L^{\mathrm{af}}_{30}
 =2(210112)+0(110202)+(211110)
 =(001001)\quad\text{en }\mathbb F_3^6.
\]
L'égalité est coordonnée par coordonnée : avant de réduire modulo trois,  
le vecteur de la somme est $(6,3,1,3,3,4)$.

\begin{proposition}[Récupération intégrale du répertoire par incidences]
\label{act21:prop:repertorio-terminal-ocho-incidencias}
Considérons le registre des incidences comparatives du
tableau~\ref{act21:tab:terminal-incidencias-completas}, augmenté du témoin
affine $(j,t,L)=(10,30,001001)$. Pour chaque position $5\leq j\leq12$,
l'ensemble des mots de sortie de ses incidences est un singleton. L'
extraction de ces huit mots et l'oubli ultérieur de leur ordre
produisent exactement
\[
 \mathcal S_{\mathrm{term}}=
 \{002001,020111,200110,121012,010122,001001,221110,222110\}.
\]
Sa codification en tant que nombres entiers en base trois est
\[
 \{28,55,98,175,437,498,687,714\}.
\]
\end{proposition}
\begin{proof}
La multiplicité d'incidences par position $j=5,6,7,8,9,10,11,12$
est, respectivement,
\[
 (4,2,3,2,1,1,2,4).
\]
À chaque position, toutes les lignes impriment le même mot. Chaque ligne se vérifie en substituant ses trois bandes et les trois coefficients dans la formule de $L$ ; le dixième cas a été calculé explicitement au moyen de $L^{\mathrm{af}}_{30}$. Par exemple, en $t=67$ et avec les coefficients $(2,1,1)$, on obtient
\[
 2(010012)+(001220)+(011120)=(002001)\pmod3.
\]
Les huit résultats sont distincts en tant que mots de longueur six. Leur codification $\sum_{k=1}^{6}w_k3^{6-k}$ donne, dans l'ordre des positions, $(55,175,498,437,98,28,714,687)$. En oubliant l'ordre, on obtient l'ensemble des entiers énoncé. Ainsi, l'extraction préserve chaque mot complet, y compris ses positions nulles, et fournit huit éléments distincts au sélecteur.
\end{proof}

La preuve précédente est une récupération exacte à partir d'incidences
identifiées. Sa donnée d'entrée comprend le registre d'incidences
avec leurs positions et leurs témoins. La sélection orbitale ultérieure reçoit
uniquement $\mathcal S_{\mathrm{term}}$ et le descripteur régional initial ;
elle parcourt les $8!$ ordonnancements et détermine leur compatibilité avec les
frontières fonctionnelles et les incidences orientées de la roue
dodécaphasique. Les deux opérations sont ainsi contiguës, avec des domaines
explicites : provenance temporelle des huit mots et détermination
constructive de leur ordre.

Le tableau comparatif contient les dix-neuf incidences que sa formule produit sur la segmentation de douze blocs aux cent temps déclarés. Cette exhaustivité finie se vérifie en parcourant $100\cdot2\cdot3\cdot2=1200$ évaluations. Le lecteur affine étend la couverture des positions terminales de sept à huit. Ce décompte caractérise le répertoire local de lecteurs qui vient d'être défini ; les autres opérations du TPK conservent les domaines et les fonctions constructives établis dans leurs dérivations respectives.

\subsection{Lecteurs comparatifs, lecteur affine et incidence temporelle}

\subsubsection{Transformation de Witt du sélecteur}

Le sélecteur utilise la coordinatisation à six coordonnées
\[
 A_{\rm ter}=
 \begin{pmatrix}
 0&1&1&1&1&1\\
 1&0&1&2&2&1\\
 1&1&0&1&2&2\\
 1&2&1&0&1&2\\
 1&2&2&1&0&1\\
 1&1&2&2&1&0
 \end{pmatrix}.
\]
On définit $W(b)=bA_{\rm ter}$ dans $\mathbb F_3^6$, ainsi que les charges
\[
 \eta(b)=\sum_{j=0}^5 b_j\pmod3,
 \qquad\eta_{\rm ter}^\vee(b)=\eta(W(b)).
\]
La coordinatisation du descripteur et celle du sélecteur sont reliées par la
permutation $\sigma=(0,1,2,4,5,3)$, écrite comme liste d'images :
\begin{equation}
 (A_{\rm ter})_{ij}=(A_{\rm reg})_{\sigma(i),\sigma(j)}.
 \label{act21:eq:kdes-cambio-carta}
\end{equation}

\begin{lemma}[Compatibilité de la fonctionnelle de charge]
Pour toute bande $b$, on a
\[
 \eta_{\rm ter}^\vee(b)=\eta^\vee(b).
\]
\end{lemma}
\begin{proof}
La somme entière de chaque ligne correspondante des deux matrices coïncide :
elle vaut cinq pour la première ligne et sept pour les cinq autres. Par
distributivité,
\[
 \sum_j\sum_i b_i(A_{\rm ter})_{ij}
 =\sum_i b_i\sum_j(A_{\rm ter})_{ij}
 =\sum_i b_i\sum_j(A_{\rm reg})_{ij}.
\]
La réduction modulo trois donne l'égalité des charges. L'identité des matrices \eqref{act21:eq:kdes-cambio-carta} se vérifie par substitution de ses trente-six entrées et conserve, en outre, le changement de coordonnées.
\end{proof}

L'égalité de la fonctionnelle totale de charge permet de comparer ses
représentations. Les transformations vectorielles restent identifiées
par leurs matrices et par la permutation explicite ; la preuve précédente
conserve cette information de coordonnées.

\subsubsection{Douze lecteurs comparatifs par instant}

On définit la phase temporelle
\[
 \theta(t)=(t+2)\bmod3.
\]
Cette écriture réalise également $(t-1)\bmod3$ pour $t=0$, où la phase
vaut deux. Pour chaque type de charge
$\nu\in\{\eta,\eta_{\rm ter}^\vee\}$, déplacement
$o\in\{0,1,2\}$ et orientation $s\in\{1,2\}$, on prend
\[
 \epsilon_\chi^{\nu,o,s}(t)=s
 \begin{cases}
 1,&\nu(B_\chi(t))=\theta(t)+o\pmod3,\\
 2,&\nu(B_\chi(t))\ne\theta(t)+o\pmod3.
 \end{cases}
\]
Le mot comparatif correspondant est
\begin{equation}
 C_{\nu,o,s}(t)=
 \sum_{\chi\in\{\pi,e,\varphi\}}
 \epsilon_\chi^{\nu,o,s}(t)B_\chi(t)
 \quad\text{en }\mathbb F_3^6.
 \label{act21:eq:kdes-comparativo}
\end{equation}
Les paramètres engendrent douze représentations étiquetées par ligne. Deux
étiquettes peuvent produire le même mot ; les témoins temporels et
les étiquettes restent disponibles dans le calendrier complet.

\subsubsection{Lecteur affine de phase et de charge}

Le lecteur affine emploie des coefficients
\[
 a_\chi^{\rm af}(t)=\eta_{\rm ter}^\vee(B_\chi(t))+\theta(t)
 \pmod3,
\]
et produit
\begin{equation}
 F(t)=\sum_{\chi\in\{\pi,e,\varphi\}}
 a_\chi^{\rm af}(t)B_\chi(t)
 \quad\text{en }\mathbb F_3^6.
 \label{act21:eq:kdes-afin}
\end{equation}
La règle comparative distingue la coïncidence et la différence de charges ; la règle affine additionne charge et phase avant de combiner les bandes. C'est cette différence opératoire qu'utilise la sélection hétérotypique.

\subsubsection{Incidence mot--résidu}

Soit $R$ la rotation cyclique des six positions d'une bande. À la
ligne temporelle $t$, on forme le répertoire résiduel
\[
 \mathcal R(t)=\{R^jv:
      v\in\{q(t),a(t),c(t),\bar w(t)\},\ 0\leq j<6\}.
\]
Les relations d'incidence temporelle sont
\begin{align}
 \mathcal A_{\rm cmp}
 &=\{(C_{\nu,o,s}(t),u):
      0\leq t<100,\ u\in\mathcal R(t),\ \nu,o,s\text{ admissibles}\},\\
 \mathcal A_{\rm af}
 &=\{(F(t),u):0\leq t<100,\ u\in\mathcal R(t)\}.
 \label{act21:eq:kdes-atlas}
\end{align}
Chaque couple relie un mot lu à un résidu de la même ligne.
Le temps se quantifie de manière commune dans les deux coordonnées du couple.
Dans le calendrier employé, les temps $0,\ldots,99$ sont distincts ;
identifier une même ligne et identifier un même temps sont, par conséquent,
des opérations équivalentes.

\begin{proposition}[Compression existentielle du répertoire temporel]
La fonction
\[
 M(v,u)=\mathbf1_{(v,u)\in\mathcal A_{\rm cmp}}
        +2\mathbf1_{(v,u)\in\mathcal A_{\rm af}}
        \in\{0,1,2,3\}
\]
conserve exactement l'existence de chaque type de lecture pour le couple  
$(v,u)$.
\end{proposition}
\begin{proof}
Le premier bit vaut un exactement lorsqu'il existe un temps, un champ
résiduel, une rotation et une étiquette comparative qui produisent le couple.
Le second bit exprime la condition correspondant au lecteur affine.
L'union par disjonction des bits conserve chaque condition d'existence
lors du parcours de toutes les lignes. Un couple peut admettre les deux types et, dans ce
cas, le masque vaut trois.
\end{proof}

La fonction $M$ est un lecteur d'existence. Les temps témoins, leurs multiplicités et les étiquettes individuelles demeurent dans \eqref{act21:eq:kdes-atlas} et dans le calendrier qui les produit. L'implémentation finie emploie $M$ pour décider de l'hétérotypie, tandis que l'archive généalogique conserve les données de provenance plus riches.

\subsection{Énumération de cylindres et candidats décimaux}

\subsubsection{Concaténation de soixante-douze et soixante-dix-huit trits}

Pour un ordonnancement $\tau=(s_1,\ldots,s_8)$ de
$\mathcal S_{\rm term}$, on définit
\[
 P_{72}(\tau)=[\Omega_{24}\Vert s_1\Vert\cdots\Vert s_8]_3.
\]
De manière récursive, on part de $p_0=[\Omega_{24}]_3$ et on calcule
$p_i=729p_{i-1}+[s_i]_3$ pour $1\leq i\leq8$. Cette récurrence
conserve les zéros initiaux de chaque bloc par la multiplication
explicite par $729$. Pour chaque position $b$ du panneau fonctionnel,
\[
 P_{78}(\tau,b)=729P_{72}(\tau)+b.
\]
Les longueurs sont respectivement $24+8\cdot6=72$ et $72+6=78$.
Le cylindre ternaire correspondant est
\begin{equation}
 I_{\tau,b}=
 \left[\frac{P_{78}(\tau,b)}{3^{78}},
       \frac{P_{78}(\tau,b)+1}{3^{78}}\right).
 \label{act21:eq:kdes-cilindro}
\end{equation}

\subsubsection{Intersection exacte avec la grille de douze triades}

La lecture décimale utilise douze triades, c'est-à-dire la grille $N/10^{36}$. On utilise les abréviations $D=10^{36}$ et $Q=3^{78}$.

\begin{lemma}[Extrémités entières du cylindre]
Pour un préfixe entier $p$ de longueur soixante-dix-huit, les numérateurs
des points de la grille qui appartiennent à son cylindre sont
\begin{equation}
 \left\{N\in\mathbb Z:
 \left\lceil\frac{pD}{Q}\right\rceil
 \leq N<
 \left\lceil\frac{(p+1)D}{Q}\right\rceil\right\}.
 \label{act21:eq:kdes-rejilla}
\end{equation}
Chaque cylindre contient au plus un point de cette grille.
\end{lemma}
\begin{proof}
La condition $N/D\in[p/Q,(p+1)/Q)$ est équivalente à
$pD/Q\leq N<(p+1)D/Q$. Pour $N$ entier, la première inégalité
est équivalente à $\lceil pD/Q\rceil\leq N$ ; la seconde est équivalente à
$N<\lceil(p+1)D/Q\rceil$. La longueur de l'intervalle en coordonnée
$N$ est $D/Q<1$, car $10^{36}<3^{78}$. Par conséquent, il peut contenir,
au maximum, un entier.
\end{proof}

Les plafonds sont évalués par division entière :
$\lceil a/Q\rceil=\lfloor(a+Q-1)/Q\rfloor$ pour $a\geq0$.
Toute l'énumération utilise des entiers exacts ; le choix des points de frontière est fixé par l'extrémité droite semi-ouverte.

On conserve d'abord la liste indexée d'incidences
\[
 \mathcal C_{\rm ind}=
 \{(\tau,j,N):N/D\in I_{\tau,\mathcal P_9(j)}\},
 \qquad1\leq j\leq9,
\]
et l'on obtient ensuite l'ensemble des numérateurs
\[
 \mathcal C=\{N:\exists\tau,j\ (\tau,j,N)\in\mathcal C_{\rm ind}\}.
\]
Le passage à $\mathcal C$ élimine des répétitions d'un numérateur, conservées
dans $\mathcal C_{\rm ind}$ avec leur origine.

\begin{lemma}[Récupération de l'ordre à partir du point du cylindre]
Si $(\tau,j,N)\in\mathcal C_{\rm ind}$, le bloc numéro $i$ du
mot de soixante-dix-huit trits se récupère par
\[
 \left\lfloor\frac{N\,729^i}{D}\right\rfloor\bmod729,
 \qquad1\leq i\leq13.
\]
En particulier, les blocs cinquième à douzième récupèrent $\tau$.
\end{lemma}
\begin{proof}
L'appartenance au cylindre signifie
$p\leq QN/D<p+1$, où $p=P_{78}(\tau,\mathcal P_9(j))$.
Par conséquent, $\lfloor QN/D\rfloor=p$. L'extraction des blocs
précédents coïncide avec la division successive du préfixe $p$ par
des puissances de $729$ : tronquer une expansion positionnelle avant d'extraire
un bloc conserve ce bloc. La formule écrite réalise
exactement cette extraction. Les quatre premiers blocs forment
$\Omega_{24}$ et les huit suivants sont l'ordonnancement $\tau$.
\end{proof}

\begin{proposition}[Recensement exact des candidats]
Pour $\Omega_{24}$, $\mathcal S_{\rm term}$ et $\mathcal P_9$
définis précédemment,
\[
 \#\operatorname{Ord}(\mathcal S_{\rm term})=40320,
 \qquad\#\mathcal C_{\rm ind}=21902,
 \qquad\#\mathcal C=19446.
\]
\end{proposition}
\begin{proof}
Les huit blocs sont distincts, donc leurs ordonnements sont les
$8!=40320$ permutations. Pour chaque permutation et chacune des
neuf positions du panneau, on calcule $p=P_{78}(\tau,b)$, on applique
les extrêmes \eqref{act21:eq:kdes-rejilla} et on ajoute son unique entier lorsque
l'intervalle entier n'est pas vide. La somme des cardinaux de ces
intervalles est $21902$. La réunion de leurs entiers contient $19446$
éléments distincts. L'évaluation exhaustive de ces opérations
est certifiée par les théorèmes
\texttt{indexed\_cylinder\_count} et \texttt{candidate\_count} du
sélecteur fini. L'algorithme réalise exactement l'énumération
décrite : permutations, neuf incidences du panneau, limites par
division entière et élimination des doublons.
\end{proof}

Les nombres $21902$ et $19446$ comptent des objets différents. Le premier
compte des incidences indexées qui contiennent un point décimal ; le second
compte des numérateurs distincts. Le nombre total de requêtes de cylindre
est $40320\cdot9=362880$, y compris les cylindres dont l'intersection avec
la grille est vide.

\subsection{Incidence exceptionnelle et condition d'hétérotypie}

\subsubsection{Face exceptionnelle produite par le panneau}

On élève une bande à douze coordonnées par le biais de
\[
 \Gamma(b)=(b,W(b))\in\mathbb F_3^{12}.
\]
Le support d'un mot est l'ensemble des positions de ses coordonnées non nulles, avec les indices humains $1,\ldots,12$. Les premières bandes régionales du panneau produisent
\begin{equation}
 \mathcal B=\operatorname{supp}\Gamma(B_\pi(0))
          \cap\operatorname{supp}\Gamma(B_e(0))
          =\{4,6,9,10\}.
 \label{act21:eq:kdes-cara}
\end{equation}
L'opération utilise les mots régionaux $103$ et $523$ et la matrice de Witt du sélecteur. La face se calcule avant de consulter les coordonnées des candidats.

Pour $N\in\mathcal C$, ses douze triades sont
\begin{equation}
 K_i(N)=\left\lfloor\frac{N}{1000^{12-i}}\right\rfloor\bmod1000,
 \qquad1\leq i\leq12.
 \label{act21:eq:kdes-triadas}
\end{equation}
Le support de couronne est
\[
 \mathcal H(N)=\{i:K_i(N)\geq729\},
 \qquad
 \mathcal C_{\mathcal B}=\{N\in\mathcal C:\mathcal H(N)=\mathcal B\}.
\]
La frontière $729=3^6$ relie la bande de six trits avec la triade décimale. L'évaluation exhaustive de l'égalité des supports donne
\begin{equation}
 \#\mathcal C_{\mathcal B}=79.
 \label{act21:eq:kdes-79}
\end{equation}
Le support conserve les positions, tandis que le registre complet
conserve aussi les douze grandeurs. L'étape suivante utilise cette
information positionnelle conjointement avec des relations entre grandeurs.

\subsubsection{Mots ternaires et différences orientées du candidat}

La $i$-ième bande ternaire du point $N/D$ s'obtient par
\begin{equation}
 T_i(N)=\left\lfloor\frac{N\,729^i}{D}\right\rfloor\bmod729,
 \qquad1\leq i\leq13.
 \label{act21:eq:kdes-tercandidato}
\end{equation}
L'arête orientée $a\to b$ a un résidu
\begin{equation}
 R_{a,b}(N)=(K_a(N)-K_b(N))\bmod729.
 \label{act21:eq:kdes-resarista}
\end{equation}
La soustraction s'effectue en $\mathbb Z$ avant la réduction. Pour cette
arête, on consulte dans les relations \eqref{act21:eq:kdes-atlas} le couple
\[
 \bigl(\operatorname{enc}_6^{-1}(T_b(N)),
  \operatorname{enc}_6^{-1}(R_{a,b}(N))\bigr).
\]
C'est donc une consultation conjointe sur le mot cible et sur la différence orientée des triades.

On écrit $\mathrm{Cmp}(a,b;N)$ lorsque le couple appartient à $\mathcal A_{\rm cmp}$ et $\mathrm{Af}(a,b;N)$ lorsqu'il appartient à $\mathcal A_{\rm af}$.

\subsubsection{Détermination de l'orbite transversale}

Le descripteur de vingt-quatre trits détermine les trois premières triades décimales avant les filtrations terminales. En effet, son cylindre satisfait l'inclusion exacte
\[
 \left[\frac{66241910521}{3^{24}},
       \frac{66241910522}{3^{24}}\right)
 \subseteq
 \left[\frac{234543140}{10^9},
       \frac{234543141}{10^9}\right).
\]
Les deux inégalités aux extrémités se vérifient par des produits croisés
entiers. Tout candidat partage donc les trois positions de
référence $1,2,3$ et leurs représentations $(234,543,140)$.

La translation de pas quatre sur $\mathbb Z/12\mathbb Z$ possède les
quatre orbites
\[
 (1,5,9),\qquad(2,6,10),\qquad(3,7,11),\qquad(4,8,12).
\]
Les trois premières contiennent respectivement l'une des positions de référence ; la quatrième est la seule disjointe de $\{1,2,3\}$. La forêt orientée des pas trois et quatre est fixée au moyen des arêtes
\[
 \begin{array}{lll}
 1\to4,&1\to5,&2\to6,\\
 3\to7,&4\to8,&5\to9,\\
 6\to10,&7\to11,&8\to12.
 \end{array}
\]
Sa première arête a un pas de trois et les autres ont un pas de quatre. Les deux
arêtes internes de l'orbite distinguée sont précisément
$4\to8$ et $8\to12$. Ainsi est fixé le lieu de la condition de
lecture par la géométrie des positions et par le préfixe initial,
avant d'examiner les candidats de la face exceptionnelle.

\subsubsection{Condition terminale d'hétérotypie}

La condition terminale est
\begin{equation}
 \begin{split}
 \mathcal T(N)={}&
 [\mathrm{Cmp}(4,8;N)\wedge\mathrm{Af}(8,12;N)]\\
 &\vee[\mathrm{Af}(4,8;N)\wedge\mathrm{Cmp}(8,12;N)].
 \end{split}
 \label{act21:eq:kdes-heterotipia}
\end{equation}
Les deux affectations de types aux deux arêtes sont admises. La
sélection exige qu'il existe une réalisation de types distincts ; les
étiquettes individuelles et leurs temps restent disponibles dans les
témoins d'incidence.

\begin{theorem}[Sélection terminale du registre de couplage]
Sur le domaine de candidats produit par
$\Omega_{24}$, $\mathcal S_{\rm term}$, $\mathcal P_9$ et le calendrier
régional de cent lignes, existe un unique numérateur qui satisfait
simultanément l'incidence exceptionnelle et l'hétérotypie :
\[
 \{N\in\mathcal C_{\mathcal B}:\mathcal T(N)\}
 =\{234543140729659824621058914794146601\}.
\]
Son registre de douze triades est
\begin{equation}
 K=(234,543,140,729,659,824,621,058,914,794,146,601).
 \label{act21:eq:kdes-kfinal}
\end{equation}
\end{theorem}
\begin{proof}
Le recensement précédent produit les $79$ éléments de
$\mathcal C_{\mathcal B}$. Pour chacun, on extrait ses douze
coordonnées par \eqref{act21:eq:kdes-triadas} ; on calcule les bandes de
destination $T_8,T_{12}$ et les résidus $R_{4,8},R_{8,12}$ ; on consulte
les deux masques du répertoire temporel et on applique la disjonction
\eqref{act21:eq:kdes-heterotipia}. La liste résultante contient exactement
l'entier indiqué, selon l'évaluation exhaustive certifiée par
\texttt{terminal\_selection}. L'égalité de liste avec un singleton
implique existence et unicité dans le domaine déclaré. L'extraction
base mille de cet entier donne \eqref{act21:eq:kdes-kfinal}.

La procédure de sélection reçoit les producteurs et les règles énumérées. L'entier final intervient dans l'énoncé de l'égalité vérifiée, après avoir exécuté la sélection. Cette position logique permet de comparer le résultat sans l'utiliser comme critère de choix.
\end{proof}

\subsubsection{Témoin arithmétique d'appartenance à un cylindre}

Une des incidences du numérateur sélectionné utilise la mise en ordre
\[
 (002001,020111,200110,121012,010122,001001,221110,222110)
\]
et la frontière $438=121020_3$. Pour elle,
\[
 P_{72}=5283881584882673136514102926090957,
\]
\[
 P_{78}=3851949675379468716518781033120308091.
\]
La formule des extrémités entières produit
\begin{align*}
 \left\lceil\frac{P_{78}10^{36}}{3^{78}}\right\rceil
 &=234543140729659824621058914794146601,\\
 \left\lceil\frac{(P_{78}+1)10^{36}}{3^{78}}\right\rceil
 &=234543140729659824621058914794146602.
\end{align*}
L'unique entier de l'intervalle est donc $N_K$. Les treize bandes ternaires de son point décimal sont
\[
 \begin{array}{rrrr}
 020022&222111&211201&021101\\
 002001&020111&200110&121012\\
 010122&001001&221110&222110\\
 121020&&&
 \end{array}
\]
Ce calcul exhibe un témoin concret d'appartenance au domaine énuméré. L'unicité du registre procède de l'examen exhaustif de tous les candidats, outre ce témoin particulier.

\subsection{Composition régionale et continuité des réalisations}

\subsubsection{Composition des générateurs régionaux}

Soit $\operatorname{Sel}(\Omega,\mathcal S,\mathcal P,\mathcal N)$ le
sélecteur défini par l'énumération et les deux filtres précédents.
Les producteurs régionaux satisfont les égalités
\[
 \Omega_{\rm reg}=\Omega_{24},\qquad
 \mathcal P_{\rm reg}=\mathcal P_9,\qquad
 \mathcal N_{\rm reg}=\mathcal N_{100},
\]
où $\mathcal N_{\rm reg}$ se calcule par extraction de bandes et
lecture de signatures. Par congruence de fonctions,
\begin{equation}
 \begin{split}
 &\operatorname{Sel}(\Omega_{\rm reg},\mathcal S_{\rm term},
                    \mathcal P_{\rm reg},\mathcal N_{\rm reg})\\
 &\hspace{10mm}=
 \operatorname{Sel}(\Omega_{24},\mathcal S_{\rm term},
                    \mathcal P_9,\mathcal N_{100})=\{N_K\}.
 \end{split}
 \label{act21:eq:kdes-composicion}
\end{equation}
L'égalité remplace trois interfaces documentaires par leur construction
régionale effective. Le répertoire terminal conserve la provenance
spécifiée dans \eqref{act21:eq:kdes-sterm}.

Le registre se forme alors par les douze divisions euclidiennes de
\eqref{act21:eq:kdes-triadas}. Chaque coordonnée appartient à
$\{0,\ldots,999\}$ et la longueur est douze par construction. L'
implémentation qui prend le premier élément de la liste sélectionnée
est justifiée par l'égalité avec le singleton : la valeur par défaut
d'une opération informatique sur des listes vides n'intervient jamais dans
ce résultat.

\subsubsection{Relation avec les incidences signées et l'inversion}

La construction présente fournit une sélection prospective dans un
domaine fini explicite. Le développement ultérieur du registre
dodécaphasique produit ses canaux orientés à partir de l'état
terminal enrichi et démontre une inversion intégrale par
Hadamard. Cette inversion récupère un registre à partir de canaux
compatibles ; ici, on sélectionne un registre à partir de ses
antécédents régionaux et incidentiels. L'égalité de ses résultats
permet de composer les deux réalisations en maintenant l'origine de chaque
opération.

L'unicité de sélection et l'unicité d'inversion sont des propriétés distinctes. La première quantifie les candidats construits dans ce chapitre. La seconde quantifie les registres qui partagent une signature compatible. Leur conjonction justifie que le registre sélectionné puisse être transporté vers la réalisation signée et être récupéré à partir d'elle, sans convertir une récupération rétrospective en producteur du registre.

\subsubsection{Relation entre les recensements terminaux}

La suite
\[
 40320\ \longrightarrow\ 21902\ \longrightarrow\ 19446
 \ \longrightarrow\ 79\ \longrightarrow\ 1
\]
enregistre successivement des ordonnancements, des incidences indexées avec la grille, des numérateurs distincts, des candidats de la face exceptionnelle et des survivants hétérotypiques. La notation à flèches indique l'ordre des opérations ; le deuxième nombre compte une classe d'objets différente de la première.

Le recensement des catalogues régionaux
$196\cdot197\cdot196\longrightarrow331\longrightarrow2
\longrightarrow1$, développé à l'endroit correspondant, part d'une
autre population : blocs régionaux de trois catalogues et conditions
terminales de Gram, d'occupation, de distance et d'orientation. Les deux recensements
conservent leurs domaines. Leur convergence vers une même représentation
s'exprime par l'identification de leurs lecteurs et de leurs sorties ;
les chiffres intermédiaires décrivent des opérations différentes et restent
documentés séparément.

\subsubsection{Composition régionale de la constante de structure fine}

La sortie $K$ intervient ensuite dans la composition régionale
ordonnée et dans la retenue en base mille. Son emploi ultérieur conserve les
douze positions et la frontière de lecture. Le couplage avec la racine
analytique et les représentations à profondeur arbitraire utilise cette
même sortie, avec les coordonnées régionales déjà engendrées et les
règles de la coordinatisation analytique. Cette dépendance prépare le chapitre
consacré aux deux réalisations corrélées de la constante de
structure fine.

La portée de ce chapitre est précise : production du descripteur et de ses antécédents régionaux, détermination du domaine fini, sélection unique sur ce domaine et composition matérielle avec les producteurs. La récurrence analytique de tous les degrés et la compatibilité des retenues infinies sont développées ensuite avec leurs démonstrations propres.

\subsection{Certification exécutable et traçabilité des opérations}

Le développement formel conserve la correspondance entre les
définitions mathématiques et leurs réalisations exécutables. Les preuves
des signatures et de l'horloge se trouvent dans
\texttt{N69RegionalSignature.lean} ; l'extraction des bandes et la
reproduction des cent lignes, dans
\texttt{GeneratedN69Rows.lean} ; le critère minimal et la construction
de $\Omega_{24}$, dans \texttt{RegionalW24.lean}. La compatibilité des
deux matrices de Witt est réunie dans
\texttt{WittReaderCompatibility.lean}.

Les lecteurs temporels, les rotations, les résidus orientés et la
condition hétérotypique sont formalisés dans
\texttt{TerminalReaders.lean}. L'énumération et les recensements sont
formalisés dans \texttt{Terminal}\allowbreak\texttt{Orbital}\allowbreak\texttt{Selection.lean}. Le panneau
régional se relie par \texttt{Regional}\allowbreak\texttt{Panel}\allowbreak\texttt{Bridge.lean}, et la
composition \eqref{act21:eq:kdes-composicion} s'exprime en
\texttt{Selected}\allowbreak\texttt{From}\allowbreak%
\texttt{Regional}\allowbreak\texttt{Inputs.lean}.

Les vérifications finies d'énumération utilisent l'évaluation native de Lean ; leurs rapports enregistrent \texttt{Lean.ofReduceBool}, ainsi que les axiomes logiques habituels de la bibliothèque. La reconstruction régionale du descripteur utilise des preuves du noyau et une réduction ordinaire. Cette distinction décrit les mécanismes de vérification réellement employés. Les reçus de compilation certifient ses modules et dépendances déclarés ; les preuves mathématiques conservent les domaines, producteurs et hypothèses qui ont été explicités dans le corps du chapitre.

\clearpage
% END NUCLEO_ADITIVO_20260921_SELECCION
\import{foundation/}{sections/registro_k.tex}
\import{foundation/}{sections/alpha.tex}
\import{foundation/}{sections/excepcional.tex}
\import{foundation/}{sections/01_hilbert_nonadico.tex}
\endgroup
\clearpage
\part{Géométrie positive, raffinement et dualité réticulaire}
\section{Orientation tritique, prolongement compatible et double projection}
\label{sec:trit-projection}

Le caractère engendré des coordonnées réelles s’exprime dans la relation
entre profondeur arithmétique et évaluation. La donnée primitive est une histoire
admissible d’opérations avec leurs restes, quotients et orientations. Une
lecture archimédienne associe à ses préfixes des intervalles de plus en plus étroits ;
une lecture incidentielle conserve des relations de frontière du même état.
La double projection du corpus \cite{hmtI,HMT-IX} réunit ces deux opérations
sans convertir une évaluation scalaire en la totalité de l’histoire.
L’article utilise cet antécédent pour étudier les réalisations ultérieures
du registre dodécaphasique.

\subsection{Cylindres positionnels et détermination de la coordonnée réelle}

Pour un bloc ternaire \(b=(b_1,\ldots,b_6)\), avec des représentants
\(b_j\in\{0,1,2\}\) de \(\mathbb F_3\), l’évaluation entière
\(\nu(b)=\sum_{j=1}^6b_j3^{6-j}\) appartient à
\(\{0,\ldots,728\}\). Si \(b_{n+1}\) est le bloc admissible suivant,
le prolongement positionnel vérifie
\[
 N_{n+1}=729N_n+\nu(b_{n+1}),\qquad
 a_n=\frac{N_n}{729^n},\qquad b_n^+=\frac{N_n+1}{729^n}.
\]
Ici, \(N_0\) est la partie entière que détermine le lecteur régional. La lettre \(b_n^+\) désigne une extrémité réelle et se distingue du bloc ternaire \(b\). Le cylindre de représentation est \([a_n,b_n^+)\) ; son adhérence est utilisée dans l'argument d'existence de la limite. La retenue conserve le passage d'un préfixe au précédent par division euclidienne.

\begin{proposition}[Limite des centres cylindriques]
Pour un prolongement compatible, les adhérences des cylindres sont
emboîtées et déterminent une unique coordonnée \(x\). Si
\(c_n=(a_n+b_n^+)/2\), alors
\[
 x=\lim_{n\to\infty}a_n=\lim_{n\to\infty}b_n^+
 =\lim_{n\to\infty}c_n,\qquad
 |x-c_n|\leq\frac1{2\,729^n}.
\]
\end{proposition}
\begin{proof}
L'inégalité \(0\leq\nu(b_{n+1})<729\) situe le cylindre successeur dans l'adhérence du prédécesseur. Ses extrémités sont rationnelles, sa largeur est \(729^{-n}\) et elles demeurent dans le premier intervalle fermé borné. L'emboîtement et la largeur décroissante produisent une unique classe de Cauchy d'extrémités. Dans la réalisation archimédienne, cette classe a pour représentant \(x\) ; la distance au point milieu est au plus égale à la moitié de la largeur. La convention de frontière attribue une écriture lorsque deux développements positionnels représentent la même limite.
\end{proof}

La moyenne arithmétique intervient donc dans une évaluation précise : le
centre de chaque intervalle qui borne la coordonnée. La limite conserve
l’exactitude de la valeur bien que la largeur disparaisse. La représentation
enrichie conserve en outre la suite des préfixes et leur incidence.
La projection numérique et la projection exceptionnelle sont deux applications
à codomaines différents ; l’une conserve la valeur limite et l’autre transporte
les relations qui permettent la construction réticulaire. Le passage à
\(\RR\) se situe ainsi dans la réalisation du lecteur, après la
génération arithmétique des données qu’il évalue.

Pour expliciter la compatibilité, soient \(p_{n+1,n}\) les restrictions de l'état enrichi et \(x_n=p_{n+1,n}(x_{n+1})\). Si \(I_n(x_n)\) est le cylindre numérique et \(\mathcal B_n(x_n)\) l'incidence de frontière, les deux transports conservent simultanément
\[
 \overline{I_{n+1}(x_{n+1})}\subseteq\overline{I_n(x_n)},\qquad
 r_{n+1,n}\mathcal B_{n+1}(x_{n+1})=\mathcal B_n(x_n).
\]
La seconde égalité utilise les applications d'incidence du corpus ; la proposition~\ref{prop:lf-incidence-naturality} explicite leur réalisation au moyen de repères et de codes transportés. Ses cinq constructions consubstantielles ---spectrale, solénoïdale, de jauge, cohomologique et déterminantale--- appartiennent à ce même système compatible. L'étude géométrique ultérieure conserve cette unité d'origine en spécifiant laquelle de ses observations elle utilise dans chaque théorème.

\subsection{Trajectoires orientées et temporalité de l’actualisation}

Soit \(x\xrightarrow{\rm adm}y\) la relation d’admissibilité du TPK dans
l’espace enrichi \(\widehat X\). Une trajectoire paramétrée par
un ensemble discrètement ordonné \(\mathcal T\) est
\[
 \gamma:\mathcal T\longrightarrow\widehat X,
 \qquad\gamma(t)\xrightarrow{\rm adm}\gamma(t^+).
\]
La relation de compatibilité détermine les transitions ; le paramétrage
permet de les ordonner. Dans l’espace \(E=\mathcal T\times\widehat X\),
avec \(p(t,x)=t\), la trajectoire détermine la section
\(\widetilde\gamma(t)=(t,\gamma(t))\), qui vérifie
\(p\circ\widetilde\gamma=\id_{\mathcal T}\).
Cette précision conserve la différence entre un état, sa trajectoire et
la coordonnée utilisée pour décrire l’évolution.

La signature tritique \(\tau\in\{+1,0,-1\}\) type les générateurs au moyen de
\[
 J_\tau^2=-\tau I.
\]
Dans le lecteur temporel du corpus, \(+1\) représente le retour et la mémoire ; \(0\), la mise à jour de seuil ; et \(-1\), l'ouverture de prolongements conjugués. Leur interprétation comme passé, présent et futur se rapporte aux fonctions de l'antécédent conservé, de la transition effective et de la continuation admissible. Les trois temps verbaux acquièrent ainsi un contenu opératoriel sur les histoires.

Si \(\mathcal P_n(x)\) est l'ensemble des continuations admissibles de longueur \(n\), la suppression de la dernière étape définit \(r_n:\mathcal P_{n+1}(x)\to\mathcal P_n(x)\). Une continuation illimitée est une collection \((\gamma_n)\) avec \(r_n(\gamma_{n+1})=\gamma_n\). La restriction reconstruit les antécédents ; l'actualisation \(U_t\) agit dans la transition ; le prolongement maintient la frontière du préfixe. L'existence de ces collections est utilisée dans les domaines survivants fixés par le théorème de prolongement du corpus.

Sur une signature finie \(\sigma=(\tau_1,\ldots,\tau_n)\), l’inversion
orientée est
\[
 C_{\rm rev}\sigma=(-\tau_n,\ldots,-\tau_1),
 \qquad C_{\rm rev}^2\sigma=\sigma.
\]
La trajectoire conjuguée conserve également les transformations de feuillet, de position et de frontière de son transport. La bidirectionnalité s'exprime donc dans l'historique complet, tandis que cette formule exprime sa signature. Le retour de phase de \(\Gamma_9\) admet un incrément de mémoire : des états de phase identique peuvent occuper des profondeurs différentes.

\subsection{Algèbres quadratiques et activité du seuil parabolique}

Dans un régime fixe, la conjugaison de \(aI+bJ_\tau\) change le signe de \(b\), et sa norme algébrique est
\[
 (aI+bJ_\tau)(aI-bJ_\tau)=(a^2+\tau b^2)I.
\]
Les types elliptique, parabolique et hyperbolique correspondent aux trois
relations quadratiques. Leurs groupes locaux sont représentés par
\[
 e^{tJ_{+1}}=\cos t\,I+\sin t\,J_{+1},\quad
 e^{tJ_0}=I+tJ_0,\quad
 e^{tJ_{-1}}=\cosh t\,I+\sinh t\,J_{-1}.
\]
La preuve sépare les termes pairs et impairs de la série exponentielle ; dans le type nilpotent, seuls demeurent les degrés zéro et un. Chaque collection satisfait la loi de composition dans son paramètre et l'inversion \(e^{-tJ_\tau}=(e^{tJ_\tau})^{-1}\).

La valeur tritique zéro permet en particulier un générateur distinct de
zéro dont le carré est nul. Pour
\(N=\left(\begin{smallmatrix}0&1\\0&0\end{smallmatrix}\right)\),
\[
 e^{tN}(q,p)=(q+tp,p).
\]
Le seuil possède une action linéaire effective sur sa fibre. Son paramètre,
son déplacement et la mémoire de la transition restent déterminés
même lorsque la signature de régime est nulle. Tel est le contenu précis qui
permet de développer dans l’exposé le présent comme mise à jour active.

La forme elliptique positive a des niveaux circulaires en coordonnées
orthonormées et elliptiques après une transformation linéaire ; la forme hyperbolique
possède deux directions isotropes et des niveaux hyperboliques ; le type parabolique
possède une forme dégénérée et un transport unipotent. Une géométrie sphérique
requiert en outre l’espace homogène ou la métrique qui la réalise. La
classification précédente fixe les générateurs qui peuvent intervenir dans cette
réalisation et maintient séparés le type algébrique et la courbure d’un
espace ultérieur.

Le calendrier des modes concrétise la connexion avec l'auto-échelle. La règle \(+1:m\mapsto\Sigma\), \(-1:m\mapsto\Pi\), \(0:m\mapsto m\) appliquée au bloc cyclique \((+1,-1,0)\) produit \((\Sigma,\Pi,\Pi)\). Compter ses trois transitions, dans l'ordre \((\Sigma,\Pi)\), détermine
\begin{equation}
 F_{\rm av}=\begin{pmatrix}0&1\\1&1\end{pmatrix}.
 \label{eq:geo-fav}
\end{equation}
Son rayon propre positif, normalisé sous la forme \((1,x)^\mathsf T\), satisfait \(x=\lambda\), \(1+x=\lambda x\), et donc \(x^2-x-1=0\). La racine positive est la coordonnée d'auto-échelle exprimée par \(\varphi\). La matrice provient du calendrier des opérations ; l'expression radicale et les identités exponentielles suivantes sont ses réalisations ultérieures.

% Articulación posterior a la generación de pi, phi y e.
% Contenido acumulativo; no reemplaza el desarrollo previo del TRIT.
\subsection{Réalisations exponentielles des coordonnées générées}
\label{geo-euler-trit-articulacion}

Les coordonnées régionales déjà déterminées permettent de reconnaître les
réalisations concrètes de la collection quadratique du TRIT. L'ordre est
essentiel : la construction des régions précède l'évaluation d'une
exponentielle aux paramètres \(\pi\) ou \(2\log\varphi\). Ces expressions
identifient la fonction des coordonnées dans des opérateurs construits à partir
d'APP et du TPK ; elles n'interviennent pas dans la sélection des préfixes qui les produisent.

\subsubsection{Cycle ternaire, transport nilpotent et incidence des feuillets}

La multiplication \(a(r)=4r\) dans \(\mathbb Z/9\mathbb Z\) est d’ordre
\(3\). Sur les unités, elle détermine les orbites \(\{1,4,7\}\) et
\(\{2,8,5\}\). Dans le module d’augmentation de l’une quelconque d’entre elles, formé
des combinaisons entières dont les coefficients ont une somme nulle, la base
\((e_r-e_{4r},e_{4r}-e_{7r})\) donne
\[
 C=\begin{pmatrix}0&-1\\1&-1\end{pmatrix},
 \qquad C^2+C+I=0.
\]
La forme restreinte a pour matrice de Gram
\(\bigl(\begin{smallmatrix}2&-1\\-1&2\end{smallmatrix}\bigr)\).
Dans les évaluations résiduelles modulo \(9\), l’action distingue les feuillets :
\(S(ax,ay)=aS(x,y)\), tandis que \(P(ax,ay)=a^{-1}P(x,y)\).
La réalisation cyclique conserve donc une orientation contragrédiente
entre somme et produit. Les quotients du relèvement entier sont conservés
dans le registre arithmétique correspondant.

Le transport parabolique élémentaire est représenté par
\[
 N=\begin{pmatrix}0&1\\0&0\end{pmatrix}.
\]
L’incidence surfacique--volumique déjà construite dans \eqref{eq:geo-fav} fournit
\[
 F_{\rm av}=\begin{pmatrix}0&1\\1&1\end{pmatrix},
 \qquad A=F_{\rm av}^2=\begin{pmatrix}1&1\\1&2\end{pmatrix}.
\]
Les trois opérateurs agissent dans des espaces réels déclarés : le module
d’augmentation tensorisé avec \(\mathbb R\), le plan du transport nilpotent et
le plan d’incidence modale. Le partage d’une dimension matricielle n’identifie
ni ces espaces ni leurs coordonnées.

\begin{proposition}[Générateurs concrets des trois types]
\label{geo-euler-trit-generadores}
Les opérateurs
\[
 J_{\mathrm e}=\frac{2C+I}{\sqrt3},\qquad
 J_{\mathrm p}=N,\qquad
 J_{\mathrm h}=\frac{2A-3I}{\sqrt5}
\]
satisfont, respectivement,
\[
 J_{\mathrm e}^{\,2}=-I,\qquad
 J_{\mathrm p}^{\,2}=0,\qquad
 J_{\mathrm h}^{\,2}=I.
\]
\end{proposition}
\begin{proof}
La relation de \(C\) implique \((2C+I)^2=-3I\).
La multiplication directe donne \(N^2=0\).
Enfin, \(A\) est de trace \(3\) et de déterminant \(1\), d’où
\(A^2-3A+I=0\) et \((2A-3I)^2=5I\).
\end{proof}

\begin{proposition}[Trois évaluations de l’identité exponentielle]
\label{geo-euler-trit-evaluaciones}
Dans les réalisations précédentes,
\begin{equation}
 \boxed{
 \exp(\pi J_{\mathrm e})+I=0,\qquad
 \exp(J_{\mathrm p})=I+J_{\mathrm p},\qquad
 \exp(2\log\varphi\,J_{\mathrm h})=F_{\rm av}^{\,2}.}
 \label{geo-euler-trit-tres-evaluaciones}
\end{equation}
\end{proposition}
\begin{proof}
La spécialisation elliptique de l’exponentielle tritique donne
\(\cos\pi\,I+\sin\pi\,J_{\mathrm e}=-I\).
La nilpotence élimine exactement les termes de degré au moins \(2\)
dans la deuxième expression. Pour la troisième, \(\varphi^2=\varphi+1\) implique
\[
 \cosh(2\log\varphi)=\frac32,\qquad
 \sinh(2\log\varphi)=\frac{\sqrt5}{2}.
\]
L’exponentielle vaut donc
\(\frac32I+\frac{\sqrt5}{2}J_{\mathrm h}=A\).
\end{proof}

La première égalité représente le demi-tour elliptique ; le tour complet
satisfait \(\exp(2\pi J_{\mathrm e})=I\). La deuxième maintient un terme
linéaire non nul : le régime parabolique exprime un transport, et non une absence
d'évolution. La troisième réalise une avancée hyperbolique unimodulaire. Ses
valeurs propres sont \(\varphi^2\) et \(\varphi^{-2}\) ; une direction se dilate et
l'autre se contracte en conservant le déterminant. Les trois évaluations
emploient des paramètres distincts d'une collection exponentielle commune.

Dans cette collection, \(e\) intervient par l'évaluation scalaire
\(\exp(I)=eI\) et la loi \(\exp(sI)\exp(tI)=\exp((s+t)I)\).
Son rôle ne se limite pas au type parabolique. La coordonnée \(e\), obtenue
auparavant à partir de sa région, est reconnue ici comme l'évaluation en l'unité
de l'exponentielle ; cette identité ne remplace pas sa généalogie coinductive.

\subsubsection{Résolution polynomiale des régimes}

Les trois réalisations admettent une présentation compacte sur la somme
directe de leurs espaces. On définit
\[
 \mathbb J=J_{\mathrm e}\oplus J_{\mathrm p}\oplus J_{\mathrm h},
 \qquad Q=\mathbb J^2.
\]
Les relations précédentes impliquent
\[
 \mathbb J^6=\mathbb J^2,\qquad Q^3=Q.
\]
L’emploi de \(Q\) maintient ce carré opératoriel distinct du registre
dodécaphasique \(K\).

\begin{proposition}[Projecteurs polynomiaux de régime]
\label{geo-euler-trit-proyectores}
Les applications
\[
 p_{\mathrm e}=\frac{Q^2-Q}{2},\qquad
 p_{\mathrm p}=I-Q^2,\qquad
 p_{\mathrm h}=\frac{Q^2+Q}{2}
\]
sont des idempotents mutuellement orthogonaux, de somme \(I\), qui restituent les trois
termes de la somme directe représentant les régimes.
\end{proposition}
\begin{proof}
Les trois polynômes prennent les valeurs \((1,0,0)\), \((0,1,0)\) et
\((0,0,1)\) aux points \(-1,0,+1\), respectivement. Les différences
entre leurs carrés et eux-mêmes, ainsi que leurs produits croisés, sont
des multiples de \(X^3-X\). Substituer \(Q\), qui annule ce polynôme, démontre
les identités. Dans les trois blocs, \(Q\) agit comme \(-I,0,I\).
\end{proof}

On obtient ainsi
\begin{align}
 \exp(t\mathbb J)
 ={}&p_{\mathrm e}(\cos t\,I+\sin t\,\mathbb J)
   +p_{\mathrm p}(I+t\mathbb J)\notag\\
   &+p_{\mathrm h}(\cosh t\,I+\sinh t\,\mathbb J).
 \label{geo-euler-trit-exponencial-total}
\end{align}
Chaque terme conserve sa relation quadratique et la conjugaison inverse
son évolution. La représentation totale réunit les trois régimes, mais ne
définit pas à elle seule un opérateur de transition entre eux. Une telle transition
requiert les applications de transport du TPK avec leurs domaines, leur orientation et
leur mémoire. L’articulation exponentielle permet de reconnaître leurs types sans
remplacer cette dynamique par une identification des fibres.

% PROCEDENCIA FOCAL (no impresa)
% Resultado recuperado: tres generadores y tres evaluaciones exponenciales.
% Propietario completo:
% BASE/manuscrito/sections/hmt/09b_algebras_cuadraticas.tex:289-516.
% Antecedente causal de C y acción contragrediente:
% BASE/manuscrito/sections/hmt/03d_esqueleto_ciclotomico_app.tex:13-47,145-246.
% Los tres tipos y la orientación están en:
% BASE/manuscrito/sections/hmt/02_trit_geometrias.tex:83-196.
% Resolución polinómica recuperada:
% output/INVESTIGACION_HOLONOMIA_NONADICA_CRISTAL_APERIODICO_20260903/
% ALGEBRA_HOLONOMICA_UNIVERSAL_Y_EULER_TRITICO.md, secciones 8-11 y 22.
% Los cuatro propietarios se leyeron completos. BASE designa:
% /Users/ruben/Documents/HMT2/
% EL_CIERRE_HOLOGRAFICO_DEL_INFINITO_SUCESOR_102_CAPITULOS_2026-08-28_EN_TRABAJO/
% No se incorpora la identidad P9(alpha)=delta del documento especializado:
% esa identificación no pertenece a este bloque.
% Tampoco se identifica una dirección nilpotente con la vacancia electrónica,
% ni se presume una representación global del transporte por la suma directa.
% COMPROBACIONES: Python estándar, enteros y Fraction.
% C^2+C+I=0; (2C+I)^2=-3I; N^2=0; (2Fav^2-3I)^2=5I.
% 81 pares de residuos verifican las dos acciones contragredientes.
% Los 3 proyectores se verificaron en Q=-1,0,+1; prueba general en texto.
% 6 labels únicos y entornos equilibrados. Sin compilación.


\section{Transport du bloc modulaire au quotient d'incidence}
\label{sec:incidence-transport-new}
Le bloc central d'APP et l'incidence des faces et des sommets fournissent deux réalisations d'une même involution. Leur relation s'exprime au moyen d'un opérateur d'entrelacement rectangulaire, avant toute interprétation métrique. La congruence diagonale adjacente
\[
 k^2\equiv(k+1)^2\pmod9
 \quad\Longleftrightarrow\quad 2k+1\equiv0\pmod9
\]
sélectionne \(k=4\) dans la fenêtre déclarée. Le bloc correspondant est
\[
 B_c=\begin{pmatrix}7&2\\2&7\end{pmatrix}
 =7I_2+2R_c=9P_++5P_-,\qquad
 R_c=\begin{pmatrix}0&1\\1&0\end{pmatrix},\quad
 P_\pm=\frac{I_2\pm R_c}{2}.
\]
La différence entre les deux valeurs propres est de quatre. Le transport
incidenciel détermine ensuite comment cette différence se conserve
lorsque les multiplicités des secteurs changent.

Les faces du cube sont indexées par \((i,\epsilon)\), avec \(i\in\{1,2,3\}\) et \(\epsilon\in\{-1,1\}\) ; ses sommets sont indexés par \(\nu\in\{-1,1\}^3\). Nous définissons
\[
 M_{(i,\epsilon),\nu}=\mathbf1_{\{\nu_i=\epsilon\}}.
\]
Soient \(R_F\) l'opposition des faces et \(R_V\) l'inversion des sommets \(\nu\mapsto-\nu\). Dans chaque entrée, on vérifie
\[
 \mathbf1_{\{\nu_i=-\epsilon\}}
 =\mathbf1_{\{(-\nu)_i=\epsilon\}},\qquad R_FM=MR_V.
\]
La matrice \(M\) transporte donc l’involution complète.

\begin{theorem}[Descente incidentielle du bloc central]
Soient \(B_F=7I_6+2R_F\) et \(B_V=7I_8+2R_V\). Alors
\[
 B_FM=MB_V.
\]
Si \(u=(1,1,1,1,1,1)^{\mathsf T}\),
\(P_u=uu^{\mathsf T}/6\) et \(P_-=(I_6-R_F)/2\), le quotient
\[
 Q_\square=\RR^6/\ker(MM^{\mathsf T})
 \simeq\RR u\oplus E_-
\]
a pour dimension \(1+3\), et l’opérateur induit est
\[
 \overline B_F=9P_u+5P_-.
\]
\end{theorem}
\begin{proof}
La première identité résulte de la multiplication de \(R_FM=MR_V\) par deux et de l'ajout de \(7M\). Chaque face contient quatre sommets ; deux faces adjacentes en partagent deux et deux faces opposées n'en partagent aucun. De ce décompte, on obtient, entrée par entrée,
\[
 MM^{\mathsf T}=12P_u+4P_-.
\]
Le premier projecteur est de rang un et le second de rang trois ; leurs images
sont orthogonales. Les valeurs propres sont \(12,4,0\), avec les multiplicités
\(1,3,2\). L’opposition agit comme \(+1\) dans \(\RR u\) et comme
\(-1\) dans \(E_-\), et conserve le noyau de la matrice de Gram. Ainsi,
\(B_F\) descend au quotient et agit par neuf et cinq dans ces deux
facteurs directs. Son polynôme caractéristique est
\((\lambda-9)(\lambda-5)^3\).
\end{proof}

L’assignation temporelle à l’axe \(\RR u\) fixe la forme signée
\(J_L=P_--P_u\). Sur le quotient, \(I_Q=P_u+P_-\), de sorte que
\[
 J_L=\frac{7I_Q-\overline B_F}{2}.
\]
Cette égalité relie la forme de signature \((3,1)\), avec la
convention temporelle indiquée, et l’opérateur modulaire transporté. La
matrice de Gram initiale conserve sa positivité ; la signature est introduite
dans la réalisation ultérieure au moyen du signe assigné à chaque secteur.
Les dilatations neuf et cinq demeurent celles du bloc : l’identité
affine détermine la forme et ne transforme pas le bloc en une isométrie.

\subsection{Agrégation arithmétique, norme réticulaire et fonctionnelle de chemin}
La conservation de l'entier précédant la réduction devient visible lorsque l'on additionne les deux feuillets. Les matrices d'agrégation modulaire sont
\[
 M_\Pi=\begin{pmatrix}4&5&6\\5&4&6\\6&6&9\end{pmatrix},\qquad
 M_\Sigma=\begin{pmatrix}5&6&4\\6&4&5\\4&5&6\end{pmatrix}.
\]
Leurs sommes sont cinquante et un et quarante-cinq. Lors de la reconstruction des
neuf éléments de chaque bloc, l’entier et sa retenue vérifient
\[
 2025-810=(459-405)+9(174-45)=54+9\cdot129.
\]
L'agrégat résiduel et le quotient constituent deux coordonnées de cette égalité. Dans la construction réticulaire ultérieure, la valeur cinquante-quatre apparaît comme la norme de la chambre marquée \(v_\alpha=4\rho_1+\rho_2+\cdots+\rho_{12}\), pour des vecteurs orthogonaux \(\rho_j^2=2\). Dans la réalisation matérielle, elle apparaît comme l'évaluation de la fonctionnelle de chemin \(D_1\) du profil \((80,54,6)\). Chaque apparition conserve son application d'origine : agrégation avec retenue, forme réticulaire et évaluation orientée de chemin. La coïncidence de leurs valeurs acquiert un contenu structurel grâce à ces applications, dont les domaines et les preuves sont développés dans leurs sections respectives.


\section{Extraction dodécaphasique et projection isotypique de la direction}
\label{sec:register}
Avant d’évaluer des nombres, les lecteurs distinguent des régions dynamiques.
La symétrie spéculaire du corpus agit sur le bloc
\(B=(u^\pi,u^e,u^\varphi)\in(\mathbb F_3^6)^3\) par
\[
 \mathfrak cB=(u^e,u^\pi,u^\varphi).
\]
Avec \(q=u^\pi+u^e+u^\varphi\), \(a=u^\pi+u^e-u^\varphi\) et \(d=u^\pi-u^e\), les formules
\[
 u^\varphi=2(q-a),\quad s=q-u^\varphi,\quad
 u^\pi=2(s+d),\quad u^e=2(s-d)
\]
sont des identités composante par composante dans \(\mathbb F_3\). Le miroir conserve l'axe et l'agrégat dans une fibre et inverse la répartition antisymétrique. Le transport nonadique fait évoluer la fibre elle-même. La cinquième phase a pour donnée de frontière \(111111\), pour signature résiduelle \(000\), une solution locale et une extension : la marque de \(\varphi\) en cinq identifie l'axe régional de cette opération. Dans les phases six, sept et huit, les agrégats régionaux sont \(012100\), \(101001\) et \(112000\), respectivement. Ces données appartiennent au domaine ternaire ; les lecteurs d'histoires compatibles produisent ensuite leurs évaluations réelles.

Le registre des douze fenêtres et les blocs régionaux conservent ainsi leurs domaines et leur lien dans le même état. La réalisation positive qui suit utilise le registre exprimé, avec son ordre et son échelle. Pour transporter le miroir régional vers une déformation géométrique, il faut d'abord composer son action sur ce registre ; l'égalité de généalogie exprime à elle seule moins que la commutativité de ce carré. Cette précision permet de conserver la dynamique régionale complète en étudiant l'une de ses réalisations finies.

Le calendrier divise les cent huit positions en douze fenêtres nonadiques.
L’extraction terminale possède la composition
\[
 x_{\rm term}\longmapsto\mathcal R_{12}(x_{\rm term})
 \longmapsto(b^{90},b^{120},Q)\longmapsto U_{\rm sgn}\longmapsto K.
\]
Les deux premières flèches conservent l'histoire orientée dans les lecteurs du corpus. Pour montrer intégralement la reconstruction finie qu'utilise cet article, leurs sorties entières sont
\begin{align*}
 b^{90}&=(-495,-116,-684,108,601,-90,-173,-88,313,560,-397,461),\\
 b^{120}&=(-425,-281,-481,671,-255,30,475,-543,680,251,6,-128),
 \qquad Q=6263.
\end{align*}
On fixe l’ordre cyclique de douze positions. Soient
\(B_r=\sum_{j=0}^3 b^{120}_{r+3j}\) et
\(d_{r,j}=b^{90}_{r+3j}\), pour \(r=1,2,3\). Alors
\[
 s_1=\frac{Q+2B_1+B_2}{3},\quad s_2=s_1-B_1,\quad s_3=s_2-B_2,
 \qquad
 U^{(r)}=(s_r,d_{r,1}-d_{r,3},d_{r,0}+d_{r,2},d_{r,0}-d_{r,2}).
\]
L’évaluation donne
\[
 U^{(1)}=(2378,-452,-668,-322),\quad
 U^{(2)}=(1406,998,-204,-28),\quad
 U^{(3)}=(2479,-551,-371,-997).
\]
La matrice
\[
 H_4=\begin{pmatrix}1&1&1&1\\1&1&-1&-1\\1&-1&1&-1\\1&-1&-1&1\end{pmatrix},
 \qquad H_4^2=4I_4,
\]
reconstruit chaque orbite au moyen de \(H_4U^{(r)}/4\). Restituer l'ordre des positions produit
\begin{equation}
 K=(234,543,140,729,659,824,621,58,914,794,146,601),
 \qquad Q=\sum_{j=1}^{12}K_j=6263.
 \label{eq:K-main}
\end{equation}
Chaque composante est positive. La division par trois provient des trois
orbites et la division par quatre de l’inverse de Hadamard. Le contrôle
direct supplémentaire est \((I-S_3)K=b^{90}\),
\((I-S_4)K=b^{120}\), où \((S_aK)_j=K_{j+a}\) avec des indices cycliques.
Ces identités rendent la reconstruction reproductible sans sélectionner
des coefficients au moyen d’une géométrie cible.

La coordinatisation incidentielle du corpus réalise les douze positions comme classes latérales de \(A_5/C_5\). Si \(\rho\) est cette représentation et \(\chi_3\) son caractère réel tridimensionnel, son projecteur est
\[
 P_3=\frac3{60}\sum_{a\in A_5}\chi_3(a^{-1})\rho(a),\qquad
 P_{11}=I-\frac1{12}\mathbf1\mathbf1^{\mathsf T}.
\]
L'énumération et l'orientation des classes latérales sont conservées dans la source de cette coordinatisation. Pour que l'extension géométrique soit directement évaluable, on écrit sa sortie exacte \(u=P_3P_{11}K\), avec \(r=\sqrt5\) :
\begin{equation}
\begin{split}
u=\bigl(&-495/4-233r/10,\quad403/4+169r/10,\\
&-403/4-169r/10,\quad495/4+233r/10,\\
&601/4+1031r/10,\quad203/4+211r/5,\\
&-203/4-211r/5,\quad-601/4-1031r/10,\\
&313/4+569r/10,\quad162+219r/20,\\
&-162-219r/20,\quad-313/4-569r/10\bigr).
\end{split}
\label{eq:u-main}
\end{equation}
On vérifie par somme et produit dans \(\QQ(\sqrt5)\) que
\[
 \mathbf1^{\mathsf T}u=0,\qquad
 \|u\|^2=\frac{6638585+2275584\sqrt5}{20}>0.
\]
En particulier, ses coordonnées ne sont pas constantes. La direction \(u\), le registre intégral \(K\) et le projecteur \(P_{10}=P_{11}-uu^{\mathsf T}/\|u\|^2\) conservent des quantités différentes d'information. Cette distinction permet d'utiliser la direction pour un flux sans la substituer au registre qu'emploient les lecteurs d'action.

\section{Reconstruction de la géométrie positive à partir de la mémoire nonadique}
\label{sec:memory}

L'extraction dodécaphasique conserve davantage d'information qu'une direction de déformation. La direction sélectionne un mouvement ; le registre complet détermine également la partition sur laquelle ce mouvement agit. La récupération du registre à partir de ses mémoires permet de formuler cette distinction comme un théorème de reconstruction. Le résultat découle de la composition des compressions nonadiques du corpus \cite{hmtX} ; son application géométrique conserve la charge uniforme et les compléments des deux étapes.

Dans cette section, les coordonnées sont numérotées de zéro à onze. Soit \(\mathsf S:\RR^{12}\to\RR^{12}\) le décalage \((\mathsf Sx)_i=x_{i+1}\), avec indices modulo douze. Dans la coordinatisation euclidienne postérieure à l'état enrichi, on définit
\[
 T_j=\frac{8I+\mathsf S^j}{9},\qquad
 D_j=\frac{\sqrt8}{9}(\mathsf S^j-I),\qquad j\in\{3,4\}.
\]
Le coefficient huit compte les positions ordinaires de la coordinatisation nonadique ; la neuvième position contient le transport. L'orthogonalité du décalage implique, par multiplication directe,
\[
 T_j^*T_j+D_j^*D_j=I.
\]
La première mémoire est le complément de la première compression ; la seconde
est prise sur l’état qui a déjà traversé cette compression. Par conséquent,
\begin{equation}
 q=\mathbf1^*K,\qquad m_0=D_3K,\qquad m_1=D_4T_3K,
 \qquad MK=(m_0,m_1).
 \label{eq:mem-data}
\end{equation}
L’ordre de la composition fait partie de l’objet. L’application
\(M:\RR^{12}\to\RR^{24}\) conserve deux observations corrélées,
produites par une même trajectoire linéaire.

\begin{theorem}[Inversion exacte de la représentation de mémoire]
\label{thm:memory-inversion}
L'application \(K\mapsto(q,MK)\) est un isomorphisme linéaire entre \(\RR^{12}\) et \(\RR\times\operatorname{im}M\). Son inverse s'obtient au moyen de
\begin{align}
 T_3^{-1}&=\frac{512I-64\mathsf S^3+8\mathsf S^6-\mathsf S^9}{455},
 \label{eq:memory-T-inverse}\\
 b&=-\frac9{\sqrt8}m_0,\qquad
 c=-\frac9{\sqrt8}T_3^{-1}m_1,\qquad w_i=c_i-b_{i+1},\notag\\
 B_0&=0,\qquad B_i=\sum_{j=0}^{i-1}w_j,\qquad
 K_i=\frac{q+\sum_{j=0}^{11}B_j}{12}-B_i.
 \label{eq:memory-inverse}
\end{align}
\end{theorem}
\begin{proof}
En écrivant \(X=\mathsf S^3\), on a \(X^4=I\) et
\((8I+X)(512I-64X+8X^2-X^3)=4095I\). La division par neuf
prouve \eqref{eq:memory-T-inverse}. Tous les opérateurs utilisés sont des
polynômes du même décalage et commutent. Ainsi,
\[
 b=(I-\mathsf S^3)K,\qquad c=(I-\mathsf S^4)K,
 \qquad w_i=K_i-K_{i+1}.
\]
La somme télescopique donne \(B_i=K_0-K_i\) ; sommer cette égalité sur
les douze positions fournit \eqref{eq:memory-inverse}. En outre,
\(MK=0\) équivaut à l’invariance sous \(\mathsf S^3\) et
\(\mathsf S^4\). Comme \(\mathsf S=\mathsf S^4(\mathsf S^3)^{-1}\),
le noyau est exactement \(\RR\mathbf1\). La charge récupère cette
composante et démontre l’injectivité. La décomposition
\(K=(q/12)\mathbf1+K^\perp\) prouve la surjectivité sur le
codomaine déclaré.
\end{proof}

La forme géométrique est déterminée par l'information qu'une lecture comprimée semblait avoir déplacée. Le complément vectoriel contient les différences entre positions ; la charge contient la composante commune. Leur réunion reconstruit les douze coordonnées avant leur normalisation en longueurs. Cette affirmation relie des opérations effectives : compression, archivage des compléments et inversion. La conservation d'une norme exprime une propriété distincte et, à elle seule, plus faible que cette récupérabilité.

\subsection{Cône de mémoire compatible et stabilité des mineurs}

Soit \(L\) la pseudo-inverse de \(M\), d’image
\(\mathbf1^\perp\). De manière équivalente,
\[
 L=\left(M^*M\big|_{\mathbf1^\perp}\right)^{-1}M^*,
 \qquad LM=P_{11},\qquad K=\frac q{12}\mathbf1+Lm.
\]
La notation d’inverse restreint s’interprète comme un opérateur à
valeurs dans \(\mathbf1^\perp\). La positivité stricte du registre
équivaut à l’appartenance au cône ouvert
\begin{equation}
 \mathcal M_+=\left\{(q,m)\in\RR\times\operatorname{im}M:
 \frac q{12}\mathbf1+Lm>0\right\}.
 \label{eq:memory-positive-cone}
\end{equation}
Les inégalités sont coordonnées. Le cône a dimension douze ; sa section
à charge positive fixée a dimension onze. En particulier, les vingt-quatre
entrées des deux mémoires satisfont des relations de compatibilité.

\begin{proposition}[Constante optimale de reconstruction en norme euclidienne]
\label{prop:memory-stability}
On a \(\|L\|=9/4\). Pour des perturbations arbitraires des observations,
\begin{equation}
 \|\delta K\|^2\leq\frac{|\delta q|^2}{12}
 +\frac{81}{16}\bigl(\|\delta m_0\|^2+\|\delta m_1\|^2\bigr).
 \label{eq:memory-stability}
\end{equation}
\end{proposition}
\begin{proof}
La base de Fourier diagonalise \(\mathsf S\). Si \(z^{12}=1\), la valeur propre de \(M^*M\) dans le mode correspondant est
\[
 \lambda(z)=\frac8{81}|z^3-1|^2
 +\frac8{81}|z^4-1|^2\left|\frac{8+z^3}{9}\right|^2.
\]
L’évaluation des douze modes donne
\[
 \operatorname{Spec}(M^*M)=
 \left\{0^{[1]},(16/81)^{[2]},(8/27)^{[2]},(32/81)^{[1]},
 (952/2187)^{[4]},(1256/2187)^{[2]}\right\}.
\]
Les exposants entre crochets indiquent la multiplicité. La plus petite valeur propre
positive est \(16/81\), d’où \(\|L\|=9/4\). La composante uniforme
\((\delta q/12)\mathbf1\) est orthogonale à \(L\delta m\) ; l’identité
de Pythagore et la norme de \(L\) produisent \eqref{eq:memory-stability}.
\end{proof}

Pour passer à la géométrie des intervalles, nous reprenons la numérotation de un à douze
et posons \(d_j=K_j/q\), \(t_0=0\),
\(t_j=\sum_{r=1}^j d_r\). Le cône \eqref{eq:memory-positive-cone}
détermine exactement \(0=t_0<t_1<\cdots<t_{12}=1\), et donc
\[
 \Delta_{ij}=t_j-t_i=\sum_{i<r\leq j}d_r>0\qquad(i<j).
\]
Ce sont les mineurs de la matrice de chemins construite dans la section \ref{app:network-cluster}. Leur dépendance à la mémoire est explicite. Les points \((t_j,t_j^2)\) sont également déterminés et produisent la réalisation amplituédrique de rang un développée ensuite.

Si la racine du membre droit de \eqref{eq:memory-stability} est inférieure à \(\min_j K_j\), la reconstruction perturbée demeure dans le cône positif. Pour le registre exprimé, \(\min_jK_j=58\) ; à charge fixe, il suffit que
\[
 \sqrt{\|\delta m_0\|^2+\|\delta m_1\|^2}<\frac{232}{9}.
\]
La borne conserve l'ordre des sommets et le signe de tous les mineurs ordonnés. Elle décrit la stabilité de cette coordinatisation mathématique, avec la normalisation indiquée. Pour des observations hors de \(\operatorname{im}M\), la reconstruction utilise la projection \(MLm\) ; le résidu \((I-ML)m\) est conservé séparément lorsqu'il importe de restituer l'observation intégrale.

\subsection{Récupération nodale et réduction à la frontière}

La partition obtenue a également une réalisation quadratique :
\[
 \mathcal E_d(f)=\sum_{j=1}^{12}\frac{|f_j-f_{j-1}|^2}{d_j}
 =f^*L_df.
\]
Le laplacien nodal récupère chaque longueur par \(d_j=-1/(L_d)_{j-1,j}\). Si seules les extrémités globales sont conservées, la minimisation élimine les nœuds intérieurs et produit, en revanche, la matrice \(\left(\begin{smallmatrix}1&-1\\-1&1\end{smallmatrix}\right)\),
puisque \(\sum_jd_j=1\). Le lecteur nodal distingue la forme des intervalles ;
le lecteur à deux extrémités n’observe que leur longueur totale.

La démonstration de cette élimination est formulée intégralement dans la section
énergétique : le prolongement affine minimisant \(Hf\) et le détail \(z\),
nul aux anciens nœuds, vérifient
\[
 \mathcal E_{\widetilde d}(Hf+z)
 =\mathcal E_d(f)+\mathcal E_{\widetilde d}(z).
\]
Le second terme est non négatif et s'annule exactement pour le détail nul. La lecture effective conserve l'énergie minimale ; la paire \((f,z)\) conserve en outre la configuration raffinée. Ainsi, positivité des mineurs, positivité énergétique et récupérabilité de la mémoire sont reliées au moyen d'applications définies sur la même partition, tout en maintenant leurs espaces d'arrivée respectifs.

\section{Reconstruction intégrale du registre et isométrie d’incidence}
\label{rev4:ki:section}

La reconstruction de la section~\ref{sec:memory} permet de conserver dans
une même composition la géométrie positive, l’incidence exceptionnelle
et la direction de déformation. L’information initiale de cette
composition est le registre produit par l’état enrichi
APP--TRIT--TPK : les deux feuillets arithmétiques conservent leurs opérations ;
l’orientation tritique détermine le régime ; le transport, la
mise à jour et le prolongement nonadique conservent les différences
entre positions et la composante uniforme. Les réalisations linéaires
qui suivent agissent sur ce registre déjà déterminé. Leur fonction est
de transporter l’information et d’expliciter ses invariants.

La distinction fondamentale concerne le domaine conservé. Le
registre \(K\) contient douze coordonnées ordonnées ; la tétrade
\(B_K\) conserve quatre positions sélectionnées par l’incidence ;
une direction géométrique conserve une droite dans un espace de
représentation. Chaque opération a une image et une fibre propres.
La construction suivante conserve d’abord les douze coordonnées et
dérive ensuite les réductions qui interviennent dans la géométrie. Ainsi,
une reconstruction exacte est disponible avant d’effectuer une
extraction de support ou une réduction dimensionnelle.

\subsection{Forme de Gram du plan dodécaphasique}
\label{rev4:ki:gram-subsection}

Nous adoptons la numérotation de un à douze et notons
\(\Omega=\{1,\ldots,12\}\) l’ensemble des positions. Le
codage Paley--Witt de l’état enrichi détermine l’ensemble
\(\mathcal H\) de ses \(132\) hexades. Une hexade est un support
de six positions d’un vecteur de poids six du code ternaire
étendu. La construction explicite du code et la démonstration du fait
que chaque quintuplet appartient à une unique hexade sont développées dans
le théorème~\ref{thm:lf-code} et la
proposition~\ref{prop:lf-witt}. Nous utilisons ici précisément cette
incidence, avec les positions transportées dans le même ordre que
les coordonnées du registre.

Soient \(V=\RR^{12}\) et \(W=\RR^{\mathcal H}\), munis de leurs produits scalaires euclidiens. La matrice d'incidence définit l'application
\begin{equation}
 \mathcal I:V\longrightarrow W,
 \qquad (\mathcal I k)_H=\sum_{p\in H}k_p,
 \qquad \mathcal I_{H,p}=\mathbf1_{\{p\in H\}}.
 \label{rev4:ki:incidence-map}
\end{equation}
Le symbole \(\mathbf1\in V\) désigne le vecteur uniforme à douze
composantes. Nous définissons
\[
 \mu(k)=\frac1{12}\mathbf1^{\mathsf T}k,\qquad
 P_0=I_{12}-\frac1{12}\mathbf1\mathbf1^{\mathsf T},\qquad
 V_0=\mathbf1^\perp.
\]
Le projecteur \(P_0\) coïncide avec \(P_{11}\) de la section~\ref{sec:memory} ; l'indice zéro indique dans cette présentation la condition de moyenne nulle. La décomposition orthogonale est \(k=\mu(k)\mathbf1+P_0k\).

\begin{lemma}[Identité de Gram et normalisation incidentielle]
\label{rev4:ki:gram}
L’incidence satisfait
\begin{equation}
 \mathcal I^{\mathsf T}\mathcal I
       =36I_{12}+30\mathbf1\mathbf1^{\mathsf T}.
 \label{rev4:ki:gram-formula}
\end{equation}
Par conséquent,
\[
 U_W=\frac16\mathcal I\big|_{V_0}:
 V_0\longrightarrow E_{\rm inc}:=\mathcal I(V_0)
\]
est une isométrie sur un sous-espace de dimension onze.
\end{lemma}
\begin{proof}
Un point étant fixé, les quintuplets qui le contiennent sont au nombre de \(\binom{11}{4}\). Chaque hexade contenant ce point en comprend \(\binom54\). L'unicité de l'extension de chaque quintuplet donne \(66\) hexades par point. Une paire de points étant fixée, le même argument donne \(\binom{10}{3}/\binom43=30\) hexades par paire. L'entrée \((p,r)\) de \(\mathcal I^{\mathsf T}\mathcal I\) compte exactement les hexades contenant simultanément \(p\) et \(r\). Ses entrées diagonales sont donc \(66\), et les autres \(30\), ce qui prouve \eqref{rev4:ki:gram-formula}.

Pour \(z\in V_0\), le terme uniforme s’annule et
\[
 \left\|\frac16\mathcal I z\right\|^2
 =\frac1{36}z^{\mathsf T}
       \mathcal I^{\mathsf T}\mathcal I z
 =\|z\|^2.
\]
L’injectivité et la dimension de l’image se déduisent de cette
égalité. Le facteur six est ainsi fixé par le gramien du plan.
\end{proof}

Le dénombrement d’incidence détermine l’échelle de l’isométrie. Les
sommes sur les hexades sont des observables linéaires des positions ; leurs
produits scalaires contiennent la multiplicité avec laquelle chaque
position et chaque couple apparaissent dans le plan. L’élimination du
mode uniforme sépare ces deux niveaux de dénombrement et laisse le facteur
quadratique \(36\). Cette opération fournit un lien explicite
entre la représentation à douze positions et sa représentation
sur les hexades.

\subsection{Conservation de la moyenne et inversion des douze coordonnées}
\label{rev4:ki:full-subsection}

La composante uniforme est conservée au moyen d’une coordonnée scalaire
supplémentaire. Nous munissons \(\RR\oplus E_{\rm inc}\) du produit
scalaire
\[
 \langle(\mu,y),(\nu,z)\rangle_*=
 12\mu\nu+\langle y,z\rangle.
\]
Le poids douze est la norme au carré de \(\mathbf1\), fixée par le
nombre de positions. Nous définissons
\begin{equation}
 F:V\longrightarrow\RR\oplus E_{\rm inc},\qquad
 F(k)=\left(\mu(k),\frac16\mathcal I P_0k\right).
 \label{rev4:ki:F}
\end{equation}

\begin{theorem}[Isomorphisme isométrique du registre complet]
\label{rev4:ki:full-isometry}
L'application \(F\) est un isomorphisme linéaire isométrique. Son inverse est
\begin{equation}
 F^{-1}(\mu,y)=\mu\mathbf1+
                  \frac16P_0\mathcal I^{\mathsf T}y,
 \qquad y\in E_{\rm inc}.
 \label{rev4:ki:F-inverse}
\end{equation}
En particulier,
\begin{equation}
 12|\mu(k)|^2+
 \left\|\frac16\mathcal I P_0k\right\|^2=\|k\|^2.
 \label{rev4:ki:F-norm}
\end{equation}
\end{theorem}
\begin{proof}
La décomposition uniforme--centrée est orthogonale et le
lemme~\ref{rev4:ki:gram} conserve la norme de son second terme.
Cela démontre \eqref{rev4:ki:F-norm} et l’injectivité.
De plus,
\[
 \frac1{36}P_0\mathcal I^{\mathsf T}\mathcal I P_0=P_0.
\]
En substituant \(F(k)\) dans \eqref{rev4:ki:F-inverse}, on obtient \(\mu(k)\mathbf1+P_0k=k\). Réciproquement, si \(y\in E_{\rm inc}\), il existe un unique \(z\in V_0\) avec \(y=\mathcal I z/6\). La formule proposée renvoie \(\mu\mathbf1+z\), dont l'image est \((\mu,y)\). La surjectivité sur le codomaine déclaré et les deux identités d'inversion sont ainsi prouvées.
\end{proof}

La restriction \(y\in E_{\rm inc}\) exprime une compatibilité
linéaire effective. L’opérateur
\begin{equation}
 \Pi_{\rm inc}=\frac1{36}\mathcal I P_0
                         \mathcal I^{\mathsf T}:W\longrightarrow W
 \label{rev4:ki:inc-projector}
\end{equation}
est le projecteur orthogonal sur \(E_{\rm inc}\). En effet, il est
symétrique et son carré se réduit à lui-même au moyen de
\eqref{rev4:ki:gram-formula} ; il agit comme l’identité dans l’image
de \(\mathcal I P_0\). Par conséquent, la compatibilité se vérifie
au moyen de \(\Pi_{\rm inc}y=y\). Pour une donnée arbitraire de
\(W\), la reconstruction suivie de l’incidence restitue
\(\Pi_{\rm inc}y\) ; la composante orthogonale
\((I_W-\Pi_{\rm inc})y\) quantifie l’information externe à
cette représentation du registre.

\begin{proposition}[Covariance de la reconstruction]
\label{rev4:ki:equivariance}
Soit \(g\) un automorphisme du plan et soient \(\rho(g)\) et
\(\sigma(g)\) ses actions par permutation sur les positions et les
hexades. Alors
\[
 F\rho(g)=(1\oplus\sigma(g))F.
\]
La moyenne est conservée et la reconstruction commute avec le transport simultané des positions et de leurs hexades.
\end{proposition}
\begin{proof}
L’équivalence \(p\in H\Longleftrightarrow g(p)\in g(H)\)
implique \(\mathcal I\rho(g)=\sigma(g)\mathcal I\).
Les permutations fixent \(\mathbf1\), conservent la moyenne et
commutent avec \(P_0\). Substituer ces identités dans
\eqref{rev4:ki:F} démontre l’affirmation. La covariance de
l’inverse s’obtient en composant avec \(F^{-1}\).
\end{proof}

La normalisation isométrique et l'inversion arithmétique remplissent des fonctions distinctes. Dans la mise en coordonnées de Hadamard du registre, \(U=\mathcal H_{12}K\) et \(K=\mathcal H_{12}U/4\), avec \(\mathcal H_{12}^{\mathsf T}=\mathcal H_{12}\) et \(\mathcal H_{12}^2=4I_{12}\). L'opérateur \(J_H=\mathcal H_{12}/2\) est orthogonal. Par conséquent, \(FJ_H\) conserve la norme du registre signé normalisé \(U/2\) et le représente dans l'espace d'incidence élargi. L'inversion entière au moyen de \(\mathcal H_{12}/4\) conserve en outre le sous-réseau de congruences correspondant. Dans les deux cas, le registre engendré est transporté, avec l'ordre de phase fixé avant la transformation \cite{hmtX}.

\subsection{Composition exacte à partir des deux mémoires nonadiques}
\label{rev4:ki:memory-subsection}

Nous reprenons \(M\), \(q\), \(m_0\) et \(m_1\)
de \eqref{eq:mem-data}, avec l’opérateur de reconstruction
\(L=(M^*M|_{V_0})^{-1}M^*\) de la section~\ref{sec:memory}.
Le domaine compatible est
\(\mathcal D=\RR\times\operatorname{im}M\), avec
\(m=(m_0,m_1)\). Le théorème~\ref{thm:memory-inversion}
établit la reconstruction
\[
 \mathcal R(q,m)=\frac q{12}\mathbf1+Lm,
 \qquad LM=P_0.
\]
Comme \(L\) prend ses valeurs dans \(V_0\), la composition avec
l’incidence possède une expression directe :
\begin{equation}
 \mathcal C=F\mathcal R,\qquad
 \mathcal C(q,m)=\left(\frac q{12},\frac16\mathcal I Lm\right).
 \label{rev4:ki:memory-incidence}
\end{equation}

\begin{corollary}[Équivalence entre mémoire compatible et incidence intégrale]
\label{rev4:ki:composition}
L’application \(\mathcal C\) est un isomorphisme entre
\(\mathcal D\) et \(\RR\oplus E_{\rm inc}\), d’inverse
\begin{equation}
 \mathcal C^{-1}(\mu,y)=
 \left(12\mu,\frac16MP_0\mathcal I^{\mathsf T}y\right).
 \label{rev4:ki:C-inverse}
\end{equation}
Sa forme quadratique conservée est
\[
 \|\mathcal C(q,m)\|_*^2=
 \frac{|q|^2}{12}+\|Lm\|^2=\|\mathcal R(q,m)\|^2.
\]
\end{corollary}
\begin{proof}
La formule \eqref{rev4:ki:memory-incidence} résulte de
\(P_0\mathbf1=0\) et de \(P_0L=L\). L’inverse s’obtient
en appliquant \(k\mapsto(\mathbf1^{\mathsf T}k,Mk)\) à
\eqref{rev4:ki:F-inverse} ; on utilise \(M\mathbf1=0\).
Les identités de composition proviennent de \(LM=P_0\) et de
\(MLm=m\) pour \(m\in\operatorname{im}M\). L’égalité des
normes est \eqref{rev4:ki:F-norm} appliquée à \(\mathcal R(q,m)\).
\end{proof}

\begin{figure}[htbp]
\centering
\begin{tikzpicture}[>=Latex, every node/.style={font=\small},
  box/.style={draw,rounded corners=2pt,align=center,
              inner sep=6pt,minimum height=12mm}]
 \node[box] (memory) at (0,0) {Charge et mémoires\\$(q,m_0,m_1)$};
 \node[box] (register) at (4.1,0) {Registre intégral\\$K$};
 \node[box] (incidence) at (8.2,0) {Moyenne et incidence\\$(\mu,y)$};
 \node[box] (direction) at (4.1,-1.75)
    {Direction géométrique\\$\ell_K=\RR P_3K$};
 \draw[<->] (memory) -- node[above] {$\mathcal R$} (register);
 \draw[<->] (register) -- node[above] {$F$} (incidence);
 \draw[->] (register) -- node[right] {projection et passage à la droite}
    (direction);
\end{tikzpicture}
\caption{La charge et les deux mémoires compatibles reconstruisent les douze
coordonnées. La moyenne et la composante incidencielle fournissent une autre
représentation réversible du même registre. La direction géométrique
est une réduction ultérieure à fibres de dimension positive.}
\label{rev4:ki:figure}
\end{figure}

Cette composition détermine également la partition positive de \ref{sec:memory}. Si \(q>0\), ses longueurs se récupèrent au moyen de
\begin{equation}
 d_j=\frac1{12}+\frac{1}{6q}
             \bigl(P_0\mathcal I^{\mathsf T}y\bigr)_j,
 \qquad y=\frac16\mathcal I P_0K.
 \label{rev4:ki:positive-lengths}
\end{equation}
Les douze conditions \(d_j>0\) sont les équations de la coordinatisation positive exprimées en coordonnées d'incidence. Avec \(t_j=\sum_{r\leq j}d_r\), les mineurs ordonnés \(t_j-t_i=\sum_{i<r\leq j}d_r\) sont reconstruits à partir de \((\mu,y)\) ou de \((q,m_0,m_1)\). La positivité des mineurs et la conservation isométrique de l'incidence sont ainsi reliées par des applications explicites. La première est une inégalité dans la coordinatisation ordonnée ; la seconde est une identité de produits scalaires.

\subsection{Récupération de la polarisation et portée des réductions}
\label{rev4:ki:direction-subsection}

La réalisation des douze positions sur les classes de
\(A_5/C_5\) détermine le projecteur tridimensionnel
\[
 P_3=\frac3{60}\sum_{a\in A_5}
               \chi_3(a^{-1})\varrho(a).
\]
Ici, \(\varrho\) est la représentation par permutations et \(\chi_3\) le caractère tridimensionnel choisi dans la coordinatisation marquée. Ses représentants, valeurs de caractère et matrice effective sont fixés dans \eqref{eq:leech-projectors} ; le lemme~\ref{lem:leech-u} démontre \(P_3=P_3^{\mathsf T}=P_3^2\), \(\operatorname{rank}P_3=3\) et \(P_3\mathbf1=0\). Dans cette coordinatisation, la direction du registre s'obtient par
\begin{equation}
 u_K=P_3P_0K=P_3Lm
     =\frac16P_3P_0\mathcal I^{\mathsf T}y.
 \label{rev4:ki:direction-reconstruction}
\end{equation}
Les égalités découlent des deux inverses déjà démontrés. Ainsi,
la même direction se récupère à partir des mémoires ou de la
composante incidentielle ; la moyenne est conservée pour reconstruire
\(K\), même si le projecteur tridimensionnel l’annule.

Pour le registre de \eqref{eq:leech-K}, l’évaluation exacte
\eqref{eq:leech-u-norm} donne
\[
 \|u_K\|^2=
       \frac{6638585+2275584\sqrt5}{20}>0.
\]
La droite \(\ell_K=\RR u_K\) et le complément transversal
\(V_{10}(K)=V_0\cap u_K^\perp\) sont ainsi déterminés.
Son projecteur est
\begin{equation}
 P_{10}(K)=P_0-
       \frac{u_Ku_K^{\mathsf T}}{\|u_K\|^2}.
 \label{rev4:ki:rank-ten}
\end{equation}
En effet, le second terme est un projecteur orthogonal de rang un contenu dans \(V_0\) ; son produit avec \(P_0\), dans l'un ou l'autre ordre, coïncide avec lui-même. En élevant \eqref{rev4:ki:rank-ten} au carré, on récupère le même opérateur, et sa trace vaut \(11-1=10\). On obtient \(V_0=\ell_K\oplus^\perp V_{10}(K)\). Cette réduction est celle qui intervient ensuite dans la polarisation des douze plans réticulaires.

L’incidence transporte également ces opérateurs. Sur
\(E_{\rm inc}\), l’opérateur
\(\widehat P_3=U_WP_3U_W^*\) vérifie
\(\widehat P_3U_W=U_WP_3\) ; l’identité découle de
\(U_W^*U_W=I_{V_0}\). Transporter de la même manière
\(P_{10}(K)\) conserve son rang et son produit scalaire.
La correspondance entre les deux représentations s’établit
donc sur les vecteurs et les opérateurs, en plus de leurs dimensions.

La nécessité de conserver les différentes composantes admet
des vérifications précises. L’observation centrée de \(k\) coïncide
avec celle de \(k+c\mathbf1\) pour tout \(c\), tandis que leurs
charges diffèrent de \(12c\). En particulier, \(k=0\) et
\(k=\mathbf1\) ont la même composante incidentielle centrée
et des moyennes distinctes. La coordonnée \(\mu\) est exactement
l’information qui sépare ces deux registres. De même,
\(u_{ak+b\mathbf1}=a u_k\) pour \(a\ne0\) ; la droite
résultante coïncide avec \(\ell_k\), même si l’amplitude
et la moyenne changent. Ces identités décrivent les fibres des
applications sur l’espace ambiant déclaré.

L’extraction \(k\mapsto\{j:k_j\geq729\}\) possède d’autres
fibres. Pour le registre concret \(K\), aussi bien \(K\) que
\(K+e_1\) ont pour support \(\{4,6,9,10\}\), où
\(e_1\) est le premier vecteur de la base canonique. Ce sont deux vecteurs
ambiants différents ayant la même tétrade. Cette vérification
caractérise la perte d’information de l’application de support ; la
sélection terminale du registre reste fixée par sa
construction préalable. Le drapeau
\(B_K\subset H^-_\alpha\) conserve en outre l’hexade déterminée
par les signes antérieurs à la retenue, dont la coïncidence avec la
sélection incidencielle est démontrée dans la
proposition~\ref{prop:lf-flag}.

Le résultat commun de ces compositions est une continuité mathématique explicite. La mémoire compatible reconstruit le registre ; sa décomposition uniforme--centrée détermine une représentation incidentielle isométrique ; la reconstruction de ses longueurs détermine la coordinatisation positive ; et la projection tridimensionnelle détermine la direction de déformation. La tétrade et le drapeau conservent les incidences utilisées par les constructions exceptionnelles, tandis que \(F\) conserve la totalité du registre dodécaphasique. L'historique enrichi maintient à son tour l'orientation, le parcours et le prolongement qui situent ce registre dans la structure discrète conjointe du continu.

% Procedencia HMT: output/EXTREMA_MEDIA_RAZON_NARRACION_Y_REVISION_20260916/
% ARTICULO/ES/source/libro/01_acoplamiento.tex, lib01:cadena-registro y lib01:K.
% Dirección: ARTICULO/ES/source/sections/k_direccion_dimensional.tex.
% Estatuto de las redes, cartas y flujos de este apéndice: FORMALIZACION_NUEVA.
% La extracción de K y el proyector P_3 son RESULTADO_RECUPERADO.
\section{Mesure planaire de frontière, coordonnées de Plücker et mutations positives}
\label{app:network-cluster}

La représentation dodécaphasique de l'état enrichi admet une réalisation positive dont le réseau, les coordonnées et les changements de coordinatisation peuvent être écrits intégralement. La composition de départ est
\[
 \mathrm{APP}\longrightarrow\mathrm{TRIT}\longrightarrow\mathrm{TPK}
 \longrightarrow x_{\rm term}\longrightarrow\mathcal R_{12}
 \longrightarrow(b^{90},b^{120},Q_{\rm TPK})
 \longrightarrow U_{\rm sgn}\longrightarrow K.
\]
Cette lecture agit sur l’état enrichi qui conserve les deux feuillets,
l’orientation, la retenue, la mémoire et la frontière. Le registre déjà engendré
est
\[
 K=(234,543,140,729,659,824,621,58,914,794,146,601),
 \qquad Q=\sum_{r=1}^{12}K_r=6263.
\]
Son extraction se récupère exactement à partir des deux canaux et de la charge totale :
\[
 b^{90}=(-495,-116,-684,108,601,-90,-173,-88,313,560,-397,461),
\]
\[
 b^{120}=(-425,-281,-481,671,-255,30,475,-543,680,251,6,-128).
\]
Sur les orbites de positions \((r,r+3,r+6,r+9)\), pour \(r=1,2,3\),
soient \(B_r=\sum_{j=0}^3b^{120}_{r+3j}\) et
\(d_{r,j}=b^{90}_{r+3j}\). On obtient
\[
 s_1=(Q+2B_1+B_2)/3,\quad s_2=s_1-B_1,\quad s_3=s_2-B_2,
\]
\[
 U^{(r)}=(s_r,d_{r,1}-d_{r,3},d_{r,0}+d_{r,2},d_{r,0}-d_{r,2}),
 \qquad K|_r=\tfrac14H_4U^{(r)},
\]
\[
 H_4=\begin{pmatrix}1&1&1&1\\1&1&-1&-1\\1&-1&1&-1\\1&-1&-1&1\end{pmatrix},
 \qquad H_4^2=4I_4.
\]
Le réseau construit ci-dessous réalise les valeurs positives produites par cette extraction. Sa forme combinatoire est une formalisation ultérieure du registre, avec une origine et une orientation déclarées.

\subsection{Réseau orienté acyclique et déterminants de chemins sans intersection}

Pour un registre positif de longueur \(L\), nous posons
\[
 d_r=K_r/Q>0,\qquad t_0=0,\qquad
 t_j=\sum_{r=1}^jd_r\quad(1\le j\le L).
\]
Ainsi \(t_L=1\). Les \(L\) séparations produisent \(L+1\) points marqués ;
dans le cas dodécaphasique, il y a treize points. La matrice qui les réalise est
\begin{equation}
 C(d)=\begin{pmatrix}1&1&\cdots&1\\t_0&t_1&\cdots&t_L\end{pmatrix}.
 \label{eq:nc-path-matrix}
\end{equation}
La grassmannienne \(\operatorname{Gr}(2,L+1)\) identifie les matrices de rang
deux par changements inversibles des deux lignes. Sa partie positive est formée
des classes dont tous les mineurs ordonnés d’ordre deux
peuvent être strictement positifs.

Nous construisons un réseau orienté de sommets \(a_j,b_j\), pour \(0\le j\le L\), et de puits \(q_j\). Il comporte des arêtes horizontales \(a_{j-1}\to a_j\) et \(b_{j-1}\to b_j\), de poids un ; des arêtes verticales \(a_j\to b_j\), de poids \(d_j\), pour \(j\ge1\) ; et des arêtes \(b_j\to q_j\), de poids un. Les sources ordonnées sont \(b_0,a_0\). Les deux lignes sont dessinées horizontalement, les arêtes verticales descendent et les puits sortent en dessous de la ligne inférieure. Ce dessin est planaire et l'orientation est acyclique. La mesure de frontière \(M_{ij}\) est la somme des produits de poids de tous les chemins allant de la source \(i\) au puits \(q_j\). Le réseau de chemins a deux sources et \(L+1\) puits, soit \(L+3\) connexions extérieures. La matrice qu'il réalise a \(L+1\) colonnes. Le germe plabique introduit ensuite possède précisément \(L+1\) sommets de frontière et constitue une autre présentation de ses coordonnées ; ces deux décomptes de frontière restent explicites.

\begin{theorem}[Mesure de frontière et mineurs du registre]
\label{thm:nc-network}
La mesure de frontière de ce réseau est \(M=C(d)\). Pour
\(0\le i<j\le L\),
\begin{equation}
 \Delta_{ij}(C(d))=t_j-t_i=\sum_{r=i+1}^{j}d_r>0.
 \label{eq:nc-minor-gap}
\end{equation}
Le même nombre est la somme des poids des paires de chemins
à sommets disjoints \(b_0\to q_i\), \(a_0\to q_j\).
\end{theorem}
\begin{proof}
La source inférieure a un chemin unique vers chaque puits, de poids un.
Un chemin de \(a_0\) à \(q_j\) descend exactement une fois, par une
arête d’indice \(1\le r\le j\) ; son poids est \(d_r\). Cela prouve la
formule de la matrice. Le chemin inférieur jusqu’à \(q_i\) occupe tous les
sommets \(b_0,\ldots,b_i\). Par conséquent, un chemin de l’autre source vers
\(q_j\) en est disjoint précisément lorsqu’il descend en
\(i<r\le j\). Sa somme de poids est le membre droit de
\eqref{eq:nc-minor-gap}. Le couple dont les destinations sont échangées
partage toujours un sommet. Le développement du déterminant annule tous les
couples qui se croisent et conserve exactement les couples disjoints. On a ici
vérifié directement, pour le réseau concret, le mécanisme déterminantal
des chemins sans intersection.
\end{proof}

L’identification des mesures de réseaux plans aux cellules de la
grassmannienne non négative fournit leur reconnaissance géométrique ultérieure ;
voir A. Postnikov,
\href{https://arxiv.org/abs/math/0609764}{\emph{Total positivity,
Grassmannians, and networks}}. La preuve précédente donne, en outre, une formule
finie pour chaque mineur sans utiliser de valeur métrologique.

La collection \(d_r>0\), \(\sum_rd_r=1\), a pour dimension \(L-1\). Son application dans la grassmannienne est injective : si les lignes de \(C(d)\) et \(C(d')\) engendrent le même plan, les secondes lignes diffèrent par une transformation affine ; les conditions \(t_0=t'_0=0\) et \(t_L=t'_L=1\) fixent cette transformation. En particulier, pour \(L=12\), on obtient une collection de dimension onze à l'intérieur de la cellule positive de dimension vingt-deux. Le point appartient à la cellule supérieure ; la collection conserve sa dimension propre. Les treize points et les douze séparations ont ainsi des rôles mathématiques distincts.

Autoriser \(d_r\ge0\) conserve la non-négativité. Lorsque \(d_r=0\), les colonnes consécutives \(r-1\) et \(r\) coïncident et forment des éléments parallèles du matroïde. Le rang reste égal à deux car \(t_0=0\) et \(t_L=1\). Les bases sont exactement les paires \(i<j\) qui traversent une séparation positive. Cette dégénérescence produit des strates d'incidence consécutive. Une colonne nulle, qui donnerait une boucle, correspond à une opération distincte de la coïncidence de deux colonnes non nulles.

\subsection{Germe de triangulation et relation d’échange}

Soit \(N=L+1\) et considérons les sommets \(0,\ldots,L\) d’un polygone
convexe dans leur ordre cyclique. À chaque côté ou diagonale \((i,j)\), on assigne
la coordonnée \(A_{ij}=\Delta_{ij}(C(d))\), avec des indices ordonnés.
La triangulation en éventail depuis le sommet zéro a pour diagonales
\((0,j)\), \(2\le j\le L-1\). Ses côtés sont des variables gelées et
ses diagonales des variables mutables. Le carquois s’obtient en orientant
cycliquement les trois côtés de chaque triangle et en annulant les couples de
flèches opposées. C’est un germe de type \(A_{N-3}\) ; pour le registre
dodécaphasique, son type est \(A_{10}\).

Le germe possède une réalisation plabique explicite. On place un sommet
noir à l’intérieur de chaque triangle ; aux sommets du polygone, on place un
sommet blanc s’il est incident à une diagonale, et noir dans le cas contraire.
On relie les centres aux trois sommets de leurs triangles, on retire les
côtés originaux et on contracte les arêtes entre sommets de même couleur.
On ajoute un feuillet de frontière à chaque sommet marqué. Cette construction
reprend l’algorithme de triangulation de Kodama et Williams,
\href{https://people.math.harvard.edu/~williams/papers/KW-ArxivVersion.pdf}
{\emph{KP solitons and total positivity for the Grassmannian}, algorithme 12.1}.
Les étiquettes de coordonnées de ce germe sont évaluées dans les mineurs déjà produits par le réseau à deux lignes. Pour une face mutable, la coordonnée de type \(\mathcal X\) est le quotient des produits des coordonnées \(\mathcal A\) incidentes prescrits par sa colonne du carquois ; dans un quadrilatère, elle sera écrite explicitement comme un birapport. Changer le graphe modifie la coordinatisation de ces mineurs, tout en maintenant leur provenance.

\begin{theorem}[Échange positif de la réalisation du registre]
\label{thm:nc-ptolemy}
Pour quatre sommets \(a<b<c<d\),
\begin{equation}
 A_{ac}A_{bd}=A_{ab}A_{cd}+A_{ad}A_{bc}.
 \label{eq:nc-ptolemy}
\end{equation}
Par conséquent, le changement de diagonale \((a,c)\leftrightarrow(b,d)\)
s’évalue au moyen de
\[
 A_{bd}=\frac{A_{ab}A_{cd}+A_{ad}A_{bc}}{A_{ac}}>0.
\]
Toutes les triangulations fournissent des coordinatisations positives compatibles sur la même réalisation \(C(d)\).
\end{theorem}
\begin{proof}
En substituant \(A_{ij}=t_j-t_i\), l’égalité se réduit à
\[
 (t_c-t_a)(t_d-t_b)
 =(t_b-t_a)(t_d-t_c)+(t_d-t_a)(t_c-t_b).
\]
Le développement des deux membres prouve l’identité. Les différences sont
positives ; chaque flip conserve donc une évaluation positive et le
quotient est défini. La connexion des triangulations par flips
permet de passer de l’une à l’autre au moyen de ces échanges. L’interprétation
en amas de la relation de Plücker est la reconnaissance ultérieure développée
par J. Scott dans
\href{https://arxiv.org/abs/math/0311148}{\emph{Grassmannians and Cluster
Algebras}}. La preuve algébrique de l’évaluation concrète est déjà contenue
dans l’identité montrée.
\end{proof}

\subsection{Raffinement des séparations et naturalité des échanges}

Soit \(B\ge2\). À la profondeur \(n\), chaque séparation \(d_r\) est divisée
en \(B^n\) séparations consécutives égales :
\[
 d^{(n)}_{(r-1)B^n+c}=d_r/B^n,
 \qquad 1\le c\le B^n.
\]
Cela définit \(LB^n+1\) points et une matrice \(C_n\). Si \(\iota_n(j)=Bj\) identifie un point à son descendant de même position, alors
\begin{equation}
 C_{n+1}|_{\iota_n(\{0,\ldots,LB^n\})}=C_n,
 \qquad
 \Delta_{\iota_n(i),\iota_n(j)}(C_{n+1})=\Delta_{ij}(C_n).
 \label{eq:nc-refinement}
\end{equation}
Le réseau raffiné remplace chaque arête verticale de poids \(d\) par une
séquence de \(B\) positions de descente, de poids \(d/B\). La somme
des chemins entre extrémités héritées reste \(d\). La composition des
raffinements conserve cette égalité car les sommes finies se regroupent
de manière associative.

\begin{proposition}[Naturalité sur les extrémités héritées]
\label{prop:nc-inherited}
Tout échange \eqref{eq:nc-ptolemy} entre quatre extrémités héritées commute exactement avec le raffinement. Toute triangulation ancienne s'étend à une triangulation du polygone raffiné en préservant ses triangles intérieurs : les polygones ajoutés entre deux extrémités consécutives sont triangulés à l'extérieur de ces triangles. Les flips entre anciennes diagonales conservent alors leurs relations d'échange.
\end{proposition}
\begin{proof}
La première conclusion résulte du remplacement de chaque mineur par son égalité dans
\eqref{eq:nc-refinement}. Pour la seconde, on conserve les cordes qui
relient les anciennes extrémités consécutives et on triangule chaque poche extérieure
formée par les nouvelles marques. Ces régions ont des intérieurs disjoints.
Les deux triangles incidents à une ancienne diagonale restent identiques,
et donc leur quadrilatère et leur identité de Ptolémée également.
\end{proof}

Cette naturalité spécifie les extrémités sur lesquelles elle agit. Un côté
ancien, lorsqu’il reçoit de nouveaux sommets sur son arc de frontière, devient une
diagonale du polygone plus grand. Une variable anciennement gelée
peut donc être mutable dans le problème étendu. La compatibilité précédente est
une compatibilité exacte du système d’évaluations et des flips
hérités. Une inclusion de motifs de germes ayant le même régime de
coefficients s’obtient en maintenant gelées ces cordes et les nouvelles
diagonales auxiliaires. Cela déclare l’opération qui conserve le germe ;
l’accroissement de profondeur et une mutation demeurent des opérations
distinctes.

\subsection{Action diagonale projective et invariance des birapports}

Si \(\lambda_j>0\), l'action par colonnes \(C\mapsto C\operatorname{diag}(\lambda_0,\ldots,\lambda_L)\) donne
\[
 A_{ij}\longmapsto\lambda_i\lambda_j A_{ij}.
\]
Elle préserve la positivité et toutes les strates définies par l’annulation de
mineurs. Le birapport associé à un quadrilatère est
\begin{equation}
 X_{abcd}=\frac{A_{ab}A_{cd}}{A_{ad}A_{bc}}>0.
 \label{eq:nc-crossratio}
\end{equation}
Les quatre facteurs \(\lambda_a\lambda_b\lambda_c\lambda_d\) se simplifient exactement. Par conséquent, toute action diagonale positive par colonnes laisse \(X_{abcd}\) fixe, même si elle dépend d'un paramètre et a un déterminant égal à un. La modification du poids des colonnes conserve le point projectif de chaque colonne ; leurs birapports expriment cette conservation. Pour que le registre change la forme projective, il doit agir sur les séparations ordonnées.

\subsection{Déformation positive des séparations dirigée par le registre}

Le corpus fixe la représentation des douze positions et son projecteur
orthogonal \(P_3\) sur le secteur tridimensionnel. Avec
\(P_{11}=I_{12}-J_{12}/12\), la direction retrouvée est
\[
 u=P_3P_{11}K=P_3K\in\mathbf1^\perp.
\]
Dans la coordinatisation déclarée, en posant \(r=\sqrt5\), son évaluation exacte est
\begin{align*}
u=(&-495/4-233r/10,\ 403/4+169r/10,\ -403/4-169r/10,\\
   &495/4+233r/10,\ 601/4+1031r/10,\ 203/4+211r/5,\\
   &-203/4-211r/5,\ -601/4-1031r/10,\ 313/4+569r/10,\\
   &162+219r/20,\ -162-219r/20,\ -313/4-569r/10).
\end{align*}
On conserve \(\sum_ru_r=0\) et
\(\|u\|^2=(6638585+2275584\sqrt5)/20>0\).

Pour \(s\in\mathbb R\), nous définissons la déformation des poids par
\begin{equation}
 d_r(s)=\frac{K_re^{su_r}}{\sum_{\ell=1}^{12}K_\ell e^{su_\ell}},
 \qquad t_j(s)=\sum_{r=1}^jd_r(s),\qquad C(s)=C(d(s)).
 \label{eq:nc-gap-flow}
\end{equation}
La formule utilise le registre et sa direction déjà engendrés. C'est une nouvelle réalisation analytique du lecteur ; le paramètre \(s\) mesure sa déformation. Lorsque la sortie interne \(\rho_A=\pi_{\rm HMT}/\varphi_{\rm HMT}^2\) est utilisée, l'écriture \(s=\tau\log\rho_A\) exprime exactement la même collection. Cette reparamétrisation agit après la génération des constantes. L'identification de \(s\) ou \(\tau\) à un temps physique requiert la loi temporelle du lecteur correspondant.

\begin{theorem}[Changement de forme, positivité et compatibilité avec la profondeur]
\label{thm:nc-shape-flow}
La collection \eqref{eq:nc-gap-flow} appartient à \(\operatorname{Gr}_{>0}(2,13)\) pour tout \(s\) réel et conserve \(t_0=0\), \(t_{12}=1\). Pour \(i<j\), soient
\[
 \mu_{ij}(s)=
 \frac{\sum_{r=i+1}^jK_ru_re^{su_r}}{\sum_{r=i+1}^jK_re^{su_r}},
 \qquad
 \mu(s)=\frac{\sum_rK_ru_re^{su_r}}{\sum_rK_re^{su_r}}.
\]
Alors
\begin{equation}
 \frac{d}{ds}\log A_{ij}(s)=\mu_{ij}(s)-\mu(s),
 \quad
 \frac{d}{ds}\log X_{abcd}(s)
 =\mu_{ab}(s)+\mu_{cd}(s)-\mu_{ad}(s)-\mu_{bc}(s).
 \label{eq:nc-crossratio-derivative}
\end{equation}
La déformation commute avec chaque raffinement qui attribue aux successeurs de la séparation \(r\) le poids \(d_r/B\) et la même vitesse \(u_r\).
\end{theorem}
\begin{proof}
Toutes les exponentielles et tous les \(K_r\) sont positifs. La normalisation
fait que la somme des séparations est un ; les mineurs sont des sommes de
séparations positives par \eqref{eq:nc-minor-gap}. Dériver le logarithme de
\[
 A_{ij}(s)=\frac{\sum_{r=i+1}^jK_re^{su_r}}
                      {\sum_rK_re^{su_r}}
\]
donne la première identité. La combinaison de quatre logarithmes dans \eqref{eq:nc-crossratio} annule les quatre moyennes globales avec leurs signes et donne la seconde. Enfin, remplacer \(K_r\) par \(B\) copies de \(K_r/B\) de même \(u_r\) conserve le numérateur et le dénominateur lorsque les successeurs sont regroupés. L'égalité vaut pour tous les paramètres, à toute profondeur et pour toute composition de raffinements.
\end{proof}

Il existe un témoin exact du changement de forme de cette collection. Pour les quatre premières extrémités,
\begin{equation}
 X_{0123}(0)=\frac{1560}{23711},\qquad
 \left.\frac{d}{ds}\log X_{0123}(s)\right|_{s=0}
 =-\frac{309899}{917}-\frac{268596}{4585}\sqrt5<0.
 \label{eq:nc-nontrivial-witness}
\end{equation}
Le calcul utilise uniquement les trois premières séparations \((234,543,140)\) et leurs vitesses déjà récupérées. L'inégalité stricte prouve un mouvement projectif effectif. Les identités de Ptolémée restent exactes au cours de ce mouvement, puisqu'elles s'appliquent aux colonnes \((1,t_j(s))\) pour chaque paramètre. Le registre contrôle ainsi une déformation positive de la forme et, simultanément, un système de coordinatisations demeurant compatible avec elle.

La portée démontrée réunit un réseau planaire de rang deux, sa mesure
explicite, un germe plabique de triangulation, toutes ses mutations par
flips, la naturalité des extrémités héritées et une déformation des
birapports dirigée par \(K\). Chaque opération a un domaine et un
résultat vérifiables. La comparaison de cette réalisation avec la
dynamique temporelle enrichie et avec la branche matérielle utilise leurs
applications de lecture respectives ; aucune identité de coordonnées ne remplace ces applications.

\section{Géométries amplituédriques de rang un et compatibilité des résidus}
\label{sec:k-poligono-amplituhedral}

Le registre dodécaphasique permet de réunir explicitement des données positives,
un domaine grassmannien, une image cinématique, un recouvrement, un degré sur les cellules
et une forme canonique. Le premier résultat concerne l’amplituèdre
bidimensionnel de paramètres \((n,k,m)=(13,1,2)\) ; la construction est ensuite
prolongée aux polytopes cycliques de rang un et de dimension supérieure.
Il s’agit d’une réalisation géométrique ultérieure du registre produit
par APP--TRIT--TPK. L’opération supplémentaire déclarée est l’élévation aux
puissances d’une partition rationnelle ordonnée. Ce choix fixe une
réalisation particulière ; les correspondances physiques ultérieures conservent
leurs lecteurs et leurs conditions propres.

La provenance du registre conserve la composition
\[
 x_{\mathrm{term}}\longmapsto\mathcal R_{12}(x_{\mathrm{term}})
 \longmapsto(b^{90},b^{120},Q_{\mathrm{TPK}})
 \longmapsto U_{\mathrm{sgn}}\longmapsto K.
\]
APP conserve les deux feuillets et leurs résidus et quotients ; TRIT oriente le régime de chaque transition ; TPK transporte et actualise l'état avec retenue, frontière, route et mémoire. La représentation \(K\) se lit dans le même système d'états enrichis qui engendre conjointement le continu. Son extraction signée et son inversion de Hadamard produisent
\begin{equation}
 K=(234,543,140,729,659,824,621,58,914,794,146,601),
 \qquad S_K=\sum_{j=1}^{12}K_j=6263.
 \label{eq:polygon-K}
\end{equation}
La preuve géométrique utilise la positivité de ces douze coordonnées, leur
ordre et leur total. Ses coefficients proviennent de cette extraction et non de
la sélection rétrospective d’une valeur physique.

\subsection{Partition rationnelle et relèvement positif}

Nous définissons treize points rationnels de partition :
\begin{equation}
 t_0=0,\qquad t_i=\frac{K_1+\cdots+K_i}{S_K}\quad(1\leq i\leq12),
 \qquad Z_i=(1,t_i,t_i^2)^\mathsf T,\qquad Z=(Z_0\ \cdots\ Z_{12}).
 \label{eq:polygon-Z}
\end{equation}
Leurs numérateurs, de dénominateur commun \(6263\), sont
\[
 (0,234,777,917,1646,2305,3129,3750,3808,4722,5516,5662,6263).
\]
Ainsi, \(0=t_0<t_1<\cdots<t_{12}=1\). La partition étiquetée et le total conservent le registre au moyen de l'inverse exact
\begin{equation}
 K_i=S_K(t_i-t_{i-1}).
 \label{eq:polygon-recover-K}
\end{equation}
La partition conserve les proportions ; retenir \(S_K\) conserve également
l’échelle entière.

\begin{proposition}[Données cinématiques strictement positives]
\label{prop:polygon-vandermonde}
La matrice \(Z\) a rang trois et tous ses mineurs ordonnés de taille
trois sont positifs. Pour \(0\leq a<b<c\leq12\),
\begin{equation}
 \langle abc\rangle:=\det(Z_a,Z_b,Z_c)
 =(t_b-t_a)(t_c-t_a)(t_c-t_b)>0.
 \label{eq:polygon-minors}
\end{equation}
\end{proposition}
\begin{proof}
On soustrait la première colonne des deux dernières et on développe selon la
première ligne. Chaque différence \(t_j^2-t_a^2\) se factorise en
\((t_j-t_a)(t_j+t_a)\), ce qui donne le produit indiqué. Les sommes partielles
strictement croissantes prouvent sa positivité. Chacun de ces
mineurs atteste le rang trois.
\end{proof}

\subsection{Image, fibres et cellules de degré un}

La grassmannienne \(\operatorname{Gr}_{\geq0}(1,13)\) est l'espace projectif des vecteurs \(c=(c_0,\ldots,c_{12})\) non négatifs et non nuls, modulo une échelle positive. La coordinatisation \(\sum c_j=1\) l'identifie au simplexe fermé de dimension douze. On définit
\begin{equation}
 \Phi_Z:\operatorname{Gr}_{\geq0}(1,13)\longrightarrow\mathbb P^2,
 \qquad[c]\longmapsto\left[\sum_{j=0}^{12}c_jZ_j\right].
 \label{eq:polygon-map}
\end{equation}
Dans la coordinatisation \(Y_0=1\), soient \(v_j=(t_j,t_j^2)\) et \(P_K=\operatorname{conv}(v_0,\ldots,v_{12})\).

\begin{theorem}[Amplituèdre polygonal du registre]
\label{thm:polygon-amplituhedron}
L’image fermée de \(\Phi_Z\) est exactement le polygone convexe
\(P_K=\mathcal A_{13,1,2}(Z)\), de dimension deux et à treize sommets. L’image
de la grassmannienne strictement positive est l’intérieur de \(P_K\).
Les onze cellules
\begin{equation}
 \mathcal S_i=\{[c]:c_0,c_i,c_{i+1}>0,\
 c_j=0\text{ si }j\notin\{0,i,i+1\}\},\qquad 1\leq i\leq11,
 \label{eq:polygon-cells}
\end{equation}
s'appliquent bijectivement, avec un degré un et une orientation positive, sur les intérieurs de \(T_i=[v_0,v_i,v_{i+1}]\). Leurs adhérences recouvrent \(P_K\) et leurs intérieurs sont deux à deux disjoints.
\end{theorem}
\begin{proof}
Dans la normalisation \(\sum c_j=1\), l’application est
\(c\mapsto\sum c_jv_j\), dont l’image est l’enveloppe convexe. Les mineurs
de \eqref{eq:polygon-minors} montrent que chaque triplet ordonné a une
orientation positive. Chaque côté consécutif laisse tous les autres sommets
strictement à sa gauche ; il en va de même du côté final allant de \(v_{12}\)
à \(v_0\). Tous les points sont donc des sommets extrémaux d’un polygone
convexe.

Toute combinaison avec \(c_j>0\) est intérieure, car pour toute droite
d’appui, l’un des sommets se trouve strictement dans le demi-plan
intérieur. Réciproquement, soient \(y\) intérieur et
\(b=13^{-1}\sum v_j\). Pour \(\varepsilon>0\) suffisamment petit,
\(y'=(y-\varepsilon b)/(1-\varepsilon)\) demeure dans le polygone.
En écrivant \(y'=\sum a_jv_j\), \(a_j\geq0\), \(\sum a_j=1\), on obtient
\(y=\sum((1-\varepsilon)a_j+\varepsilon/13)v_j\), avec tous les
coefficients positifs.

Les diagonales issues de \(v_0\) divisent le polygone en onze triangles. Sur chaque support triple, les coordonnées barycentriques sont uniques, car \(\langle0,i,i+1\rangle>0\) ; leur inverse explicite apparaît dans \eqref{eq:polygon-barycentric}. Chaque support définit une cellule positroïdale de rang un de coordonnées \([1:a:b]\), \(a,b>0\), à ces positions. Un réseau orienté planaire comportant une source et trois trajectoires de poids \(1,a,b\) vers les trois destinations réalise ces coordonnées par mesure de frontière. Le déterminant positif prouve l'orientation et le degré un.
\end{proof}

La fibre générique de l’application complète est de dimension dix :
\(\sum c_j=1\), \(\sum c_jt_j=x\) et \(\sum c_jt_j^2=y\) sont trois
équations indépendantes en treize coordonnées. Pour \(y\) intérieur, la
fibre contient un point strictement positif, de sorte que son intersection
avec le simplexe est de dimension \(13-3=10\). Le degré un concerne les
cellules triangulaires ; la forme bidimensionnelle s’obtient au moyen de ces
cellules.

\subsection{Forme canonique et résidus orientés}

Pour \(Y=(1,x,y)^\mathsf T\), nous définissons
\begin{equation}
 \ell_{ab}(Y)=\det(Z_a,Z_b,Y)
 =(t_b-t_a)\{y-(t_a+t_b)x+t_at_b\}.
 \label{eq:polygon-lines}
\end{equation}
L’expression vaut également lorsque \(b<a\). Pour \(a<b<c\), soient
\(D_{abc}=\langle abc\rangle\) et
\begin{equation}
 \lambda_a=\frac{\ell_{bc}}{D_{abc}},\qquad
 \lambda_b=\frac{\ell_{ca}}{D_{abc}},\qquad
 \lambda_c=\frac{\ell_{ab}}{D_{abc}},\qquad
 \lambda_a+\lambda_b+\lambda_c=1.
 \label{eq:polygon-barycentric}
\end{equation}
L’orientation affine est \(dx\wedge dy\). On fixe la convention de frontière
\(\operatorname{Res}^{\partial}(\eta\wedge du/u)=\eta|_{u=0}\) :
la différentielle normale occupe la dernière position. En dimension deux, le
résidu de Poincaré qui place \(du/u\) en premier a le signe opposé.
Les annulations suivantes utilisent uniformément la convention déclarée.

\begin{theorem}[Forme rationnelle et annulation des diagonales]
\label{thm:polygon-form}
La forme canonique du triangle orienté \([v_a,v_b,v_c]\) est
\begin{equation}
 \Omega_{abc}=
 \frac{D_{abc}^2\,dx\wedge dy}
 {\ell_{ab}(Y)\ell_{bc}(Y)\ell_{ca}(Y)}.
 \label{eq:polygon-triangle-form}
\end{equation}
La forme canonique du polygone est
\begin{equation}
 \boxed{\displaystyle
 \Omega_{P_K}(x,y)=
 \sum_{i=1}^{11}
 \frac{t_it_{i+1}(t_{i+1}-t_i)\,dx\wedge dy}
 {(y-t_ix)\{y-(t_i+t_{i+1})x+t_it_{i+1}\}(t_{i+1}x-y)}.}
 \label{eq:polygon-full-form}
\end{equation}
Les pôles des diagonales intérieures s’annulent. Les seuls diviseurs
polaires sont les treize droites des côtés extérieurs. Chaque côté étant paramétré
de son sommet initial à son sommet final par \(s\in[0,1]\), son résidu de
frontière est \(ds/[s(1-s)]\).
\end{theorem}
\begin{proof}
La forme du simplexe barycentrique est \(d\lambda_b\wedge d\lambda_c/(\lambda_a\lambda_b\lambda_c)\). Le jacobien satisfait \(dx\wedge dy=D_{abc}\,d\lambda_b\wedge d\lambda_c\). Substituer \eqref{eq:polygon-barycentric} donne \eqref{eq:polygon-triangle-form}. Pour \(a=0\), simplifier les trois facteurs scalaires des droites donne chaque terme de \eqref{eq:polygon-full-form}.

Sur le côté \(ab\), posons \(\lambda_c=u\),
\(\lambda_b=(1-u)s\), \(\lambda_a=(1-u)(1-s)\). Alors
\[
 \Omega_{abc}=\frac{ds}{s(1-s)}\wedge\frac{du}{u(1-u)}.
\]
Le résidu est la forme de l’intervalle orienté. Deux triangles contigus
induisent des orientations contraires sur leur diagonale et produisent des résidus
opposés. Tous les pôles étant simples, annuler le résidu élimine le
pôle de la diagonale comme diviseur. Chaque côté extérieur apparaît une fois.
La formule homogène est de degré zéro comme forme projective et n’a pas
d’autres diviseurs polaires. Si deux formes rationnelles projectives avaient les
mêmes résidus et aucun autre pôle, leur différence serait une deux-forme
holomorphe sur \(\mathbb P^2\), qui est nécessairement nulle. La
caractérisation et l’unicité sont ainsi prouvées.
\end{proof}

On emploie la notion conventionnelle de géométrie positive au moyen de pôles logarithmiques et de résidus canoniques de frontière, développée par N.~Arkani-Hamed, Y.~Bai et T.~Lam, \emph{Positive Geometries and Canonical Forms}, JHEP 11 (2017), 039, \url{https://arxiv.org/abs/1703.04541}, sections 2, 3, 5.2, 6.1 et 7.2. La sélection des coefficients du présent polygone précède cette reconnaissance et provient de \(K\).

\subsection{Raffinement par retenue sur la géométrie fixe}

Dans un triangle \(T=[A,B,C]\) d’orientation positive, nous écrivons
\begin{equation}
 Y(u,r)=(1-u)A+u\{(1-r)B+rC\},\qquad 0<u,r<1.
 \label{eq:polygon-blowdown}
\end{equation}
Le sommet \(A\) correspond à \(u=0\) et le côté opposé à \(u=1\).
La forme devient
\begin{equation}
 \Omega_T=\frac{du}{u(1-u)}\wedge\frac{dr}{r(1-r)}.
 \label{eq:polygon-product-form}
\end{equation}
Soit \(B_0\geq2\) une base entière déclarée, distincte du sommet \(B\).
La coordonnée \(r\) admet la mise à jour par résidu et retenue
\begin{equation}
 c=\lfloor B_0r\rfloor,\qquad r'=B_0r-c,\qquad
 r=\frac{c+r'}{B_0}.
 \label{eq:polygon-carry}
\end{equation}
Les extrémités conservent leurs marques de frontière. Le chiffre \(c\) demeure comme mémoire et permet d'inverser l'actualisation. Il convient de préciser les coordinatisations des deux transports affines. Le zoom de la sous-cellule agit sur \(r\in[c/B_0,(c+1)/B_0]\) au moyen de \(r'=B_0r-c\). En revanche, le transport de la paire relevée \(L_c(x,y)=(B_0x-c,B_0y+c)\), avec \(x+y=M\), change la masse totale en \(B_0M\) et exprime la coordonnée
\[
 \frac{B_0x-c}{B_0M}=\frac{x}{M}-\frac{c}{B_0M}.
\]
Ce sont des applications de domaines et de normalisations différents. Toutes deux conservent la retenue dans l'état enrichi et permettent une récupération avec leurs données d'échelle. La démonstration suivante utilise le zoom de sous-cellule et la partition géométrique qu'il identifie.

\begin{theorem}[Invariance exacte sous subdivision avec mémoire]
\label{thm:polygon-carry}
Soient \(s_j=j/B_0\) et \(W_j=(1-s_j)B+s_jC\).
Les triangles \(T_j=[A,W_j,W_{j+1}]\), \(0\leq j<B_0\), subdivisent
\(T\) et satisfont
\begin{equation}
 \sum_{j=0}^{B_0-1}\Omega_{T_j}=\Omega_T.
 \label{eq:polygon-carry-additivity}
\end{equation}
L'égalité reste valable à toute profondeur finie et pour toute partition rationnelle ordonnée du côté. Subdiviser les onze triangles de \(P_K\) laisse \(\Omega_{P_K}\) exactement invariante.
\end{theorem}
\begin{proof}
Les régions \(r\in[s_j,s_{j+1}]\) couvrent \(T\) avec des intérieurs disjoints.
Dans les coordonnées communes,
\[
 \Omega_{T_j}=\frac{du}{u(1-u)}\wedge
 \frac{(s_{j+1}-s_j)\,dr}{(r-s_j)(s_{j+1}-r)}.
\]
L’identité rationnelle
\[
 \frac{s_{j+1}-s_j}{(r-s_j)(s_{j+1}-r)}
 =\frac1{r-s_j}+\frac1{s_{j+1}-r}
\]
se télescope : chaque terme intérieur annule celui du sous-intervalle contigu.
Il reste \(1/r+1/(1-r)=1/[r(1-r)]\). Cela prouve l’égalité des formes
rationnelles ; son itération fournit toute profondeur finie.
La preuve utilise seulement \(s_0=0<s_1<\cdots<s_N=1\), de sorte qu’elle couvre
les partitions non uniformes. Le changement
\(r'=(r-s_j)/(s_{j+1}-s_j)\) restitue localement la forme unitaire.
Le mot des retenues identifie la sous-cellule et son inverse récupère la
coordonnée antérieure. Les nouveaux pôles intérieurs s’annulent au moyen de
la même identité télescopique.
\end{proof}

La somme précédente réunit la partition géométrique complète. Lorsqu’un lecteur
de survie sélectionne seulement certaines histoires, son image est l’union
des sous-cellules sélectionnées et sa forme est la somme correspondante ;
elle coïncide avec la forme du polygone complet précisément lorsque cette union
le recouvre avec une multiplicité orientée de un. Le critère de
recouvrement pour une sélection concrète d’histoires s’exprime donc par
une condition vérifiable sur sa chaîne de cellules.

Ce raffinement conserve le polygone fixe. Insérer un nouveau point \((t,t^2)\) de la parabole entre deux paramètres consécutifs \(a<t<b\) produit une autre géométrie : la hauteur de la corde moins \(t^2\) est \((t-a)(b-t)>0\), et le nouveau point se trouve à l'extérieur du polygone précédent. L'invariance démontrée correspond à la subdivision rectiligne des cellules, dont la frontière extérieure reste fixe.

\subsection{Changements de triangulation et reparamétrage positif}

Pour quatre sommets ordonnés \(A,B,C,D\), on a
\begin{equation}
 \Omega_{ABC}+\Omega_{ACD}=\Omega_{ABD}+\Omega_{BCD}.
 \label{eq:polygon-flip}
\end{equation}
Les deux membres ont les mêmes résidus extérieurs et annulent leur diagonale intérieure ; leur égalité découle de l'unicité démontrée. Elle se vérifie aussi en multipliant par leurs dénominateurs linéaires. Tout changement entre triangulations d'un polygone convexe se décompose en ces changements de diagonale, si bien que la forme est indépendante de la triangulation. Une identification avec des coordonnées cluster doit en outre conserver la variété et ses relations de changement.

Si \(D=\operatorname{diag}(d_0,\ldots,d_{12})\), \(d_j>0\), alors
\begin{equation}
 \mathcal A_{13,1,2}(ZD)=\mathcal A_{13,1,2}(Z).
 \label{eq:polygon-rescale}
\end{equation}
L'automorphisme \([c]\mapsto[Dc]\) reparamétrise l'application cinématique. Dans la coordinatisation affine, \(c_j\) devient \(c_jd_j/\sum_\ell c_\ell d_\ell\), qui parcourt de nouveau le simplexe. L'image et sa forme canonique demeurent identiques. Le déplacement d'un historique paramétré et le changement de l'ensemble géométrique complet sont ainsi des opérations distinctes.

\subsection{Extension de rang un aux dimensions supérieures}

Pour \(2\leq m\leq12\), nous définissons
\begin{equation}
 Z_i^{(m)}=(1,t_i,t_i^2,\ldots,t_i^m)^\mathsf T,\qquad
 P_K^{(m)}=\operatorname{conv}
 \{(t_i,t_i^2,\ldots,t_i^m):0\leq i\leq12\}.
 \label{eq:polygon-higher-lift}
\end{equation}
Les mineurs maximaux ordonnés sont les produits de Vandermonde
\(\prod_{a<b}(t_{i_b}-t_{i_a})>0\). La matrice a rang \(m+1\).
L’image de \(\operatorname{Gr}_{\geq0}(1,13)\) est le polytope cyclique
\(P_K^{(m)}=\mathcal A_{13,1,m}(Z^{(m)})\), de dimension \(m\) ; sa fibre
générique a dimension \(12-m\). On utilise ici l’extension mathématique
de la définition amplituédrique à tout \(m\) ; les applications
habituelles aux amplitudes imposent leurs propres valeurs de \(m\).

La triangulation est déterminée par l’ordre fixe des étiquettes. On choisit
le plus petit sommet d’un polytope, on triangule récursivement toutes ses
facettes qui ne le contiennent pas et on prend les cônes depuis ce sommet. La
récurrence commence par les points. Une droite allant du sommet initial
à un point intérieur coupe la frontière opposée dans une facette ; cela
prouve le recouvrement. L’unicité du point d’intersection et la
compatibilité des restrictions induites par le même ordre prouvent
que les intérieurs sont disjoints et que les intersections sont des faces
communes. À chaque étape, la dimension diminue, de sorte que le procédé
s’achève. Il s’agit d’une triangulation par traction déterminée par les
étiquettes.

Pour un simplexe orienté \(T=[p_0,\ldots,p_m]\) de cette triangulation, soient
\[
 A_T=(p_1-p_0\ \cdots\ p_m-p_0),\qquad
 (\lambda_1,\ldots,\lambda_m)^\mathsf T=A_T^{-1}(x-p_0),
 \qquad\lambda_0=1-\sum_{j=1}^m\lambda_j.
\]
L’orientation de ses sommets est choisie de sorte que \(\det A_T>0\).
Sa forme canonique est
\begin{equation}
 \Omega_T(x)=
 \frac{dx_1\wedge\cdots\wedge dx_m}
 {\det(A_T)\prod_{j=0}^m\lambda_j(x)},\qquad
 \Omega_{P_K^{(m)}}=\sum_{T}\Omega_T.
 \label{eq:polygon-higher-form}
\end{equation}
Le changement barycentrique prouve la formule et le fait que la restriction cinématique sur le support de \(T\) est de degré un. Les résidus sont les formes des faces avec les signes induits par l'orientation. Sur une face intérieure, les deux orientations sont opposées et le résidu total s'annule. Sur une face extérieure reste la somme des formes de la triangulation de cette face. La récurrence sur la dimension démontre que cette somme est sa forme canonique. Répéter le résidu le long d'un drapeau de faces atteint les sommets, où la normalisation est \(\pm1\). On obtient ainsi la correspondance des résidus dans toutes les codimensions pour cette collection de rang un.

La même démonstration couvre toute subdivision simpliciale compatible du polytope fixé : ses faces intérieures s'annulent et les faces extérieures conservent les résidus. L'absence de formes holomorphes de degré maximal sur \(\mathbb P^m\) assure l'unicité de la forme rationnelle. En particulier, la stabilité par raffinement, la couverture et les résidus sont établis dans cette collection de dimensions \(2,\ldots,12\). Les degrés grassmanniens \(k>1\) conservent leurs applications et leurs obligations démonstratives spécifiques.

\section{Matrices de moments et réalisations amplituédriques de rang supérieur}
\label{sec:higher-rank}
La partition du registre dodécaphasique permet de prolonger la construction de rang un à des points explicites de rang supérieur. Les mêmes treize nœuds ordonnés fournissent simultanément des mineurs de Plücker positifs et une forme de moments définie positive. Ces deux positivités sont obtenues à partir d'une seule liste d'intervalles ; la seconde prouve que la composition conserve le rang requis.

Soit \(d_j=K_j/Q\), où \(Q=\sum_{j=1}^{12}K_j=6263\), et définissons
\[
 t_0=0,\qquad t_j=\sum_{r=1}^jd_r.
\]
La positivité des douze blocs donne
\(0=t_0<t_1<\cdots<t_{12}=1\). Pour chaque \(2\le k\le9\), nous posons
\[
 C^{(k)}_{ai}=t_i^a\quad(0\le a<k),\qquad
 Z^{(k)}_{bi}=t_i^b\quad(0\le b<k+4),\qquad 0\le i\le12.
\]
La convention \(t_i^0=1\) inclut le nœud zéro. Ainsi, \(C^{(k)}\)
est une matrice de taille \(k\times13\), tandis que \(Z^{(k)}\) a pour
taille \((k+4)\times13\).

\begin{theorem}[Positivité et rang de la réalisation des moments]
La matrice \(C^{(k)}\) représente un point de \(\Gr_{>0}(k,13)\), et tous les mineurs maximaux ordonnés de \(Z^{(k)}\) sont positifs. La matrice
\[
 Y^{(k)}=C^{(k)}(Z^{(k)})^{\mathsf T},\qquad
 Y^{(k)}_{ab}=\sum_{i=0}^{12}t_i^{a+b},
\]
est de rang \(k\). En particulier, pour la donnée externe positive
\(Z^{(k)}\) déjà construite à partir du registre,
\[
 [Y^{(k)}]\in\mathcal A_{13,k,4}(Z^{(k)}),\qquad 2\le k\le9,
\]
où
\[
 \mathcal A_{13,k,4}(Z^{(k)})
 =\{[C(Z^{(k)})^{\mathsf T}]:[C]\in\Gr_{\ge0}(k,13)\}.
\]
\end{theorem}
\begin{proof}
Un mineur maximal de l’une ou l’autre des deux matrices initiales a la forme
\[
 \det(t_{i_j}^{a})_{a=0}^{r-1}
 =\prod_{p<q}(t_{i_q}-t_{i_p})>0,
\]
pour des indices croissants. L'inégalité prouve les deux affirmations sur Plücker. Le bloc initial \(k\times k\) de \(Y^{(k)}\) est la matrice de Hankel \(H^{(k)}\). Pour \(x\ne0\),
\[
 x^{\mathsf T}H^{(k)}x
 =\sum_{i=0}^{12}\left(\sum_{a=0}^{k-1}x_at_i^a\right)^2>0.
\]
En effet, un polynôme non nul de degré inférieur à \(k\le9\) ne
s’annule pas en treize points distincts. Donc \(H^{(k)}\) est définie positive,
son déterminant est strictement positif et \(Y^{(k)}\) est de rang
\(k\). L’appartenance amplituédrique est alors la composition des
deux matrices positives précédentes, dont le rang est déjà vérifié.
\end{proof}

En rang deux, la preuve prend la forme particulièrement explicite
\[
 \det H^{(2)}=13\sum_i t_i^2-\left(\sum_i t_i\right)^2
 =\sum_{i<j}(t_j-t_i)^2>0.
\]
La séparation des nœuds, produite par la positivité du registre, est ainsi convertie en coercivité d'une forme quadratique. À chaque registre fixé correspond un point déterminé pour chaque rang. La couverture complète, les triangulations et les résidus démontrés dans la section précédente conservent leur domaine de rang un ; le présent théorème ajoute la collection explicite de points de rangs deux à neuf et leurs matrices de moments. La distinction permet de réutiliser chaque démonstration avec sa portée exacte.

% Sección aditiva. No modifica la edición de 211 páginas.
% Procedencia y comprobaciones: documentos internos de este mismo directorio.
\section{Raffinement nonadique, drapeaux positifs et compatibilité des moments}
\label{sec:dim-cofinal}

La longueur d'une représentation finie détermine le nombre de mineurs pouvant être formés avec ses colonnes. Le raffinement conserve les colonnes héritées et en ajoute d'autres dont la provenance est connue. La dimension d'une réalisation matricielle peut donc croître sans remplacer le registre qui fixe ses séparations. La construction suivante réunit cette croissance avec deux notions précises de compatibilité : la restriction des évaluations aux marques antérieures et la moyenne des observables sur les cellules descendantes. La première préserve les drapeaux polynomiaux ; la seconde conserve une mesure et sépare l'incrément d'information en détails orthogonaux.

La donnée de départ est la sortie dodécaphasique déjà produite par APP--TRIT--TPK. Les deux feuillets d'APP conservent leurs divisions complètes ; TRIT oriente le régime et l'opérateur \(U_t=\operatorname{Upd}_t\circ\operatorname{Tra}_t\circ
\operatorname{Sel}_t\) transporte l'état enrichi. Le prolongement \(w_6\to w_{12}\to w_{18}\to w_{24}\to w_{30}\to R_{36}\to G_9\) conserve le retour de phase et l'avancement de mémoire. La structure discrète conjointe du continu conserve corrélativement les représentations de cet état, avec ses feuillets, son orientation, sa mémoire et ses frontières. De ce même état proviennent les représentations régionales et le registre signé dont l'inversion donne
\[
 K=(234,543,140,729,659,824,621,58,914,794,146,601),
 \qquad Q=\sum_{j=1}^{12}K_j=6263.
\]
Cette sortie est utilisée avec ses marques et son orientation ; les matrices de
moments sont des réalisations ultérieures. L’indice de raffinement introduit
ci-dessous compte les subdivisions du lecteur géométrique. Son identification
à une durée physique ou à un nombre de microtransitions du TPK
exige de conserver l’horloge de cette autre réalisation.

\subsection{Partitions marquées et croissance cofinale}

Écrivons \(d_j=K_j/Q\), \(a_0=0\) et
\(a_j=\sum_{h=1}^j d_h\). Pour \(r\geq0\), nous définissons
\[
 L_r=12\,9^r,\qquad N_r=L_r+1,\qquad
 t_{r,(j-1)9^r+c}=a_{j-1}+\frac{c}{9^r}d_j,
 \quad 0\leq c\leq9^r.
\]
Les extrémités communes ne sont comptées qu'une fois. Les \(L_r\) cellules sont indexées par \((j,c)\), avec \(1\leq j\leq12\) et \(0\leq c<9^r\), et ont pour longueur \(\ell_{r,j,c}=d_j/9^r\). L'écriture de \(c\) en \(r\) chiffres nonadiques est également conservée, y compris les zéros initiaux. Un successeur a pour indice \(9c+h\), avec \(0\leq h\leq8\) ; la division euclidienne restitue son prédécesseur et le dernier chiffre. Les extrémités partagées retiennent leur marque de frontière ; la convention d'intervalles semi-ouverts ne fait qu'organiser les cellules et ne supprime aucune marque.

\begin{proposition}[Cofinalité et récupération du registre]
\label{prop:dim-grid}
Les marques satisfont
\[
 0=t_{r,0}<t_{r,1}<\cdots<t_{r,L_r}=1,
 \qquad t_{r+1,9i}=t_{r,i}.
\]
Le maximum des longueurs est
\(h_r=(\max_jK_j)/(Q9^r)\), qui tend vers zéro. Pour toute paire
d’entiers \(k\geq1\), \(m\geq1\), il existe \(r_0\) tel que
\(k+m\leq N_r\) pour \(r\geq r_0\). Le registre se retrouve à
tout niveau par
\[
 K_j=Q\sum_{c=0}^{9^r-1}\ell_{r,j,c}.
\]
\end{proposition}
\begin{proof}
Chaque différence consécutive est \(d_j/9^r>0\). En remplaçant \(c\) par \(9c\), la formule du niveau suivant donne exactement la marque antérieure. Les neuf longueurs des successeurs ont pour somme la longueur du prédécesseur ; l'itération prouve la récupération. Enfin, \(N_r=12\,9^r+1\) croît sans borne. Un choix explicite est tout \(r_0\geq0\) tel que \(9^{r_0}\geq(k+m-1)/12\). L'argument entier conserve \(Q\), l'indice de bloc et les chiffres de descendance.
\end{proof}

La normalisation détermine les rapports \(K_j/Q\). La récupération des entiers utilise en outre la charge \(Q\) conservée par le registre. On distingue ainsi la géométrie normalisée et l'information d'échelle accompagnant sa représentation.

\subsection{Drapeaux d'évaluations et réalisations en toute dimension finie}

Pour \(1\leq q\leq N_r\), soit
\[
 V_r^{(q)}=(t_{r,i}^{\,a})_{0\leq a<q,\,0\leq i<N_r},
 \qquad F_r^q=\operatorname{fila}V_r^{(q)}.
\]
La puissance d’exposant zéro vaut un également au nœud zéro.

\begin{theorem}[Drapeau positif et composition de rang arbitraire]
\label{thm:dim-all-ranks}
Les espaces \(F_r^q\) forment un drapeau \(F_r^1\subset F_r^2\subset\cdots\subset F_r^{N_r}\), avec \(\dim F_r^q=q\), et chaque matrice \(V_r^{(q)}\) a tous ses mineurs maximaux ordonnés strictement positifs. Pour \(k\geq1\), \(m\geq1\), \(k+m\leq N_r\), les matrices
\[
 C_r=V_r^{(k)},\quad Z_r=V_r^{(k+m)},\quad
 Y_r=C_rZ_r^{\mathsf T}
\]
satisfont
\[
 (Y_r)_{ab}=\sum_{i=0}^{L_r}t_{r,i}^{a+b},\qquad
 \operatorname{rank}Y_r=k,\qquad
 [Y_r]\in\mathcal A_{N_r,k,m}(Z_r).
\]
La construction atteint toute paire finie \((k,m)\) par un niveau
suffisamment profond de la même partition marquée.
\end{theorem}
\begin{proof}
Pour \(i_1<\cdots<i_q\), le mineur correspondant est
\[
 \det(t_{r,i_b}^{a-1})_{a,b=1}^{q}
   =\prod_{a<b}(t_{r,i_b}-t_{r,i_a})>0.
\]
Cela prouve le rang, la positivité et les inclusions du drapeau. Le bloc
initial \(k\times k\) de \(Y_r\) est le gramien des évaluations
de \(1,t,\ldots,t^{k-1}\). Pour \(x\neq0\),
\[
 x^{\mathsf T}(Y_r)_{[0,k-1],[0,k-1]}x
 =\sum_i\left(\sum_{a=0}^{k-1}x_at_{r,i}^{a}\right)^2>0,
\]
car un polynôme non nul de degré inférieur à \(k\) ne s’annule pas
aux \(N_r\geq k+m\) nœuds distincts. Le bloc est inversible,
et \(Y_r\), qui a \(k\) lignes, est de rang \(k\).
L’appartenance indiquée est l’application de la matrice externe positive
\(Z_r\) au point positif représenté par \(C_r\).
\end{proof}

On emploie la convention mathématique
\[
 \mathcal A_{N,k,m}(Z)=
 \{[CZ^{\mathsf T}]:[C]\in\operatorname{Gr}_{\geq0}(k,N)\}.
\]
La définition pour \(m\) général peut être consultée dans S.~N.~Karp, L.~K.~Williams et Y.~X.~Zhang,
\href{https://people.math.harvard.edu/~williams/papers/amplituhedron-plus.pdf}
{\emph{Decompositions of amplituhedra}}, définition~1.1.
Ici, la donnée externe provient déjà de la partition exprimée par HMT. La conclusion construit une collection déterminée de points et de drapeaux ; les dimensions de la grassmannienne ambiante et d'un atlas de son image sont des objets supplémentaires aux rangs démontrés.

Soit \(E_r:\mathbb R^{N_{r+1}}\to\mathbb R^{N_r}\) la restriction aux positions \(9i\). Pour tout \(q\leq N_r\),
\[
 E_r\bigl((V_{r+1}^{(q)})^{\mathsf T}x\bigr)
       =(V_r^{(q)})^{\mathsf T}x.
\]
Par conséquent, \(E_r:F_{r+1}^q\to F_r^q\) est un isomorphisme sur ces espaces d'évaluations : ses deux réalisations identifient le même polynôme de degré inférieur à \(q\). La composition des restrictions coïncide avec la restriction directe. Les mineurs sur les indices hérités conservent exactement leur valeur.

\begin{proposition}[Récupération à partir d’un drapeau calibré]
Si l’on conserve l’ordre des colonnes, les marques extrêmes et la
calibration \(t_{r,0}=0\), \(t_{r,L_r}=1\), le plan \(F_r^2\)
détermine les nœuds et les séparations. Avec \(Q\) et les
marques de descendance, il détermine \(K\).
\end{proposition}
\begin{proof}
Le plan contient le vecteur constant et le vecteur des nœuds. Toute
autre coordonnée qui, avec le vecteur constant, engendre le même
plan est \(a+bt\), avec \(b\neq0\). Les deux calibrages fixent
\(a=0\), \(b=1\). Les différences consécutives récupèrent les
longueurs, et la proposition~\ref{prop:dim-grid} récupère le registre.
\end{proof}

\subsection{Corps convexes de rang un et croissance du domaine}

Pour \(m\geq1\) et \(m+1\leq N_r\), nous définissons
\[
 P_r^{(m)}=\operatorname{conv}
 \{(t_{r,i},t_{r,i}^2,\ldots,t_{r,i}^m):0\leq i\leq L_r\}.
\]
La normalisation d'une ligne positive de \(C\) en coefficients barycentriques prouve \(\mathcal A_{N_r,1,m}(Z_r)=P_r^{(m)}\). Le déterminant de Vandermonde fournit \(m+1\) points affinement indépendants, de sorte que le corps a pour dimension \(m\).

\begin{proposition}[Recouvrement simplicial en dimension finie et limite convexe]
\label{prop:dim-polytopes}
Chaque \(P_r^{(m)}\) admet une triangulation par ses sommets, de
degré un sur l’intérieur de chaque cellule simpliciale. La somme de ses
formes canoniques a ses résidus intérieurs annulés. De plus,
\[
 P_r^{(m)}\subseteq P_{r+1}^{(m)},\qquad
 \overline{\bigcup_{r\geq r_0}P_r^{(m)}}=
 \operatorname{conv}\{(t,t^2,\ldots,t^m):0\leq t\leq1\}.
\]
\end{proposition}
\begin{proof}
On ordonne les sommets selon leurs marques. Dans chaque face, on choisit son plus petit
sommet ; on triangule récursivement les facettes qui ne le contiennent pas
et on forme les cônes correspondants. La récurrence sur la dimension
prouve la couverture et la disjonction des intérieurs, et la même règle sur les faces
communes garantit la compatibilité. Dans chaque simplexe orienté
\([v_0,\ldots,v_m]\), de matrice
\(A=(v_1-v_0,\ldots,v_m-v_0)\) à déterminant positif et de
coordonnées barycentriques \(\lambda_0,\ldots,\lambda_m\),
l’application barycentrique est affine et inversible. Sa forme est
\[
 \Omega_{[v_0,\ldots,v_m]}=
 \frac{dx_1\wedge\cdots\wedge dx_m}
 {\det A\prod_{j=0}^{m}\lambda_j}.
\]
La restriction résiduelle à chaque face est sa forme orientée. Deux simplexes adjacents induisent des orientations opposées sur leur face commune, de sorte que les résidus intérieurs s'annulent. C'est la même opération de triangulation et de résidu appliquée dans la réalisation de rang un précédente, le nombre de sommets étant maintenant augmenté par le raffinement.

L’inclusion des nœuds hérités prouve l’inclusion des corps. Le maillage
tend vers zéro, de sorte que les courbes échantillonnées sont denses dans la courbe
compacte des moments. Toute combinaison convexe finie de points de la
courbe s’approche par des combinaisons de nœuds d’un niveau commun.
Réciproquement, tous les corps sont contenus dans l’enveloppe
convexe compacte de la courbe. Les deux inclusions prouvent l’égalité.
\end{proof}

L'annulation des résidus opère sur une triangulation ou une subdivision d'un corps fixe. Le passage de \(P_r^{(m)}\) à \(P_{r+1}^{(m)}\) ajoute des sommets et peut agrandir le corps. Leurs formes appartiennent à ces domaines respectifs. La limite convexe précédente spécifie la convergence des corps, sans imposer d'identité entre leurs formes.

\subsection{Mesure du calendrier et convergence des moments nodaux}

Soit \(F_K:[0,1]\to[0,1]\) la fonction croissante affine par morceaux
qui envoie \((j-1)/12\) sur \(a_{j-1}\) et \(j/12\) sur \(a_j\).
Sur l’intervalle correspondant, sa pente est \(12d_j>0\).
La mesure du calendrier transporté est
\[
 \mu_K=(F_K)_*(ds),\qquad
 \frac{d\mu_K}{dt}=\frac1{12d_j}\quad
 (a_{j-1}<t<a_j).
\]
Chaque intervalle initial a pour masse \(1/12\), et chaque cellule de niveau
\(r\) a pour masse \(1/L_r\). Ce choix attribue le même poids
aux positions du calendrier avant de réaliser leurs longueurs.
La mesure de longueur \(dt\) est un autre lecteur ; les deux coïncident lorsque
toutes les séparations initiales sont égales.

Ses moments sont rationnels et déterminés par \(K\) :
\begin{equation}
 M_p=\int_0^1t^p\,d\mu_K(t)
   =\frac1{12(p+1)}\sum_{j=1}^{12}
       \frac{a_j^{p+1}-a_{j-1}^{p+1}}{d_j},\qquad p\geq0.
 \label{eq:dim-moments}
\end{equation}
\begin{proposition}[Récupération par moments et quantiles]
\label{prop:dim-moment-recovery}
La suite complète \((M_p)_{p\geq0}\) détermine \(\mu_K\). Si \(H_K(t)=\mu_K([0,t])\) est sa fonction de répartition, alors
\[
 a_j=H_K^{-1}(j/12),\qquad
 K_j=Q\bigl(H_K^{-1}(j/12)-H_K^{-1}((j-1)/12)\bigr).
\]
Par conséquent, les moments complets, la charge \(Q\) et l’ordre
du calendrier récupèrent le registre initial.
\end{proposition}
\begin{proof}
Deux probabilités sur \([0,1]\) ayant les mêmes moments ont les
mêmes intégrales pour tous les polynômes. L’approximation uniforme
des fonctions continues par des polynômes et la masse finie étendent cette
égalité à toutes les fonctions continues ; celles-ci déterminent la mesure
de Borel. La densité de \(\mu_K\) est positive et constante sur chaque
intervalle initial. Sa fonction de répartition est donc continue et
strictement croissante, et vérifie \(H_K(a_j)=j/12\). L’inversion
de cette égalité et la différence des extrémités prouvent les formules.
\end{proof}

La probabilité nodale correspondant au théorème \ref{thm:dim-all-ranks} est \(\nu_r=N_r^{-1}\sum_{i=0}^{L_r}\delta_{t_{r,i}}\).

\begin{theorem}[Limite de la représentation nodale]
\label{thm:dim-nodal-limit}
On a \(\nu_r\rightharpoonup\mu_K\). Si \(f\) est continue,
\[
 \left|\int f\,d\nu_r-\int f\,d\mu_K\right|
 \leq \omega_f(h_r)+\frac{\operatorname{osc}(f)}{N_r},
\]
où \(\omega_f\) est son module de continuité. Pour \(p\geq1\), l'erreur du moment d'ordre \(p\) est bornée par \(p h_r+1/N_r\). Pour chaque paire fixée \((k,m)\),
\[
 \frac1{N_r}Y_r\longrightarrow (M_{a+b})_{
     0\leq a<k,\,0\leq b<k+m}=:Y_\infty^{(k,m)},
 \qquad\operatorname{rank}Y_\infty^{(k,m)}=k.
\]
\end{theorem}
\begin{proof}
La mesure qui assigne une masse \(1/L_r\) à chaque extrémité gauche de
cellule approche son intégrale avec une erreur au plus égale à \(\omega_f(h_r)\).
En effet, chaque cellule a cette masse et un diamètre au plus égal à \(h_r\).
La mesure \(\nu_r\) est \(L_r/N_r\) fois cette mesure plus
\(\delta_1/N_r\), ce qui ajoute au plus
\(\operatorname{osc}(f)/N_r\). Pour \(f(t)=t^p\), la constante
de Lipschitz est \(p\). La convergence de chaque entrée découle de
ces bornes. Le bloc initial de la matrice limite est défini positif :
l’intégrale du carré d’un polynôme non nul est positive, car
\(\mu_K\) a une densité strictement positive sur tout intervalle
ouvert de \([0,1]\). La limite conserve ainsi le rang \(k\).
\end{proof}

Les moments nodaux varient exactement par l'incorporation de nœuds. Si \(\mathcal J_r\) est l'ensemble des nouvelles positions et \(z_q(t)=(1,t,\ldots,t^{q-1})^{\mathsf T}\), alors
\[
 V_{r+1}^{(q)}(V_{r+1}^{(q)})^{\mathsf T}
 =V_r^{(q)}(V_r^{(q)})^{\mathsf T}
  +\sum_{i\in\mathcal J_r}z_q(t_{r+1,i})z_q(t_{r+1,i})^{\mathsf T}.
\]
Chaque terme est positif semi-défini. Pour les matrices rectangulaires \(Y_r\), la même actualisation vaut avec \(z_kz_{k+m}^{\mathsf T}\). L'égalité exprime le changement de la mesure de comptage ; la normalisation et le théorème précédent décrivent sa limite. L'objet limite de rang \(k\) possède des coordonnées dans une grassmannienne finie, tandis que les données externes \(Z_r\) appartiennent aux différentes représentations de chaque niveau.

\subsection{Moyennes cellulaires, opérateurs compatibles et détail orthogonal}

Pour une cellule \(I_{r,i}=[b_{r,i},b_{r,i}+\ell_{r,i}]\), nous posons
\[
 a_{r,i}^{(p)}=
 \frac{(b_{r,i}+\ell_{r,i})^{p+1}-b_{r,i}^{p+1}}
 {(p+1)\ell_{r,i}},\qquad
 A_r^{(q)}=(a_{r,i}^{(p)})_{0\leq i<L_r,\,0\leq p<q}.
\]
Ce sont les moyennes des monômes par rapport à la probabilité conditionnelle
de \(\mu_K\) dans chaque cellule. Cette probabilité est uniforme dans la
cellule, bien que \(\mu_K\) ne soit pas la mesure de longueur globale.

Soit \(\mathcal H_r=\mathbb R^{L_r}\) muni du produit \(\langle x,y\rangle_r=L_r^{-1}\sum_i x_i y_i\). L'application \(J_r:\mathcal H_r\to\mathcal H_{r+1}\) répète la valeur d'un prédécesseur dans ses neuf successeurs. Son adjointe \(R_r=J_r^*\) prend la moyenne de ces neuf valeurs. On a \(R_rJ_r=I\).

\begin{theorem}[Compatibilité exacte et bilan des moments]
\label{thm:dim-cell-balance}
Pour tout \(p\geq0\),
\[
 R_r a_{r+1}^{(p)}=a_r^{(p)},\qquad
 \frac1{L_r}\sum_i a_{r,i}^{(p)}=M_p.
\]
Définissons
\[
 D_r^{(q)}=A_{r+1}^{(q)}-J_rA_r^{(q)},\quad
 G_r^{(q)}=\frac1{L_r}(A_r^{(q)})^{\mathsf T}A_r^{(q)}.
\]
Alors
\[
 R_rD_r^{(q)}=0,\qquad
 G_{r+1}^{(q)}=G_r^{(q)}+
       \frac1{L_{r+1}}(D_r^{(q)})^{\mathsf T}D_r^{(q)},
\]
et
\[
 G_r^{(q)}\longrightarrow H_K^{(q)}=(M_{a+b})_{0\leq a,b<q}.
\]
Les détails et le niveau initial reconstruisent tous les niveaux. Pour
\(q\leq L_r\), la matrice \(A_r^{(q)}\) a rang \(q\) et
tous les mineurs ordonnés d’ordre \(q\) sont positifs.
\end{theorem}
\begin{proof}
Les masses des neuf successeurs sont égales et leurs intervalles partitionnent l'intervalle prédécesseur. L'additivité de l'intégrale donne la première identité ; la somme sur toutes les cellules donne la seconde. La décomposition \(A_{r+1}=J_rA_r+D_r\) est orthogonale dans chaque colonne puisque \(R_rD_r=0\). En développant son Gram, les termes croisés s'annulent, et \(J_r\) est une isométrie. Cela prouve le bilan matriciel.

L'observable constant par cellules qui prend la valeur \(a_{r,i}^{(p)}\) approche uniformément \(t^p\), avec une erreur ne dépassant pas \(p h_r\) pour \(p\geq1\). Il converge donc dans \(L^2(\mu_K)\), et les produits scalaires convergent vers \(M_{a+b}\). La reconstruction s'obtient en itérant \(A_{r+1}=J_rA_r+D_r\).

Pour prouver les mineurs, choisissons \(q\) cellules ordonnées. Par
multilinéarité et Fubini, le déterminant de leurs moyennes est la moyenne
sur le produit de ces cellules de
\(\prod_{a<b}(x_b-x_a)\). Cet intégrande est positif à
l’intérieur du produit car les cellules sont disjointes et ordonnées.
Leurs frontières sont de mesure nulle. La moyenne est strictement positive.
\end{proof}

L'identité se prolonge à la paire de degrés \(k\) et \(k+m\) :
\[
 \overline Y_r=\frac1{L_r}(A_r^{(k)})^{\mathsf T}A_r^{(k+m)},
 \qquad
 \overline Y_{r+1}-\overline Y_r
 =\frac1{L_{r+1}}(D_r^{(k)})^{\mathsf T}D_r^{(k+m)}.
\]
Pour \(k+m\leq L_r\), les matrices de moyennes définissent une
seconde réalisation amplituédrique positive, à \(L_r\) colonnes,
et \(\overline Y_r\) est de rang \(k\). Ce sont des moyennes de cellules,
tandis que \(Y_r\) utilise des évaluations en \(N_r=L_r+1\) nœuds. Les deux
réalisations convergent vers \(Y_\infty^{(k,m)}\) ; leur coïncidence
à la limite découle des démonstrations précédentes.

La version opératorielle est également explicite. Pour une fonction continue \(f\), soit \(\Psi_r(f)\) l'opérateur diagonal de ses moyennes conditionnelles dans \(\mathcal H_r\). Il est linéaire, positif et unital, et
\[
 R_r\Psi_{r+1}(f)J_r=\Psi_r(f).
\]
Cette égalité se prouve en appliquant les deux membres à chaque vecteur de base.
Pour \(X(t)=t\), les deux premiers opérateurs déterminent, cellule par
cellule,
\[
 \Psi_r(X)_i=b_{r,i}+\frac{\ell_{r,i}}2,\qquad
 \Psi_r(X^2)_i-\Psi_r(X)_i^2=\frac{\ell_{r,i}^2}{12}.
\]
Si \(v_{r,i}\) désigne cette variance, la formule \(\ell_{r,i}=\sqrt{12v_{r,i}}\) récupère la longueur ; la première identité récupère l'extrémité initiale. Avec l'orientation, les marques et \(Q\), on récupère à nouveau \(K\). La moyenne est une compression positive : la variance écrite quantifie pourquoi \(\Psi_r(X^2)\) et \(\Psi_r(X)^2\) sont des opérations distinctes.

\subsection{Transport de la mesure et portée de la croissance dimensionnelle}

Deux listes positives de séparations \(d\) et \(d'\), avec les
mêmes marques, déterminent un homéomorphisme croissant affine par intervalles
\(F=F_{K'}\circ F_K^{-1}\). Il envoie chaque cellule sur sa cellule homologue
et vérifie \(F_*\mu_K=\mu_{K'}\). C’est pourquoi
\[
 \mathcal U_F:L^2(\mu_K)\longrightarrow L^2(\mu_{K'}),
 \qquad\mathcal U_F f=f\circ F^{-1}
\]
est une isométrie surjective. Si \(P_r\) et \(P'_r\) sont les projections sur les fonctions constantes par cellules, alors \(P'_r\mathcal U_F=\mathcal U_FP_r\) : les moyennes conditionnelles sont transportées parce que les masses des cellules correspondantes coïncident. Cela s'applique en particulier au flot positif de séparations dirigé par \(K\), en conservant la direction de chaque bloc dans ses successeurs. Les moments de la coordonnée spatiale peuvent changer ; l'isométrie transporte l'observable avec la mesure et ses projections.

La croissance démontrée est cofinale en toute dimension externe finie. Elle conserve un système de drapeaux positifs, une collection explicite de points amplituèdraux, des corps convexes complets en rang un et une mesure dont les opérateurs de moyenne sont exactement compatibles. La mémoire des détails permet d'inverser la compression entre niveaux. Les preuves de couverture et de résidus de rang un conservent leur domaine ; pour les rangs supérieurs, le présent résultat fixe les matrices, la positivité, la récupération, la compatibilité et la limite avec des types explicites. Cette séparation situe la portée de chaque construction dans le même développement sans confondre croissance de dimension, paramétrisation d'un point et classification de toutes les cellules d'une image.

\section{Moments généalogiques, opérateurs de Jacobi et résidus spectraux}
\label{sec:jacobi-generated}

La mesure induite par le calendrier fournit un lien entre la
géométrie des cellules et une réalisation spectrale positive. Ce lien
utilise toute la partition marquée : chacun de ses douze intervalles conserve
la même masse chronologique, tandis que sa longueur provient de la coordonnée
correspondante du registre. La densité compense exactement cette
longueur. Par conséquent, la distribution spatiale de la mesure retient
l’anisotropie du registre bien que la distribution entre fenêtres soit
uniforme.

Soient \(K_j>0\), \(Q=\sum_{j=1}^{12}K_j\), \(d_j=K_j/Q\) et \(t_j=\sum_{i\leq j}d_i\), avec \(t_0=0\). La mesure de probabilité utilisée dans cette section est
\[
 \dd\mu_K(t)=\sum_{j=1}^{12}
 \frac{\mathbf1_{[t_{j-1},t_j]}(t)}{12d_j}\,\dd t.
\]
Les extrémités partagées sont de mesure nulle. Ses moments sont rationnels pour tout registre entier :
\begin{equation}
 m_r=\int_0^1t^r\,\dd\mu_K(t)
 =\frac1{12(r+1)}\sum_{j=1}^{12}
   \frac{t_j^{r+1}-t_{j-1}^{r+1}}{d_j},\qquad r\geq0.
 \label{eq:jacobi-moments}
\end{equation}
L'égalité \(m_0=1\) exprime la conservation de la masse totale. Pour chaque \(n\geq1\), la matrice de Hankel \(H_n=(m_{a+b})_{0\leq a,b<n}\) est définie positive : sa forme quadratique est l'intégrale du carré d'un polynôme non nul sur un intervalle où la densité est positive. Cette observation permet de passer des moments à un opérateur sans choisir un spectre cible.

\begin{theorem}[Récurrence tridiagonale déterminée par la partition]
\label{thm:jacobi-recurrence}
Il existe une unique suite de polynômes unitaires \(p_n\), de degré \(n\),
orthogonaux pour \(\mu_K\). Si
\[
 h_n=\int p_n^2\,\dd\mu_K>0,\qquad
 a_n=\frac{\int t p_n^2\,\dd\mu_K}{h_n},\qquad
 \beta_n=\frac{h_n}{h_{n-1}}>0\quad(n\geq1),
\]
alors \(p_0=1\), \(p_1=t-a_0\) et
\begin{equation}
 p_{n+1}(t)=(t-a_n)p_n(t)-\beta_np_{n-1}(t),\qquad n\geq1.
 \label{eq:jacobi-recurrence}
\end{equation}
Tous les coefficients \(a_n,\beta_n\) sont rationnels pour le registre
entier déclaré. Dans la base orthonormée \(\psi_n=p_n/\sqrt{h_n}\),
la compression de la multiplication par \(t\) aux polynômes de degré
inférieur à \(n\) a pour matrice symétrique
\[
 J_n=\begin{pmatrix}
 a_0&\sqrt{\beta_1}&&0\\
 \sqrt{\beta_1}&a_1&\ddots&\\
 &\ddots&\ddots&\sqrt{\beta_{n-1}}\\
 0&&\sqrt{\beta_{n-1}}&a_{n-1}
 \end{pmatrix}.
\]
\end{theorem}
\begin{proof}
La positivité de \(H_n\) garantit que le système linéaire qui impose
l’orthogonalité à un polynôme unitaire a une solution unique. Comme ses
coefficients sont des moments rationnels, la solution est rationnelle. On
décompose \(tp_n\) dans la base unitaire jusqu’au degré \(n+1\). Pour
\(j\leq n-2\), la symétrie du produit scalaire donne
\(\langle tp_n,p_j\rangle=\langle p_n,tp_j\rangle=0\).
Seuls subsistent les termes de degrés \(n+1,n,n-1\). Leurs
coefficients sont, respectivement, \(1,a_n,h_n/h_{n-1}\), ce qui
démontre la récurrence. La normalisation par \(\sqrt{h_n}\)
transforme les deux coefficients adjacents en \(\sqrt{\beta_n}\).
\end{proof}

Pour le registre exprimé dans la section~\ref{sec:register}, la première entrée diagonale et le carré du premier couplage sont
\[
 a_0=m_1=\frac{71195}{150312},\qquad
 \beta_1=m_2-m_1^2=\frac{2206226167}{22593697344}.
\]
La première égalité calcule le centre de masse de la partition chronologique ;
la seconde calcule sa dispersion et détermine le premier lien de la
matrice tridiagonale. Les deux nombres s’obtiennent en intégrant les longueurs
engendrées, et leur signification structurelle précède toute évaluation
décimale.

\begin{theorem}[Spectre simple, quadrature et résidus positifs]
\label{thm:jacobi-residues}
Les valeurs propres \(\lambda_1<\cdots<\lambda_n\) de \(J_n\) sont simples et appartiennent à \((0,1)\). Il existe des poids uniques \(w_i>0\), de somme un, tels que
\begin{equation}
 \int f(t)\,\dd\mu_K(t)=\sum_{i=1}^{n}w_i f(\lambda_i)
 \qquad(\deg f\leq2n-1).
 \label{eq:jacobi-quadrature}
\end{equation}
La résolvante observée dans la constante unitaire \(e_0\) satisfait
\begin{equation}
 g_n(z)=\langle e_0,(zI-J_n)^{-1}e_0\rangle
       =\sum_{i=1}^{n}\frac{w_i}{z-\lambda_i}
       =\cfrac1{z-a_0-\cfrac{\beta_1}{z-a_1-
          \cfrac{\beta_2}{\ddots-\cfrac{\beta_{n-1}}{z-a_{n-1}}}}}.
 \label{eq:jacobi-resolvent}
\end{equation}
En particulier, ses pôles et ses résidus s’obtiennent à partir du registre par les
opérations de moment, d’orthogonalisation et de compression.
\end{theorem}
\begin{proof}
Pour tout polynôme non nul \(p\) de degré inférieur à \(n\),
\[
 0<\int_0^1 t|p(t)|^2\,\dd\mu_K(t)
   <\int_0^1 |p(t)|^2\,\dd\mu_K(t).
\]

La caractérisation variationnelle place le spectre de \(J_n\) dans
\((0,1)\). Chaque sous-diagonale est positive. Les équations d’un vecteur propre
déterminent récursivement toutes ses coordonnées à partir de la première ;
si celle-ci s’annulait, le vecteur entier s’annulerait. Chaque espace propre est
ainsi de dimension un. La symétrie de la matrice donne la diagonalisation orthogonale
et la simplicité du spectre. Pour des vecteurs propres unitaires \(v_i\), on prend
\(w_i=|\langle e_0,v_i\rangle|^2>0\). La complétude de la base
propre donne \(\sum_iw_i=1\) et la première expression de la résolvante.

La récurrence tridiagonale montre par développement du déterminant que \(\det(tI-J_n)=p_n(t)\). Pour \(\deg f\leq2n-1\), la division euclidienne s'écrit \(f=qp_n+r\), avec \(\deg q,\deg r<n\). L'orthogonalité annule l'intégrale de \(qp_n\), et son évaluation aux \(\lambda_i\) s'annule également. Pour \(r\) de degré inférieur à \(n\), les multiplications successives de la constante par \(t\) restent dans l'espace des polynômes jusqu'au degré requis ; par conséquent, \(\int r\,\dd\mu_K=\langle e_0,r(J_n)e_0\rangle\). La résolution spectrale prouve la formule de quadrature. Son unicité se déduit de l'inversibilité de la matrice de Vandermonde évaluée aux racines distinctes. Enfin, le complément de Schur de la première ligne et de la première colonne de \(zI-J_n\), itéré sur les sous-matrices terminales, produit la fraction continue.
\end{proof}

\begin{figure}[htbp]
\centering
\begin{tikzpicture}[>=Latex, node distance=8mm and 8mm,
 every node/.style={font=\small}, box/.style={draw,rounded corners=1pt,
 text width=37mm,minimum height=13mm,align=center}]
\node[box] (k) {Registre marqué\\\(K,Q\)};
\node[box,right=of k] (m) {Mesure chronologique\\\(\mu_K\)};
\node[box,right=of m] (h) {Moments positifs\\\(H_n\)};
\node[box,below=of h] (j) {Opérateur tridiagonal\\\(J_n\)};
\node[box,left=of j] (g) {Résolvante\\\(g_n(z)\)};
\node[box,left=of g] (w) {Pôles et résidus\\\(\lambda_i,w_i>0\)};
\draw[->] (k)--(m);\draw[->] (m)--(h);\draw[->] (h)--(j);
\draw[->] (j)--(g);\draw[->] (g)--(w);
\end{tikzpicture}
\caption{Réalisation spectrale de la partition HMT. La géométrie des
cellules détermine la mesure ; la positivité de ses moments produit
l’opérateur et les résidus positifs de sa résolvante.}
\label{fig:jacobi-generated}
\end{figure}

La formule de la résolvante apporte une lecture spectrale de la mémoire
spatiale. Les poids positifs décrivent combien la coordonnée
constante observe de chaque mode ; les pôles sont les fréquences algébriques de
l’opérateur de multiplication comprimé. Les moments d’ordre croissant
récupèrent des détails successifs de la mesure, et les coefficients de la
fraction continue conservent cette information au moyen d’une récurrence
locale à trois termes. L’apparence classique d’une matrice de Jacobi
se reconnaît après avoir construit ses coefficients à partir des
cellules HMT. La chaîne causale est fixée par
\[
 K\longmapsto(d,t)\longmapsto\mu_K\longmapsto(m_r)_{r\geq0}
 \longmapsto\bigl((a_n)_{n\geq0},(\beta_n)_{n\geq1}\bigr)
 \longmapsto(J_n,g_n).
\]

\begin{proposition}[Compatibilité de dimension et de raffinement]
\label{prop:jacobi-refinement}
Les matrices \(J_n\) sont des compressions principales compatibles de la
même multiplication bornée par \(t\) dans \(L^2(\mu_K)\). Leurs
mesures spectrales \(\sum_iw_i\delta_{\lambda_i}\) convergent
faiblement vers \(\mu_K\). Si chaque sous-intervalle nonadique est représenté
par son milieu et par sa masse \(1/(12\cdot9^r)\), la mesure
atomique résultante \(\mu_{K,r}\) approche les moments de
\(\mu_K\). Pour un degré fixé \(n\), la construction à partir de ces
moments retrouve les coefficients \(a_j,\beta_j\) de la même
matrice lorsque le raffinement tend vers une profondeur infinie.
\end{proposition}
\begin{proof}
L'orthogonalisation de la suite complète utilise les mêmes formules lorsque \(n\) augmente, de sorte que chaque bloc principal reste fixe. La quadrature reproduit exactement tout polynôme lorsque \(2n-1\) dépasse son degré. L'approximation uniforme de fonctions continues sur \([0,1]\) par des polynômes, avec la masse totale égale à un, prouve la convergence faible. Les sommes de points milieux convergent dans chacun des douze intervalles avec sa densité constante, et approchent donc tous les moments. En dimension fixée, les coefficients d'orthogonalisation sont des fonctions rationnelles d'un nombre fini de moments, dont les dénominateurs sont formés de déterminants de Gram positifs. La convergence de ces moments et la positivité de la limite assurent la convergence des coefficients.
\end{proof}

\begin{falsifier}
Une longueur nulle invalide la densité déclarée et peut détruire le
rang des moments. Le remplacement de la mesure fixée par des dénombrements
non normalisés sur des maillages distincts change l’opérateur. Les résidus
de \(g_n\) sont des poids spectraux de la résolvante, tandis que les résidus
d’une forme canonique se définissent sur des diviseurs géométriques : une
identification entre les deux notions exige de spécifier la forme et
l’application correspondante. De même, \(J_n\) est adimensionnel ; son
identification à un hamiltonien physique requiert la représentation
et l’échelle d’action et d’horloge utilisées par ce secteur.
\end{falsifier}

\begin{theorem}[Fidélité de la réalisation spectrale marquée]
\label{thm:jacobi-marked-faithfulness}
L’opérateur infini de Jacobi \(J\), avec son vecteur cyclique
unitaire \(e_0\) et la charge \(Q\), détermine le registre \(K\).
Deux réalisations de ce type sont unitairement équivalentes par
une application qui conserve le vecteur marqué et la charge si et seulement
si leurs registres coïncident. Le spectre comme ensemble est, en revanche,
\([0,1]\) pour tout registre positif du domaine.
\end{theorem}
\begin{proof}
Les polynômes sont denses dans \(L^2(\mu_K)\) : on approche d’abord des
fonctions continues, puis on utilise leur densité pour la mesure finie.
La base orthonormale construite définit une application unitaire qui envoie
\(e_n\) sur \(\psi_n\). Par cette application, \(J\) est la multiplication par
\(t\), un opérateur auto-adjoint borné, et \(e_0\) représente la
fonction constante un. Par conséquent,
\[
 \langle e_0,J^re_0\rangle=m_r\qquad(r\geq0).
\]
Ces moments déterminent la mesure : deux mesures ayant les mêmes
moments coïncident sur les polynômes, puis sur les fonctions continues par
approximation uniforme et enfin comme mesures de Borel.
Sa fonction de répartition \(F(t)=\mu_K([0,t])\) est continue et strictement
croissante. Chaque intervalle initial a une masse \(1/12\), de sorte que
\[
 t_j=F^{-1}(j/12),\qquad
 K_j=Q\bigl(F^{-1}(j/12)-F^{-1}((j-1)/12)\bigr).
\]
Une équivalence unitaire qui conserve \(e_0\) conserve les moments et donc le registre. La réciproque découle de la construction unique.

Pour le dernier énoncé, la densité de \(\mu_K\) est positive dans tout sous-intervalle. Chaque \(\lambda\in[0,1]\) admet des fonctions unitaires concentrées dans des intervalles de plus en plus petits autour de \(\lambda\), sur lesquelles la norme de \((J-\lambda I)f\) tend vers zéro. Ainsi, tout l'intervalle appartient au spectre. En dehors de celui-ci, la multiplication par \(1/(t-\lambda)\) est bornée et définit la résolvante. Cela prouve l'égalité du spectre.
\end{proof}

Le théorème distingue l'étendue du spectre et l'information qu'il conserve relativement à une observation marquée. Le même intervalle spectral admet les différentes distributions de la collection ; la mesure observée depuis l'unité restitue les séparations originales. Le lien avec la représentation grassmannienne est désormais réversible : les mineurs marqués reconstruisent \(K\), le registre construit \((J,e_0,Q)\), et ses moments et quantiles restituent le même registre. Chaque lecteur qui se factorise par \((K,\xi)\) possède donc une expression spectrale marquée dès lors que les autres coordonnées \(\xi\) sont conservées. La réalisation spectrale contient une structure, en plus des positions des pôles ou des valeurs propres.

\begin{figure}[htbp]
\centering
\begin{tikzpicture}[x=1.9cm,y=1.05cm,>=Stealth,font=\small]
\foreach \j/\lab in {0/0,1/1,2/2,3/3,5/12}{
  \node[circle,fill=black,inner sep=1.6pt,label=above:{$a_{\lab}$}] (a\j) at (\j,1.4) {};
  \node[circle,fill=black,inner sep=1.6pt,label=below right:{$b_{\lab}$}] (b\j) at (\j,0) {};
  \node (q\j) at (\j,-1) {$q_{\lab}$};
  \draw[->] (b\j) -- (q\j);
}
\foreach \j/\prev in {1/0,2/1,3/2}{
 \draw[->] (a\prev)--(a\j);\draw[->] (b\prev)--(b\j);
 \draw[->,blue!65!black,thick] (a\j)--node[right]{$d_{\j}$}(b\j);
}
\draw[->,blue!65!black,thick] (a5)--node[right]{$d_{12}$}(b5);
\draw[dashed,->] (a3)--(a5);\draw[dashed,->] (b3)--(b5);
\node at (4,.7) {$\cdots$};
\node[left] at (a0) {$s_2$};\node[left] at (b0) {$s_1$};
\end{tikzpicture}
\caption{Réseau planaire à deux lignes. Les arêtes horizontales et de sortie ont un poids égal à un. Chaque chemin depuis la source supérieure descend une seule fois. Pour deux destinations \(i<j\), les paires disjointes descendent après \(i\) et avant \(j\) ; leur poids total est \(\Delta_{ij}=d_{i+1}+\cdots+d_j\). Les traits discontinus conservent les huit positions intermédiaires omises uniquement dans le dessin.}
\end{figure}

\begin{figure}[htbp]
\centering
\begin{tikzpicture}[x=9cm,y=6cm,font=\small]
\foreach \j/\num in {0/0,1/234,2/777,3/917,4/1646,5/2305,6/3129,7/3750,8/3808,9/4722,10/5516,11/5662,12/6263}{
 \pgfmathsetmacro{\tx}{\num/6263}
 \coordinate (v\j) at (\tx,\tx*\tx);
}
\fill[blue!5] (v0)--(v1)--(v2)--(v3)--(v4)--(v5)--(v6)--(v7)--(v8)--(v9)--(v10)--(v11)--(v12)--cycle;
\foreach \j in {2,...,11}{\draw[blue!45,thin] (v0)--(v\j);}
\draw[blue!80!black,thick] (v0)--(v1)--(v2)--(v3)--(v4)--(v5)--(v6)--(v7)--(v8)--(v9)--(v10)--(v11)--(v12)--cycle;
\foreach \j in {0,...,12}{\fill (v\j) circle[radius=.035cm];}
\node[below left] at (v0) {$v_0$};
\node[below right] at (v3) {$v_3$};
\node[below right] at (v6) {$v_6$};
\node[below right] at (v9) {$v_9$};
\node[above] at (v12) {$v_{12}$};
\end{tikzpicture}
\caption{Polygone rationnel déterminé par les treize paramètres
\(t_j\) de \eqref{eq:polygon-Z}, dessiné avec ses coordonnées effectives
\((t_j,t_j^2)\). La triangulation à partir de \(v_0\) comporte onze cellules.
Chaque diagonale intérieure reçoit deux résidus opposés. Cinq
sommets sont étiquetés pour faciliter la lecture ; les treize sont représentés.}
\end{figure}

\section{Flux de poids, géométrie de l'information et dualité de Legendre}
\label{sec:convex-refinement}

La déformation des intervalles admet une description variationnelle
exacte. Soient \(d_j=K_j/Q>0\), \(\sum_jd_j=1\), et \(u_j\) les
coordonnées de la direction \eqref{eq:u-main}. Pour \(s\in\RR\),
on définit
\begin{equation}
 Z(s)=\sum_{j=1}^{12}d_j e^{su_j},\qquad
 \psi(s)=\log Z(s),\qquad
 d_j(s)=\frac{d_j e^{su_j}}{Z(s)}.
 \label{eq:convex-partition}
\end{equation}
Tous les arguments de l'exponentielle sont sans dimension. Le paramètre \(s\) est la coordonnée de cette déformation géométrique. L'attribution d'une durée physique relève d'un lecteur temporel déclaré.

\begin{theorem}[Convexité stricte et coordonnée duale]
La fonction \(\psi\) est analytique et strictement convexe sur \(\RR\).
Si \(u_- =\min_j u_j\) et \(u_+=\max_j u_j\), l’application
\(v=\psi'(s)\) est un difféomorphisme de \(\RR\) sur \((u_-,u_+)\).
Pour chaque \(v\) dans cet intervalle, il existe une unique coordonnée \(s(v)\), et
\[
 \psi^*(v)=s(v)v-\psi(s(v)),\qquad
 (\psi^*)'(v)=s(v),\qquad
 (\psi^*)''(v)=\frac1{\psi''(s(v))}>0.
\]
\end{theorem}
\begin{proof}
Les sommes finies permettent de dériver terme à terme. On obtient
\begin{align}
 \psi'(s)&=\sum_jd_j(s)u_j=:\overline u(s),\notag\\
 \psi''(s)&=\sum_jd_j(s)(u_j-\overline u(s))^2.
 \label{eq:variance}
\end{align}
Les poids sont strictement positifs et les \(u_j\) ne sont pas tous égaux ; la variance est donc strictement positive. Lorsque \(s\) tend vers \(+\infty\), la division du numérateur et du dénominateur de \(\psi'\) par \(e^{su_+}\) montre que seuls subsistent les indices tels que \(u_j=u_+\). Ainsi \(\psi'(s)\to u_+\). L'argument avec \(u_-\) donne la limite opposée. La monotonie et le théorème d'inversion locale fournissent une bijection lisse. La fonction \(sv-\psi(s)\) a son unique maximum en \(s(v)\) ; la dérivation de la valeur maximale démontre les deux identités duales.
\end{proof}

L’évolution de chaque intervalle est
\begin{equation}
 \dot d_j(s)=d_j(s)\bigl(u_j-\overline u(s)\bigr).
 \label{eq:replicator}
\end{equation}
La somme des dérivées est nulle. La longueur totale demeure donc
unitaire tandis que les longueurs internes se redistribuent. Cette conservation
de l’échelle globale permet d’attribuer le changement des birapports à une
déformation des positions et non à une dilatation commune.

Dans le simplexe ouvert des poids, la métrique
\(g_d(\xi,\zeta)=\sum_j\xi_j\zeta_j/d_j\), sur les vecteurs tangents
de somme nulle, est positive. La fonction linéaire
\(F(d)=\sum_jd_ju_j\) a pour gradient relativement à cette métrique
précisément le champ de \eqref{eq:replicator}, car
\[
 g_d\bigl(d_j(u_j-\overline u),\xi\bigr)
 =\sum_j(u_j-\overline u)\xi_j=\sum_ju_j\xi_j=dF(\xi).
\]
Par conséquent, \(dF/ds=\psi''(s)>0\) le long du flot. Le
caractère ascendant correspond à cette fonction et à cette métrique précises.

\begin{theorem}[Invariance variationnelle par subdivision héritée]
Pour chaque intervalle \(j\), soient \(r_{j,c}>0\), \(\sum_c r_{j,c}=1\), les poids relatifs d'une subdivision fixe. Attribuons aux successeurs
\[
 d_{j,c}=r_{j,c}d_j,\qquad u_{j,c}=u_j.
\]
La fonction de partition, son logarithme, toutes ses dérivées et la transformée de Legendre coïncident exactement avec ceux de la partition initiale. De plus, le flot commute avec la somme des successeurs :
\[
 \sum_c d_{j,c}(s)=d_j(s),\qquad
 d_{j,c}(s)=r_{j,c}d_j(s).
\]
\end{theorem}
\begin{proof}
Par regroupement de la somme finie,
\[
 Z_{\rm ref}(s)=\sum_{j,c}r_{j,c}d_j e^{su_j}
 =\sum_j d_j e^{su_j}=Z(s).
\]
Toutes les conclusions se déduisent de cette identité et de la division de chaque poids par le même \(Z(s)\). Le même argument vaut après tout nombre fini de subdivisions. Pour une subdivision dénombrable, la somme de termes positifs converge vers la même valeur par regroupement selon les prédécesseurs.
\end{proof}

L'hypothèse de direction héritée a un contenu vérifiable. Si un intervalle de direction \(a\) est divisé en deux moitiés de directions \(a+b\) et \(a-b\), sa contribution passe de \(d e^{sa}\) à \(d e^{sa}\cosh(sb)\). Pour \(b\ne0\), l'égalité des fonctions de partition échoue bien que le poids total en \(s=0\) soit le même. Ainsi, la conservation géométrique et variationnelle exige de conserver également l'information directionnelle pertinente. L'identité démontrée s'applique à la subdivision héritée ; son transport à une mise à jour TPK qui change ces directions se décide en composant le lecteur de cette mise à jour.

Il existe une formulation hamiltonienne régulière de la coordonnée duale. Dans le
plan cotangent de coordonnées \((q,p)\), prenons
\(H(q,p)=\psi(p)\). Les équations sont
\(\dot q=\psi'(p)\), \(\dot p=0\). Pour des vitesses
\(\dot q\in(u_-,u_+)\), la transformation de Legendre donne
\[
 L(q,\dot q)=\psi^*(\dot q),\qquad
 \frac{\partial^2L}{\partial\dot q^2}>0.
\]
Ce système réalise variationnellement le potentiel du flot des intervalles. Son équation \(\dot p=0\) montre que le temps hamiltonien est distinct du paramètre \(s\) parcourant la collection \eqref{eq:convex-partition}. L'évaluation de \(H\) et \(L\) est conservée sous la subdivision du théorème, puisque \(\psi\) et \(\psi^*\) demeurent identiques.

Par ailleurs, tout flot \(\dot d=F(d)\) admet, dans une coordinatisation locale du simplexe, un relèvement cotangent avec \(\mathscr H(d,p)=p\cdot F(d)\). Cette seconde construction reproduit le flot \eqref{eq:replicator} lui-même sur sa base, mais son hessien par rapport à \(p\) est nul. Les deux réalisations répondent à des questions distinctes : l'une produit une dualité convexe régulière et l'autre transporte une dynamique de base donnée. Cette distinction préserve la signification de l'hamiltonien et du lagrangien lors de la comparaison de cette géométrie avec les lecteurs physiques du corpus.

% Desarrollo acotado: lector energético intervalar del registro K.
% Estatuto: formalización añadida sobre K y sobre las identidades de Schur
% ya presentes en el corpus. No se atribuye novedad histórica a Schur.
\section{Formes de Dirichlet et équivalence variationnelle du raffinement}
\label{sec:energia-refinamiento-k}

La représentation positive du registre dodécaphasique définit une partition ordonnée au moyen de
\[
 K=(234,543,140,729,659,824,621,58,914,794,146,601),
 \qquad d_j=K_j/6263,
 \qquad t_i=\sum_{j=1}^id_j.
\]
Le registre est reçu de la composition
\(x_{\rm term}\mapsto\mathcal R_{12}\mapsto
(b^{90},b^{120},Q_{\rm TPK})\mapsto U_{\rm sgn}\mapsto K\),
postérieure à APP--TRIT--TPK et à son état enrichi, conformément aux lecteurs
reconstruits dans \cite{hmtX}. Cette section applique
un lecteur quadratique à la partition résultante. La comparaison entre
profondeurs s’effectuera par restrictions aux extrémités héritées, de sorte
que chaque coefficient de l’énergie conserve son origine dans une séparation
positive déjà produite.

Le corpus contient une énergie pondérée de différences pour le graphe des
deux mémoires et une réduction de Schur des formes complètes. On
réunit ici ces opérations sur le réseau d’intervalles du registre. La
spécialisation aux conductances \(1/d_j\), l’opérateur d’interpolation
et sa compatibilité avec la déformation des séparations constituent
la formalisation développée dans cette section.

\subsection{Reconstruction des séparations à partir de la matrice de rigidité}

Soit \(d=(d_1,\ldots,d_m)\), avec \(d_j>0\), et soient
\(f=(f_0,\ldots,f_m)\in\mathbb C^{m+1}\) les valeurs d’un observable
aux extrémités de la partition. Nous définissons
\begin{equation}
 \mathcal E_d(f)=\sum_{j=1}^{m}\frac{|f_j-f_{j-1}|^2}{d_j},
 \qquad L_d=D^*\operatorname{diag}(d_1^{-1},\ldots,d_m^{-1})D,
 \label{eq:energia-intervalar}
\end{equation}
où \(D:\mathbb C^{m+1}\to\mathbb C^m\) est l'opérateur d'incidence orientée, \((Df)_j=f_j-f_{j-1}\), et l'astérisque désigne l'adjonction hermitienne. Ainsi \(\mathcal E_d(f)=f^*L_df\). La matrice de rigidité a pour entrées
\[
 (L_d)_{00}=d_1^{-1},\quad (L_d)_{mm}=d_m^{-1},\quad
 (L_d)_{ii}=d_i^{-1}+d_{i+1}^{-1}\quad(1\le i<m),
\]
\[
 (L_d)_{i-1,i}=(L_d)_{i,i-1}=-d_i^{-1},
\]
et les autres sont nuls. Elle est hermitienne positive semi-définie, son noyau
est constitué des fonctions constantes et sa forme est définie positive
sur le quotient par ce noyau. La matrice complète retrouve les
séparations par
\begin{equation}
 d_j=-\frac1{(L_d)_{j-1,j}}.
 \label{eq:recuperacion-d-rigidez}
\end{equation}

\begin{proposition}[Équivalence entre séparations et lecteur nodal complet]
\label{prop:energia-bidireccional}
L'application \(d\mapsto L_d\) est une bijection entre \((0,\infty)^m\) et l'ensemble des matrices réelles symétriques tridiagonales dont les entrées immédiatement adjacentes à la diagonale sont négatives et dont la somme de chaque ligne est nulle. Son inverse est \eqref{eq:recuperacion-d-rigidez}. La normalisation \(\sum_jd_j=1\) équivaut à \(\sum_{j=1}^m-1/(L_d)_{j-1,j}=1\).
\end{proposition}
\begin{proof}
Les formules explicites de \(L_d\) vérifient toutes les conditions.
Réciproquement, les entrées adjacentes négatives déterminent des séparations
positives uniques par \eqref{eq:recuperacion-d-rigidez}. La somme nulle de
chaque ligne détermine ses entrées diagonales comme la somme des
conductances incidentes. La tridiagonalité fixe à zéro les entrées
restantes. On récupère exactement la matrice de
\eqref{eq:energia-intervalar} ; sa semi-définie positivité découle alors
de la factorisation \(D^*\operatorname{diag}(1/d_j)D\).
\end{proof}

La dérivée discrète pondérée \(J_j=(f_j-f_{j-1})/d_j\) est le flux orienté de l'arête \(j\). Aux sommets intérieurs, \((L_df)_i=J_i-J_{i+1}\) ; aux extrémités, \((L_df)_0=-J_1\) et \((L_df)_m=J_m\). L'équation intérieure \(L_df=0\) exprime que le flux a la même valeur sur chaque arête. Cette définition est celle d'un observable quadratique du lecteur d'intervalles ; sa réalisation comme grandeur physique requiert la coordinatisation dimensionnelle pertinente.

\subsection{Interpolation harmonique, complément de Schur et reconstruction du détail}

Raffinons chaque séparation sous la forme \(d_j=\sum_{c=1}^{r_j}\delta_{j,c}\), avec \(\delta_{j,c}>0\). Soit \(\widetilde d\) la concaténation de ces séparations, \(R\) la restriction aux anciennes extrémités et \(H\) l'interpolation définie par
\begin{equation}
 (Hf)_{j,c}=(1-\theta_{j,c})f_{j-1}+\theta_{j,c}f_j,
 \qquad
 \theta_{j,c}=\frac{\sum_{a=1}^c\delta_{j,a}}{d_j}.
 \label{eq:interpolacion-armonica-k}
\end{equation}
Les extrémités partagées ne sont comptées qu’une fois. On a \(RH=I\).

\begin{theorem}[Raffinement isométrique de l’énergie]
\label{thm:energia-schur-refinamiento}
Pour tout vecteur de valeurs anciennes \(f\),
\begin{equation}
 H^*L_{\widetilde d}H=L_d,
 \qquad L_{\widetilde d}H=R^*L_d,
 \qquad
 \min_{Rg=f}\mathcal E_{\widetilde d}(g)=\mathcal E_d(f).
 \label{eq:isometria-energia-k}
\end{equation}
Le minimiseur est unique et vaut \(Hf\). Tout \(g\) admet la
décomposition unique
\begin{equation}
 g=H(Rg)+z,\qquad Rz=0,
 \qquad
 \mathcal E_{\widetilde d}(g)
 =\mathcal E_d(Rg)+\mathcal E_{\widetilde d}(z).
 \label{eq:detalle-energia-k}
\end{equation}
\end{theorem}
\begin{proof}
Sur une ancienne arête, écrivons \(h_c=g_{j,c}-g_{j,c-1}\), de sorte que \(\sum_ch_c=f_j-f_{j-1}=a\). Si \(h_c=\delta_{j,c}a/d_j+r_c\), alors \(\sum_cr_c=0\) et
\[
 \sum_c\frac{|h_c|^2}{\delta_{j,c}}
 =\frac{|a|^2}{d_j}+\sum_c\frac{|r_c|^2}{\delta_{j,c}}.
\]
L’égalité s’obtient en développant ; le terme croisé est proportionnel à
\(\sum_cr_c\) et s’annule. Le minimum est atteint exactement lorsque
tous les \(r_c\) sont nuls, condition qui détermine
\eqref{eq:interpolacion-armonica-k}. En sommant les arêtes, on obtient
le minimum et la décomposition énergétique. Le flux de \(Hf\) est
constant à l’intérieur de chaque ancienne arête, de sorte que ses composantes
intérieures de \(L_{\widetilde d}Hf\) sont nulles et ses composantes héritées
coïncident avec \(L_df\). Cela prouve la seconde identité matricielle ; en
multipliant par \(H^*\), en utilisant \(RH=I\), on obtient la première.
\end{proof}

Le terme \(\mathcal E_{\widetilde d}(z)\) est non négatif et s’annule
exactement pour \(z=0\) : un détail d’énergie nulle serait constant et,
en s’annulant aux extrémités héritées, il s’annulerait sur tout le réseau.
La paire \((Rg,z)\) conserve le vecteur raffiné complet et permet
de le retrouver par \(g=H(Rg)+z\). La minimisation sélectionne sa
composante harmonique ; l’archivage du détail \(z\) maintient
l’information que la seule forme réduite ne distingue pas. La conservation
de l’énergie observée et la conservation de l’état raffiné sont
ainsi exprimées par des applications différentes et explicites.

En ordonnant les sommets comme extrémités héritées \(E\) et sommets intérieurs \(I\), nous écrivons
\[
 L_{\widetilde d}=\begin{pmatrix}A&B^*\\B&C\end{pmatrix}.
\]
Pour un raffinement comportant des sommets intérieurs, le bloc \(C\) est défini positif : une fonction nulle aux extrémités dont l'énergie est nulle est constante sur chaque sous-intervalle et, par conséquent, identiquement nulle. On obtient
\begin{equation}
 H=\begin{pmatrix}I\\-C^{-1}B\end{pmatrix},\qquad
 \operatorname{Schur}_{I}(L_{\widetilde d})
 =A-B^*C^{-1}B=L_d.
 \label{eq:schur-intervalar-k}
\end{equation}
Si le raffinement conserve tous les sommets et n'en introduit aucun, la même affirmation se réduit à \(H=I\) et \(L_{\widetilde d}=L_d\).

\begin{proposition}[Naturalité des prolongements et de l'élimination]
\label{prop:energia-naturalidad-sucesiva}
Pour trois partitions emboîtées, les interpolations harmoniques et les compléments de Schur satisfont
\begin{equation}
 H_{n+2\leftarrow n}
 =H_{n+2\leftarrow n+1}H_{n+1\leftarrow n},
 \qquad
 \operatorname{Schur}_{n+1\to n}
 \bigl(\operatorname{Schur}_{n+2\to n+1}(L_{n+2})\bigr)=L_n.
 \label{eq:energia-naturalidad-sucesiva}
\end{equation}
\end{proposition}
\begin{proof}
L’interpolation de valeurs affines sur un intervalle demeure la
même fonction affine après toute subdivision de cet intervalle.
Cela prouve la première identité. Minimiser d’abord sur les sommets
incorporés au dernier niveau puis sur ceux du niveau intermédiaire
revient à minimiser sur leur totalité, avec les mêmes extrémités héritées fixées.
L’identité variationnelle de \eqref{eq:isometria-energia-k} et l’unicité
du minimiseur prouvent la seconde identité. La restriction du minimum
se réalise donc au moyen de la même composition que le prolongement
harmonique des données.
\end{proof}

\subsection{Réponse de Dirichlet à Neumann et défaut énergétique d’une extension}

Si seules les extrémités initiale et finale sont conservées, soit \(D_0=\sum_{j=1}^md_j\). Pour des valeurs \((a,b)\), la solution harmonique est
\[
 f_i=a+(b-a)\frac{\sum_{j=1}^id_j}{D_0},\qquad
 J_j=(b-a)/D_0.
\]
Son opérateur de Dirichlet à Neumann, qui transforme les valeurs de frontière en
flux sortants, est
\begin{equation}
 \Lambda_d=\frac1{D_0}
 \begin{pmatrix}1&-1\\-1&1\end{pmatrix},\qquad
 \min\mathcal E_d=\frac{|b-a|^2}{D_0}.
 \label{eq:dtn-intervalar-k}
\end{equation}
Lorsque la valeur est prescrite à tous les sommets hérités, leur opérateur de réponse multiponctuelle est, par le même argument, \(\Lambda=L_d\).

Une extension approchée \(g=(f,y)\), où \(y\) réunit les valeurs
aux sommets intérieurs, a pour résidu intérieur
\(r=Bf+Cy\). L’identité exacte
\begin{equation}
 \mathcal E_{\widetilde d}(g)-f^*\Lambda f
 =r^*C^{-1}r\ge0
 \label{eq:residual-dtn-k}
\end{equation}
se déduit de \(y+C^{-1}Bf=C^{-1}r\) et de la décomposition
\[
 g^*L_{\widetilde d}g
 =f^*(A-B^*C^{-1}B)f
 +(y+C^{-1}Bf)^*C(y+C^{-1}Bf).
\]
Ainsi, la valeur effective et le détail maintiennent une décomposition
hermitienne complète. En particulier, si une
extension linéaire est \(\widetilde H=\begin{pmatrix}I\\Z\end{pmatrix}\), sa matrice de résidus est \(\mathscr R=B+CZ\) et
\[
 \widetilde H^*L_{\widetilde d}\widetilde H-\Lambda
 =\mathscr R^*C^{-1}\mathscr R\succeq0.
\]
Le résidu nul caractérise l'interpolation harmonique exacte. Ce résidu énergétique se situe dans l'espace des sommets intérieurs ; l'annulation des résidus des formes méromorphes appartient à une autre application de frontière et conserve sa propre définition.

\subsection{Commutation du raffinement et du flux des séparations}

Soit \(u=P_3P_{11}K\) la direction algébrique déjà retrouvée à partir des
lecteurs du registre \cite{hmtX}. Pour le paramètre réel \(s\), définissons
\[
 d_j(s)=\frac{K_je^{su_j}}{\sum_\ell K_\ell e^{su_\ell}}.
\]
Étant fixées des fractions positives \(\vartheta_{j,c}\) de somme un pour chaque
\(j\), le raffinement hérité est
\[
 \delta_{j,c}(s)=\vartheta_{j,c}d_j(s),\qquad
 u_{j,c}=u_j.
\]
Produire d'abord les poids raffinés \(K_j\vartheta_{j,c}\) et les déformer donne la même collection : les successeurs ont une exponentielle identique et leurs coefficients ont pour somme \(K_j\). Pour tout \(s\),
\begin{equation}
 \operatorname{Schur}_{I}(L_{\widetilde d(s)})=L_{d(s)},
 \qquad H^*L_{\widetilde d(s)}H=L_{d(s)}.
 \label{eq:schur-flujo-k}
\end{equation}
Dans cette comparaison parent--enfants, \(H\) est constant : ses coefficients
sont les sommes partielles de \(\vartheta_{j,c}\). En revanche,
l’interpolation à partir des deux extrémités globales utilise
\(t_i(s)=\sum_{j\le i}d_j(s)\) et change avec le paramètre. Les deux
affirmations correspondent à des restrictions différentes du même réseau.

La normalisation \(\sum_jd_j(s)=1\) fait que l'opérateur global à deux extrémités reste égal à \(\left(\begin{smallmatrix}1&-1\\-1&1\end{smallmatrix}\right)\). La matrice de tous les nœuds \(L_{d(s)}\) conserve, elle, la déformation et la récupère au moyen de \eqref{eq:recuperacion-d-rigidez}. La perte d'information dans le lecteur à deux extrémités est identifiée exactement : il retient la longueur totale et ses flux, tandis que le lecteur nodal retient toutes les séparations. L'archive des détails fournit le complément nécessaire pour reconstruire l'observable raffiné.

\section{Reconstruction analytique de l’énergie et mémoire du raffinement}
\label{sec:limite-energetico-k}

Les séparations positives construites à partir du registre dodécaphasique déterminent une partition ordonnée et, d'après la section précédente, une forme de Dirichlet sur ses valeurs nodales. Son raffinement possède deux opérations complémentaires : l'interpolation harmonique conserve l'énergie des données héritées, tandis que chaque détail enregistre la variation incorporée entre ces données. Le passage à une profondeur illimitée consiste à réunir cette collection compatible au moyen d'une norme contrôlant la somme de ses détails. Le résultat concerne la réalisation analytique unidimensionnelle de l'énergie ; il conserve comme antécédent la construction HMT du registre et de la structure discrète conjointe du continu.

\subsection{Partitions emboîtées et espaces d’interpolation}

Fixons la partition initiale
\(P_0=\{0=t_0<t_1<\cdots<t_{12}=1\}\), dont les longueurs sont
\(d_j=t_j-t_{j-1}=K_j/\sum_iK_i>0\). Soit
\((P_n)_{n\geq0}\) une suite de partitions finies emboîtées de
\([0,1]\), obtenues par subdivision des intervalles précédents.
Nous écrirons
\[
 |P_n|=\max\{b-a:[a,b]\text{ est un intervalle de }P_n\},
 \qquad \lim_{n\to\infty}|P_n|=0.
\]
La subdivision répétée de chaque intervalle en un nombre fixe de successeurs égaux, au moins deux, satisfait cette condition. Pour d'autres règles, la contraction du maillage est une hypothèse explicite de ce lecteur. Les étiquettes des successeurs et leurs données de retenue accompagnent la subdivision et conservent la provenance du raffinement.

Soit \(V_n\) l’espace complexe des fonctions continues sur \([0,1]\)
qui sont affines sur chaque intervalle de \(P_n\). L’énergie et la norme
énergétique sont définies par
\begin{equation}
 \mathcal E_n(f)
 =\sum_{[a,b]\in P_n}\frac{|f(b)-f(a)|^2}{b-a}
 =\int_0^1|f'(x)|^2\,dx,
 \qquad
 \|f\|_{\mathcal E}^2=|f(0)|^2+\mathcal E_n(f).
 \label{eq:energia-afines-k}
\end{equation}
L’égalité intégrale résulte du fait que la dérivée d’une fonction affine est
constante sur chaque intervalle. L’énergie coïncide ainsi exactement avec
la forme nodale déjà construite. L’inclusion \(V_n\subset V_{n+1}\)
interprète la même fonction affine dans une partition plus fine ; elle conserve
aussi bien \(f(0)\) que l’énergie et est isométrique pour
\(\|\cdot\|_{\mathcal E}\).

La valeur initiale a une fonction précise : elle fixe la composante constante
que la forme de Dirichlet identifie dans son noyau. L’énergie contrôle
les différences, et le couple formé par la valeur initiale et ces différences
détermine l’observable. Dans l’espace de Sobolev
\(H^1([0,1];\mathbb C)\), dont les éléments ont un représentant
absolument continu de dérivée faible dans \(L^2\), cette norme est
équivalente à la norme usuelle. En effet,
\(f(x)=f(0)+\int_0^xf'(t)\,dt\) contrôle la norme de \(f\) à
partir de \(|f(0)|\) et de \(\|f'\|_{L^2}\) ; l’intégration de la
même identité contrôle à son tour \(|f(0)|\) au moyen de
\(\|f\|_{L^2}+\|f'\|_{L^2}\).

\subsection{Reconstruction d’une fonction à partir de ses valeurs compatibles}

\begin{theorem}[Reconstruction énergétique et complétion des interpolations]
\label{thm:limite-energia-h1-k}
Soit \(f_n\in V_n\) une collection compatible par restriction aux nœuds hérités :
\[
 f_{n+1}(a)=f_n(a)
 \quad\text{pour tout nœud }a\text{ de }P_n.
\]
Si \(\sup_n\mathcal E_n(f_n)<\infty\), il existe une unique fonction
\(f\in H^1([0,1];\mathbb C)\), considérée par son représentant
continu, dont les valeurs en tous ces nœuds sont celles prescrites. De plus,
\(f_n\to f\) uniformément et dans \(H^1\), et
\begin{equation}
 \int_0^1|f'(x)|^2\,dx
 =\lim_{n\to\infty}\mathcal E_n(f_n).
 \label{eq:limite-energia-h1-k}
\end{equation}
Réciproquement, l'interpolation nodale de tout élément de \(H^1\) produit une collection de cette classe. Par conséquent, la complétion de \(\bigcup_n V_n\) pour la norme énergétique est \(H^1([0,1];\mathbb C)\).
\end{theorem}

\begin{proof}
Posons \(g_n=f_n'\in L^2([0,1];\mathbb C)\). Pour
\(m\geq n\), la compatibilité des valeurs aux extrémités de chaque
intervalle \([a,b]\) de \(P_n\) fournit
\[
 \frac1{b-a}\int_a^bg_m(x)\,dx
 =\frac{f_m(b)-f_m(a)}{b-a}
 =\frac{f_n(b)-f_n(a)}{b-a}
 =g_n\big|_{(a,b)}.
\]
Soit \(W_n\) l’espace des fonctions de \(L^2\) constantes sur chaque
intervalle de \(P_n\), et soit \(Q_n:L^2\to W_n\) la projection
orthogonale donnée par les moyennes sur ces intervalles. L’identité
précédente affirme \(Q_ng_m=g_n\). Par conséquent,
\(g_m-g_n\) est orthogonal à \(g_n\) et
\begin{equation}
 \|g_m-g_n\|_{L^2}^2
 =\mathcal E_m(f_m)-\mathcal E_n(f_n).
 \label{eq:pitagoras-gradientes-k}
\end{equation}
Les énergies sont croissantes et bornées. L'identité de Pythagore prouve que \((g_n)\) est une suite de Cauchy dans \(L^2\), de limite \(g\). Définissons
\[
 f(x)=f_0(0)+\int_0^xg(t)\,dt.
\]
Cette fonction appartient à \(H^1\), a pour dérivée faible \(g\) et
vérifie, par l’inégalité de Cauchy--Schwarz,
\begin{equation}
 \sup_{x\in[0,1]}|f_n(x)-f(x)|
 \leq\|g_n-g\|_{L^2}\longrightarrow0.
 \label{eq:convergencia-uniforme-energia-k}
\end{equation}
La convergence uniforme conserve chaque valeur nodale prescrite ; jointe à
la convergence des dérivées dans \(L^2\), elle donne la convergence dans
\(H^1\). La convergence des normes des dérivées prouve
\eqref{eq:limite-energia-h1-k}. L’union des nœuds est dense car
\(|P_n|\to0\) ; deux représentants continus ayant les mêmes valeurs
sur cette union coïncident. On obtient l’unicité.

Pour la réciproque, prenons \(f\in H^1\) dans son représentant continu et soit \(f_n\) son interpolation aux nœuds de \(P_n\). Sa dérivée est \(g_n=Q_nf'\). Les projections \(Q_n\) sont contractantes dans \(L^2\). Pour une fonction continue \(v\), l'erreur \(|Q_nv-v|\) est bornée par son module de continuité évalué en \(|P_n|\), qui tend vers zéro. La densité des fonctions continues dans \(L^2\) et la contractivité impliquent \(Q_ng\to g\) pour tout \(g\in L^2\) : étant donné un approximant continu \(v\), on utilise
\[
 \|Q_ng-g\|_{L^2}
 \leq 2\|g-v\|_{L^2}+\|Q_nv-v\|_{L^2}.
\]
En appliquant cette convergence à \(g=f'\), on obtient la convergence
des dérivées et, au moyen de \eqref{eq:convergencia-uniforme-energia-k},
celle des fonctions. La contractivité borne leurs énergies par
\(\|f'\|_{L^2}^2\), et la convergence des normes prouve de
nouveau l’identité de la limite. Toute fonction de \(H^1\) est donc
limite énergétique d’éléments de \(\bigcup_nV_n\).
L’équivalence des normes et la complétude de \(H^1\) identifient
la complétion énoncée.
\end{proof}

L’argument décrit concrètement ce qui est conservé lors du passage du réseau
à la fonction. Chaque dérivée finie enregistre la moyenne de toutes les
dérivées ultérieures dans son intervalle ; lors du raffinement, on incorpore des
composantes orthogonales qui maintiennent cette moyenne. L’énergie bornée
contrôle la somme totale de ces incorporations. La valeur continue se
retrouve alors en intégrant la dérivée limite et en conservant la valeur
initiale. La réalisation analytique réunit, par cette voie explicite,
l’information compatible des niveaux finis.

\subsection{Décomposition orthogonale de la mémoire par échelles}

Soit \(H_n\) l’interpolation harmonique des valeurs de \(P_n\)
sur les nœuds de \(P_{n+1}\), interprétée comme fonction de
\(V_{n+1}\), et définissons le détail
\[
 z_{n+1}=f_{n+1}-H_nf_n.
\]
D'après \eqref{eq:detalle-energia-k}, le détail s'annule aux nœuds hérités. Sa dérivée est de moyenne nulle sur chaque ancien intervalle et est orthogonale à \(W_n\). Pour des niveaux distincts, la dérivée d'un détail antérieur appartient à cet espace de fonctions constantes ; les dérivées des détails sont donc orthogonales entre elles. La même orthogonalité vaut par rapport à \(f_0'\).

\begin{proposition}[Identité d’énergie et récupération des détails]
\label{prop:memoria-energetica-escalas-k}
Pour toute profondeur finie \(N\),
\begin{equation}
 \mathcal E_N(f_N)
 =\mathcal E_0(f_0)
 +\sum_{n=0}^{N-1}\mathcal E_{n+1}(z_{n+1}).
 \label{eq:energia-telescopica-detalles-k}
\end{equation}
Si la collection satisfait les hypothèses du théorème~\ref{thm:limite-energia-h1-k}, la série du membre droit converge et sa somme est \(\int_0^1|f'|^2\). Les données \(f_0\) et tous les détails déterminent la collection compatible complète et sa réalisation analytique.
\end{proposition}
\begin{proof}
L’identité finie résulte de la somme de
\(\mathcal E_{n+1}(f_{n+1})
=\mathcal E_n(f_n)+\mathcal E_{n+1}(z_{n+1})\), qui est la
décomposition orthogonale de la section précédente. Les termes sont
non négatifs et les sommes partielles convergent par
\eqref{eq:limite-energia-h1-k}. La reconstruction finie s’effectue
récursivement au moyen de \(f_{n+1}=H_nf_n+z_{n+1}\) ; le théorème
précédent détermine ensuite la fonction limite. Réciproquement, une
fonction de \(H^1\) détermine ses interpolations et, par différence,
tous ces détails. Chaque flèche conserve ainsi ses données propres.
\end{proof}

L'élimination variationnelle des nœuds intérieurs sélectionne le détail nul au niveau éliminé. La conservation des détails permet de récupérer toute configuration compatible d'énergie finie. La réponse effective et la configuration complète possèdent ainsi une relation précise : la première enregistre le minimum sur les variables intérieures ; la seconde ajoute les composantes orthogonales et leur énergie.

\subsection{Transport de la réalisation analytique sous le flot positif}

La direction \(u=P_3P_{11}K\), récupérée avant l'évaluation énergétique \cite{hmtX}, détermine la collection de séparations de \eqref{eq:schur-flujo-k}. Son paramètre \(s\in\mathbb R\) décrit cette déformation analytique. Écrivons
\[
 d_j(s)=\frac{K_je^{su_j}}{\sum_\ell K_\ell e^{su_\ell}},
 \qquad
 t_i(s)=\sum_{j=1}^id_j(s),
 \qquad
 a_j(s)=\frac{d_j(s)}{d_j(0)}>0.
\]
Soit \(F_s:[0,1]\to[0,1]\) l’application qui envoie chaque
\(t_i(0)\) sur \(t_i(s)\) et qui est affine sur les intervalles initiaux.
C’est un homéomorphisme croissant, de pente \(a_j(s)\) sur
l’intervalle \([t_{j-1}(0),t_j(0)]\). Son inverse est également affine
par intervalles et les deux applications sont lipschitziennes.

Supposons que chaque subdivision hérite des fractions de longueur et de la direction de son intervalle prédécesseur, comme dans la section précédente. Tous les nœuds raffinés sont alors transportés par le même \(F_s\). Notons leurs partitions \(P_n(s)=F_s(P_n(0))\) et leurs espaces d'interpolation \(V_n(s)\). L'application
\[
 U_s f=f\circ F_s^{-1}
\]
transporte \(V_n(0)\) vers \(V_n(s)\), préserve les valeurs nodales correspondantes et prolonge cette action à \(H^1\).

\begin{theorem}[Équivalence énergétique et commutation avec le raffinement]
\label{thm:transporte-h1-flujo-k}
Pour tout \(f\in H^1([0,1];\mathbb C)\),
\begin{equation}
 \int_0^1|(U_sf)'(y)|^2\,dy
 =\sum_{j=1}^{12}\frac1{a_j(s)}
 \int_{t_{j-1}(0)}^{t_j(0)}|f'(x)|^2\,dx.
 \label{eq:transporte-energia-h1-k}
\end{equation}
Pour \(s\) dans tout intervalle compact, les normes énergétiques avant et après le transport sont uniformément équivalentes et \(|P_n(s)|\to0\) uniformément sur ce compact. L'interpolation harmonique et la restriction aux nœuds hérités commutent avec \(U_s\) ; la reconstruction et la minimisation des détails se transportent aussi bien à profondeur finie que dans la réalisation \(H^1\).
\end{theorem}

\begin{proof}
Sur l’intervalle initial \(j\), la règle de dérivation en chaîne et le changement
de variable \(y=F_s(x)\) fournissent
\((U_sf)'(F_s(x))=f'(x)/a_j(s)\) et \(dy=a_j(s)\,dx\).
L’intégration et la somme sur les intervalles prouvent
\eqref{eq:transporte-energia-h1-k}. La formule vaut pour les
représentants absolument continus, et son membre droit est fini.

Un intervalle compact de paramètres étant fixé, la continuité et la positivité
des douze facteurs donnent des constantes \(0<a_-\leq a_+<\infty\)
telles que \(a_-\leq a_j(s)\leq a_+\) pour tout \(j\) et
tout \(s\) de ce compact. Par conséquent,
\[
 \frac1{a_+}\mathcal E(f)
 \leq\mathcal E(U_sf)
 \leq\frac1{a_-}\mathcal E(f),
 \qquad
 |P_n(s)|\leq a_+|P_n(0)|.
\]
Comme \((U_sf)(0)=f(0)\), on obtient en outre
\[
 \min\{1,a_+^{-1}\}\|f\|_{\mathcal E}^2
 \leq\|U_sf\|_{\mathcal E}^2
 \leq\max\{1,a_-^{-1}\}\|f\|_{\mathcal E}^2.
\]
Ces bornes prouvent l’équivalence uniforme et la conservation de la
condition de maillage décroissant.

Sur chaque arête prédécesseur, les coefficients d'interpolation sont les sommes partielles des fractions héritées, indépendantes de \(s\). Si \(H_n(s)\) désigne son interpolation, on a \(U_sH_n(0)=H_n(s)U_s\) sur les espaces correspondants. La restriction commute parce que les nœuds sont transportés avec leurs valeurs. La décomposition en fonction harmonique et détail est conservée, et l'identité de Schur de \eqref{eq:schur-flujo-k} préserve sa caractérisation variationnelle. Enfin, la continuité bornée de \(U_s\) permet de passer à la limite \(H^1\). En particulier, une collection reconstruite avant le transport converge ensuite vers \(U_sf\), avec les mêmes valeurs nodales transportées.
\end{proof}

La loi \eqref{eq:transporte-energia-h1-k} identifie les facteurs métriques exacts de la déformation : le changement de longueur de chaque intervalle modifie inversement sa contribution énergétique. L'équivalence des normes conserve les configurations d'énergie finie, et l'archive des détails conserve leur reconstruction. On obtient ainsi une réalisation analytique complète de ce lecteur du registre positif. Sa compatibilité avec les formes canoniques s'articule sur les subdivisions d'un domaine fixe : la somme des formes conserve l'objet méromorphe, tandis que la réduction variationnelle et les détails conservent la réponse énergétique et l'observable raffiné. L'identification avec des formes amplituédriques de rang supérieur conserve les conditions spécifiques de couverture, de degré et de correspondance des résidus établies pour cette réalisation.

\section{Transports nonadiques et conservation de la mémoire}
\label{sec:nonadic}

La reconstruction de la section~\ref{sec:memory} fournit un lien effectif entre la mémoire du transport et la géométrie du registre dodécaphasique. Le registre engendré par APP--TRIT--TPK détermine ses lectures d'incidence ; les deux compressions conservent, dans des composantes séparées, la représentation terminale et ses compléments ; la réunion des mémoires avec la charge restitue les douze coordonnées. La direction géométrique et le projecteur dimensionnel sont obtenus après cette restitution. Le développement suivant établit le bilan opératoriel soutenant cette chaîne et son comportement sous raffinement, réalisation spectrale et changement de métrique \cite{HMT-VIII,hmtX}.

\subsection{Calendrier à neuf phases et décomposition orthogonale}

Soit \(\mathcal H\) un espace de Hilbert complexe et soit
\(C\in\mathcal B(\mathcal H)\). La somme orthogonale de neuf copies,
\(\mathcal H^{\oplus9}\), conserve séparément les neuf positions
du calendrier. Nous définissons
\begin{equation}
 \begin{split}
 S_C(v_0,\ldots,v_8)&=(Cv_8,v_0,\ldots,v_7),\\
 J:\mathcal H\longrightarrow\mathcal H^{\oplus9},
 \qquad Jv&=\frac13(v,\ldots,v).
 \end{split}
 \label{nt:calendar}
\end{equation}
La normalisation \(1/3=9^{-1/2}\) donne \(J^*J=I\). L'opérateur \(P=JJ^*\) est la projection orthogonale sur la section uniforme ; son action remplace les neuf composantes par leur moyenne. Chaque composante traverse une fois la jonction qui contient \(C\) pendant un tour, de sorte que
\begin{equation}
 S_C^9=I_9\otimes C.
 \label{nt:full-return}
\end{equation}
Le retour de la position du calendrier réalise ainsi une action de
\(C\) sur la fibre de mémoire.

\begin{definition}
La compression uniforme du transport et son complément sont
\begin{equation}
 T_C=J^*S_CJ=\frac{8I+C}{9},
 \qquad
 \eta_C=(I-P)S_CJ.
 \label{nt:compression}
\end{equation}
La première application a pour codomaine \(\mathcal H\) ; la seconde
a pour codomaine \(J(\mathcal H)^\perp\subset\mathcal H^{\oplus9}\).
\end{definition}

Pour \(z=(C-I)v\), le complément s'écrit
\begin{equation}
 \eta_Cv=\frac1{27}(8z,-z,-z,-z,-z,-z,-z,-z,-z).
 \label{nt:complement}
\end{equation}
La somme de ses composantes est nulle. Les coefficients de la
décomposition proviennent donc de la moyenne de huit
continuations identiques et d’une continuation transportée.

\begin{proposition}[Bilan de la compression et de la jonction de mémoire]
Pour tout \(C\in\mathcal B(\mathcal H)\),
\begin{align}
 \eta_C^*\eta_C
   &=\frac8{81}(I-C)^*(I-C),\label{nt:eta-gram}\\
 I-T_C^*T_C
   &=\eta_C^*\eta_C+\frac19(I-C^*C).\label{nt:general-balance}
\end{align}
Si \(C^*C=I\), l'application \(v\mapsto(T_Cv,\eta_Cv)\) est isométrique.
\end{proposition}

\begin{proof}
L’expression~\eqref{nt:complement} donne
\[
 \|\eta_Cv\|^2
 =\frac{8^2+8}{27^2}\|(C-I)v\|^2
 =\frac8{81}\|(C-I)v\|^2.
\]
La polarisation prouve~\eqref{nt:eta-gram}. D’autre part,
\[
 T_C^*T_C
 =\frac{64I+8C+8C^*+C^*C}{81}.
\]
En additionnant \eqref{nt:eta-gram}, on obtient
\[
 T_C^*T_C+\eta_C^*\eta_C=\frac{8I+C^*C}{9}.
\]
Soustraire ce résultat de \(I\) donne~\eqref{nt:general-balance}. Lorsque \(C^*C=I\), la somme des deux produits de Gram est l'identité, ce qui conserve tous les produits scalaires.
\end{proof}

Le terme \(\eta_C^*\eta_C\) mesure la dispersion entre les neuf
composantes. Le terme \((I-C^*C)/9\) enregistre, séparément,
la variation de la norme dans la réunion de mémoire. La réalisation
unitaire situe toute la variation visible dans le premier terme et
conserve la norme dans la somme orthogonale complète.

\begin{proposition}[Conservation d’un observable]
Soient \(C\) unitaire et \(A=A^*\in\mathcal B(\mathcal H)\) tels que
\(C^*AC=A\). Alors
\begin{equation}
 A=T_C^*AT_C+
   \eta_C^*(I_9\otimes A)\eta_C
  =T_C^*AT_C+\frac8{81}(I-C)^*A(I-C).
 \label{nt:observable-balance}
\end{equation}
Pour \(A\succeq0\), les deux termes de la somme sont positifs.
\end{proposition}

\begin{proof}
La formule du complément, appliquée au produit sesquilinéaire
défini par \(A\), prouve la seconde égalité. Le développement des
deux termes donne
\[
 \frac1{81}\bigl(
 64A+8AC+8C^*A+C^*AC
 +8A-8AC-8C^*A+8C^*AC\bigr)
 =\frac{8A+C^*AC}{9}=A.
\]
La positivité est conservée sous une application de la forme \(B\mapsto L^*BL\).
\end{proof}

\subsection{Itération de la compression et conservation du secteur fixe}

On fixe dans cette sous-section un opérateur unitaire \(C\) et l'on écrit \(T=T_C\), \(\eta=\eta_C\). Le transport \(S_C^N\) maintient toutes les phases pendant \(N\) étapes. La puissance \(T^N\), en revanche, applique successivement la compression uniforme ; à chaque étape, \(\eta T^jv\) conserve le complément de la représentation précédente. La distinction est exprimée par les opérateurs suivants.

\begin{proposition}[Registre isométrique à horizon fini]
Pour tout entier \(N\geq1\), l’application
\begin{equation}
 V_Nv=(T^Nv,\eta v,\eta Tv,\ldots,\eta T^{N-1}v)
 \label{nt:finite-record}
\end{equation}
de \(\mathcal H\) dans
\(\mathcal H\oplus(\mathcal H^{\oplus9})^{\oplus N}\)
est isométrique. Si \(A\) satisfait les hypothèses de
\eqref{nt:observable-balance}, alors
\begin{equation}
 A=T^{*N}AT^N+
 \sum_{j=0}^{N-1}T^{*j}\eta^*(I_9\otimes A)\eta T^j.
 \label{nt:telescopy}
\end{equation}
\end{proposition}

\begin{proof}
Multiplier~\eqref{nt:observable-balance} par \(T^{*j}\) et
\(T^j\) donne
\[
 T^{*j}AT^j-T^{*(j+1)}AT^{j+1}
 =T^{*j}\eta^*(I_9\otimes A)\eta T^j.
\]
La somme télescopique prouve~\eqref{nt:telescopy}. Pour \(A=I\),
le membre de droite est le produit de Gram de \(V_N\).
Le passage de \(N\) à \(N+1\) remplace uniquement \(T^Nv\) par
\((T^{N+1}v,\eta T^Nv)\), en conservant toutes les composantes
de mémoire précédemment produites.
\end{proof}

\begin{theorem}[Limite de la compression et mémoire illimitée]
Soit \(P_{\mathrm{fix}}\) la projection orthogonale sur
\(\ker(I-C)\). Alors
\begin{align}
 \operatorname*{s\!-\!lim}_{N\to\infty}T^N
   &=P_{\mathrm{fix}},\label{nt:strong-limit}\\
 I&=P_{\mathrm{fix}}+
 \sum_{j=0}^{\infty}T^{*j}\eta^*\eta T^j.\label{nt:infinite-record}
\end{align}
La série converge fortement. Pour tout observable borné
autoadjoint qui satisfait \(C^*AC=A\), on conserve en outre
\begin{equation}
 A=P_{\mathrm{fix}}AP_{\mathrm{fix}}+
 \sum_{j=0}^{\infty}
 T^{*j}\eta^*(I_9\otimes A)\eta T^j.
 \label{nt:infinite-observable}
\end{equation}
\end{theorem}

\begin{proof}
Le théorème spectral de l’opérateur unitaire réalise \(C\) comme
la fonction coordonnée \(z\) sur le cercle unité.
La compression correspond à \(t(z)=(8+z)/9\), et
\begin{equation}
 1-|t(z)|^2=\frac8{81}|1-z|^2,
 \qquad |z|=1.
 \label{nt:circular-defect}
\end{equation}
Ainsi, \(t(z)^N\) converge vers zéro en chaque \(z\ne1\) et conserve
la valeur un en \(z=1\). La borne \(|t(z)^N|\leq1\) permet
d’appliquer la convergence dominée à chaque mesure spectrale vectorielle.
On obtient~\eqref{nt:strong-limit} et, par le même argument,
\(T^{*N}\to P_{\mathrm{fix}}\) fortement. Le caractère borné de
\(A\), avec \(\|T^N\|\leq1\), donne
\(T^{*N}AT^N\to P_{\mathrm{fix}}AP_{\mathrm{fix}}\).
La limite de~\eqref{nt:telescopy} prouve les deux identités.
\end{proof}

Le registre
\[
 V_\infty v=(P_{\mathrm{fix}}v,\eta v,\eta Tv,\eta T^2v,\ldots)
\]
est, par conséquent, une isométrie. Le secteur fixe conserve la partie qui persiste dans la représentation terminale ; les autres composantes demeurent dans la mémoire des compléments.

Pour le décalage bilatéral
\(C e_n=e_{n+1}\) dans \(\ell^2(\mathbb Z)\), une suite fixe
est constante. La condition de sommabilité des carrés impose son
annulation, de sorte que \(P_{\mathrm{fix}}=0\). Toute la norme
est alors représentée dans les compléments. Le spectre de
ce décalage est le cercle complet, et le maximum
spectral de \(|t(z)|^N\) est un ; d’où
\(\|T^N\|=1\) pour chaque \(N\). La convergence forte conserve
ce comportement non uniforme.

Pour une permutation cyclique de \(d\) positions, le secteur fixe est la droite des vecteurs uniformes. Si la permutation possède plusieurs orbites, chaque orbite apporte sa propre direction uniforme. Cette dimension détermine la partie terminale qui doit être conservée.

\begin{proposition}[Compression d’un transport d’ordre trois]
Soit \(C\) unitaire et \(C^3=I\). Alors
\begin{equation}
 P_{\mathrm{fix}}=\frac{I+C+C^2}{3},
 \qquad
 T_C^*T_C=P_{\mathrm{fix}}+
 \frac{19}{27}(I-P_{\mathrm{fix}}).
 \label{nt:order-three}
\end{equation}
\end{proposition}

\begin{proof}
La moyenne des trois puissances est le projecteur orthogonal
sur les vecteurs fixes. Comme \(C^*=C^2\),
\[
 T_C^*T_C=\frac{65I+8(C+C^2)}{81}
 =\frac{57I+24P_{\mathrm{fix}}}{81},
\]
qui coïncide avec~\eqref{nt:order-three}.
\end{proof}

Dans la réalisation à douze positions de \(\mathbb R^{12}\), \((Sx)_i=x_{i+1}\) avec indices modulo douze, \(S^3\) est d'ordre quatre et \(S^4\) est d'ordre trois. Les quatre orbites de \(S^4\) donnent \(\operatorname{rank}P_{\mathrm{fix}}=4\). La réalisation exceptionnelle \(g_c\) de l'article~X satisfait \(g_c^3=I\) et \(\ker(I-g_c)=0\) dans son propre espace porteur ; dans celle-ci,~\eqref{nt:order-three} se réduit à \(T_{g_c}^*T_{g_c}=19I/27\) \cite{hmtX}. Le coefficient et le secteur fixe sont transportés conjointement avec le domaine.

La contraction radiale \(M=qC\), \(0<q<1\), possède un autre registre
exact :
\begin{equation}
 L_Nv=\bigl(\sqrt{1-q^2}\,v,
 \sqrt{1-q^2}\,qCv,\ldots,
 \sqrt{1-q^2}\,q^{N-1}C^{N-1}v,q^NC^Nv\bigr).
 \label{nt:radial-record}
\end{equation}
Son produit de Gram est l’identité car
\((1-q^2)\sum_{j=0}^{N-1}q^{2j}+q^{2N}=1\).
Le terminal a pour norme \(q^N\|v\|\), et le registre infini
conserve le produit scalaire. Le facteur \(q=5/9\) du mode
différentiel central de l’article~VIII s’applique à un tour
complet ; sa répartition uniforme entre neuf étapes utilise
\(q^{1/9}S_C\), dont la neuvième puissance est
\(q(I_9\otimes C)\) \cite{HMT-VIII}.

\subsection{Réduction du registre et conservation de la trace}

La direction fixe du complément dans les neuf composantes
permet une réalisation minimale en deux termes. Définissons
\[
 D_C=\frac{\sqrt8}{9}(I-C).
\]
Les identités précédentes donnent
\(T_C^*T_C+D_C^*D_C=I\) pour \(C\) unitaire.
Le changement de signe par rapport à l’expression de
\(\eta_C\) correspond à une orientation de la composante
complémentaire et conserve son produit de Gram.

\begin{proposition}[Canal positif de la compression avec mémoire]
Pour \(C\) unitaire, l’application
\begin{equation}
 \mathcal Q_C(X)=T_CXT_C^*+D_CXD_C^*
              =\frac{8X+CXC^*}{9}
 \label{nt:channel}
\end{equation}
est complètement positive sur \(\mathcal B(\mathcal H)\) et conserve l'identité. Sa restriction aux opérateurs à trace conserve la trace.
\end{proposition}

\begin{proof}
En développant les deux termes de la première expression,
les produits \(CX\) et \(XC^*\) s’annulent. Les coefficients
restants sont \(72/81=8/9\) pour \(X\) et \(9/81=1/9\)
pour \(CXC^*\), ce qui prouve~\eqref{nt:channel}.
Pour tout entier \(m\geq1\), l’amplification aux matrices
d’ordre \(m\) est de la forme
\[
 X\longmapsto
 (I_m\otimes T_C)X(I_m\otimes T_C)^*
 +(I_m\otimes D_C)X(I_m\otimes D_C)^*.
\]
Chaque terme conserve la positivité ; l’application est donc
complètement positive. La seconde expression de
\eqref{nt:channel} donne \(\mathcal Q_C(I)=I\).
Si \(X\) est à trace, l’invariance de la trace sous
conjugaison unitaire implique
\(\operatorname{Tr}\mathcal Q_C(X)=\operatorname{Tr}X\).
\end{proof}

L'application \eqref{nt:channel} s'obtient également en préparant l'isométrie \(v\mapsto T_Cv\otimes e_0+D_Cv\otimes e_1\) et en effectuant la trace partielle sur \(\mathbb C^2\). La conservation du vecteur complet appartient à l'isométrie ; la conservation de la trace de l'état réduit appartient à \(\mathcal Q_C\). Les deux affirmations découlent du même bilan orthogonal.

Dans le prolongement nonadique d’algèbres finies apparaît une
conservation complémentaire. Pour
\(\mathcal A_d=M_{9^d}(\mathbb C)\), avec
\(\tau_d(a)=9^{-d}\operatorname{Tr}a\), les applications
\[
 \iota_d(a)=a\otimes I_9,\qquad
 \jmath_{d,j}(a)=a\otimes p_j,
 \quad p_j=|e_j\rangle\langle e_j|
\]
satisfont
\begin{align}
 \tau_{d+1}(\iota_d(a))&=\tau_d(a),
 &\iota_d(I)&=I,\label{nt:unital-trace}\\
 \tau_{d+1}(\jmath_{d,j}(a))&=\frac{\tau_d(a)}9,
 &\jmath_{d,j}(I)&=I\otimes p_j.\label{nt:corner-trace}
\end{align}
La démonstration utilise
\(\operatorname{Tr}(a\otimes b)=
\operatorname{Tr}(a)\operatorname{Tr}(b)\),
\(\operatorname{Tr}I_9=9\) et \(\operatorname{Tr}p_j=1\).
La première inclusion conserve toute la fibre et l’unité ;
la seconde conserve une continuation marquée et son support.

Cette distinction se manifeste également dans l'arbre des historiques enrichis. Soient \(\mathcal H_n\) les préfixes survivants de longueur \(n\), et soit \(d(h)\geq1\) le nombre d'émissions sortantes conservées de \(h\). Chaque émission avec multiplicité est représentée par des successeurs marqués distincts ; la somme suivante parcourt cet ensemble de \(d(h)\) successeurs. La loi des poids
\[
 \mu_{n+1}(h\epsilon)=\mu_n(h)\,p_h(\epsilon),
 \qquad p_h(\epsilon)=d(h)^{-1}
\]
conserve \(\sum_{\epsilon}\mu_{n+1}(h\epsilon)=\mu_n(h)\). Par conséquent, dans \(\mathscr L_n=L^2(\mathcal H_n,\mu_n)\), l'image réciproque
\[
 R_nf(h\epsilon)=f(h)
\]
est isométrique et conserve la fonction unité. Son adjoint est
\[
 E_ng(h)=\sum_{\epsilon}p_h(\epsilon)g(h\epsilon),
 \qquad E_nR_n=I,\quad E_n1=1.
\]
En effet, regrouper la somme de \(\|R_nf\|^2\) par prédécesseur produit \(\sum_h\mu_n(h)|f(h)|^2\) ; le même regroupement dans le produit scalaire détermine \(E_n\). La mémoire et les multiplicités demeurent présentes dans le domaine des historiques. Ces applications constituent la réalisation fonctionnelle du raffinement décrit dans l'article~X \cite{hmtX}.

\subsection{Coordinatisation spectrale et réciprocité du rayon}

Le transport précédent admet des réalisations spectrales possédant leurs propres domaines. La coordinatisation de Möbius de VIII agit, en coordonnées finies, par
\begin{equation}
 \tau(z)=\frac{z-1}{z},
 \qquad z\in\mathbb C\setminus\{0\}.
 \label{nt:mobius}
\end{equation}
L'inverse est \(z=(1-\tau)^{-1}\), pour \(\tau\ne1\). L'extension à la sphère de Riemann envoie \(0\) sur \(\infty\) et \(\infty\) sur \(1\).

\begin{proposition}[Miroir spectral et rayon réciproque]
Soit \(z^\sharp=1-\overline z\). Pour
\(z\notin\{0,1\}\),
\begin{equation}
 \tau(z^\sharp)=\frac1{\overline{\tau(z)}}.
 \label{nt:mirror}
\end{equation}
Si \(z=\sigma+it\), alors
\begin{equation}
 |\tau(z)|^2
 =\frac{(\sigma-1)^2+t^2}{\sigma^2+t^2};
 \qquad
 |\tau(z)|=1\ \Longleftrightarrow\ \sigma=\frac12.
 \label{nt:critical-circle}
\end{equation}
\end{proposition}

\begin{proof}
La substitution directe donne
\[
 \tau(1-\overline z)
 =\frac{-\overline z}{1-\overline z}
 =\frac{\overline z}{\overline z-1}
 =\frac1{\overline{(z-1)/z}}.
\]
Le quotient des carrés des modules produit
\eqref{nt:critical-circle} ; l’égalité entre le numérateur et le
dénominateur équivaut à \(-2\sigma+1=0\).
\end{proof}

En écrivant \(\tau(z)=r e^{i\theta}\), avec \(r>0\), le miroir préserve \(\theta\) et transforme \(r\) en \(r^{-1}\). Son ensemble fixe dans cette coordinatisation est le cercle unité. La conjugaison \(z\mapsto\overline z\), quant à elle, produit \(\tau\mapsto\overline\tau\) et change l'orientation angulaire. Les deux involutions conservent respectivement la coordonnée angulaire et la coordonnée radiale.

Sur \(z=1/2+i\gamma\), avec \(\gamma\in\mathbb R\),
\begin{equation}
 \tau_\gamma=\frac{\gamma+i/2}{\gamma-i/2},
 \qquad
 \gamma=\frac{i}{2}\frac{\tau_\gamma+1}{\tau_\gamma-1}
       =\frac12\cot\frac{\theta}{2},
 \quad \tau_\gamma=e^{i\theta}.
 \label{nt:spectral-phase}
\end{equation}
La compactification conserve la hauteur au moyen de l'inverse ; le point \(\tau=1\) correspond à l'infini.

Soit \(\mathcal D=C_c^\infty(\mathbb R;\mathbb C)\), avec sa topologie
usuelle de fonctions tests. Pour la réalisation positive de la forme
de Weil, on utilise expressément sa semi-définie positivité, son
invariance par translation et sa continuité dans cette topologie, sur le
domaine démontré par le lecteur spectral correspondant. Si
\(\mathscr W(f,g)=\langle Ef,Eg\rangle\), son noyau est
\(\mathcal N=\{f\in\mathcal D:\mathscr W(f,f)=0\}\), et l’on définit
\[
 \mathcal H_{\mathscr W}
 =\overline{\mathcal D/\mathcal N}.
\]
Les translations induisent un groupe unitaire fortement continu \(U_t[f]=[f(\,\cdot-t)]\) : l'invariance préserve la norme, et la continuité des translations dans \(\mathcal D\) et de la forme prouve la continuité forte sur un sous-espace dense, puis dans tout le complété. Avec la convention \(U_t=e^{itH}\), son générateur autoadjoint agit sur les classes de fonctions tests par \(H[f]=[if']\). La transformée de Cayley de cette réalisation est l'opérateur borné
\begin{equation}
 C_{\mathscr W}
 =(H+iI/2)(H-iI/2)^{-1}
 =I+i(H-iI/2)^{-1}.
 \label{nt:cayley-hilbert}
\end{equation}
Le calcul fonctionnel le rend unitaire, car le quotient
dans~\eqref{nt:spectral-phase} est de module un pour
\(\gamma\) réel. L’identité de compression s’applique
à \(C_{\mathscr W}\) dans \(\mathcal H_{\mathscr W}\).
L’espace, le domaine du générateur et la positivité
qui permet sa construction sont les prémisses de cette
réalisation spectrale \cite{HMT-VIII}.

La mémoire bilatérale agit, en revanche, dans
\(\ell^2(\mathbb Z)\), et les permutations \(S^3,S^4\)
agissent dans \(\mathbb R^{12}\). Chaque réalisation détermine
son spectre, ses secteurs fixes et l’action d’un tour.
Les applications de transport respectives spécifient les
compositions entre ces espaces.

\subsection{Cayley sur les fonctions tests et conservation de la forme arithmétique}

La même représentation arithmétique admet un lecteur de Cayley défini directement sur les fonctions tests. Fixons \(a>b>1/2\) et
\[
 \mathscr E_b=
 \left\{f\in C^\infty(\mathbb R):
   \sup_x e^{b|x|}|f^{(j)}(x)|<\infty
   \text{ pour tout }j\geq0\right\}.
\]
Les fonctions lisses à support compact appartiennent à \(\mathscr E_b\). Nous définissons
\begin{align}
 (R_a^+f)(x)&=\int_0^\infty e^{-at}f(x+t)\,dt,
 &(R_a^-f)(x)&=\int_0^\infty e^{-at}f(x-t)\,dt,\label{nt:resolvents}\\
 C_af&=f-2aR_a^+f,
 &C_a^{-1}f&=f-2aR_a^-f.\label{nt:test-cayley}
\end{align}
Chaque dérivée de \(R_a^\pm f\) satisfait la borne de la semi-norme correspondante de \(f\), multipliée par \((a-b)^{-1}\). Cela découle de \(e^{b|x|}|f(x\pm t)|\leq M_0e^{bt}\). Les deux applications préservent \(\mathscr E_b\).

Avec la transformée
\(\widehat f(\xi)=\int_{\mathbb R}f(x)e^{-i\xi x}\,dx\),
l’application de Fubini donne
\begin{equation}
 \widehat{C_af}(\xi)
 =c_a(\xi)\widehat f(\xi),\qquad
 c_a(\xi)=\frac{\xi-ia}{\xi+ia}.
 \label{nt:cayley-multiplier}
\end{equation}
Le multiplicateur de la seconde formule
de~\eqref{nt:test-cayley} est \(c_a^{-1}\), ce qui prouve
l’inverse. Pour \(\xi\in\mathbb R\),
\(|c_a(\xi)|=1\) ; Plancherel fournit son extension
unitaire à \(L^2(\mathbb R)\).

Pour préciser la conservation arithmétique, soient
\[
 \mathcal C_{f,g}(u)
 =\int_{\mathbb R}f(x)\overline{g(x-u)}\,dx,
 \qquad
 P_f^\pm=\int_{\mathbb R}f(x)e^{\pm x/2}\,dx.
\]
Les longueurs \(\ell_{p,k}=k\log p\) et les poids \(w_{p,k}=(\log p)p^{-k/2}\) sont les représentations premier--puissance du lecteur arithmétique. La forme complète de l'article~VIII s'écrit
\begin{equation}
 \begin{split}
 \mathscr W(f,g)={}&c_\Gamma\mathcal C_{f,g}(0)\\
 &+\int_0^\infty\frac{e^{-u/2}}{1-e^{-2u}}
 \bigl[2\mathcal C_{f,g}(0)
       -\mathcal C_{f,g}(u)-\mathcal C_{f,g}(-u)\bigr]\,du\\
 &-\sum_{p,k}w_{p,k}
   \bigl[\mathcal C_{f,g}(\ell_{p,k})
        +\mathcal C_{f,g}(-\ell_{p,k})\bigr]
 +P_f^+\overline{P_g^-}+P_f^-\overline{P_g^+},
 \end{split}
 \label{nt:weil-complete}
\end{equation}
où \(c_\Gamma=\psi(1/4)-\log\pi\), et \(\psi=\Gamma'/\Gamma\) est la dérivée logarithmique de la fonction gamma. Ces fonctions et constantes interviennent dans la réalisation analytique ultérieure du lecteur arithmétique du corpus.

La borne
\[
 |\mathcal C_{f,g}(u)|
 \leq M_fM_g\bigl(|u|+b^{-1}\bigr)e^{-b|u|}
\]
contrôle la somme sur les nombres premiers par une série convergente du type
\(\sum_{n\geq2}(\log n)(1+\log n)n^{-b-1/2}\).
La différence symétrique de corrélations est \(O(u^2)\)
à l’origine, tandis que le poids gamma est \(O(u^{-1})\).
À l’infini, ce poids décroît exponentiellement.
La marge \(b>1/2\) assure également la convergence des
deux lecteurs polaires. Le domaine
de tous les termes de~\eqref{nt:weil-complete} est ainsi fixé.

\begin{proposition}[Invariance arithmétique et bilan nonadique]
Pour \(f,g\in\mathscr E_b\),
\begin{equation}
 \mathscr W(C_af,C_ag)=\mathscr W(f,g).
 \label{nt:weil-invariance}
\end{equation}
Avec \(\mathcal T_a=(8I+C_a)/9\), on a
\begin{equation}
 \mathscr W(f,g)
 =\mathscr W(\mathcal T_af,\mathcal T_ag)
 +\frac8{81}\mathscr W((I-C_a)f,(I-C_a)g).
 \label{nt:weil-balance}
\end{equation}
\end{proposition}

\begin{proof}
Le multiplicateur \(c_a\) est unitaire sur la droite réelle et
commute avec les translations. Par conséquent,
\(\mathcal C_{C_af,C_ag}(u)=\mathcal C_{f,g}(u)\)
pour tout \(u\). Les lecteurs polaires se transforment
par Fubini selon
\[
 P_{C_af}^+
 =-\frac{a-1/2}{a+1/2}P_f^+,\qquad
 P_{C_af}^-
 =-\frac{a+1/2}{a-1/2}P_f^-.
\]
Leurs facteurs sont réels et réciproques, de sorte que chaque produit polaire croisé reste inchangé. La réunion des termes de \eqref{nt:weil-complete} prouve~\eqref{nt:weil-invariance}. Le développement sesquilinéaire du second membre de \eqref{nt:weil-balance} annule les termes mixtes et produit \(\bigl(8\mathscr W(f,g)+
\mathscr W(C_af,C_ag)\bigr)/9\), égal à \(\mathscr W(f,g)\).
\end{proof}

L’itération finie conserve la forme complète :
\[
 \mathscr W(f,g)
 =\mathscr W(\mathcal T_a^Nf,\mathcal T_a^Ng)
 +\frac8{81}\sum_{j=0}^{N-1}
 \mathscr W((I-C_a)\mathcal T_a^jf,
            (I-C_a)\mathcal T_a^jg).
\]
Cette identité sesquilinéaire conserve simultanément
gamma, nombres premiers et termes polaires. Son évaluation dans une
réalisation positive permet de l’interpréter comme répartition
d’énergie ; sa validité algébrique précède cette
interprétation. La marge \(a>b>1/2\) appartient au
domaine exponentiel des fonctions tests de cette construction.
Le Cayley~\eqref{nt:cayley-hilbert} de paramètre \(1/2\)
appartient au domaine hilbertien de son générateur.

\subsection{Covariance métrique du flux dodécaphasique}

La reconstruction de la section~\ref{sec:memory} détermine \(u_K\in\mathbb R^{12}\), avec \(\mathbf1^{\mathsf T}u_K=0\) et \(u_K\ne0\). Nous écrivons \(A_K=\operatorname{diag}(u_K)\), avec la direction et la normalisation déclarées dans \eqref{eq:u-main}, et
\[
 g_s=e^{sA_K}\in\operatorname{GL}(12,\mathbb R).
\]
La trace nulle de \(A_K\) implique \(\det g_s=1\). Chaque \(g_s\) transporte les poids et les coordonnées géométriques selon la réalisation correspondante. Pour conserver également le bilan de mémoire entre coordinatisations, il faut transporter le produit scalaire.

\begin{proposition}[Transport métrique de la compression et du complément]
Soient \(V\) un espace hermitien fini, \(C_0\) unitaire et \(g:V\to V\) inversible. Nous écrivons \(g^{-*}=(g^{-1})^*\) et définissons
\begin{equation}
 C_g=gC_0g^{-1},\qquad
 G_g=g^{-*}g^{-1},\qquad
 \langle v,w\rangle_g=\langle G_gv,w\rangle.
 \label{nt:metric-chart}
\end{equation}
Alors \(C_g^*G_gC_g=G_g\), et
\begin{align}
 T_{C_g}g&=gT_{C_0},\label{nt:metric-naturality}\\
 \eta_{C_g}g&=g^{\oplus9}\eta_{C_0},\label{nt:eta-naturality}\\
 G_g&=T_{C_g}^*G_gT_{C_g}
 +\frac8{81}(I-C_g)^*G_g(I-C_g).\label{nt:metric-balance}
\end{align}
\end{proposition}

\begin{proof}
La positivité de \(G_g\) s’obtient à partir de
\(\langle G_gv,v\rangle=\|g^{-1}v\|^2\).
En substituant~\eqref{nt:metric-chart},
\[
 C_g^*G_gC_g
 =g^{-*}C_0^*g^*g^{-*}g^{-1}gC_0g^{-1}
 =g^{-*}C_0^*C_0g^{-1}=G_g.
\]
L’identité \(C_gg=gC_0\) donne
\eqref{nt:metric-naturality} en ajoutant \(8g\) et en divisant
par neuf. Elle donne aussi
\((C_g-I)g=g(C_0-I)\) ; sa substitution dans
\eqref{nt:complement} prouve~\eqref{nt:eta-naturality}.
Enfin, le développement de
\eqref{nt:observable-balance}, avec \(A=G_g\) et
\(C_g^*G_gC_g=G_g\), prouve
\eqref{nt:metric-balance}.
\end{proof}

En appliquant la proposition à \(g=g_s\) et aux opérateurs \(C_0=S^3,S^4\) de la réalisation dodécaphasique, les compressions, les mémoires et leurs produits scalaires sont transportés par des carrés commutatifs explicites. L'identification \(g_s:(V,\langle\cdot,\cdot\rangle)
\to(V,\langle\cdot,\cdot\rangle_{g_s})\) est isométrique. Par conséquent, la formule transporte également les secteurs fixes et les registres à toute profondeur :
\[
 P_{\mathrm{fix},s}
 =g_sP_{\mathrm{fix},0}g_s^{-1}.
\]
Le projecteur du membre de droite est orthogonal relativement
à \(G_{g_s}\).

La condition de volume et la condition métrique ont
des effets distincts. Pour la matrice rationnelle
\[
 g=\begin{pmatrix}2&0\\0&1/2\end{pmatrix},
 \qquad \det g=1,
\]
la substitution directe dans la compression donne
\[
 \frac{8I+g}{9}
 =\begin{pmatrix}10/9&0\\0&17/18\end{pmatrix},
 \qquad
 T_g^*T_g+\frac8{81}(I-g)^*(I-g)
 =\frac{8I+g^*g}{9}.
\]
La première coordonnée se dilate dans la métrique
euclidienne. Le bilan~\eqref{nt:metric-balance}
conserve en revanche l’unité métrique lorsqu’on
transporte un opérateur unitaire \(C_0\) avec
son produit scalaire. Le déterminant contrôle
le volume ; \(G_g\) contrôle les normes et les
produits croisés.

Le flot satisfait \(g_{-s}=g_s^{-1}\). La réciprocité du rayon de \eqref{nt:mirror} relève de la coordinatisation spectrale, et l'inversion de \(g_s\) relève de la réalisation dimensionnelle du registre. Le lien constructif utilisé ici conserve les applications explicites de la mémoire vers \(K\) et de \(K\) vers sa direction géométrique, suivies du transport métrique de chaque réalisation. Cette composition maintient, dans un même développement, la reconstruction exacte, le raffinement positif, le registre des frontières et la conservation des produits scalaires.

\section{Incidence ternaire, recollement des discriminants et construction de Leech}
\label{sec:lf-foundations}

La construction réticulaire utilise l’information d’incidence qui
accompagne le registre dodécaphasique. Son domaine conserve les mots
ternaires, le repère dans lequel ils sont représentés, l’orientation préalable à la
retenue et le drapeau qui distingue l’une des douze positions. Le
passage à un réseau transmet ces données par des opérations
successives : codage linéaire, sélection incidencielle, recollement de
discriminants et formation d’un voisin. La classification réticulaire
intervient après la preuve de la positivité, de la parité, de l’autodualité
et de l’absence de racines de la construction obtenue.

Ces opérations poursuivent la même génération
APP--TRIT--TPK et le même raffinement compatible du continu
conjoint qui précèdent la lecture numérique \cite{HMT-IV}.
La matrice finie, le code et la métrique définis
ci-dessous sont des réalisations aux domaines précis. La preuve
réticulaire reçoit l’état marqué ; le scalaire final d’une
constante ne joue le rôle ni de code de recollement ni de longueur.

\subsection{Coordinatisation elliptique et transport de l'incidence}
\label{subsec:lf-incidence}

Les six unités \(\{1,2,4,5,7,8\}\) modulo neuf, conservées par le transport modulaire, admettent une coordinatisation à six positions. On l'identifie à \(\mathbb P^1(\mathbb F_5)=\{\infty,0,1,2,3,4\}\), dans cet ordre. Cette identification fixe les coordonnées permettant de réaliser la relation quadratique du régime elliptique du TRIT.

Soit \(\chi:\mathbb F_5\to\{0,1,-1\}\) le caractère quadratique :
\(\chi(0)=0\), \(\chi(1)=\chi(4)=1\) et
\(\chi(2)=\chi(3)=-1\). Nous définissons la matrice entière
\(C_{\mathrm P}\) à diagonale nulle,
\((C_{\mathrm P})_{\infty,a}=(C_{\mathrm P})_{a,\infty}=1\),
et d’entrées finies \((C_{\mathrm P})_{a,b}=\chi(b-a)\).
Son carré est \(5I_6\). En effet, chaque entrée diagonale du
carré additionne cinq unités. Les entrées qui font intervenir
\(\infty\) hors de la diagonale s’annulent car la somme de
\(\chi\) est nulle. Pour deux indices finis distincts, la
contribution un de la position \(\infty\) s’annule avec
\(\sum_c\chi(c-a)\chi(c-b)=-1\). Cette dernière identité se vérifie
directement sur les cinq résidus, ou au moyen de la substitution
affine qui envoie \((a,b)\) sur \((0,1)\).

La réduction modulo trois produit la matrice symétrique
\begin{equation}
 A_W=
 \begin{pmatrix}
 0&1&1&1&1&1\\
 1&0&1&2&2&1\\
 1&1&0&1&2&2\\
 1&2&1&0&1&2\\
 1&2&2&1&0&1\\
 1&1&2&2&1&0
 \end{pmatrix},
 \qquad A_W^2=2I_6=-I_6
 \quad\text{en }\mathbb F_3.
 \label{eq:lf-aw}
\end{equation}
L’égalité conserve le type elliptique dans ce codomaine fini.
Le corps fini et la réalisation archimédienne ultérieure de
l’unité imaginaire ont des domaines différents.

Pour une matrice \(A\in M_6(\mathbb F_3)\) et un mot ligne
\(w\in\mathbb F_3^6\), nous écrivons
\[
 \gamma_A(w)=(w,wA)\in\mathbb F_3^{12}.
\]
Si \(f\in\operatorname{GL}_6(\mathbb F_3)\), alors
\begin{equation}
 \gamma_{fAf^{-1}}(wf^{-1})
   =\gamma_A(w)(f^{-1}\oplus f^{-1}).
 \label{eq:lf-frame-covariance}
\end{equation}
La première moitié est \(wf^{-1}\) ; la seconde est \(wf^{-1}fAf^{-1}=wAf^{-1}\), ce qui prouve l'identité.

Soit \(X_n^{\mathrm{mar}}\) l’espace des préfixes enrichis de
longueur \(n\), avec un repère initial déclaré, et soient
\(p_{n,m}:X_n^{\mathrm{mar}}\to X_m^{\mathrm{mar}}\) leurs
troncatures. À l’étape \(j\), le préfixe contient un mot
visible \(w_j\) et un repère \(f_j\). La mise à jour du repère
est de la forme \(f_{j+1}=g_j(x)f_j\), où \(g_j(x)\) dépend
uniquement de l’information déjà contenue dans le préfixe. Un
prolongement compatible conserve donc tous les repères
antérieurs.

\begin{proposition}[Naturalité du transport d'incidence]
\label{prop:lf-incidence-naturality}
L'application
\begin{equation}
 \mathcal Q_n(x)=
 \bigl((f_j,\gamma_{f_jA_Wf_j^{-1}}(w_j))\bigr)_{0\leq j<n}
 \label{eq:lf-Qn}
\end{equation}
satisfait \(r_{n,m}\mathcal Q_n=\mathcal Q_mp_{n,m}\), où \(r_{n,m}\) tronque aux \(m\) premières composantes codées. Elle induit donc une application entre les limites projectives des systèmes de préfixes et d'incidences marquées.
\end{proposition}
\begin{proof}
Deux préfixes compatibles ont les mêmes mots et les mêmes opérateurs de transport jusqu'à l'étape \(m-1\). À partir du même repère initial, la relation récursive donne par récurrence les mêmes \(f_j\) jusqu'à cette étape. Les \(m\) premières composantes de \eqref{eq:lf-Qn} coïncident alors. La propriété universelle de la limite projective produit l'application de prolongement.
\end{proof}

Dans chaque repère, la première projection de \(\gamma_A\) restitue exactement le mot visible. Cette portée est locale : la mémoire profonde est conservée dans l'état enrichi et sa récupération exige les coordonnées correspondantes. Le passage des mots aux supports retient encore moins d'information. De même, un changement général de base dans \(\operatorname{GL}_6(\mathbb F_3)\) transporte également la métrique de Hamming et l'incidence de la coordinatisation. Les changements monomiaux les préservent directement. Les supports utilisés dans cette section sont calculés dans la coordinatisation explicite \eqref{eq:lf-aw}.

La naturalité précédente concrétise la lecture incidencielle qui
accompagne les cylindres numériques. L’emboîtement des cylindres
et la restriction de \eqref{eq:lf-Qn} agissent sur un même
préfixe compatible, en conservant respectivement son évaluation
positionnelle et ses relations marquées. L’espace ambiant
\(\mathbb F_3^6\), de \(729\) mots, fournit le domaine
linéaire du codage ; il conserve sa distinction par rapport aux
\(243\) mots visibles et aux \(468\) émissions du recensement
admissible d’APP--TRIT--TPK.

\subsection{Code ternaire autodual et énumération de ses poids}
\label{subsec:lf-code}

\begin{definition}[Code de la coordinatisation dodécaphasique]
Le code linéaire associé à \eqref{eq:lf-aw} est
\[
 C_W=\{(w,wA_W):w\in\mathbb F_3^6\}\subset\mathbb F_3^{12}.
\]
Le poids \(\operatorname{wt}(c)\) compte les coordonnées non nulles de \(c\), et \(\operatorname{supp}(c)\subset\{1,\ldots,12\}\) est l'ensemble de ses positions non nulles. Une hexade est le support d'un mot de poids six ; leur collection est notée \(\mathcal H\).
\end{definition}

\begin{theorem}[Autodualité, distance minimale et énumérateur]
\label{thm:lf-code}
Le code \(C_W\) est autodual de paramètres \([12,6,6]_3\) et
a pour énumérateur
\begin{equation}
 W_{C_W}(z)=1+264z^6+440z^9+24z^{12}.
 \label{eq:lf-enumerator}
\end{equation}
\end{theorem}
\begin{proof}
La matrice génératrice \([I_6\mid A_W]\) a rang six.
Par la symétrie de \(A_W\) et de \eqref{eq:lf-aw},
\[
 [I_6\mid A_W][I_6\mid A_W]^{\mathsf T}
       =I_6+A_W^2=0.
\]
Ainsi, le code est auto-orthogonal de dimension égale à la moitié de sa longueur et, par conséquent, autodual.

L’énumération restante est une opération entière finie sur
la matrice explicite. Pour \(r\in\{0,\ldots,6\}\), nous regroupons
les mots par \(\operatorname{wt}(w)=r\). Le tableau suivant
donne le nombre de mots de chaque poids total :
\[
\begin{array}{c|rrrr|r}
r&0&6&9&12&\text{total}\\ \hline
0&1&0&0&0&1\\
1&0&12&0&0&12\\
2&0&60&0&0&60\\
3&0&120&40&0&160\\
4&0&60&180&0&240\\
5&0&12&180&0&192\\
6&0&0&40&24&64\\ \hline
\text{total}&1&264&440&24&729
\end{array}
\]
Chaque ligne contient \(\binom6r2^r\) mots. Pour vérifier
toutes ses entrées, soit \(q\) un entier de \(0\) à \(728\), et
définissons ses six chiffres ternaires
\[
 d_i(q)=\left\lfloor\frac{q}{3^{6-i}}\right\rfloor\bmod3,
 \qquad
 e_j(q)=\left(\sum_{i=1}^{6}d_i(q)(A_W)_{ij}\right)\bmod3.
\]
La ligne et la colonne de la contribution de \(q\) sont, respectivement,
\[
 r(q)=\sum_i\mathbf1_{\{d_i(q)\ne0\}},\qquad
 h(q)=r(q)+\sum_j\mathbf1_{\{e_j(q)\ne0\}}.
\]
La récurrence d'accumulation commence par \(T^{(0)}_{r,h}=0\) et applique
\[
 T^{(q+1)}_{r,h}
 =T^{(q)}_{r,h}
   +\mathbf1_{\{r=r(q)\}}\mathbf1_{\{h=h(q)\}}
 \quad(0\leq q<729).
\]
Il produit exactement le tableau montré. Le développement ternaire
donne une bijection entre ces \(729\) entiers et tous les mots
de \(\mathbb F_3^6\), de sorte que l’opération épuise le code.
La somme de ses lignes prouve \eqref{eq:lf-enumerator} ; le premier
poids non nul est six.
\end{proof}

\begin{remark}[Évaluation reproductible de la récurrence finie]
La boucle suivante implémente littéralement les sommes de la preuve. Ses seules données sont les entiers de la matrice \eqref{eq:lf-aw} ; elle calcule l'énumérateur avant de le comparer à l'égalité obtenue.
\end{remark}
{\small
\begin{verbatim}
from collections import Counter
from itertools import product
A = ((0,1,1,1,1,1), (1,0,1,2,2,1),
     (1,1,0,1,2,2), (1,2,1,0,1,2),
     (1,2,2,1,0,1), (1,1,2,2,1,0))
rows = [Counter() for _ in range(7)]
for w in product(range(3), repeat=6):
    v = tuple(sum(w[i]*A[i][j] for i in range(6)) % 3
              for j in range(6))
    r = sum(x != 0 for x in w)
    h = r + sum(x != 0 for x in v)
    rows[r][h] += 1
total = sum(rows, Counter())
if dict(total) != {0:1, 6:264, 9:440, 12:24}:
    raise ArithmeticError("Enumerador incompatible")
\end{verbatim}
}

\begin{proposition}[Hexades et plan de Witt]
\label{prop:lf-witt}
La collection \(\mathcal H\) contient \(132\) hexades et chaque quintuplet de positions appartient à une unique hexade.
\end{proposition}
\begin{proof}
Deux mots de poids six ayant le même support diffèrent par
un scalaire. En effet, parmi leurs six quotients de coordonnées
non nulles, l’un des signes \(1,-1\) apparaît au moins trois
fois. En soustrayant le multiple correspondant, si les mots
n’étaient pas proportionnels, on obtiendrait un mot non nul
de poids au plus trois, en contradiction avec
le théorème~\ref{thm:lf-code}. Chaque support correspond ainsi à
la paire \(\{c,-c\}\), et il y a \(264/2=132\) supports.

Si deux supports distincts partageaient cinq positions, dans cette intersection l'un des deux quotients apparaîtrait au moins trois fois. La combinaison linéaire qui annule ces coordonnées aurait un support dans une réunion de taille sept et un poids au plus égal à quatre. La distance minimale l'interdit. Un quintuplet appartient donc à au plus une hexade. Comme
\[
 132\binom65=792=\binom{12}{5},
\]
appartient à exactement une.
\end{proof}

L’objet construit est le code ternaire étendu de Golay
avec son petit plan de Witt. La reconnaissance s’effectue
après l’encodage, l’autodualité et le calcul de la
distance. Pour la construction réticulaire, seules sont nécessaires
les propriétés déjà démontrées et le drapeau marqué obtenu
ci-dessous.

\subsection{Drapeau dodécaphasique et sélection de l'origine}
\label{subsec:lf-flag}

Dans le repère du registre engendré, on conserve les mots
\[
 w_\pi=010211,\qquad
 w_e=201101,\qquad
 w_\varphi=121200,\qquad
 w_{\mathrm{APP}}=022020.
\]
Les indices \(\pi,e,\varphi\) identifient les régions correspondantes de lecture des constantes ; \(w_e\) conserve ici la région d'Euler. Appliquer \(\gamma_{A_W}\) et prendre les supports donne
\[
\begin{aligned}
 P_\pi&=\{2,4,5,6,7,8,9,10,11\},\\
 H_e&=\{1,3,4,6,9,10\},\\
 H_\varphi&=\{1,2,3,4,7,9\},\\
 H_{\mathrm{APP}}&=\{2,3,5,10,11,12\}.
\end{aligned}
\]
Le premier support est de poids neuf ; les trois autres sont des
hexades. Le registre
\[
 K=(234,543,140,729,659,824,621,58,914,794,146,601)
\]
sélectionne
\[
 B_K=\{j:K_j\geq729\}=\{4,6,9,10\}=H_e\cap P_\pi.
\]
Le seuil \(729\) appartient au format de blocs de la
lecture déclarée. Dans ce registre concret, les seuils
entiers de \(660\) à \(729\) produisent le même support, de
sorte que le support conserve l’incidence sélectionnée,
tandis que le registre complet conserve aussi les blocs
et le seuil de sa lecture.

Le registre signé antérieur à la retenue est
\[
\begin{split}
 Z_0={}&(7,297,353,-430,-715,-1199,\\
       &\hspace{9mm}-3,286,-894,-528,380,663).
\end{split}
\]
Son support négatif est l’hexade
\[
 H_\alpha^-=\{4,5,6,7,9,10\}.
\]
Les signes sont reçus de l'état antérieur à la retenue et ne sont pas extraits du développement du scalaire final.

\begin{proposition}[Coïncidence de la sélection signée et de la sélection incidencielle]
\label{prop:lf-flag}
L’hexade \(H_\alpha^-\) est l’unique \(H\in\mathcal H\) qui vérifie
\[
\begin{gathered}
 B_K\subset H\subset P_\pi,\\
 (|H\cap H_e|,|H\cap H_\varphi|,
                   |H\cap H_{\mathrm{APP}}|)=(4,3,2).
\end{gathered}
\]
\end{proposition}
\begin{proof}
Le codage des mots de poids six, ou la récurrence finie précédente appliquée à leurs supports, donne les quatre hexades sur \(B_K\). Leurs paires complémentaires sont
\[
 \{1,3\},\qquad\{2,11\},\qquad\{5,7\},\qquad\{8,12\}.
\]
Seuls le deuxième et le troisième sont contenus dans \(P_\pi\).
L’hexade du deuxième a des intersections de taille trois
avec \(H_\varphi\) et avec \(H_{\mathrm{APP}}\), et ne satisfait donc pas
la dernière condition. Celle du troisième a pour
intersections \(4,3,2\), et coïncide avec \(H_\alpha^-\).
\end{proof}

\begin{lemma}[Origine du drapeau marqué]
\label{lem:lf-origin}
L’origine sélectionnée par l’incidence est
\[
 o_*=(H_e\cap H_\varphi)
                 \setminus(H_\alpha^-\cup H_{\mathrm{APP}})
     =\{1\}.
\]
La sélection commute avec toute bijection appliquée simultanément aux supports.
\end{lemma}
\begin{proof}
L’intersection \(H_e\cap H_\varphi\) est
\(\{1,3,4,9\}\). L’union soustraite contient \(3,4,9\)
et exclut \(1\). L’équivariance provient du fait que les
bijections commutent avec l’union, l’intersection et la différence.
\end{proof}

Le drapeau distingue une composante parmi douze avant l'étape métrique. Le vecteur radial est déterminé ensuite, dans une chambre positive déclarée. La sélection combinatoire et l'opération réticulaire sont ainsi séparées et articulées.

\subsection{Recollement ternaire des douze composantes}
\label{subsec:lf-glue}

Soit \(A_2=\mathbb Za_1\oplus\mathbb Za_2\), avec une forme
bilinéaire de matrice
\[
 B_2=\begin{pmatrix}2&-1\\-1&2\end{pmatrix},
 \qquad \rho=a_1+a_2,\qquad \rho^2=2.
\]
Pour \(x=ua_1+va_2\),
\[
 x^2=2(u^2-uv+v^2).
\]
L’identité
\(4(u^2-uv+v^2)=(2u-v)^2+3v^2\) montre
la positivité. Le déterminant est trois. Les trois
classes de \(A_2^*/A_2\) sont représentées par
\[
 \varpi_0=0,\qquad
 \varpi_1=\frac{2a_1+a_2}{3},\qquad
 \varpi_2=\frac{a_1+2a_2}{3}.
\]
L’appariement avec \(a_1,a_2\) vérifie que ces
vecteurs appartiennent à \(A_2^*\), et
\(2\varpi_1-\varpi_2=a_1\). C’est pourquoi
\(j\mapsto\varpi_j+A_2\) identifie additivement
\(\mathbb F_3\) au discriminant.

Pour \(c\in C_W\), nous définissons
\(\phi(c)=(\varpi_{c_1},\ldots,\varpi_{c_{12}})\)
et
\begin{equation}
 N=\bigcup_{c\in C_W}\bigl(A_2^{12}+\phi(c)\bigr).
 \label{eq:lf-glue}
\end{equation}

\begin{theorem}[Réseau de recollement pair et unimodulaire]
\label{thm:lf-niemeier}
L’ensemble \(N\) est un réseau positif, pair et
unimodulaire de rang \(24\). Ses vecteurs de norme deux
sont exactement les racines de \(A_2^{12}\).
\end{theorem}
\begin{proof}
La linéarité du code et l’additivité du discriminant
font de l’union un sous-groupe de
\((A_2^*)^{12}\) contenant \(A_2^{12}\).
C’est donc un réseau de rang \(24\).

L'appariement des classes discriminantes est \(2cd/3\) modulo \(\mathbb Z\) dans chaque composante. L'auto-orthogonalité de \(C_W\) annule leur somme et prouve l'intégralité de \(N\). La norme d'une classe représentée par \(\phi(c)\) est \(2\operatorname{wt}(c)/3\) modulo \(2\mathbb Z\). Tous les poids de \eqref{eq:lf-enumerator} sont des multiples de trois, donc \(N\) est pair.

Son indice sur \(A_2^{12}\) est \(|C_W|=3^6\).
Avec le déterminant de la matrice de Gram,
\[
 \det N=\frac{\det(A_2^{12})}{[N:A_2^{12}]^2}
       =\frac{3^{12}}{(3^6)^2}=1.
\]
La forme ambiante est la somme orthogonale de douze formes
positives ; la positivité est démontrée.
L’intégralité et le déterminant un impliquent
\(N^*=N\).

Dans une classe non nulle du discriminant local, un
vecteur a la forme \((ua_1+va_2)/3\), avec des résidus
\((u,v)\equiv(2,1)\) ou \((1,2)\pmod3\).
L’entier \(u^2-uv+v^2\) est positif et divisible
par trois, il vaut donc au moins trois. Les représentants
\(\varpi_1,\varpi_2\) atteignent cette valeur : la norme
minimale de chaque classe non nulle est \(2/3\).
Un mot non nul du code a au moins
six coordonnées non nulles. Tout vecteur d’une classe
de recollement non triviale a donc une norme
d’au moins \(6(2/3)=4\).

Dans \(A_2^{12}\), une norme deux ne peut provenir que d'une unique composante, car chaque composante a une norme paire positive minimale égale à deux. Les solutions de \(u^2-uv+v^2=1\) sont \((\pm1,0),(0,\pm1),(1,1),(-1,-1)\). Ce sont les six racines de \(A_2\), et elles complètent le recensement des racines de \(N\).
\end{proof}

Les propriétés précédentes identifient \(N\) comme le
réseau de Niemeier de système de racines \(A_2^{12}\).
La description de toutes ses racines est nécessaire pour
contrôler lesquelles survivent au passage au voisin.

\subsection{Chambre radiale et voisin sans racines}
\label{subsec:lf-neighbour}

On fixe la chambre radiale
\[
 v(a)=\sum_{j=1}^{12}a_j\rho_j,\qquad
 a_j=1+3b_j,\qquad b_j\in\mathbb Z_{\geq0},
\]
où \(\rho_j\) est la copie de \(\rho\) dans la composante \(j\). La parité d'un voisin d'indice trois requiert \(v(a)^2\equiv0\pmod{18}\). Comme
\[
 v(a)^2=2\sum_j(1+3b_j)^2
       =24+12\sum_jb_j+18\sum_jb_j^2,
\]
la congruence équivaut à \(\sum_jb_j\equiv1\pmod3\).
La norme minimale à l’intérieur de la chambre est atteinte avec
un unique \(b_j=1\) et les autres nuls, et vaut \(54\).
L’origine \(o_*=1\) du lemme~\ref{lem:lf-origin}
sélectionne
\begin{equation}
 v_\alpha=4\rho_1+\rho_2+\cdots+\rho_{12},
 \qquad v_\alpha^2=54.
 \label{eq:lf-valpha}
\end{equation}
Le qualificatif de minimal se rapporte à la chambre déclarée. Dans la classe résiduelle complète, \((a_1-2a_2,\rho,\ldots,\rho)\) représente la même classe modulo trois et a pour norme \(14+11\cdot2=36\). Cette distinction préserve la condition qui détermine le coefficient quatre.

On définit
\begin{equation}
\begin{split}
 N_0&=\{x\in N:\langle x,v_\alpha\rangle\equiv0\pmod3\},\\
 \Lambda&=N_0+\mathbb Z\,v_\alpha/3.
\end{split}
\label{eq:lf-lambda}
\end{equation}

\begin{theorem}[Voisin pair autodual sans racines]
\label{thm:lf-leech}
La construction \eqref{eq:lf-lambda} est un
réseau positif, pair et autodual de rang \(24\)
sans vecteurs de norme deux. Par conséquent, elle est
isométrique au réseau de Leech.
\end{theorem}
\begin{proof}
Le caractère \(x\mapsto\langle x,v_\alpha\rangle\bmod3\)
est non nul. Une racine simple de la composante \(j\)
a pour appariement \(a_j\equiv1\pmod3\) avec
\(a_j\rho_j\). Par conséquent, \([N:N_0]=3\), d’où
\(\det N_0=9\).

Pour \(x\in N_0\), l’appariement
\(\langle x,v_\alpha/3\rangle\) est entier, et ainsi
\(v_\alpha/3\in N_0^*\). En outre,
\(\langle v_\alpha,v_\alpha\rangle=54\) implique
\(v_\alpha\in N_0\). Si \(v_\alpha/3\in N\),
l’intégralité de \(N\) imposerait que tous les
appariements de \(N\) avec \(v_\alpha\) soient
divisibles par trois, contredisant le caractère
non nul. La classe de \(v_\alpha/3\) sur \(N_0\)
a alors un ordre exactement égal à trois. Il en résulte
\[
 [\Lambda:N_0]=3,\qquad
 \det\Lambda=\frac9{3^2}=1.
\]
Les appariements de \(N_0\) avec le nouveau
générateur sont entiers et la norme de ce générateur
est six. L’extension est entière. Pour
\(x\in N_0\) et \(m\in\mathbb Z\),
\[
 \left(x+m\frac{v_\alpha}{3}\right)^2
 =x^2+2m\frac{\langle x,v_\alpha\rangle}{3}
       +6m^2\in2\mathbb Z.
\]
Par conséquent, \(\Lambda\) est pair ; son intégralité et son déterminant égal à un prouvent l'autodualité. La positivité et le rang sont hérités de l'espace euclidien ambiant.

Pour exclure les racines, nous considérons d'abord la classe \(N_0\). Toutes les racines de \(N\) appartiennent à \(A_2^{12}\) d'après le théorème~\ref{thm:lf-niemeier}. Les appariements de ses six racines locales avec \(a_j\rho_j\) sont \(\pm a_j,\pm2a_j\), congrus à \(\pm1,\pm2\pmod3\). Aucune racine n'appartient à \(N_0\).

Les deux autres classes de \(\Lambda/N_0\) peuvent
être représentées avec \(\pm v_\alpha/3\).
Dans chaque composante, leurs vecteurs appartiennent à
l’une des classes
\[
 A_2+\varpi_j\pm\rho/3,\qquad j=0,1,2.
\]
La composante distinguée possède la même
description car \(4\rho/3\equiv\rho/3\pmod{A_2}\).
Pour les numérateurs \((u,v)\) dans
\((ua_1+va_2)/3\), les six résidus sont
\[
 (1,1),(0,2),(2,0),(2,2),(1,0),(0,1)\pmod3.
\]
Tous sont non nuls. La positivité de
\(u^2-uv+v^2\) impose qu’il vaille au moins un,
et les représentants
\((1,1),(0,-1),(-1,0),(-1,-1),(1,0),(0,1)\)
atteignent ce minimum dans les six classes.
Chaque composante a une norme d’au moins \(2/9\).
Un vecteur de l’une ou l’autre des deux classes
nouvelles a une norme d’au moins
\[
 12\cdot\frac29=\frac83>2.
\]
Les racines ont été exclues dans les trois classes. Le théorème de classification des réseaux positifs pairs unimodulaires de rang vingt-quatre identifie l'unique cas sans racines au réseau de Leech \cite{ConwaySloane}.
\end{proof}

L’identification finale utilise une classification
externe sur un objet dont les propriétés ont déjà
été démontrées. La sélection de la composante
distinguée conserve le drapeau provenant du
registre signé ; l’indice \(\alpha\) exprime
cette provenance. La norme et les coefficients du
vecteur radial s’obtiennent à partir de la métrique et de
la chambre au moyen de la congruence indiquée.

Le réseau \(\Lambda\), sa forme bilinéaire, ses
douze plans ambiants et le vecteur \(v_\alpha/3\)
sont ainsi déterminés avec toutes les propriétés
utilisées par la déformation ultérieure. Le
code contient également le mot
\[
 c_*=(1,1,1,1,1,1,2,1,1,1,1,1),
\]
puisque \(q_*=(1,1,1,1,1,1)\) satisfait \(q_*A_W=(2,1,1,1,1,1)\). C'est la donnée de support complet qui permet de construire et de prouver la stabilité de l'isométrie d'ordre trois dans la section~\ref{sec:leech-dualidad}.

\section{Polarisation dodécaphasique, dualité réticulaire et réciprocité thêta}
\label{sec:leech-dualidad}

L'incidence exceptionnelle et la polarisation du registre dodécaphasique proviennent de deux réalisations du même état enrichi. La première conserve les relations de frontière, le code de recollement et l'origine marquée qui déterminent le réseau ; la seconde fait agir une représentation sur les douze coordonnées entières déjà obtenues par les lecteurs du Topological Phase Kernel. Leur composition permet d'étudier une déformation métrique dont le paramètre modifie les longueurs relatives de douze plans et conserve simultanément le volume total et une symétrie ternaire explicite.

Le point de départ causal demeure la construction APP--TRIT--TPK de l'état enrichi et du continu conjoint. Le registre \(K\) est reçu après le registre dodécaphasique et ses opérations d'incidence et d'orientation. La représentation employée agit sur cette sortie et conserve ses douze valeurs comme des données déjà déterminées. La réalisation réticulaire est reçue avec sa forme bilinéaire, ses classes de recollement et son vecteur marqué. Ces structures, développées respectivement dans \cite{hmtX,HMT-IV}, sont maintenues lors de leur composition. La collection construite ci-dessous est une réalisation géométrique de ces données avec une coordinatisation commune à douze positions.

\subsection{Polarisation du registre et marquage des douze plans}
\label{subsec:leech-polarizacion}

Le registre dodécaphasique et sa somme sont
\begin{equation}
 \begin{split}
 K={}&(234,543,140,729,659,824,\\
     &\hspace{13mm}621,58,914,794,146,601),\\
 \sum_{j=1}^{12}K_j={}&6263.
 \end{split}
 \label{eq:leech-K}
\end{equation}
Soit \(A_5\) le groupe des permutations paires de cinq éléments et soit \(C_5=\langle(1\,2\,3\,4\,5)\rangle\). Ses douze classes à gauche déterminent une représentation par permutations \(\varrho:A_5\to O(\mathbb R^{12})\). La coordinatisation marquée est fixée au moyen des représentants suivants, écrits comme listes d'images :
\[
\begin{array}{c|c@{\qquad}c|c}
j&r_j&j&r_j\\ \hline
1&(4,3,5,2,1)&7&(5,4,3,2,1)\\
2&(4,2,1,3,5)&8&(1,3,4,2,5)\\
3&(1,2,4,5,3)&9&(4,2,3,5,1)\\
4&(2,5,3,4,1)&10&(3,2,5,4,1)\\
5&(4,3,1,5,2)&11&(4,5,2,3,1)\\
6&(5,1,2,3,4)&12&(5,3,2,4,1).
\end{array}
\]
Concrètement, \(\varrho(a)e_j=e_\ell\) si \(ar_jC_5=r_\ell C_5\). Le caractère tridimensionnel \(\chi_3\) prend les valeurs
\[
\begin{array}{c|ccccc}
\text{classe}&1&2^2&3&5_a&5_b\\ \hline
\chi_3&3&-1&0&(1+\sqrt5)/2&(1-\sqrt5)/2.
\end{array}
\]
La classe \(5_a\) contient le générateur de \(C_5\) indiqué et \(5_b\) contient son carré. Ce choix distingue les deux caractères tridimensionnels conjugués. Si \(I_{12}\) est l'identité et \(\mathbf1\) est le vecteur de douze uns, on définit
\begin{equation}
 P_3=\frac3{60}\sum_{a\in A_5}\chi_3(a^{-1})\varrho(a),
 \qquad
 P_{11}=I_{12}-\frac1{12}\mathbf1\mathbf1^{\mathsf T}.
 \label{eq:leech-projectors}
\end{equation}
Le corps \(\mathbb Q(\sqrt5)\) apparaît ici dans la réalisation de la
représentation qui agit sur le registre engendré.

\begin{lemma}[Polarisation de somme nulle]
\label{lem:leech-u}
Les opérateurs de \eqref{eq:leech-projectors} satisfont
\[
 P_3=P_3^{\mathsf T}=P_3^2,\qquad
 \operatorname{rank}P_3=3,\qquad P_3\mathbf1=0,\qquad
 P_3P_{11}=P_{11}P_3=P_3.
\]
Par conséquent,
\begin{equation}
 u=P_3P_{11}K=P_3K,\qquad
 \sum_{j=1}^{12}u_j=0,\qquad
 \|u\|^2=\frac{6638585+2275584\sqrt5}{20}>0.
 \label{eq:leech-u-norm}
\end{equation}
\end{lemma}
\begin{proof}
La somme du caractère est le projecteur isotypique de la représentation.
Le caractère réel et la substitution \(a\mapsto a^{-1}\)
donnent sa symétrie. La multiplicité du caractère tridimensionnel dans la
représentation de \(A_5/C_5\) est la dimension de ses invariants par
\(C_5\), qui vaut
\[
 \frac15\left(
 3+2\frac{1+\sqrt5}{2}+2\frac{1-\sqrt5}{2}
 \right)=1.
\]
Le projecteur a donc rang trois. Son orthogonalité au caractère
trivial implique \(P_3\mathbf1=0\) et les identités avec \(P_{11}\).
L’application de la somme finie au registre
\eqref{eq:leech-K} produit les douze coordonnées
\begin{equation}
\begin{aligned}
u={}&(-495/4-233\sqrt5/10,\quad403/4+169\sqrt5/10,\\
&-403/4-169\sqrt5/10,\quad495/4+233\sqrt5/10,\\
&601/4+1031\sqrt5/10,\quad203/4+211\sqrt5/5,\\
&-203/4-211\sqrt5/5,\quad-601/4-1031\sqrt5/10,\\
&313/4+569\sqrt5/10,\quad162+219\sqrt5/20,\\
&-162-219\sqrt5/20,\quad-313/4-569\sqrt5/10).
\end{aligned}
\label{eq:leech-u-coordinates}
\end{equation}
Les additionner et les élever au carré dans \(\mathbb Q(\sqrt5)\) donne \eqref{eq:leech-u-norm}. De manière équivalente,
\[
 \|u\|^2=K^{\mathsf T}P_3K
 =\frac1{20}\sum_{a\in A_5}
          \chi_3(a^{-1})K^{\mathsf T}\varrho(a)K.
\]
Les coefficients positifs de son évaluation prouvent que \(u\ne0\).
\end{proof}

La somme nulle a une fonction géométrique concrète : les expansions de certains plans pourront compenser les contractions des autres dans le déterminant total. La non-nullité de \(u\), en revanche, garantit que cette compensation contient une déformation effective.

L’espace euclidien ambiant de la réalisation réticulaire est
\begin{equation}
 E=\bigoplus_{j=1}^{12}E_j,\qquad
 B=\operatorname{diag}(B_2,\ldots,B_2),\qquad
 B_2=\begin{pmatrix}2&-1\\-1&2\end{pmatrix}.
 \label{eq:leech-ambient}
\end{equation}
Chaque \(E_j\) est le plan réel engendré par le réseau de racines
\(A_2\), exprimé dans sa base de racines simples. La décomposition
orthogonale appartient à l’espace ambiant. Le réseau de Leech
s’obtient par recollement et passage à un voisin, opérations qui relient
les douze facteurs.

Soit \(C_W\subset\mathbb F_3^{12}\) le code ternaire de dimension six
de la construction de Paley--Witt. En identifiant le discriminant de chaque
facteur \(A_2\) à \(\mathbb F_3\), le recollement est
\[
 N=N(C_W)=
 \bigcup_{c\in C_W}\bigl(A_2^{12}+\phi(c)\bigr),
\]
où \(\phi(c)\) choisit des représentants des classes discriminantes.
L’auto-orthogonalité et l’autodualité du code transmettent parité et
intégralité à \(N\) ; son déterminant est
\[
 \det N=\frac{3^{12}}{(3^6)^2}=1.
\]
La construction exceptionnelle fixe l'origine marquée \(j=1\) et le vecteur
\begin{equation}
 v_\alpha=4\rho_1+\rho_2+\cdots+\rho_{12},
 \qquad \rho_j=(1,1)_j,\qquad
 \|v_\alpha\|_B^2=54.
 \label{eq:leech-marked-vector}
\end{equation}
Avec lui, on forme
\begin{equation}
\begin{split}
 N_{\alpha,0}
  &=\{x\in N:\langle x,v_\alpha\rangle_B\equiv0\pmod3\},\\
 \Lambda&=N_{\alpha,0}+\mathbb Z\,v_\alpha/3.
\end{split}
\label{eq:leech-neighbour}
\end{equation}
La réalisation exceptionnelle établit que \(\Lambda\) est un réseau
positif, pair, unimodulaire, sans racines et de rang \(24\), de sorte qu’il est
isométrique au réseau de Leech \cite{HMT-IV}. En particulier,
\(\Lambda^*=\Lambda\), avec le dual relativement à \(B\).
Ces propriétés et l’appartenance \(v_\alpha/3\in\Lambda\) sont les
hypothèses réticulaires précises utilisées dans les théorèmes suivants.

\subsection{Isométrie ternaire et stabilité du voisin}
\label{subsec:leech-c3}

Nous définissons les matrices bidimensionnelles
\[
 C=\begin{pmatrix}0&-1\\1&-1\end{pmatrix},
 \qquad
 R=\begin{pmatrix}0&1\\1&0\end{pmatrix},
 \qquad
 \varpi=\frac13\begin{pmatrix}2\\1\end{pmatrix}.
\]
La multiplication directe vérifie
\[
 C^3=I_2,\quad C^2+C+I_2=0,\quad
 C^{\mathsf T}B_2C=B_2,\quad
 RCR=C^{-1},\quad R^{\mathsf T}B_2R=B_2.
\]
En outre, \(C\varpi-\varpi=(-1,0)^{\mathsf T}\), de sorte que \(C\)
agit trivialement sur \(A_2^*/A_2\). La réflexion \(R\) change le signe
de la classe discriminante. Le code \(C_W\) contient le mot de
support complet
\[
 c_*=(1,1,1,1,1,1,2,1,1,1,1,1).
\]
Dans la présentation \(C_W=\{(q,qA_W):q\in\mathbb F_3^6\}\), ce mot provient de \(q_*=(1,1,1,1,1,1)\) et de \(q_*A_W=(2,1,1,1,1,1)\).

Pour chaque plan, on fixe le repère
\begin{equation}
 Q_j=\begin{cases}
 R,&(c_*)_j=1,\\
 I_2,&(c_*)_j=2,
 \end{cases}
 \qquad
 g_c=\bigoplus_{j=1}^{12}Q_jCQ_j^{-1}.
 \label{eq:leech-c3}
\end{equation}
Les symboles \(Q_j\) désignent des repères de plans ; le symbole \(P_3\) reste réservé au projecteur sur le secteur tridimensionnel.

\begin{proposition}[Symétrie d’ordre trois du réseau marqué]
\label{prop:leech-c3}
La transformation \(g_c\) est une isométrie d’ordre trois de \(E\),
a pour unique vecteur fixe le vecteur nul et satisfait
\[
 g_cN=N,\qquad g_cN_{\alpha,0}=N_{\alpha,0},
 \qquad g_c\Lambda=\Lambda.
\]
\end{proposition}
\begin{proof}
Chaque bloc est \(C\) ou \(C^{-1}\), dont le polynôme minimal est \(X^2+X+1\). Cela prouve l'ordre et l'affirmation sur les vecteurs fixes, ainsi que la conservation de \(B\). Chaque bloc agit trivialement sur le discriminant et conserve donc toutes les classes de recollement de \(N\).

Écrivons \(v=v_\alpha=\sum_j a_j\rho_j\), avec \(a_1=4\) et
\(a_j=1\) pour \(j>1\). Tous les \(a_j\) sont congrus à un
modulo trois. Les identités
\[
 (C-I_2)\rho=(-2,-1)^{\mathsf T},
 \qquad (C^{-1}-I_2)\rho=(-1,-2)^{\mathsf T}
\]
montrent que
\[
 y=\frac{g_cv-v}{3},\qquad z=\frac{g_c^{-1}v-v}{3}
\]
ont pour classes discriminantes \(c_*\) et \(-c_*\), respectivement.
Les deux appartiennent à \(C_W\), donc \(y,z\in N\). Dans chaque facteur,
\[
 \left\langle
 \frac{(C^{\pm1}-I_2)a_j\rho}{3},a_j\rho
 \right\rangle_{B_2}=-a_j^2.
\]
Par sommation, \(\langle y,v\rangle_B=\langle z,v\rangle_B=-27\) ;
c’est pourquoi \(y,z\in N_{\alpha,0}\).
Si \(x\in N_{\alpha,0}\), alors
\[
 \langle g_cx,v\rangle_B
  =\langle x,g_c^{-1}v\rangle_B
  =\langle x,v\rangle_B+3\langle x,z\rangle_B
  \equiv0\pmod3.
\]
La même opération avec \(g_c^{-1}\) démontre l’égalité
\(g_cN_{\alpha,0}=N_{\alpha,0}\). Enfin,
\(g_c(v/3)=v/3+y\) conserve la classe qui engendre
\(\Lambda=N_{\alpha,0}+\mathbb Z(v/3)\).
\end{proof}

La symétrie ternaire agit donc sur le réseau
recollé complet. Sa conservation utilise le mot du code et
la congruence du vecteur marqué, en plus de la rotation de chaque
plan. Cette précision permet de transporter la symétrie au flot sans
perdre l’information arithmétique qui a rendu sa construction possible.

\subsection{Flux anisotrope et conservation du covolume}
\label{subsec:leech-flow}

\begin{definition}[Déformation déterminée par la polarisation]
Pour \(s\in\mathbb R\), paramètre sans dimension, on définit
\begin{equation}
 H_u=\bigoplus_{j=1}^{12}u_jI_2,\qquad
 D_s=\exp(sH_u)=\bigoplus_{j=1}^{12}e^{su_j}I_2,\qquad
 \Lambda_s=D_s\Lambda.
 \label{eq:leech-flow}
\end{equation}
La coordonnée \(u_j\) est attribuée au plan \(E_j\) de la même coordinatisation marquée.
\end{definition}

Le facteur positif \(e^{su_j}\) conserve l'orientation de chaque plan. Les douze étiquettes transmettent la correspondance entre la polarisation et la réalisation réticulaire. Cette correspondance fait partie de la structure de départ ; une permutation des étiquettes est traitée par leur transport simultané dans les deux objets.

\begin{theorem}[Covolume et symétrie du flux]
\label{thm:leech-flow}
Pour tous \(s,t\in\mathbb R\), on a
\[
\begin{gathered}
 D_{s+t}=D_sD_t,\qquad D_0=I,\qquad D_s^{-1}=D_{-s},\\
 D_s^\dagger=D_s,\qquad D_sg_c=g_cD_s,\qquad \det D_s=1,
\end{gathered}
\]
où \(A^\dagger=B^{-1}A^{\mathsf T}B\).
Chaque \(\Lambda_s\) est un réseau discret de rang \(24\) et de covolume
un, préservé par l’isométrie \(g_c\) d’ordre trois.
\end{theorem}
\begin{proof}
Les identités de groupe se vérifient dans chaque bloc exponentiel.
Les blocs scalaires sont autoadjoints pour \(B_2\) et commutent
avec les blocs de \(g_c\). La somme nulle de
\eqref{eq:leech-u-norm} donne
\[
 \det D_s=\prod_{j=1}^{12}e^{2su_j}
 =\exp\left(2s\sum_{j=1}^{12}u_j\right)=1.
\]
L’image d’un réseau par un opérateur inversible est un réseau de
même rang ; son covolume est multiplié par la valeur absolue du
déterminant. Le covolume de \(\Lambda\) est un, de sorte que celui
de \(\Lambda_s\) l’est également. La commutation et
la proposition~\ref{prop:leech-c3} donnent
\[
 g_c\Lambda_s=g_cD_s\Lambda=D_sg_c\Lambda=D_s\Lambda.
\]
L'ordre et l'espace des points fixes de \(g_c\) demeurent inchangés.
\end{proof}

\begin{remark}[Paramètre et symétrie conservée]
La direction \(u\) est déterminée par le registre ; \(s\) parcourt une collection de réalisations. Remplacer \(u\) par \(u/\|u\|\) reparamétrise la même collection. Une interprétation physique sélectionnant une valeur particulière de \(s\) ajoute sa règle de réalisation. Le théorème conserve la symétrie explicite \(g_c\) ; le transport d'autres automorphismes requiert leurs relations correspondantes avec \(H_u\).
\end{remark}

\subsection{Dual euclidien et déformation effective}
\label{subsec:leech-dual}

\begin{definition}[Dual euclidien]
Pour un réseau \(L\subset E\), son dual par rapport à \(B\) est
\[
 L^*=\{y\in E:\langle y,x\rangle_B\in\mathbb Z
                         \text{ pour tout }x\in L\}.
\]
\end{definition}

\begin{theorem}[Réciprocité du dual]
\label{thm:leech-dual}
La collection \eqref{eq:leech-flow} satisfait
\begin{equation}
 \Lambda_s^*=\Lambda_{-s}\qquad(s\in\mathbb R).
 \label{eq:leech-dual}
\end{equation}
Son autodualité en un paramètre \(s\) équivaut à
\(D_{2s}\Lambda=\Lambda\).
\end{theorem}
\begin{proof}
Pour un opérateur inversible \(A\),
\[
 (AL)^*=A^{-\dagger}L^*.
\]
En effet, \(\langle y,Ax\rangle_B=\langle A^\dagger y,x\rangle_B\)
est entier pour tout \(x\in L\) précisément lorsque
\(A^\dagger y\in L^*\). En prenant \(A=D_s\), l’autoadjonction,
l’identité \(D_s^{-1}=D_{-s}\) et \(\Lambda^*=\Lambda\) prouvent
\eqref{eq:leech-dual}. L’égalité
\(D_s\Lambda=D_{-s}\Lambda\) équivaut, en multipliant par \(D_s\),
à \(D_{2s}\Lambda=\Lambda\).
\end{proof}

Le covolume contrôle le volume d'un domaine fondamental ; l'autodualité contrôle tous les produits entiers avec le réseau. La première propriété est continue le long de cette collection. La seconde impose la condition arithmétique exacte du théorème précédent. On emploiera donc l'expression \emph{covolume un} pour le flot euclidien et l'on spécifiera l'intégralité lorsque l'unimodularité arithmétique intervient.

\begin{proposition}[Témoin de perte locale d'intégralité]
\label{prop:leech-nonintegral}
Il existe \(\delta>0\) tel que, si \(0<|s|<\delta\), le réseau \(\Lambda_s\) n'est pas entier et n'est pas isométrique au réseau de Leech.
\end{proposition}
\begin{proof}
Le vecteur \(v_\alpha/3\) appartient à \(\Lambda\). Sa norme
transportée est
\begin{equation}
 f(s)=\left\|D_s\frac{v_\alpha}{3}\right\|_B^2
 =\frac29\left(16e^{2su_1}+
                       \sum_{j=2}^{12}e^{2su_j}\right).
 \label{eq:leech-witness}
\end{equation}
Par \(\sum_j u_j=0\) et par la première coordonnée de
\eqref{eq:leech-u-coordinates},
\begin{equation}
\begin{split}
 f(0)&=6,\\
 f'(0)&=\frac49\left(16u_1+\sum_{j=2}^{12}u_j\right)
       =\frac{20}{3}u_1
       =-825-\frac{466}{3}\sqrt5\ne0.
\end{split}
\label{eq:leech-witness-derivative}
\end{equation}
La continuité de \(f'\) permet de choisir un intervalle \((-\delta,\delta)\) dans lequel \(f\) est strictement monotone. En réduisant \(\delta\), on obtient en outre \(11/2<f(s)<13/2\). Pour \(0<|s|<\delta\), il en résulte \(f(s)\ne6\) ; dans cet intervalle, six est le seul entier. Par conséquent, \(\Lambda_s\) contient un vecteur de norme non entière. Un réseau entier a une norme entière pour tous ses vecteurs. Puisque Leech est entier et que la norme est conservée par les isométries, la seconde affirmation est également démontrée.
\end{proof}

La déformation préserve volume et symétrie ternaire tout en
modifiant effectivement la structure métrique arithmétique. Les couples
opposés de coordonnées de \(u\) définissent une permutation de l’espace
ambiant qui échange dilatations et contractions. Sa promotion
en un automorphisme du réseau recollé exigerait de conserver le code
et le voisin. La réciprocité du dual déjà démontrée utilise
uniquement l’autoadjonction du flot et l’autodualité initiale.

\subsection{Convergence et réciprocité des fonctions thêta}
\label{subsec:leech-theta}

Pour \(s\in\mathbb R\) et \(t>0\), nous définissons
\begin{equation}
 \Theta_s(t)=\sum_{\lambda\in\Lambda}
                \exp\bigl(-\pi t\|D_s\lambda\|_B^2\bigr).
 \label{eq:leech-theta-real}
\end{equation}
C’est la somme gaussienne de toutes les longueurs du réseau déformé,
avec leurs multiplicités. Si \(M=\max_j|u_j|\), on a
\begin{equation}
 e^{-2|s|M}\|x\|_B^2
 \leq\|D_sx\|_B^2
 \leq e^{2|s|M}\|x\|_B^2.
 \label{eq:leech-quadratic-bounds}
\end{equation}
Sur un compact de \(\mathbb R\times(0,\infty)\), la borne
inférieure fournit une gaussienne dominante sur le réseau fixe.
La série converge absolument et uniformément. Les dérivées
d’ordre fini ajoutent des facteurs polynomiaux dans les coordonnées de
\(\lambda\) ; ces facteurs restent dominés par une autre gaussienne.
Les dérivations terme à terme sont, par conséquent,
légitimes sur ces compacts.

\begin{theorem}[Reciprocidad gaussiana]
\label{thm:leech-theta}
Avec le volume euclidien associé à \(B\),
\begin{equation}
 \Theta_s(t)=t^{-12}\Theta_{-s}(t^{-1}),
 \qquad s\in\mathbb R,\quad t>0.
 \label{eq:leech-theta-reciprocity}
\end{equation}
En particulier, \(\Theta_s(1)=\Theta_{-s}(1)\).
\end{theorem}
\begin{proof}
Nous fixons la transformée de Fourier
\[
 \widehat f(y)=\int_E
      f(x)e^{-2\pi i\langle x,y\rangle_B}\,dx.
\]
Pour un réseau \(L\subset E\) et
\(f_t(x)=e^{-\pi t\|x\|_B^2}\), nous périodisons
\[
 F_t(x)=\sum_{\ell\in L}f_t(x+\ell).
\]
La convergence normale de la série et de ses dérivées donne une fonction lisse sur \(E/L\). Si \(\xi\in L^*\), son coefficient de Fourier est
\[
\begin{split}
 c_\xi
 &=\frac1{\operatorname{covol}(L)}
   \int_{E/L}F_t(x)e^{-2\pi i\langle x,\xi\rangle_B}\,dx\\
 &=\frac{\widehat f_t(\xi)}{\operatorname{covol}(L)}.
\end{split}
\]
La seconde égalité décompose \(E\) en translations d’un domaine
fondamental ; le caractère vaut un sur \(L\). Le produit
des intégrales gaussiennes unidimensionnelles, dans une base
orthonormée pour \(B\), donne
\[
 \widehat f_t(\xi)=
 t^{-12}\exp\bigl(-\pi\|\xi\|_B^2/t\bigr).
\]
La série de Fourier est absolument convergente. En l’évaluant en
\(x=0\), nous obtenons la sommation de Poisson
\[
 \sum_{\ell\in L}e^{-\pi t\|\ell\|_B^2}
 =\frac{t^{-12}}{\operatorname{covol}(L)}
      \sum_{\xi\in L^*}e^{-\pi\|\xi\|_B^2/t}.
\]
À présent, \(L=\Lambda_s\) a un covolume égal à un et \(L^*=\Lambda_{-s}\) d'après le théorème~\ref{thm:leech-dual}. La substitution prouve \eqref{eq:leech-theta-reciprocity}. Le cas \(t=1\) est immédiat.
\end{proof}

La périodisation précédente est la démonstration gaussienne de Poisson ;
elle peut être comparée à \cite[teorema 2, pp.~10--11]{Elkies}.
La convention de signe positif dans Fourier conduit à la même
transformée de la gaussienne, du fait de sa parité. L’égalité des thêta
porte sur les séries complètes et conserve le domaine analytique de
convergence.

\begin{proposition}[Extension holomorphe]
\label{prop:leech-theta-complex}
Pour \(\tau\) dans le demi-plan supérieur
\(\mathbb H=\{\tau\in\mathbb C:\operatorname{Im}\tau>0\}\), la série
\begin{equation}
 \theta_s(\tau)=\sum_{\lambda\in\Lambda}
             e^{\pi i\tau\|D_s\lambda\|_B^2}
 \label{eq:leech-theta-complex}
\end{equation}
converge normalement et satisfait
\begin{equation}
 \theta_s(-1/\tau)=(-i\tau)^{12}\theta_{-s}(\tau).
 \label{eq:leech-theta-modular}
\end{equation}
\end{proposition}
\begin{proof}
La valeur absolue de chaque terme est \(e^{-\pi\operatorname{Im}\tau\|D_s\lambda\|_B^2}\). Sur les compacts du demi-plan supérieur, la partie imaginaire possède une borne inférieure positive, si bien que \eqref{eq:leech-quadratic-bounds} prouve la convergence normale et l'holomorphie. Pour \(\tau=it\), avec \(t>0\), \eqref{eq:leech-theta-reciprocity} appliquée en \(t^{-1}\) donne \(\theta_s(i/t)=t^{12}\theta_{-s}(it)\). Les deux membres de \eqref{eq:leech-theta-modular} sont holomorphes dans \(\mathbb H\) ; le théorème d'identité prolonge l'égalité depuis l'axe imaginaire positif.
\end{proof}

L'inversion modulaire relie les paramètres \(s\) et \(-s\). La périodicité supplémentaire sous \(\tau\mapsto\tau+1\) utilise la parité des normes. La proposition~\ref{prop:leech-nonintegral} présente des paramètres arbitrairement proches de zéro pour lesquels l'intégralité euclidienne est perdue. La collection conserve donc la réciprocité analytique démontrée, tandis qu'une réalisation par des formes modulaires scalaires ou des algèbres vertex de réseaux incorpore en outre les hypothèses arithmétiques correspondantes. Dans la réalisation exceptionnelle initiale, ces hypothèses soutiennent le prolongement vers les algèbres vertex et le Monstrous Moonshine \cite{HMT-IV}.

\subsection{Forme déployée et autodualité en dimension quarante-huit}
\label{subsec:leech-split}

La perte d’intégralité d’une projection euclidienne conduit à
considérer simultanément les deux projections. Dans l’espace
\(E\oplus E\), on introduit la forme bilinéaire
\begin{equation}
 \mathcal B((x,p),(x',p'))=
       \langle x,p'\rangle_B+\langle p,x'\rangle_B.
 \label{eq:leech-split-form}
\end{equation}
Leurs produits apparient les deux coordonnées. Nous définissons
\begin{equation}
\begin{split}
 \mathcal D_s&=D_s\oplus D_{-s},\\
 \Gamma_s&=\mathcal D_s(\Lambda\oplus\Lambda)\\
 &=\{(D_s\lambda,D_{-s}\mu):\lambda,\mu\in\Lambda\}.
\end{split}
\label{eq:leech-split-lattice}
\end{equation}

\begin{theorem}[Autodualité entière pour la forme déployée]
\label{thm:leech-split}
La forme \(\mathcal B\) a pour signature \((24,24)\).
Pour tout \(s\in\mathbb R\), \(\mathcal D_s\) est une isométrie de
\(\mathcal B\), et \(\Gamma_s\) est un réseau entier, pair et autodual
relativement à cette forme.
\end{theorem}
\begin{proof}
En introduisant
\[
 p_L=\frac{x+p}{\sqrt2},\qquad
 p_R=\frac{x-p}{\sqrt2},
\]
on obtient
\[
 \mathcal B((x,p),(x,p))=\|p_L\|_B^2-\|p_R\|_B^2.
\]
Comme \(B\) est positive en dimension \(24\), la signature est \((24,24)\). Le caractère autoadjoint et la réciprocité des blocs donnent
\[
 \langle D_sx,D_{-s}p'\rangle_B=\langle x,p'\rangle_B,
\]
et de manière analogue pour le second produit croisé. Par conséquent,
\(\mathcal B(\mathcal D_s\xi,\mathcal D_s\eta)=\mathcal B(\xi,\eta)\).

Dans \(\Gamma_0=\Lambda\oplus\Lambda\), tous les produits croisés sont entiers et
\[
 \mathcal B((\lambda,\mu),(\lambda,\mu))
   =2\langle\lambda,\mu\rangle_B\in2\mathbb Z.
\]
Cela prouve l’intégralité et la parité. Si \((x,p)\) appartient au dual
de \(\Gamma_0\) pour \(\mathcal B\), l’appariement avec
\((\lambda,0)\) et \((0,\mu)\) exige respectivement
\(p,x\in\Lambda^*=\Lambda\). Le dual est donc \(\Gamma_0\).
L’isométrie \(\mathcal D_s\) transporte ces trois propriétés
à \(\Gamma_s\).
\end{proof}

La métrique spécifie le contenu de la conclusion. La collection \(\Gamma_s\) conserve l'intégralité pour les produits croisés de \(\mathcal B\), même lorsque \(\Lambda_s\) perd l'intégralité pour \(B\). En tant que réseaux abstraits à forme déployée, les \(\Gamma_s\) sont isométriques. Leurs deux projections euclidiennes, considérées avec les métriques positives respectives, varient avec le paramètre.

\begin{proposition}[Échange des projections réciproques]
\label{prop:leech-split-involution}
L'involution \(\mathcal J(x,p)=(p,x)\) satisfait
\[
 \mathcal J\Gamma_s=\Gamma_{-s},\qquad
 \mathcal J\mathcal D_s\mathcal J=\mathcal D_{-s}.
\]
Dans les coordonnées \((p_L,p_R)\), elle fixe \(p_L\) et change le signe de \(p_R\).
\end{proposition}
\begin{proof}
L'échange transforme \((D_s\lambda,D_{-s}\mu)\) en \((D_{-s}\mu,D_s\lambda)\), qui parcourt \(\Gamma_{-s}\). La conjugaison des deux blocs donne la seconde identité. Enfin, \(x+p\) est conservé et \(x-p\) change de signe.
\end{proof}

L’involution exprime une dualité de la construction réticulaire
complète. Son éventuelle représentation physique conserve, en plus de
la forme déployée, les opérateurs et les données d’identification
propres à cette réalisation.

\subsection{Énergie réticulaire et équivalence unitaire}
\label{subsec:leech-energy}

L'énergie de rayon réciproque \(m^2/R^2+w^2R^2\) échange ses deux termes sous \((m,w,R)\mapsto(w,m,R^{-1})\), avec une conjugaison unitaire qui préserve les domaines opératoriels \cite{HMT-IV}. La réalisation sur les douze plans marqués emploie
\begin{equation}
 \mathcal E_s(\lambda,\mu)=
       \|D_s\lambda\|_B^2+\|D_{-s}\mu\|_B^2,
 \qquad \lambda,\mu\in\Lambda.
 \label{eq:leech-energy}
\end{equation}
L’énergie est positive et vérifie
\(\mathcal E_s(\lambda,\mu)=\mathcal E_{-s}(\mu,\lambda)\).

\begin{theorem}[Opérateur positif et conjugaison des domaines]
\label{thm:leech-energy}
Dans l’espace de Hilbert \(\ell^2(\Lambda\oplus\Lambda)\), soit
\[
 (H_s\psi)(\lambda,\mu)=
       \mathcal E_s(\lambda,\mu)\psi(\lambda,\mu)
\]
de domaine
\begin{equation}
 \mathcal D(H_s)=
 \left\{\psi\in\ell^2(\Lambda\oplus\Lambda):
 \sum_{\lambda,\mu}\mathcal E_s(\lambda,\mu)^2
             |\psi(\lambda,\mu)|^2<\infty\right\}.
 \label{eq:leech-energy-domain}
\end{equation}
Alors \(H_s\) est positif, auto-adjoint et à résolvante compacte.
L’application \(U\psi(\lambda,\mu)=\psi(\mu,\lambda)\) est unitaire et
\begin{equation}
 U\mathcal D(H_s)=\mathcal D(H_{-s}),\qquad
 UH_sU^{-1}=H_{-s}.
 \label{eq:leech-unitary}
\end{equation}
\end{theorem}
\begin{proof}
La multiplication par une fonction réelle sur un ensemble discret,
avec le domaine maximal \eqref{eq:leech-energy-domain}, est
autoadjointe. On peut le vérifier directement au moyen de la base
orthonormée des vecteurs à support unitaire : l’appartenance
au domaine de l’adjoint exige que la suite obtenue en
multipliant par \(\mathcal E_s\) soit de carré sommable, soit exactement
la même condition. En outre,
\[
 \langle\psi,H_s\psi\rangle
   =\sum_{\lambda,\mu}\mathcal E_s(\lambda,\mu)
                         |\psi(\lambda,\mu)|^2\geq0.
\]
Par \eqref{eq:leech-quadratic-bounds},
\[
 \mathcal E_s(\lambda,\mu)\geq
 e^{-2|s|M}\bigl(\|\lambda\|_B^2+\|\mu\|_B^2\bigr).
\]
Un réseau a un nombre fini de points dans tout ensemble borné ; chaque sous-niveau de \(\mathcal E_s\) est donc fini. Les coefficients diagonaux \((1+\mathcal E_s)^{-1}\) tendent vers zéro en dehors d'ensembles finis, de sorte que \((H_s+I)^{-1}\) est limite en norme d'opérateurs de rang fini et est compact.

L’échange des deux indices conserve la norme de
\(\ell^2(\Lambda\oplus\Lambda)\), de sorte que \(U\) est unitaire
et \(U^{-1}=U\). L’égalité des énergies lors de l’inversion de \(s\) transforme
la condition de domaine de \(H_s\) en celle de \(H_{-s}\), et
l’évaluation sur chaque paire \((\lambda,\mu)\) prouve
\eqref{eq:leech-unitary}.
\end{proof}

Trois conservations de domaines distincts sont ainsi articulées. Dans les plans euclidiens, le flux maintient le covolume et la symétrie ternaire et relie chaque réseau au dual du paramètre opposé. Dans l'espace doublé, les produits croisés conservent un réseau entier, pair et autodual. Dans la réalisation opératorielle, l'échange réciproque conserve l'énergie au moyen d'une équivalence unitaire d'opérateurs avec leurs domaines. Les formules thêta expriment la même réciprocité dans la somme convergente des longueurs de tout le réseau.

% Procedencia primaria (lectura de 18-09-2026):
% XROOT = output/EXTREMA_MEDIA_RAZON_NARRACION_Y_REVISION_20260916/ARTICULO/ES/source
% XROOT/libro/01_acoplamiento.tex:51-147,727-749.
% XROOT/sections/k_direccion_dimensional.tex:54-155.
% XROOT/nuclear/12.tex:9-88.
% VIROOT = output/AMPLIACION_FORMAL_LEAN_SERIE_20260916/delivery/USB_ES_EN_STAGE10_20260917/ES/PAQUETES/06_PARTICULAS_Y_MASAS_ES/documentacion_original/payload/source_es
% VIROOT/sections/05_registro_operador_masa.tex:11-119,193-294.
% VIROOT/sections/02_particula_persistencia.tex:210-235.
% XROOT/nuclear/07.tex:5-25,74-154: Hamiltoniano de transporte eliptico.
% XROOT/nuclear/11.tex:878-979: elevacion cotangente y accion discreta.
% Estatutos: cadena K->alpha->accion->masa = RESULTADO_RECUPERADO;
% carta 1+1+10 de eventos y criterio de entrelazamiento = FORMALIZACION_NUEVA;
% congruencia, espectro generalizado y control racional = pruebas completas aqui.

\section{Congruences positives, lecteurs d’action et dynamique hamiltonienne}
\label{sec:k-mass-bridge}

La relation entre le registre dodécaphasique et la Loi générale des masses possède une composition effective. APP produit les deux feuillets avec leurs résidus et quotients ; TRIT oriente la transition et détermine son régime ; TPK compose sélection, transport et actualisation en conservant route, retenue, frontière, survie et mémoire. Les représentations régionales et le registre à douze positions proviennent conjointement de l'état enrichi et de son prolongement compatible. La clôture numérique et la lecture matérielle partagent cette origine et conservent leurs opérations spécifiques.

L'origine de phase, l'orientation et la base mille étant déclarées, le registre est
\[
 K=(234,543,140,729,659,824,621,58,914,794,146,601).
\]
Le lecteur positionnel utilise
\begin{equation}
 \kappa_{\rm ag}(K)=\frac{\sum_{j=1}^{12}K_j1000^{12-j}}{1000^{12}-1},
 \qquad
 \alpha=\pi_{\rm HMT}+e_{\rm HMT}-\varphi_{\rm HMT}-4-
              \kappa_{\rm ag}(K).
 \label{eq:mass-k-alpha}
\end{equation}
Les valeurs du membre droit sont des représentations internes du même processus coinductif. La formule exprime le lecteur de clôture infinie compatible, avec la frontière de retenue qui le définit.

L'action reçoit ces représentations au moyen des lecteurs
\begin{align}
 H_5={}&\frac12\sqrt{1000\alpha/\varphi}-\alpha+
       \frac9{16}\alpha^2-\frac59\alpha^3+
       \frac7{48}\alpha^4-\frac1{54}\alpha^5,\notag\\
 D_A={}&\exp(-100\pi\alpha/9),\qquad
 \mathcal C_\pi=169\alpha^6+\frac{D_A\alpha^7}{1-D_A\alpha},\notag\\
 \eta_{\rm ret}={}&H_5-\frac{90}{\pi}\mathcal C_\pi,
 \qquad R_{\rm act}=\frac{H_5}{\eta_{\rm ret}}>0.
 \label{eq:mass-k-action}
\end{align}
Ici et dans les formules suivantes, \(\pi,\varphi,e\) abrègent les représentations HMT déjà construites. La norme déterminantale locale \(\Sigma_H=\varphi\alpha^{16}\) fixe la décade \(-34\), et la section \(\hbar_{\rm ret}=\eta_{\rm ret}10^{-34}\mathcal U_S\) appartient à la droite d'action déclarée. Le rapport \(R_{\rm act}\) conserve la relation entre les deux sections d'action lors du changement de leur base dimensionnelle commune.

La route électronique centrale fournit, par ses lecteurs de
déplacement et d’autocorrélation, le triplet \((80,54,6)\). Avec
\[
 \Delta_4=\frac{\pi-e\log\pi}{270}\left(1+\frac\pi{729}\right),
 \qquad
 b_e^{\rm alg}=\frac{\sqrt3}{4}
       \bigl[\exp(\varphi/\pi^2)-22\alpha^3\bigr](1+15\Delta_4),
\]
sa représentation matérielle contient
\begin{equation}
 \kappa_e=80-54\alpha+6\Delta_4,
 \qquad \varepsilon_e^\beta=b_e^{\rm alg}R_{\rm act}^{\kappa_e}.
 \label{eq:mass-k-electron}
\end{equation}
La composition \eqref{eq:mass-k-alpha}--\eqref{eq:mass-k-electron} identifie l'intervention concrète de \(K\) : sa lecture de clôture participe à \(\alpha\), et cette sortie agit ensuite dans le lecteur d'action et dans l'exposant de retour. Le choix central de l'électron relève de l'involution de la coordinatisation APP, dont l'unique point fixe est \((5,5)\). La centralité de cette cellule et son évaluation scalaire sont deux résultats reliés par la route et son registre.

\subsection{Coordonnées réversibles du registre d’événements matériels}

Le registre matériel assigne à une route \(\gamma\) douze événements
orientés \(\mathbf e(\gamma)=(e_0,\ldots,e_{11})\), un pour chaque
fenêtre nonadique du calendrier de cent huit positions. Le même
ordre temporel permet de réaliser ces événements dans le porteur à douze
coordonnées utilisé par \(K\). Cette identification requiert de conserver
l’origine et l’orientation du calendrier.

Soient \(\mathbf1\) le vecteur uniforme, \(P_{11}=I-\mathbf1\mathbf1^*/12\), et \(P_3\) le projecteur tridimensionnel de la coordinatisation incidentielle déclarée. La construction géométrique du registre détermine
\[
 u_K=P_3P_{11}K,\qquad
 \|u_K\|^2=\frac{6638585+2275584\sqrt5}{20}>0,
 \qquad u_K\perp\mathbf1,
\]
et le projecteur transversal
\[
 P_{10}(K)=P_{11}-\frac{u_Ku_K^*}{\|u_K\|^2}.
\]

\begin{proposition}[Décomposition matérielle relative à la direction de \(K\)]
Pour tout vecteur d'événements \(\mathbf e\), l'application
\begin{equation}
 \mathbf e\longmapsto
 \left(s_C=\mathbf1^*\mathbf e,\quad
       a_K=\frac{u_K^*\mathbf e}{\|u_K\|^2},\quad
       v_K=P_{10}(K)\mathbf e\right)
 \label{eq:mass-k-event-chart}
\end{equation}
est inversible sur son image et satisfait
\begin{equation}
 \mathbf e=\frac{s_C}{12}\mathbf1+a_Ku_K+v_K,
 \qquad
 \|\mathbf e\|^2=\frac{|s_C|^2}{12}
          +|a_K|^2\|u_K\|^2+\|v_K\|^2.
 \label{eq:mass-k-event-inverse}
\end{equation}
\end{proposition}
\begin{proof}
Les projecteurs \(\mathbf1\mathbf1^*/12\), \(u_Ku_K^*/\|u_K\|^2\) et \(P_{10}(K)\) sont deux à deux orthogonaux et leur somme est l'identité. Appliquer cette égalité à \(\mathbf e\) démontre l'inversion. L'orthogonalité des trois termes donne l'identité des normes.
\end{proof}

Cette coordinatisation conserve le vecteur complet des événements. Elle conserve donc également les fonctionnelles dépendant de leurs signes, de leurs sommes et de leurs corrélations, une fois effectuée la reconstruction \eqref{eq:mass-k-event-inverse}. Le registre des prédécesseurs des neuf, le décompte d'action et les coordonnées de tour demeurent dans la trajectoire enrichie : la coordinatisation précédente ne modifie que les coordonnées des douze événements. Le caractère est évalué sur la même signature récupérée. La construction fournit ainsi une représentation géométrique de l'information utilisée par la masse, avec une direction distinguée par \(K\) et dix composantes transversales.

La quantité d'information retenue par chaque représentation peut être démontrée exactement. Sur le domaine algébrique des registres, \(K'=K+\mathbf1\) vérifie
\[
 u_{K'}=u_K,\qquad P_{10}(K')=P_{10}(K),\qquad
 \kappa_{\rm ag}(K')-\kappa_{\rm ag}(K)=\frac1{999}.
\]
La dernière identité résulte de la sommation de la progression géométrique des douze poids positionnels. Si les représentations régionales sont maintenues fixes, la lecture de \(\alpha\) varie de \(-1/999\). Ce contrôle algébrique détermine la portée exacte de la direction : son projecteur conserve une polarisation, tandis que la clôture scalaire utilise en outre l'information uniforme et positionnelle du registre. L'admissibilité de \(K'\) comme autre historique terminal requiert sa propre construction APP--TRIT--TPK.

\subsection{Congruence positive et spectre de l'opérateur matériel}

Pour une multisection \(M\), soit \(\mathscr H_M\) l'espace fini muni d'une base orthonormale étiquetée par ses chemins. Le lecteur de retour \(r\in\{D,\Omega\}\) produit un opérateur autoadjoint \(B_M^r e_\gamma=\kappa_\gamma^r e_\gamma\). Le caractère raffiné, évalué sur les différences par rapport à l'électron, donne l'opérateur basal
\[
 M_M^{(0)}e_\gamma=
 m_e^{(0)}\chi_{\rm ref}(\sigma_\gamma-\sigma_e,
                  \eta_\gamma-\eta_e)e_\gamma.
\]
Les coefficients de tour et leur convergence relèvent du lecteur de chaque route. Leur représentation bêta est la congruence
\begin{equation}
 D_M^r=R_{\rm act}^{B_M^r/2},\qquad
 M_M^{\beta,r}=D_M^rM_M^{(0)}D_M^r.
 \label{eq:mass-k-congruence}
\end{equation}

\begin{proposition}[Invariants exacts de la congruence matérielle]
L’opérateur \(D_M^r\) est positif et inversible. La congruence
\eqref{eq:mass-k-congruence} conserve le rang et l’inertie ; en particulier,
elle conserve la positivité. Si \(M_M^{(0)}\) et \(B_M^r\) sont diagonaux
dans la base des routes, alors
\[
 m_\gamma^\beta=m_\gamma^{(0)}R_{\rm act}^{\kappa_\gamma^r}.
\]
Pour deux routes dans la même coordinatisation dimensionnelle, on obtient
\[
 \frac{m_Y^\beta}{m_X^\beta}=
 \chi_{\rm ref}(\sigma_Y-\sigma_X,\eta_Y-\eta_X)
 R_{\rm act}^{\kappa_Y^r-\kappa_X^r}.
\]
\end{proposition}
\begin{proof}
Le calcul fonctionnel donne \(D_M^r=\exp((\log R_{\rm act})B_M^r/2)\), de spectre strictement positif. Son inverse est l'exponentielle de l'opposé. Pour tout \(v\), la forme quadratique transformée vaut \((D_M^rv)^*M_M^{(0)}(D_M^rv)\). Le changement inversible de variable préserve les dimensions maximales des sous-espaces positifs, négatifs et nuls ; celles-ci déterminent l'inertie et le rang. Dans une base diagonale commune, les deux facteurs apportent conjointement \(R_{\rm act}^{\kappa_\gamma^r}\). Diviser deux valeurs annule la section dimensionnelle commune et utilise la loi multiplicative du caractère.
\end{proof}

La distinction entre transport d’une forme et modification d’un
opérateur relativement à une métrique fixe peut être vérifiée sans approximations.
Pour
\[
 M_0=\begin{pmatrix}1&0\\0&4\end{pmatrix},\qquad
 D=\begin{pmatrix}2&0\\0&1/2\end{pmatrix},\qquad
 M'=D^*M_0D=\begin{pmatrix}4&0\\0&1\end{pmatrix},
\]
le déterminant de \(D\) est un et les valeurs associées aux deux
routes s’échangent. Même le spectre non ordonné change si
\(M_0=I\) : il se transforme de \(\{1,1\}\) en \(\{4,1/4\}\).
Les deux contrôles préservent positivité, rang et inertie.

\begin{proposition}[Invariance spectrale avec métrique transportée]
Soient \(G_0>0\) une métrique et \(D\) inversible. Si
\[
 M'=D^*M_0D,\qquad G'=D^*G_0D,
\]
les paires \((M',G')\) et \((M_0,G_0)\) ont le même spectre
généralisé, avec les mêmes multiplicités.
\end{proposition}
\begin{proof}
On a
\[
 \det(M'-\lambda G')=|\det D|^2\det(M_0-\lambda G_0),
 \qquad (G')^{-1}M'=D^{-1}G_0^{-1}M_0D.
\]
La première égalité préserve les racines et leurs multiplicités ; la
seconde exhibe une similitude des opérateurs associés.
\end{proof}

Le flot diagonal positif induit par \(K\) possède donc deux réalisations mathématiques précises sur une forme matérielle. Lorsqu'il transporte simultanément la forme et sa métrique, il décrit le même objet dans une autre coordinatisation. Lorsqu'il agit sur la forme en conservant la métrique des routes, il produit une collection opératorielle pouvant modifier les masses relatives. Son identification physique est déterminée par le lecteur des routes et par l'action déclarée.

\subsection{Compatibilité du flux dodécaphasique avec le retour matériel}

Soit
\[
 A_K=(\log\rho_A)\operatorname{diag}(\widehat u_{K,1},\ldots,
       \widehat u_{K,12}),\quad
 \widehat u_K=u_K/\|u_K\|,\quad \rho_A=\pi/\varphi^2,
\]
de sorte que \(g_K(t)=e^{tA_K}\). La collection de retour matériel dans la multisection \(M\) a pour générateur
\[
 C_M^r=\tfrac12(\log R_{\rm act})B_M^r,\qquad
 S_M^r(t)=e^{tC_M^r}.
\]

\begin{proposition}[Transport de la congruence par un lecteur d'entrelacement]
Une application linéaire déclarée \(J_M:\mathbb C^{12}\to\mathscr H_M\) qui satisfait
\begin{equation}
 C_M^rJ_M=J_MA_K
 \label{eq:mass-k-intertwiner}
\end{equation}
satisfait
\[
 S_M^r(t)J_M=J_Mg_K(t),
\]
et, pour \(M_M(t)=S_M^r(t)M_M^{(0)}S_M^r(t)\),
\begin{equation}
 J_M^*M_M(t)J_M
 =g_K(t)^*\bigl(J_M^*M_M^{(0)}J_M\bigr)g_K(t).
 \label{eq:mass-k-pullback}
\end{equation}
\end{proposition}
\begin{proof}
L’identité \eqref{eq:mass-k-intertwiner} s’itère à chaque puissance
et se somme dans la série exponentielle convergente. Prendre les adjoints et
substituer dans la forme matérielle démontre
\eqref{eq:mass-k-pullback}.
\end{proof}

Le critère est fini et admet une épreuve de résistance directement opératorielle. Dans les bases diagonales déclarées, si \(a_j\) et \(c_\gamma\) sont les valeurs propres de \(A_K\) et \(C_M^r\), respectivement, on doit avoir
\[
 (c_\gamma-a_j)(J_M)_{\gamma j}=0
 \qquad\text{pour chaque }(\gamma,j).
\]
Chaque coefficient non nul du lecteur exige l'égalité correspondante des valeurs propres. La composition récupérée de \eqref{eq:mass-k-alpha}--\eqref{eq:mass-k-electron} et la coordinatisation réversible \eqref{eq:mass-k-event-chart} sont des résultats effectifs. L'identification supplémentaire de leurs deux flots requiert un lecteur \(J_M\) possédant la propriété \eqref{eq:mass-k-intertwiner} ; cette proposition détermine exactement ce que doit vérifier une réalisation qui la soutient.

\subsection{Transport elliptique et transformation régulière de Legendre}

La réalisation variationnelle du transport possède un antécédent explicite dans l'ellipse d'action. La paire angulaire déjà exprimée détermine \(\eta=\operatorname{artanh}(C^*/A)\), dans la chambre \(A\ne0\), \(|C^*/A|<1\). Les coordonnées \((Q,P)\), toutes deux de dimension racine d'action, réalisent la forme
\[
 \mathcal I_\eta(Q,P)=\frac12\bigl(e^{-\eta}Q^2+e^\eta P^2\bigr),
 \qquad \lambda=P\,dQ,\quad \Omega=dQ\wedge dP.
\]
Le transport entre formes, avec \(M_\eta=\operatorname{diag}(e^{\eta/2},e^{-\eta/2})\) et \(J_2=\left(\begin{smallmatrix}0&-1\\1&0\end{smallmatrix}\right)\), est \(M_{\eta'}e^{-\vartheta J_2}M_\eta^{-1}\). Une interpolation \(\eta(t)\), \(\vartheta(t)\), avec \(\dot\vartheta=\omega(t)>0\), possède l'hamiltonien récupéré de l'article X :
\begin{equation}
 H_{\rm tr}(Q,P,t)=\omega(t)\mathcal I_{\eta(t)}(Q,P)
                 +\frac{\dot\eta(t)}2QP.
 \label{eq:mass-k-hamiltonian}
\end{equation}
Le second terme provient de \(\dot M_\eta M_\eta^{-1}\). Les équations
\[
 \dot Q=\omega e^\eta P+\frac{\dot\eta}2Q,
 \qquad
 \dot P=-\omega e^{-\eta}Q-\frac{\dot\eta}2P
\]
conservent \(\mathcal I_{\eta(t)}(Q(t),P(t))\) : la variation explicite de \(\eta\) et la contribution du terme de connexion s'annulent. Ce résultat et l'identité \(P\dot Q-H_{\rm tr}=p\dot q-\omega(q^2+p^2)/2\), pour \(Q=e^{\eta/2}q\), \(P=e^{-\eta/2}p\), sont démontrés dans la section de conservation variationnelle de l'ouvrage sur l'extrême et moyenne raison holographique.

\begin{proposition}[Lagrangien régulier du transport elliptique]
Supposons \(\eta\in C^2\), \(\omega\in C^1\) et \(\omega>0\).
La transformation de Legendre par rapport à \(P\) est inversible et
produit
\begin{equation}
 L_{\rm tr}(Q,\dot Q,t)=
 \frac{\bigl(\dot Q-\dot\eta Q/2\bigr)^2}{2\omega e^\eta}
 -\frac{\omega e^{-\eta}}2Q^2.
 \label{eq:mass-k-lagrangian}
\end{equation}
\end{proposition}
\begin{proof}
Écrivons \(a=\dot\eta/2\), \(b=\omega e^\eta>0\) et
\(c=\omega e^{-\eta}\). Alors
\(H_{\rm tr}=bP^2/2+aQP+cQ^2/2\), et
\[
 \dot Q=\partial_PH_{\rm tr}=bP+aQ,
 \qquad P=\frac{\dot Q-aQ}{b}.
\]
En substituant cette expression dans \(P\dot Q-H_{\rm tr}\), les
termes contenant \(P\) se regroupent en
\(P(\dot Q-aQ)-bP^2/2=(\dot Q-aQ)^2/(2b)\), ce qui
démontre \eqref{eq:mass-k-lagrangian}. De plus,
\[
 \partial_{\dot Q}L_{\rm tr}=P,
 \qquad
 \partial_P^2H_{\rm tr}=\omega e^\eta>0,
 \qquad
 \partial_{\dot Q}^2L_{\rm tr}=\frac1{\omega e^\eta}>0.
\]
La transformation est régulière et ses équations d’Euler--Lagrange
équivalent aux équations hamiltoniennes précédentes.
\end{proof}

Cette élimination algébrique réunit le Hamiltonien déjà construit et
sa forme lagrangienne du second ordre. La régularité provient du
terme quadratique de la réalisation elliptique. Pour un relèvement
cotangent purement linéaire \(H_K(q,p)=p^{\mathsf T}A_Kq\), l’équation
\(\dot q=A_Kq\) est indépendante de \(p\) et le Hessien
\(\partial_p^2H_K\) est nul. Sa formulation variationnelle utilise
l’action du premier ordre \(\int p^{\mathsf T}(\dot q-A_Kq)\,dt\).
Le corpus contient également son antécédent discret
\(L_n=p_{n+1}^{\mathsf T}(q_{n+1}-U_nq_n)\), dont les équations
produisent le transport \(q_{n+1}=U_nq_n\) et le transport dual
\(p_{n+1}=U_n^{-\mathsf T}p_n\). Les deux formulations possèdent
des domaines et des opérateurs déterminés. L’interpolation temporelle et la
section d’action sont conservées comme parties de la réalisation physique
déclarée.

% Incorporación aditiva. No modifica las fuentes ni el PDF de 211 páginas.
% Procedencia, dependencias y alcance en preservation_map.json.
\section{Représentation grassmannienne des lecteurs et formes à valeurs dans la droite d’action}
\label{sec:compat-lectores-formas}

Les réalisations positives conservent une information qui peut être réutilisée
par les lecteurs physiques. Cette conservation possède une formulation
précise : une représentation marquée du registre permet de reconstruire ses
coordonnées et de transporter, par composition, chaque fonction qui les utilise.
La géométrie représente un résultat de la construction ; son choix de
coordonnées conserve la provenance de ce résultat.

Le domaine de départ est formé des états produits par APP, TRIT et TPK dans les sections précédentes. APP conserve ses feuillets additif et multiplicatif au moyen du résidu et du quotient ; TRIT sélectionne le régime local avec son orientation ; la composition \(U_t=\operatorname{Upd}_t\circ\operatorname{Tra}_t\circ
\operatorname{Sel}_t\) met à jour la position, le registre et la mémoire. Le prolongement
\[
 w_6\longrightarrow w_{12}\longrightarrow w_{18}\longrightarrow
 w_{24}\longrightarrow w_{30}\longrightarrow R_{36}\longrightarrow G_9
\]
conserve le préfixe et distingue retour de phase et retour de l'état. Les réalisations spectrale, solénoïdale, de jauge, cohomologique et déterminantale demeurent des opérations corrélatives de la structure discrète conjointe du continu. L'extraction du registre et les représentations régionales s'effectuent sur ce même domaine. Dans ce qui suit, ce sont toutes des sorties HMT précédemment construites.

\subsection{La collection marquée et son inverse}

Soit \(L\geq2\), et soit \(K=(K_1,\ldots,K_L)\) un registre positif de charge \(Q=\sum_jK_j\). Dans l'application dodécaphasique, \(L=12\) et \(Q=6263\). Nous définissons
\[
 d_j=K_j/Q,\qquad t_0=0,\qquad
 t_j=\sum_{a=1}^{j}d_a,\qquad
 B(K)=\begin{pmatrix}1&\cdots&1\\t_0&\cdots&t_L\end{pmatrix}.
\]
La collection \(\mathcal G_L^{\rm aff}\) est l'image de ces plans de lignes dans \(\operatorname{Gr}_{>0}(2,L+1)\), en conservant les colonnes ordonnées \(0,\ldots,L\), les marques extrêmes et la charge \(Q\). La normalisation affine fixe \(t_0=0,t_L=1\). On note \(\Delta_{ij}\) le mineur des colonnes \(i,j\) de tout représentant du plan.

\begin{proposition}[Inverse de la réalisation marquée]
\label{prop:compat-inversa-marcada}
L'application \(K\mapsto([B(K)],Q)\) est injective. Son inverse sur l'image est donné par
\begin{equation}
 K_j=Q\frac{\Delta_{j-1,j}}{\Delta_{0L}},
 \qquad
 t_j=\frac{\Delta_{0j}}{\Delta_{0L}},
 \quad 1\leq j\leq L.
 \label{eq:compat-inversa-plucker}
\end{equation}
Les quotients sont invariants sous tout changement inversible de base
des deux lignes.
\end{proposition}
\begin{proof}
Dans le représentant déclaré, \(\Delta_{ij}=t_j-t_i\), et \(\Delta_{0L}=1\). Les formules récupèrent les différences consécutives et leurs sommes. Si l'on remplace \(B\) par \(AB\), avec \(A\) inversible, tous les mineurs sont multipliés par \(\det A\) ; leurs quotients restent égaux. La charge récupère l'échelle antérieure à la normalisation. Par conséquent, deux registres ayant les mêmes données marquées ont les mêmes coordonnées.
\end{proof}

La condition d'appartenance à \(\mathcal G_L^{\rm aff}\) importe. Les quotients consécutifs d'un plan positif arbitraire peuvent ne pas satisfaire l'identité \(\sum_j\Delta_{j-1,j}=\Delta_{0L}\). Dans la collection construite, cette identité est télescopique. L'inverse agit sur l'image marquée qui vient d'être définie.

\subsection{Algèbre des lecteurs et changements de coordinatisation}

Soit \(\mathcal D\) l'image d'une représentation de l'état enrichi de la forme \(x\mapsto(K(x),\xi(x))\). Le symbole \(\xi\) réunit exactement les autres coordonnées utilisées par le lecteur : représentations régionales, origine et orientation, route, signature, conditions de frontière et sections dimensionnelles, selon l'objet considéré. On ne suppose pas ces coordonnées indépendantes. On définit
\[
 \mathcal F:\mathcal D\longrightarrow\widetilde{\mathcal G},
 \qquad (K,\xi)\longmapsto([B(K)],Q,\xi),
 \qquad \widetilde{\mathcal G}:=\mathcal F(\mathcal D).
\]
La proposition précédente fournit l’inverse \(\mathcal I\) de
\(\mathcal F\) sur son image.

\begin{theorem}[Représentation fidèle de l'algèbre des lecteurs retenus]
\label{thm:compat-algebra-lectores}
Pour chaque lecteur \(\mathcal R:\mathcal D\to V\), de codomaine
déclaré \(V\), il existe un unique lecteur géométrique
\(\mathcal R^{\rm geom}:\widetilde{\mathcal G}\to V\) tel que
\[
 \mathcal R^{\rm geom}=\mathcal R\circ\mathcal I,
 \qquad \mathcal R^{\rm geom}\circ\mathcal F=\mathcal R.
\]
Pour des lecteurs scalaires, cette correspondance préserve les sommes, les produits, l'unité et la conjugaison. C'est un isomorphisme de l'algèbre des fonctions sur \(\mathcal D\) avec l'algèbre correspondante sur \(\widetilde{\mathcal G}\).
\end{theorem}
\begin{proof}
La composition est définie parce que \(\mathcal I\mathcal F\) est
l’identité. Si deux lecteurs géométriques ont la même composition
avec \(\mathcal F\), ils coïncident sur son image, qui est tout
\(\widetilde{\mathcal G}\). Les opérations algébriques se vérifient
point par point ; par exemple,
\((\mathcal R_1\mathcal R_2)\circ\mathcal I=
(\mathcal R_1\circ\mathcal I)(\mathcal R_2\circ\mathcal I)\).
La composition inverse est la restriction par \(\mathcal F\).
\end{proof}

La fidélité concerne la représentation \((K,\xi)\). Pour un observable \(\mathcal O\) de l'état enrichi complet, le critère de descente est que \(\mathcal O(x)=\mathcal O(x')\) dès que \((K(x),\xi(x))=(K(x'),\xi(x'))\). La nécessité s'obtient par composition ; la suffisance permet de définir la valeur sur une fibre à partir de n'importe lequel de ses représentants. Ce critère conserve explicitement l'information qui doit accompagner le registre.

Soit maintenant \(\mathcal T\) une triangulation du polygone étiqueté déterminant une coordinatisation de Plücker positive. Ses coordonnées de frontière et une normalisation homogène sont incluses, par exemple \(\Delta_{0L}=1\). Si l'on échange la diagonale \(ac\) contre \(bd\), avec \(a<b<c<d\), le changement de coordinatisation est
\begin{equation}
 \Delta_{bd}=
 \frac{\Delta_{ab}\Delta_{cd}+\Delta_{ad}\Delta_{bc}}{\Delta_{ac}}.
 \label{eq:compat-cambio-carta}
\end{equation}
Le dénominateur est positif et la relation de Plücker prouve que les deux coordinatisations évaluent le même point.

\begin{corollary}[Compatibilité des lecteurs avec les mutations]
\label{cor:compat-mutaciones}
Soient \(r_{\mathcal T}\) et \(r_{\mathcal T'}\) les expressions de \(\mathcal R^{\rm geom}\) dans deux coordinatisations conservant les marques, et \(\mu_{\mathcal T\mathcal T'}\) leur changement positif. Alors
\[
 r_{\mathcal T'}\circ\mu_{\mathcal T\mathcal T'}
 =r_{\mathcal T}
\]
sur le domaine commun. Une suite fermée de changements de coordinatisation qui conserve le point marqué conserve tous ces lecteurs.
\end{corollary}
\begin{proof}
Les deux reconstructions au moyen de \eqref{eq:compat-inversa-plucker} récupèrent le même \(K\). Les coordonnées \(\xi\) sont transportées avec leurs marques. Appliquer \(\mathcal R\) prouve l'égalité ; l'itérer prouve l'affirmation concernant une suite fermée.
\end{proof}

Les lecteurs peuvent s'exprimer dans des coordinatisations distinctes tout en restant des fonctions du même objet. Une permutation des étiquettes agit en outre sur la marque positionnelle ; une transformation qui change le point agit sur l'argument du lecteur. Les deux opérations possèdent leur loi de transport spécifique.

L’action diagonale par colonnes fournit un contrôle explicite.
Pour \(\lambda_i>0\),
\(\Delta_{ij}\mapsto\lambda_i\lambda_j\Delta_{ij}\).
Le birapport \(\Delta_{ab}\Delta_{cd}/
(\Delta_{ad}\Delta_{bc})\) est conservé, tandis que
\(\Delta_{j-1,j}/\Delta_{0L}\) reçoit le facteur
\(\lambda_{j-1}\lambda_j/(\lambda_0\lambda_L)\).
Avec \(t=(0,1/4,1/2,1)\) et \(\lambda=(1,2,1,1)\), les
séparations reconstruites seraient \((1/2,1/2,1/2)\), dont la somme
est \(3/2\). Les seuls birapports conservent un quotient différent
de celui exigé par l’inverse marqué.

\subsection{Composition explicite avec l'action et la réponse constitutive}

Dans la coordinatisation dodécaphasique, on écrit \(b=1000\) et
\begin{equation}
 \kappa^{\rm geom}=
 \frac{Q}{b^{12}-1}\sum_{j=1}^{12}
 b^{12-j}\frac{\Delta_{j-1,j}}{\Delta_{0,12}},
 \qquad
 \alpha^{\rm geom}=\pi_{\rm HMT}+e_{\rm HMT}
 -\varphi_{\rm HMT}-4-\kappa^{\rm geom}.
 \label{eq:compat-alpha-geometrica}
\end{equation}
Ce sont les expressions géométriques du lecteur de clôture \eqref{eq:mass-k-alpha}. Les trois représentations régionales et \(K\) partagent une génération coinductive ; l'expression utilise leur ordre constructif interne. Le lecteur d'action \eqref{eq:mass-k-action}, évalué en \(\alpha^{\rm geom}\), détermine \(\eta_{\rm ret}^{\rm geom}\), et donc
\[
 \hbar_{\rm ret}^{\rm geom}
 =\eta_{\rm ret}^{\rm geom}10^{-34}\mathcal U_S
 =\hbar_{\rm ret}.
\]
L’égalité a lieu dans la droite d’action \(\mathcal L_S\)
avec sa section \(\mathcal U_S\) ; elle conserve la décade provenant
du lecteur déterminantal \(\varphi_{\rm HMT}\alpha^{16}\).
La composition ne modifie pas la section et ne l’applique pas une seconde fois.

Dans la coordinatisation angulaire relevée qui conserve l'origine et la branche réelle de la représentation, la paire angulaire ultérieure détermine
\[
 x=\frac{\pi_{\rm HMT}}{180}(1000\alpha),\qquad
 y=\frac{\pi_{\rm HMT}}{90}(\eta_{\rm ret}+\alpha),\qquad
 q_\pm=e^{-x\mp y}.
\]
La branche fait partie des marques retenues : la réduction d'un angle modulo \(360\) à elle seule conserve une autre information. Les égalités qui suivent emploient le relèvement déclaré, dans lequel la première coordonnée est \(1000\alpha\).
Dans la chambre \(x>|y|\), avec des projecteurs de feuillet fixés \(P_\pm\), soit \(T=q_+P_++q_-P_-\). La réponse du corpus compare les incidences \(90\) et \(120\) au moyen de
\[
 \mathcal V=(I-T^{90})(I-T^{120})^{-1}
 =r_+P_++r_-P_-,\qquad
 r_\pm=f(q_\pm^{30}),\quad
 f(s)=\frac{1+s+s^2}{1+s+s^2+s^3}.
\]
La règle vérifie
\(\mathcal V\Sigma_4(T^{30})=\Sigma_3(T^{30})\), où
\(\Sigma_j(S)=\sum_{a=0}^{j-1}S^a\). Son unicité découle
de l’inversibilité de \(\Sigma_4(T^{30})\).
Les quatre lectures normalisées sont
\[
 \widehat\varepsilon=r_+^2,\quad
 \widehat\mu=r_-^2,\quad
 \widehat Z=r_-/r_+,\quad
 \widehat c=(r_+r_-)^{-1}.
\]
Dans les droites dimensionnelles d’action, de charge et de vitesse, elles se réalisent comme
\[
 Z=\widehat Z\,u_Su_Q^{-2},\quad c=\widehat c\,u_C,
 \qquad
 \mu=\widehat\mu\,u_Su_C^{-1}u_Q^{-2},\quad
 \varepsilon=\widehat\varepsilon\,u_S^{-1}u_C^{-1}u_Q^2.
\]
Les identités \(\mu\varepsilon=c^{-2}\) et \(\mu/\varepsilon=Z^2\) s'obtiennent par produit et quotient. Le corollaire des changements de coordinatisation conserve cette réponse complète, avec les marques de feuillet et les sections dimensionnelles incluses dans \(\xi\). Le même argument s'applique à un lecteur de masse dont la route, la signature et l'opérateur de retour sont retenus, sur le domaine de la section~\ref{sec:k-mass-bridge}.

\begin{theorem}[Récupération du registre à partir de la réponse étiquetée]
\label{thm:compat-recuperacion-constitutiva}
Les représentations régionales, l'orientation des feuillets, la base \(b=1000\), l'origine, la longueur douze et le relèvement angulaire déclaré étant fixés, l'application \(K\mapsto(\widehat\varepsilon,\widehat\mu)\) est injective sur les registres entiers de son image physique satisfaisant \(0\leq K_j<b\) et la chambre contractive indiquée. Son inverse exact est la composition des opérations suivantes :
\begin{align*}
 r_+&=\sqrt{\widehat\varepsilon},&
 r_-&=\sqrt{\widehat\mu},\\
 r_\pm s_\pm^3+(r_\pm-1)(1+s_\pm+s_\pm^2)&=0,
 &0<s_\pm<1,\\
 x&=-\frac1{60}\log(s_+s_-),&
 \alpha&=\frac{180x}{1000\pi_{\rm HMT}},\\
 \kappa&=\pi_{\rm HMT}+e_{\rm HMT}-\varphi_{\rm HMT}-4-\alpha,
 &N&=(b^{12}-1)\kappa,\\
 K_j&=\left\lfloor N/b^{12-j}\right\rfloor\bmod b.
\end{align*}
\end{theorem}
\begin{proof}
Les racines positives récupèrent les coefficients étiquetés. La dérivée
\[
 f'(s)=-\frac{s^2(3+2s+s^2)}{(1+s+s^2+s^3)^2}<0
 \qquad(0<s<1)
\]
et ses limites \(1\) et \(3/4\) prouvent que l'équation cubique possède exactement une racine dans cet intervalle. Puisque \(s_\pm=e^{-30x\mp30y}\), leur produit est \(e^{-60x}\). On récupère \(x\), \(\alpha\) et \(\kappa\). Dans l'image considérée, \(N=\sum_jK_jb^{12-j}\) est un entier compris entre zéro et \(b^{12}-1\). La division euclidienne successive récupère de manière unique ses douze chiffres, y compris les zéros initiaux.
\end{proof}

La précision nécessaire à une récupération numérique peut également être exprimée. Avec des représentations régionales exactes, une erreur \(|\widetilde\alpha-\alpha|<1/[2(b^{12}-1)]\) implique \(|\widetilde N-N|<1/2\) ; l'arrondi à l'entier le plus proche restitue \(N\). Si les représentations régionales comportent des erreurs, la même condition s'applique à la somme des erreurs de \(\pi+e-\varphi-\alpha\). L'injectivité exacte et la résolution métrologique sont ainsi des propriétés distinctes.

Le domaine entier est essentiel. Dans un prolongement analytique des registres positifs, la perturbation de trois positions consécutives proportionnelle à \((1,-(b+1),b)\) conserve à la fois la somme et la lecture positionnelle, puisque \(1-(b+1)+b=0\) et \(b^2-b(b+1)+b=0\). Pour des perturbations suffisamment petites, elle conserve la positivité et change le point grassmannien. L'injectivité précédente utilise la condition de douze chiffres entiers ; la collection analytique continue possède d'autres fibres.

\subsection{Raffinement et lecteurs hérités}

Si chaque séparation \(d_j\) est subdivisée en séparations positives \(d_{j,a}\) avec \(\sum_ad_{j,a}=d_j\), les extrémités originales conservent leur position. Soient \(\rho\) l'opération consistant à sommer les successeurs de chaque intervalle et \(H\) l'inclusion de ses extrémités dans le nouvel ensemble ordonné. Alors
\[
 \mathcal I_{\rm heredada}(B(\widetilde d),Q)
 =\bigl(Q\sum_ad_{j,a}\bigr)_j=K,
 \qquad
 \mathcal R_{\rm heredada}=\mathcal R(K,\xi).
\]
L’égalité prouve la commutativité entre reconstruction héritée et
subdivision. Dans une suite de subdivisions, les sommes s’associent ;
le résultat persiste dans chaque composition finie. Cela conserve le
lecteur de longueur douze et ses marques. Un lecteur défini sur les
nouveaux blocs aura son propre domaine et sa loi de comparaison.

\subsection{Formes canoniques à valeurs dans une droite dimensionnelle}

Soit \(P\) l’un des polytopes de rang un construits dans le
manuscrit et \(P=\bigcup_\sigma P_\sigma\) une subdivision
orientée du même support, de multiplicité un et à faces compatibles.
Nous écrivons \(\Omega_\sigma\) et \(\Omega_P\) pour leurs formes
canoniques. Leurs identités sont
\(\Omega_P=\sum_\sigma\Omega_\sigma\) et
\(\operatorname{Res}_F\Omega_P=\Omega_F\), avec la convention
d’orientation fixée dans le développement amplituédrique.
La droite d’action \(\mathcal L_S\) est considérée comme constante sur
ce polytope ; sa section provient de l’état HMT fixé.

\begin{theorem}[Transport de la forme par une section d'action]
\label{thm:compat-forma-accion}
Pour \(a\in\mathcal L_S\), constant par rapport aux coordonnées du polytope,
\[
 \boldsymbol\Omega_P:=a\otimes\Omega_P
 =\sum_\sigma a\otimes\Omega_\sigma,
 \qquad
 \operatorname{Res}_F\boldsymbol\Omega_P=a\otimes\Omega_F.
\]
Les identités se conservent sous des subdivisions successives et des résidus
itérés sur des drapeaux de faces. Elles s’appliquent en particulier à
\(a=\hbar_{\rm ret}\).
\end{theorem}
\begin{proof}
Le produit tensoriel avec une droite est linéaire. Le résidu agit sur le facteur différentiel et commute avec le coefficient constant. La compatibilité des subdivisions et des résidus itérés est transportée par les mêmes opérations linéaires. Dans un changement de base dimensionnelle \(u'_S=\lambda\,u_S\), avec \(\lambda>0\), le coefficient numérique change selon \(\lambda^{-1}\) ; le tenseur et ses résidus restent égaux.
\end{proof}

\begin{theorem}[Condition de recollement des coefficients de frontière]
\label{thm:compat-pegado}
Soient \(a_\sigma\) des sections de la même droite, régulières dans un
voisinage des faces et méromorphes dans une réalisation ambiante commune.
Supposons que leurs produits avec les formes canoniques aient, le
long des faces considérées, des pôles au plus simples. Dans
l’intérieur relatif d’une face commune \(F\) de codimension un de deux cellules
\(\sigma,\tau\), d’orientations opposées,
\begin{equation}
 \operatorname{Res}_F\left(a_\sigma\Omega_\sigma+a_\tau\Omega_\tau\right)
 =\left(a_\sigma|_F-a_\tau|_F\right)\Omega_F.
 \label{eq:compat-defecto-frontera}
\end{equation}
Par conséquent, le pôle de cette paire sur le diviseur de \(F\) disparaît
exactement lorsque les deux traces coïncident. Si les coefficients sont
constants et si le graphe d’adjacence des cellules est connexe, l’annulation
de chaque paire de cellules à travers leur face commune équivaut à l’existence
d’un unique coefficient commun \(a\).
\end{theorem}
\begin{proof}
Avec une coordonnée locale transverse \(z\), une forme logarithmique s'écrit \(dz/z\wedge\omega+\eta\), où \(\eta\) est régulière par rapport à ce diviseur. La régularité de \(a_\sigma\) donne \(\operatorname{Res}_F(a_\sigma\Omega_\sigma)
=a_\sigma|_F\operatorname{Res}_F\Omega_\sigma\). Les deux orientations produisent \(\Omega_F\) et \(-\Omega_F\). La forme \(\Omega_F\) est non nulle dans l'intérieur relatif de la face ; l'égalité des résidus équivaut à l'égalité des traces. Pour un pôle simple, l'annulation du résidu élimine le pôle. Dans le cas constant, l'égalité se propage le long de chaque arête du graphe connexe d'adjacence.
\end{proof}

Pour la somme globale, il faut compter tous les termes ayant un pôle sur le
même diviseur ambiant. Si \(H\) est l’hyperplan qui contient \(F\),
et si \(I_H\) réunit ces termes, l’identité générale est
\[
 \operatorname{Res}_{H}\left(\sum_\nu a_\nu\Omega_\nu\right)
 =\sum_{\nu\in I_H}a_\nu|_H\operatorname{Res}_H\Omega_\nu,
\]
pourvu que les coefficients soient réguliers sur \(H\). Une facette
coplanaire, même si elle n’est pas adjacente à la face réelle \(F\), participe à cette
somme. La formule à deux termes représente le résidu global lorsque
les autres termes sont réguliers sur ce diviseur. L’annulation
par paires de toutes les facettes intérieures élimine leurs contributions ;
les facettes extérieures conservent les résidus de frontière correspondants.

Sur une face de codimension supérieure, cette formule est appliquée successivement avec les orientations induites. Les signes sont ceux de l'incidence cellulaire et satisfont l'annulation du bord du bord. La collection de lecteurs locaux définit un recollement compatible deux à deux lorsque leurs traces coïncident après transport vers la même droite de coefficients. Le résidu global est obtenu en sommant les contributions de chaque diviseur. Cet énoncé fournit un critère effectif pour répartir une lecture physique sur des coordinatisations ou des cellules sans introduire de discontinuités artificielles.

\subsection{Annulation des pôles et normalisation globale}

L'annulation intérieure détermine une condition locale. L'identification à une forme canonique globale requiert également le diviseur des pôles et les résidus extérieurs normalisés. Dans l'espace projectif, deux formes rationnelles de degré maximal ayant les mêmes pôles logarithmiques et les mêmes résidus diffèrent d'une forme holomorphe de degré maximal ; \(H^0(\mathbb P^m,\Omega^m)=0\) démontre leur égalité. Telle est la condition globale utilisée dans la preuve du polytope \cite{positive}.

Un exemple par intervalles montre pourquoi les coefficients variables
exigent cette condition. Pour \(0<c<1\), soient
\[
 \Omega_{0c}=\frac{c\,dx}{x(c-x)},\qquad
 \Omega_{c1}=\frac{(1-c)\,dx}{(x-c)(1-x)},\qquad
 \Omega_{01}=\frac{dx}{x(1-x)}.
\]
Avec \(a_1(x)=1+x(c-x)\) et \(a_2(x)=1\), leurs traces coïncident
à toutes les extrémités pertinentes, et pourtant
\[
 a_1\Omega_{0c}+a_2\Omega_{c1}=\Omega_{01}+c\,dx.
\]
Le terme supplémentaire est régulier dans la coordinatisation affine et possède un pôle à l'infini projectif. Les résidus aux extrémités finies et l'annulation en \(c\) sont conservés, tandis que la forme résultante possède un autre comportement global. Multiplier l'exemple par une section de \(\mathcal L_S\) produit le même contrôle dimensionnel.

La conséquence d'ensemble est précise : la représentation grassmannienne marquée transporte les lecteurs retenus ; les mutations conservent leurs valeurs lorsqu'elles ne changent que la coordinatisation ; le raffinement conserve leurs évaluations héritées ; et les formes à valeurs dans la droite d'action s'assemblent par égalité des traces et contrôle des pôles. Chaque égalité exprime l'opération qui conserve l'objet correspondant.

\section{Incidence hexadique--octadique, opérateur d’aire et information sectorielle}
\label{rev4:ae:seccion}

La récupération géométrique du registre acquiert un contenu physique lorsqu’elle est
composée avec des observables dont la construction conserve les incidences et les
conditions d’état utilisées. Le passage d’une hexade à son octade marquée
fournit un cas précis : l’élargissement du support transforme deux
espaces d’incidence, détermine leurs multiplicités et permet de conserver
cette distinction dans un opérateur d’aire et dans le logarithme d’un état
réduit. L’exposé suivant réunit cette construction du deuxième
article de la série~\cite{hmtII} et la représentation grassmannienne
marquée de la section précédente.

Le point de départ demeure dans l’état enrichi produit par
APP--TRIT--TPK. Les deux opérations d’APP conservent leurs divisions
complètes, TRIT détermine régime et orientation, et la composition
\(U_t=\operatorname{Upd}_t\circ\operatorname{Tra}_t\circ
\operatorname{Sel}_t\) transporte et actualise l’information de position,
de retenue et de mémoire. Le prolongement compatible et l’holonomie nonadique
conservent conjointement les réalisations arithmétique et incidentielle de la
structure discrète du continu. Le registre \(K\), les coordonnées
régionales \(\pi_{\rm HMT},\varphi_{\rm HMT},e_{\rm HMT}\), la
coordonnée \(\alpha_{\rm HMT}\) et la configuration exceptionnelle proviennent
de cette construction commune. On utilise ici ses résultats et les marques
qui permettent de les prolonger vers les observables de frontière.

\subsection{Élévation d'incidence et différence de multiplicités}

Soit \(\mathcal H\) le système d'hexades de Witt sur les douze positions du registre. Fixons la coordinatisation d'incidence \(\ell:[12]\to\mathscr D\), où \(\mathscr D\) est une dodécade du code binaire de Golay étendu \(G_{24}\), de paramètres \([24,12,8]_2\). La coordinatisation transporte les hexades de \(\mathcal H\) vers le plan résiduel sur \(\mathscr D\). Cette donnée de réalisation conserve la configuration marquée déterminée dans le développement exceptionnel. Les supports de poids huit de \(G_{24}\) constituent les octades du système \(S(5,8,24)\).

\begin{proposition}[Relèvement unique d’une hexade résiduelle]
\label{rev4:ae:elevacion}
Les ensembles \(O\cap\mathscr D\), avec \(O\) une octade et
\(|O\cap\mathscr D|=6\), forment un système \(S(5,6,12)\).
Chaque hexade de ce système est l’intersection de \(\mathscr D\) avec
une unique octade. En particulier, l’hexade marquée \(H\subset\mathscr D\)
possède un relèvement déterminé \(O_H\) et une inclusion \(H\subset O_H\).
\end{proposition}
\begin{proof}
Les poids du code sont \(0,8,12,16,24\). Si
\(j=|O\cap\mathscr D|\), la somme binaire des vecteurs de support
\(O\) et \(\mathscr D\) a pour poids \(20-2j\) ; par conséquent,
\(j\in\{2,4,6\}\). Tout quintuplet \(T\subset\mathscr D\)
appartient à une unique octade. Comme son intersection avec \(\mathscr D\)
contient \(T\), cette intersection a six éléments. On obtient ainsi
une unique hexade résiduelle sur chaque quintuplet. Si deux octades relevaient
la même hexade, toutes deux contiendraient n’importe lequel de ses quintuplets et
coïncideraient par la propriété de Steiner.
\end{proof}

Dans la suite, \(\binom H2\) désigne l’ensemble des sous-ensembles à deux
éléments de \(H\). Nous conservons ce second facteur en élargissant le premier :
\begin{equation}
 \mathcal F_6=H\times\binom H2,
 \qquad
 \mathcal F_8=O_H\times\binom H2,
 \qquad
 \jmath(h,P)=(h,P).
 \label{rev4:ae:incidencias}
\end{equation}
Un élément de ces espaces est une incidence point--paire. Le point distingué parcourt d'abord l'hexade, puis son octade, tandis que la paire continue d'appartenir à l'hexade d'origine. Cette spécification fixe l'opération responsable de l'accroissement.

\begin{proposition}[Élargissement exact de l’espace des incidences]
\label{rev4:ae:incremento}
L’application \(\jmath\) est injective et son complément est
\((O_H\setminus H)\times\binom H2\). En conséquence,
\begin{equation}
 |\mathcal F_6|=6\binom62=90,
 \qquad |\mathcal F_8|=8\binom62=120,
 \qquad |\mathcal F_8\setminus\jmath(\mathcal F_6)|=30.
 \label{rev4:ae:cardinales}
\end{equation}
\end{proposition}
\begin{proof}
L’inclusion \(H\subset O_H\) conserve les deux coordonnées et prouve
l’injectivité. La partition disjointe \(O_H=H\sqcup(O_H\setminus H)\)
induit la partition du produit cartésien. Comme
\(|O_H\setminus H|=2\) et \(\binom62=15\), le complément a
\(2\cdot15=30\) éléments.
\end{proof}

Les deux multiplicités sont obtenues avant toute distribution
probabiliste sur la frontière. La différence normalisée
\((90-120)/(24\binom62)=-1/12\) conserve le signe de la comparaison
orientée. Son interprétation dépend de l’espace des incidences dans lequel
l’opération a été effectuée.

\subsection{Transport collectif et paramètre de Barbero--Immirzi}

La représentation de frontière retient les étiquettes des deux secteurs. Pour expliciter son support, écrivons \(x_j=r_{2j-1}+3r_{2j}\), avec \(r_k\in\{0,1,2\}\), en trois coordonnées de \(\mathbb Z/9\mathbb Z\). Le regroupement positionnel
\begin{equation}
 \beta_{729}(x_1,x_2,x_3)
 =\bigl(r_1+3r_2+9r_3,\ r_4+3r_5+9r_6\bigr)
 \label{rev4:ae:reagrupamiento}
\end{equation}
est une bijection vers \((\mathbb Z/27\mathbb Z)^2\). L’unicité du
développement ternaire prouve que son inverse extrait les six chiffres et
restitue les trois couples initiaux. L’ordre positionnel est préservé
avec son information de retenue ; les additions respectives restent
définies dans leurs propres groupes.

Dans le bloc \(B_\partial=\{0,\ldots,8\}^2\) du tore d'ordre vingt-sept, on distingue les liaisons orientées
\begin{align}
 E_{90}&=\{((-1,j),(0,j)),\ ((8,j),(9,j)):0\le j\le8\},\\
 E_{120}&=\{((i,-1),(i,0)),\ ((i,8),(i,9)):0\le i\le8\}.
 \label{rev4:ae:enlaces}
\end{align}
Les coordonnées sont interprétées modulo vingt-sept. Les deux ensembles
sont disjoints et contiennent dix-huit liens chacun : deux côtés de
neuf liens par direction. Les étiquettes \(90\) et \(120\) proviennent
des espaces d’incidence ; les cardinaux des ensembles de
liens sont \(18\) et \(18\).

Soient \(u_m\) les vecteurs uniformes normalisés de \(\ell^2(\mathcal F_6)\oplus\ell^2(\mathcal F_8)\), avec \(m=90,120\), et soient \(v_m\) les vecteurs uniformes normalisés correspondants de \(\ell^2(E_{90})\oplus\ell^2(E_{120})\). Par exemple, \(u_{90}=90^{-1/2}\sum_{f\in\mathcal F_6}(\delta_f,0)\) et \(v_{90}=18^{-1/2}\sum_{e\in E_{90}}(\delta_e,0)\). Chaque paire est orthonormale.

\begin{proposition}[Conservation de la structure spectrale collective]
\label{rev4:ae:colectivo}
L’application linéaire \(\mathcal J u_m=v_m\) est une isométrie
surjective entre les deux plans collectifs. Si
\[
 D_{\rm inc}=90|u_{90}\rangle\langle u_{90}|
             +120|u_{120}\rangle\langle u_{120}|,
 \qquad
 D_\partial=90|v_{90}\rangle\langle v_{90}|
             +120|v_{120}\rangle\langle v_{120}|,
\]
alors \(\mathcal J D_{\rm inc}=D_\partial\mathcal J\). Les projecteurs sectoriels et leurs combinaisons orientées sont transportés au moyen de la même application.
\end{proposition}
\begin{proof}
L’application envoie une base orthonormale sur une autre et est déterminée
par ces images. L’égalité des opérateurs vaut sur \(u_m\),
puisque les deux membres ont pour valeur \(m v_m\). Les projecteurs sont
\((120I-D)/30\) et \((D-90I)/30\) dans chaque plan, de sorte que
leurs combinaisons sont conservées par composition linéaire.
\end{proof}

Le transport précédent agit sur les modes uniformes. Les fluctuations
intrinsèques de chaque espace d’incidences demeurent dans leurs compléments
orthogonaux. La pondération aréale utilise précisément les deux étiquettes
que cette isométrie conserve.

Le paramètre de Barbero--Immirzi est déterminé à partir de ces données et de l'orientation angulaire préalablement construite. Soient \(A\) et \(C^*\) les coordonnées angulaires en degrés, dans le relèvement fixé dans la section précédente, et posons
\[
 x=\frac{\pi_{\rm HMT}A}{180},\qquad
 y=\frac{\pi_{\rm HMT}C^*}{180},\qquad
 q_\pm=e^{-x\mp y},\qquad
 S_m(q)=\frac{q^m}{1-q^{3m}}.
\]
On conserve la chambre \(x>|y|\), de sorte que \(0<q_\pm<1\).
La chaîne orientée d’un tour assigne douze secteurs à
\(\mathcal F_6\) et une couture de retour d’orientation contraire
à \(\mathcal F_8\). Son évaluation est \(12S_{90}-S_{120}\).
Le coefficient de \((1-q)^{-1}\) est déterminé par
\[
 \lim_{q\uparrow1}(1-q)(12S_{90}(q)-S_{120}(q))
 =\frac{12}{270}-\frac1{360}=\frac1{24},
\]
car \(1-q^{3m}=(1-q)\sum_{j=0}^{3m-1}q^j\).
L’évaluation conjuguée de la chaîne, jointe à la composante angulaire,
définit le paramètre structurel
\begin{equation}
 \gamma_{\rm HMT}
 =\frac{\pi_{\rm HMT}AC^*}{180}
 +12\bigl[S_{90}(q_-)-S_{90}(q_+)\bigr]
 -\bigl[S_{120}(q_-)-S_{120}(q_+)\bigr].
 \label{rev4:ae:gamma}
\end{equation}
L'origine des coefficients \(12:-1\) réside dans la chaîne orientée. La factorisation précédente détermine son coefficient de pôle. L'inversion \(C^*\mapsto-C^*\) échange les canaux et change le signe de \(\gamma_{\rm HMT}\). Nous conservons ci-dessous la branche \(C^*>0\), \(\gamma_{\rm HMT}>0\), avec la normalisation surfacique fixée par cette même orientation.

\subsection{Retour longitudinal et détermination de l’échelle aréale}
\label{rev5:longitud}

L'unité dimensionnelle de l'aire admet une construction explicite à partir du même état enrichi. Nous écrirons \(L_\alpha\) pour la longueur physique obtenue ; l'entier \(L\) qui désigne la longueur d'un registre dans la section~\ref{sec:compat-lectores-formas} conserve son sens combinatoire. La distinction sépare le cardinal d'une représentation, l'étendue d'une arête et le rayon de sa réalisation.

La composition utilise les coordonnées régionales et le registre déjà
engendrés par APP--TRIT--TPK. APP conserve les deux feuillets avec leurs restes
et quotients ; TRIT oriente le régime local ; TPK prolonge l’état,
en conservant retenue, frontière, route et mémoire. La chaîne
\(w_6\to w_{12}\to w_{18}\to w_{24}\to w_{30}\to R_{36}\to G_9\)
et l’ordre \(U_t\to\Gamma_9\to K_{\rm ph}\) maintiennent le
retour de phase avec accroissement de mémoire. La construction conjointe
du continu conserve ses aspects spectral, solénoïdal, de jauge,
cohomologique et déterminantal. La longueur qui suit est une lecture
dimensionnelle ultérieure de cette structure et de ses sorties corrélées
\(\pi_{\rm HMT},\varphi_{\rm HMT},e_{\rm HMT},\alpha_{\rm HMT}\).

Fixons le registre marqué
\[
 \Omega=\{(j,k,q,u,v):j,k\in\mathbb Z/12\mathbb Z,\quad
 k-j\in\{0,3,6,9\},\quad1\le q\le8,\quad
 u<v,\quad u,v\in\{1,2,4,5,7,8\}\}.
\]
La soustraction est calculée dans ce groupe cyclique. Son cardinal est
\(N_\Omega=12\cdot4\cdot8\binom62=5760\).
Le facteur \(120\) provient de l’espace octadique des incidences
point--paire déjà construit ; les douze positions et les quatre
positions de chaque orbite conservent leurs marques respectives. Cette
résolution élargit le porteur du retour de Hadamard. Dans les espaces
hermitiens de départ et d’arrivée, de dimension \(N_\Omega\), soit
\(B\) l’inverseur transporté, avec
\(B^*B=BB^*=I/4\). L’identité locale
\(H_4^*H_4=4I\) produit cette normalisation pour l’inverseur ;
son extension orthogonale et le transport des marques la conservent.

Dans l'espace d'arrivée, soient \(u_\Omega=N_\Omega^{-1/2}\mathbf1\), \(P=u_\Omega u_\Omega^*\), \(C=(I-P)B\) et \(M=PB\). Le transport complet \(C+M=B\) est conservé. Les projecteurs atomiques \(P_i=e_ie_i^*\), déterminés par l'inscription, définissent l'espérance diagonale \(\Delta_\Omega(T)=\sum_iP_iTP_i\).

\begin{proposition}[Réponse enregistrée et complément conservé]
\label{rev5:retorno}
Parmi les applications linéaires vers l’algèbre diagonale qui fixent
les \(P_i\) et sont bimodulaires par rapport à cette algèbre,
\(\Delta_\Omega\) est unique. En outre,
\begin{equation}
 \Delta_\Omega(CC^*)=r_\Omega I,
 \qquad r_\Omega=\frac{N_\Omega-1}{4N_\Omega}
 =\frac{5759}{23040},
 \qquad \Delta_\Omega(MM^*)=\frac{I}{4N_\Omega}.
 \label{rev5:retorno-escalar}
\end{equation}
Les deux réponses ont pour somme \(I/4\).
\end{proposition}
\begin{proof}
Pour l’unité matricielle \(E_{ij}\), la bimodularité impose
\(E(E_{ij})=P_iE(E_{ij})P_j\). Comme l’image est diagonale,
ce terme vaut zéro si \(i\ne j\) ; si \(i=j\), fixer
\(P_i\) détermine son image. Les unités matricielles constituent
une base et déterminent l’application. D’autre part,
\(CC^*=(I-P)/4\), \(MM^*=P/4\), et chaque coefficient diagonal
de \(P\) est \(1/N_\Omega\). Appliquer l’espérance démontre
les trois identités.
\end{proof}

L'inscription isométrique \(We_i=e_i\otimes e_i\) réalise \(W^*(T\otimes I)W=\Delta_\Omega(T)\). L'inscription complète conserve l'information du registre ; l'espérance extrait sa réponse diagonale. Le complément \(M\) demeure dans le transport et possède la réponse spécifique de \eqref{rev5:retorno-escalar}. Cette séparation identifie l'observable qui intervient dans la longueur et la composante complémentaire conservée.

La forme normale de Smith de \(H_4\) est
\(\operatorname{diag}(1,2,2,4)\) ; son conoyau a pour ordre
seize. Le produit du scalaire engendré
\(\alpha=\alpha_{\rm HMT}\), indexé par les classes de ce
conoyau, définit le lecteur multiplicatif fini
\[
 n_H:=\prod_{\bar v\in\operatorname{coker}_{\mathbb Z}(H_4)}
 \alpha=\alpha^{16}.
\]
Soit \(L_*>0\) une section de référence
radiale de la droite dimensionnelle de longueur. La composition longitudinale
agit linéairement sur une cochaîne orientée \(\beta_*\) :
\begin{equation}
 \mathscr T_L\beta_*
 =\alpha^{16}(\Delta_\Omega(CC^*)\otimes I)\beta_*
 =q_\alpha\beta_*,\qquad
 q_\alpha=\alpha^{16}r_\Omega,\qquad
 L_\alpha=q_\alpha L_*.
 \label{rev5:longitud-generada}
\end{equation}
La formule fixe une composition concrète de lecteurs. La linéarité conserve le type de longueur et son orientation ; le coefficient provient du retour enregistré et de la norme locale. Par exemple, la réponse diagonale d'une source \(a e_i\) est \(a r_\Omega\), tandis que la norme \(\|C^*ae_i\|=|a|\sqrt{r_\Omega}\) constitue une autre lecture du même transport.

Si \(\beta_*(\bar a)=-U(a)\beta_*(a)\), multiplier par
\(q_\alpha\) conserve l’inversion orientée. Si une arête est
subdivisée en neuf tronçons compatibles,
\(\beta_*(a)=\sum_jU_j^{-1}\beta_*(a_j)\), la même
multiplication conserve cette identité et la concaténation de mémoire.
Dans un élargissement auxiliaire du registre, on transporte
\(P\), \(CC^*\) et \(\Delta_\Omega\) par tensorisation
avec l’identité. Ainsi, \(r_\Omega\) demeure ; le cardinal d’un
maillage raffiné et le cardinal du registre marqué ont des fonctions
distinctes.

\subsection{Réciprocité radiale et reconnaissance de la longueur de Planck}
\label{rev5:radial}

La lecture circulaire du calendrier de cent huit avancées conserve
\(108\ell_0=2\pi_{\rm HMT}L_\alpha\), avec
\(\ell_0=ct_0\). Par conséquent,
\begin{equation}
 t_0=\frac{\pi_{\rm HMT}L_\alpha}{54c},\qquad
 \omega=\frac{c}{L_\alpha},\qquad
 Q_{\rm clk}=\hbar_{\rm ret}\omega
 =\frac{\hbar_{\rm ret}c}{L_\alpha}.
 \label{rev5:reloj}
\end{equation}
L’action retournée est reçue de son équation HMT complète dans
\eqref{eq:mass-k-action} ; la vitesse est reçue des deux canaux
constitutifs de la section~\ref{sec:compat-lectores-formas}.
La relation circulaire convertit l’avancée relevée en rayon. Si
\(c=\widehat c\,u_C\), \(t_0=u_T\) et
\(u_L=u_Cu_T\), il en résulte
\(L_*/u_L=54\widehat c/(\pi_{\rm HMT}q_\alpha)\).
La référence radiale \(L_*\) et la base d’arête \(u_L\)
sont reliées avec leurs types dimensionnels explicites.

Dans une fibre finie, soit \(X=X^*>0\) un caractère positif inversible, \(U=2B\) et \(Y=UXU^*\). Le transport longitudinal est \(D_L=L_\alpha U\). Sa réponse radiale positive est définie au moyen de la métrique d'arrivée, \(\|R_Xv\|^2=\|XD_L^*v\|^2\) pour tout \(v\). L'énergie et la longueur conjuguée conservent le même caractère :
\[
 H_X=Q_{\rm clk}Y,\qquad M_X=H_X/c^2,\qquad
 \Lambda_X=L_\alpha Y^{-1}.
\]

\begin{theorem}[Coefficient radial commun et échelle aréale]
\label{rev5:rigidez-radial}
Les conditions précédentes déterminent une unique réponse positive
\(R_X=L_\alpha Y\). On a
\begin{equation}
 R_X\Lambda_X=L_\alpha^2I,\qquad
 H_X\Lambda_X=\hbar_{\rm ret}cI,\qquad
 R_X=\frac{L_\alpha^2}{\hbar_{\rm ret}c}H_X.
\end{equation}
Dans la réalisation physique qui identifie \(R_X\) au rayon
gravitationnel réduit \(R_X=(G/c^2)M_X\), le coefficient est
unique et commun à tous les caractères positifs admissibles :
\begin{equation}
 G=\frac{c^3L_\alpha^2}{\hbar_{\rm ret}},\qquad
 \frac{\hbar_{\rm ret}G}{c^3}=L_\alpha^2.
 \label{rev5:G-area}
\end{equation}
\end{theorem}
\begin{proof}
La polarisation de l'égalité des normes détermine \(R_X^2=D_LX^2D_L^*=L_\alpha^2UX^2U^*\). L'opérateur \(L_\alpha UXU^*\) est positif et son carré est celui indiqué. L'unicité de la racine positive prouve la première affirmation. Les deux identités de produits découlent d'une substitution. La convention physique donne \(L_\alpha Y=(GQ_{\rm clk}/c^4)Y\). L'inversibilité de \(Y\) permet de le simplifier, et l'expression de \(Q_{\rm clk}\) fournit \eqref{rev5:G-area}. Si un autre coefficient conserve les mêmes données, leur différence multipliée par \(M_X\) vaut zéro et est donc nulle.
\end{proof}

Le domaine de la preuve est la fibre positive finie ; il admet aussi
des opérateurs bornés uniformément positifs avec les mêmes identités.
La lecture du rayon gravitationnel réduit constitue la convention
physique déclarée de cette réalisation. La reconnaissance conventionnelle
ultérieure \(\ell_P^2=\hbar_{\rm ret}G/c^3\) identifie
\(\ell_P=L_\alpha\). La longueur s’obtient d’abord au moyen de
\eqref{rev5:longitud-generada} ; cette identification conserve son
constructeur. La formule circulaire équivalente est
\[
 G=\left(\frac{54}{\pi_{\rm HMT}}\right)^2
 \frac{c^5t_0^2}{\hbar_{\rm ret}}.
\]
Le facteur circulaire est déjà inclus lorsque l’on utilise \(L_\alpha\).
La longueur et la section d’action déterminent également
\[
 t_P=\frac{L_\alpha}{c},\qquad
 m_P=\frac{\hbar_{\rm ret}}{cL_\alpha},\qquad
 E_P=\frac{\hbar_{\rm ret}c}{L_\alpha}.
\]
Ces identités découlent de la substitution de \eqref{rev5:G-area} dans les
expressions conventionnelles des trois échelles ; la positivité
détermine leurs racines. Pour un mode de caractère \(y>0\), les
lectures sont \(m(y)=m_Py\), \(r_g(y)=L_\alpha y\) et
\(\bar\lambda(y)=L_\alpha/y\). Leur produit longitudinal
vaut \(r_g(y)\bar\lambda(y)=L_\alpha^2\). L’égalité des
deux longueurs équivaut à \(y=y^{-1}\) et, par positivité,
caractérise exactement \(y=1\). Le calendrier conserve
\(t_0=(\pi_{\rm HMT}/54)t_P\) et
\(108t_0=2\pi_{\rm HMT}t_P\).

Un rayon d’horizon \(R_h=L_\alpha\zeta_h\) conserve son
facteur propre \(\zeta_h=R_h/L_\alpha\). Son aire spatiale
appartient à l’horizon déclaré ; l’échelle élémentaire de l’opérateur
sectoriel suivant est \(L_\alpha^2\). Dans la spécialisation
de Schwarzschild pour la même masse, le rayon est \(2r_g(y)\)
et, par conséquent, \(\zeta_h=2y\) ; un horizon d’une autre réalisation
conserve sa condition géométrique de formation.


\subsection{Occupations de frontière et reconstruction sectorielle}

Chaque lien de \(E_{90}\sqcup E_{120}\) détermine un facteur
\(\mathbb C^3\), de base \(|0\rangle,|1\rangle,|2\rangle\)
et d’opérateur nombre \(N_e=\operatorname{diag}(0,1,2)\). L’espace des
états de frontière est
\(\mathscr H_A=(\mathbb C^3)^{\otimes18}\otimes
(\mathbb C^3)^{\otimes18}\). Nous définissons \(N_m\) comme la somme
des opérateurs nombre des dix-huit liens du secteur \(m\).
Les deux opérateurs commutent et chacun a un spectre entier entre
zéro et trente-six.

L’échelle angulaire provient de l’horloge et du caractère normalisé du
triplet conjugué :
\begin{equation}
 \theta_{\rm clk}=\frac{C^*}{6}\frac{\pi_{\rm HMT}}{180},
 \quad
 \chi_{10}=\frac13\Tr\operatorname{diag}
 (1,e^{i\pi_{\rm HMT}/18},e^{-i\pi_{\rm HMT}/18})
 =\frac{1+2\cos(\pi_{\rm HMT}/18)}3,
 \quad a_0=\chi_{10}\theta_{\rm clk}.
 \label{rev4:ae:escala}
\end{equation}
L'identité du caractère s'obtient en additionnant les phases conjuguées. Dans la branche considérée, \(a_0>0\). L'unité d'aire est \(L_\alpha^2=\ell_P^2\), construite et reconnue dans \eqref{rev5:longitud-generada}--\eqref{rev5:G-area}. Écrivons \(a_\partial=\gamma_{\rm HMT}a_0>0\). L'opérateur d'aire est
\begin{equation}
 \widehat{\mathcal A}=L_\alpha^2a_\partial N,
 \qquad N=N_{90}+N_{120},
 \qquad D=90N_{90}+120N_{120}.
 \label{rev4:ae:area}
\end{equation}
Son spectre est \(\{nL_\alpha^2a_\partial:0\le n\le72\}\).
La multiplicité de \(n\) est le coefficient de \(z^n\) dans
\((1+z+z^2)^{36}\), car chaque facteur tensoriel apporte l’une des
trois occupations.

\begin{theorem}[Reconstruction des deux occupations sectorielles]
\label{rev4:ae:recuperacion}
La paire d’observables \((N,D)\) détermine
\begin{equation}
 N_{90}=4N-\frac{D}{30},
 \qquad N_{120}=\frac{D}{30}-3N.
 \label{rev4:ae:inversa}
\end{equation}
Dans le spectre conjoint, ses valeurs \((n,d)\) satisfont \(n\in\mathbb Z\), \(d\in30\mathbb Z\), \(90n\le d\le120n\) et \(0\le4n-d/30\le36\), \(0\le d/30-3n\le36\). Ces conditions caractérisent l'image des occupations entières.
\end{theorem}
\begin{proof}
La transformation de \((N_{90},N_{120})\) en \((N,D)\) a pour matrice
\(\left(\begin{smallmatrix}1&1\\90&120\end{smallmatrix}\right)\),
dont le déterminant est \(30\). Son inverse donne les identités
d’opérateurs. La commutativité permet de les appliquer dans chaque espace propre
commun. Les inégalités expriment la non-négativité et les capacités
maximales des secteurs ; la divisibilité rend les deux occupations entières.
Réciproquement, chaque entier entre zéro et trente-six s’exprime comme
somme de dix-huit termes de \(\{0,1,2\}\), et donc toute paire
satisfaisant ces conditions appartient au spectre conjoint.
\end{proof}

L’aire conserve la somme des occupations ; l’observable pondéré conserve
leur composition sectorielle. Ainsi, \((N_{90},N_{120})=(1,0)\) et \((0,1)\)
ont la même aire \(L_\alpha^2a_\partial\) et des poids \(90\) et \(120\).
Dans le sens complémentaire, \((4,0)\) et \((0,3)\) ont le même poids
\(D=360\), mais des aires différentes. Le couple sépare les deux situations.
À l’intérieur de chaque espace propre conjoint, les dégénérescences
microscopiques correspondant aux distributions entre liens sont conservées.

\begin{figure}[htbp]
\centering
\begin{tikzpicture}[>=Latex,every node/.style={font=\small},node distance=15mm]
\node[draw,rounded corners,align=center] (a) {Occupation sectorielle\\$(1,0)$};
\node[draw,rounded corners,align=center,right=54mm of a] (b) {Occupation sectorielle\\$(0,1)$};
\node[draw,rounded corners,align=center] (n) at ($(a)!0.5!(b)+(0,-1.8)$)
 {Aire commune\\$L_\alpha^2a_\partial$};
\draw[->] (a)--node[left] {$N=1$}(n);
\draw[->] (b)--node[right] {$N=1$}(n);
\node[above=5mm of a] {$D=90$};
\node[above=5mm of b] {$D=120$};
\end{tikzpicture}
\caption{L'égalité d'aire conserve l'occupation totale. La pondération d'incidence distingue les deux compositions ; leur lecture conjointe permet de reconstruire les occupations des deux secteurs.}
\label{rev4:ae:figura}
\end{figure}

\subsection{État réduit, purification et entropie d'intrication}

L'interprétation informationnelle requiert de fixer l'état. Pour \(\Delta\in\mathbb R\) fini, avec la normalisation du paramètre modulaire de la construction sectorielle, nous définissons
\begin{equation}
 x_m=e^{-m\Delta},\qquad z_m=1+x_m+x_m^2,
 \qquad Z=z_{90}^{18}z_{120}^{18},
 \qquad \rho_A(\Delta)=Z^{-1}e^{-\Delta D}.
 \label{rev4:ae:rho}
\end{equation}
La somme des occupations commutantes et la décomposition tensorielle prouvent
l’expression de \(Z\). Toutes les valeurs propres de \(\rho_A\) sont
strictement positives pour \(\Delta\) réel fini. Son hamiltonien
modulaire, noté \(K_A\) pour le distinguer du registre arithmétique
\(K\), est
\begin{equation}
 K_A=-\log\rho_A=(\log Z)I+\Delta D.
 \label{rev4:ae:modular}
\end{equation}
Cet opérateur caractérise la distribution spectrale de l’état réduit.
Le générateur d’évolution temporelle est défini dans le développement dynamique
avec ses propres paramètres et son propre domaine. Pour \(\Delta\ne0\), lorsque
\(\Delta\) et la normalisation sont connus, la paire
\((\widehat{\mathcal A},K_A)\) retrouve \((N,D)\) et, au moyen de
\eqref{rev4:ae:inversa}, les deux occupations. Pour \(\Delta=0\),
l’état est uniforme et \(K_A=36\log3\,I\).

Nous construisons explicitement la bipartition. Soit \(\mathscr H_{\bar A}\) une copie de \(\mathscr H_A\), avec les bases conjuguées fixées. Pour une liaison de poids \(w\), le vecteur
\[
 |\operatorname{TFD}(w)\rangle
 =\frac1{\sqrt{1+e^{-w}+e^{-2w}}}
 \sum_{n=0}^{2}e^{-wn/2}|n\rangle_A\otimes|n\rangle_{\bar A}
\]
a une norme égale à un. L'état pur
\begin{equation}
 |\Psi_\Delta\rangle
 =|\operatorname{TFD}(90\Delta)\rangle^{\otimes18}
  \otimes|\operatorname{TFD}(120\Delta)\rangle^{\otimes18}
 \label{rev4:ae:tfd}
\end{equation}
détermine la bipartition \(A\mid\bar A\).

\begin{proposition}[Réduction sectorielle et entropie de la purification]
\label{rev4:ae:entropia}
La trace partielle de \(|\Psi_\Delta\rangle\langle\Psi_\Delta|\)
sur \(\mathscr H_{\bar A}\) est \(\rho_A(\Delta)\). Étant définis
\[
 f_m=\frac{x_m+2x_m^2}{z_m},
 \qquad s(\Delta)=-\Tr(\rho_A(\Delta)\log\rho_A(\Delta)),
\]
on a
\begin{align}
 \langle N\rangle_\Delta&=18(f_{90}+f_{120}),\\
 \langle D\rangle_\Delta&=18(90f_{90}+120f_{120}),\\
 \langle\widehat{\mathcal A}\rangle_\Delta
 &=18L_\alpha^2a_\partial(f_{90}+f_{120}),\\
 s(\Delta)&=18\log z_{90}+18\log z_{120}
             +\Delta\langle D\rangle_\Delta.
 \label{rev4:ae:entropia-formula}
\end{align}
La dernière expression est l'entropie d'intrication de la purification \eqref{rev4:ae:tfd}, en unités sans dimension.
\end{proposition}
\begin{proof}
Dans la trace partielle d’un lien, l’orthogonalité de
\(|n\rangle_{\bar A}\) élimine les termes d’indices distincts et
laisse les poids \((1,x_m,x_m^2)/z_m\). La trace partielle d’un produit
tensoriel est le produit des traces partielles ; on obtient
\eqref{rev4:ae:rho}. Le premier moment de chaque lien est \(f_m\).
L’additivité des opérateurs nombre donne les trois premières identités.
La dernière résulte de la substitution de \eqref{rev4:ae:modular} dans
\(s=\Tr(\rho_AK_A)\). Les coefficients de la construction précédente
sont des coefficients de Schmidt de la bipartition déclarée, de sorte que
cette entropie coïncide avec son entropie d’intrication.
\end{proof}

L'entropie maximale de cet espace est \(36\log3\), atteinte en \(\Delta=0\). La quantité \(36\log3-s(\Delta)\) est l'entropie relative de \(\rho_A(\Delta)\) par rapport à \(I/3^{36}\). L'incidence fixe les secteurs et leurs multiplicités ; le paramètre de l'état fixe les probabilités ; la coordinatisation dimensionnelle convertit l'occupation en aire. L'identification physique de \(\mathscr H_{\bar A}\) conserve la bipartition et l'algèbre des observables qui viennent d'être spécifiées.

\subsection{Identité modulaire finie et transport de l'information géométrique}

Fixons l'état de référence \(\rho_0=\rho_A(\Delta_0)\), avec \(\Delta_0\) réel fini, et maintenons constants \(\Delta_0,a_\partial,L_\alpha\). Les aires sectorielles normalisées \(\mathcal A_m^{\rm n}=a_\partial N_m\) satisfont
\begin{equation}
 K_0=-\log\rho_0=c_0I+
 \beta_{90}\mathcal A_{90}^{\rm n}
 +\beta_{120}\mathcal A_{120}^{\rm n},
 \qquad c_0=\log Z(\Delta_0),\qquad
 \beta_m=\frac{m\Delta_0}{a_\partial}.
 \label{rev4:ae:betas}
\end{equation}
La relation \(\beta_{120}=\tfrac43\beta_{90}\) conserve le rapport des incidences. Pour \(\Delta_0\ne0\), le rapport des deux coefficients est défini et vaut \(4/3\).

\begin{theorem}[Identité finie d'aire et d'information relative]
\label{rev4:ae:ley-finita}
Pour tout état normalisé \(\sigma\) de \(\mathscr H_A\), soit \(\delta_\sigma\langle B\rangle=\Tr[(\sigma-\rho_0)B]\). Alors
\begin{equation}
 s(\sigma)-s(\rho_0)
 =\sum_{m\in\{90,120\}}\beta_m
       \delta_\sigma\langle\mathcal A_m^{\rm n}\rangle
   -D(\sigma\Vert\rho_0),
 \label{rev4:ae:ley-finita-formula}
\end{equation}
où \(D(\sigma\Vert\rho_0)=
\Tr[\sigma(\log\sigma-\log\rho_0)]\) est finie et non négative,
avec la convention \(0\log0=0\).
\end{theorem}
\begin{proof}
La définition donne \(D(\sigma\Vert\rho_0)=-s(\sigma)+\Tr(\sigma K_0)\), tandis que \(s(\rho_0)=\Tr(\rho_0K_0)\). Soustraire les identités et substituer \eqref{rev4:ae:betas} prouve \eqref{rev4:ae:ley-finita-formula} ; le terme scalaire disparaît du fait de la normalisation des traces. Pour la positivité, soient \(p_i\) et \(u_i\) les valeurs propres et les vecteurs propres de \(\sigma\), et \(q_j>0,v_j\) ceux de \(\rho_0\). Avec \(r_i=\langle u_i,\rho_0u_i\rangle\), la concavité de \(\log\) donne \(\langle u_i,\log\rho_0u_i\rangle\le\log r_i\). Ainsi, \(D(\sigma\Vert\rho_0)\ge\sum_i p_i\log(p_i/r_i)\ge0\), par convexité de \(t\log t\), puisque \(\sum_i r_i=1\). Le rang plein de \(\rho_0\) garantit la finitude.
\end{proof}

Sa version infinitésimale maintient l’état de référence fixe. Si
\(\sigma(\varepsilon)\) est différentiable, normalisé et
\(\sigma(0)=\rho_0\), alors
\[
 \left.\frac{\dd s(\sigma(\varepsilon))}{\dd\varepsilon}\right|_0
 =\sum_m\beta_m
  \left.\frac{\dd\langle\mathcal A_m^{\rm n}\rangle_{
                          \sigma(\varepsilon)}}{\dd\varepsilon}\right|_0.
\]
En effet, la dérivée de la trace de \(-\sigma\log\sigma\) est \(-\Tr[\dot\sigma(\log\rho_0+I)]\) ; la normalisation élimine \(\Tr\dot\sigma\) et la substitution de \(K_0\) donne la formule. L'identité finie ajoute à cette réponse tangente la correction exacte d'entropie relative. Les coefficients des aires physiques sont \(\beta_m/L_\alpha^2\), en conservant la même coordinatisation dimensionnelle.

La composition avec la représentation grassmannienne précise enfin quelle information est conservée. Dans \(\mathcal F(K,\xi)=([B(K)],Q,\xi)\), les coordonnées supplémentaires \(\xi\) retiennent les grandeurs régionales, l'orientation angulaire, l'inclusion marquée, l'identification des secteurs, la base des occupations, le paramètre \(\Delta\) et l'unité aréale pertinente. L'inverse de la proposition~\ref{prop:compat-inversa-marcada} restitue \(K\) ; les autres conditions sont conservées par leurs marques. D'après le théorème~\ref{thm:compat-algebra-lectores}, la composition avec cet inverse transporte \(\gamma_{\rm HMT}\), \(\widehat{\mathcal A}\), \(\rho_A\), \(K_A\) et leurs évaluations sur le domaine déclaré. Les changements de coordinatisation du même point marqué conservent ces objets, et les raffinements conservent les évaluations héritées lorsqu'ils maintiennent les mêmes données sectorielles.

Cette articulation conserve trois niveaux distinguables d’information.
La représentation positive marquée récupère le registre avec les données
retenues ; le couple aire--hamiltonien modulaire récupère les deux
occupations sous les conditions de \eqref{rev4:ae:modular} ; la
purification détermine l’intrication d’une bipartition concrète.
La perte d’information dans une projection se calcule par ses fibres :
l’aire identifie les occupations de même somme, tandis que la lecture
conjointe sépare leur composition et conserve la dégénérescence microscopique
restante. Le lien entre géométrie et information est ainsi exprimé
par des applications et des identités vérifiables sur la même construction.
\subsection{Naturalité de l’échelle et formes canoniques opératorielles}
\label{rev5:naturalidad-areal}

La longueur engendrée permet de préciser la compatibilité entre la
réalisation positive, l’opérateur d’aire et l’information sectorielle.
La représentation marquée \(\mathcal F(K,\xi)\) conserve,
parmi les données pertinentes de \(\xi\), la résolution
\(\Omega\), ses projecteurs, le retour transporté et la section
\(L_*\). Les coordonnées régionales engendrées, l’action et les canaux
constitutifs maintiennent leur généalogie antérieure. L’inverse marqué
restitue ainsi tous les arguments de
\eqref{rev5:longitud-generada} et de \eqref{rev5:G-area}.
En particulier, le lecteur géométrique déjà construit de
\(\alpha^{\mathrm{geom}}\) détermine les expressions
\[
 L_\alpha^{\mathrm{geom}}
 =L_*(\alpha^{\mathrm{geom}})^{16}\frac{5759}{23040},
 \qquad
 G^{\mathrm{geom}}
 =\frac{(c^{\mathrm{geom}})^3
 (L_\alpha^{\mathrm{geom}})^2}{\hbar_{\rm ret}^{\mathrm{geom}}}.
\]
Les lecteurs de vitesse et d’action s’obtiennent par la même
composition avec l’inverse marqué. Ces expressions conservent
les données du retour et la section radiale de référence.

\begin{corollary}[Conservation par changements de coordinatisation et raffinement hérité]
\label{rev5:conservacion}
Les changements positifs de coordinatisation conservant le point marqué et les raffinements restituant le registre initial conservent \(L_\alpha\), \(G\), \(\gamma_{\rm HMT}\) et \(\widehat{\mathcal A}=L_\alpha^2a_\partial N\) lorsqu'ils transportent également les données sectorielles et dimensionnelles déclarées. Le spectre normalisé \(\widehat{\mathcal A}/L_\alpha^2=a_\partial N\), ses multiplicités et l'état \(\rho_A(\Delta)\) demeurent identiques.
\end{corollary}
\begin{proof}
Le théorème~\ref{thm:compat-algebra-lectores} transporte chaque lecteur
par composition avec l’inverse. Les mutations du même point
reconstruisent le même registre. Dans le raffinement hérité, la somme
des séparations subdivisées reconstruit chaque coordonnée initiale.
La tensorisation de l’inscription conserve \(r_\Omega\),
et ses autres arguments restent fixés par les marques. Par
substitution, la longueur, le coefficient gravitationnel et
l’opérateur d’aire sont conservés. Les projecteurs d’occupation conservent le
polynôme de multiplicités \((1+z+z^2)^{36}\) et l’expression
normalisée de l’état.
\end{proof}

Le raffinement hérité considéré ici conserve le même espace
de trente-six liens de frontière. Une extension qui introduit
un espace physique de frontière différent requiert des opérateurs d’entrelacement
spécifiés pour \(N_{90}\) et \(N_{120}\), ainsi que le
transport des données sectorielles déclarées.

Soit maintenant \(P\) un polytope de rang un de ce manuscrit, avec une subdivision orientée \(P=\bigcup_\sigma P_\sigma\) du même support, de multiplicité un et à faces compatibles. On fixe la même fibre sectorielle \(\mathscr H_A\) sur toutes ses cellules. La forme de degré maximal à valeurs opératorielles et de dimension d'aire est définie par
\[
 \boldsymbol\Omega^{\mathcal A}_P
 =\widehat{\mathcal A}\otimes\Omega_P.
\]
Ici \(\widehat{\mathcal A}\) est constant par rapport aux
coordonnées différentielles du polytope : il provient de l’état HMT
fixé, tandis que \(\Omega_P\) décrit son support géométrique.
Le coefficient \(\widehat{\mathcal A}\) est semi-défini positif.
Sa composante dans le secteur du vide \(N=0\) est nulle. Sur un
espace propre d’occupation \(n>0\), diviser par
\(nL_\alpha^2a_\partial\) récupère la forme canonique
\(\Omega_P\) multipliée par l’identité de cet espace propre.

\begin{proposition}[Additivité aréale et résidus de frontière]
\label{rev5:forma-areal}
Dans la trivialisation sectorielle commune, on a
\begin{equation}
 \boldsymbol\Omega^{\mathcal A}_P
 =\sum_\sigma\widehat{\mathcal A}\otimes\Omega_{P_\sigma},
 \qquad
 \operatorname{Res}_F\boldsymbol\Omega^{\mathcal A}_P
 =\widehat{\mathcal A}\otimes\Omega_F.
\end{equation}
Les identités persistent sous subdivisions successives et résidus itérés. Pour un état fixé \(\rho\), prendre la trace donne \(\Tr(\rho\widehat{\mathcal A})\Omega_P\) et conserve les mêmes identités scalaires.
\end{proposition}
\begin{proof}
L’espace d’opérateurs de \(\mathscr H_A\) est de dimension
finie. Dans toute base, l’assertion est l’identité
d’additivité des formes canoniques multipliée par chaque coefficient
constant de \(\widehat{\mathcal A}\). Le résidu et la trace
sont linéaires. L’annulation des facettes intérieures conserve leurs
orientations et le même coefficient opératoriel des deux côtés.
L’itération applique ces opérations à chaque subdivision et à chaque
drapeau de faces.
\end{proof}

Si différentes cellules portent des coefficients opératoriels réguliers \(A_\sigma\), le défaut d'une paire sur sa face commune est \((A_\sigma|_F-A_\tau|_F)\otimes\Omega_F\). L'égalité des restrictions opératorielles à la frontière fournit le recollement, sous les conditions de pôles simples du théorème~\ref{thm:compat-pegado}. La somme globale conserve en outre toutes les contributions sur chaque diviseur ambiant et le contrôle des pôles extérieurs. La normalisation dimensionnelle, l'annulation des résidus et la couverture du support ont ainsi des critères séparés et compatibles.

L’identité modulaire finie conserve aussi une expression directe
en aires physiques \(\mathcal A_m=L_\alpha^2a_\partial N_m\) :
\begin{equation}
 s(\sigma)-s(\rho_0)
 =\sum_{m=90,120}\frac{m\Delta_0}{L_\alpha^2a_\partial}
   \Tr[(\sigma-\rho_0)\mathcal A_m]
  -D(\sigma\Vert\rho_0).
 \label{rev5:modular-fisica}
\end{equation}
C’est \eqref{rev4:ae:ley-finita-formula} après insertion de la
longueur construite. Le facteur angulaire \(a_0\), contenu dans
\(a_\partial=\gamma_{\rm HMT}a_0\), demeure intégral.
Si l’on compare des sections dimensionnelles avec \(L'_\alpha=sL_\alpha\)
et si l’on conserve le même état et les mêmes données adimensionnelles,
les aires sont multipliées par \(s^2\) et leurs coefficients modulaires
par \(s^{-2}\). L’entropie et l’entropie relative demeurent
égales. Cette comparaison de sections ne constitue pas une autre solution
pour les mêmes données longitudinales fixées.

Pour un horizon sphérique de rayon \(R_h=L_\alpha\zeta_h\),
l’aire géométrique vaut \(4\pi_{\rm HMT}L_\alpha^2\zeta_h^2\).
Son identification à l’espérance de l’opérateur sectoriel, si elle est
requise pour une même coupe, exige la condition explicite
\[
 a_\partial\Tr(\rho N)=4\pi_{\rm HMT}\zeta_h^2.
\]
La géométrie de l'horizon et l'état de frontière doivent satisfaire cette relation avec leurs propres constructions. Le facteur \(\zeta_h\) caractérise cet horizon ; l'échelle longitudinale du couplage reste \(L_\alpha\).


\begin{figure}[htbp]
\centering
\begin{tikzpicture}[
 box/.style={draw=black!70,rounded corners=1pt,align=center,
             font=\small,inner sep=5pt,text width=40mm},
 arrow/.style={-{Stealth[length=2mm]},thick}]
\node[box] (memory) at (0,0) {Charge et mémoires\\\((q,m_0,m_1)\)};
\node[box] (register) at (0,-1.65) {Registre reconstruit\\\(K=(q/12)\mathbf1+Lm\)};
\node[box] (weights) at (-3.8,-3.55) {Partition ordonnée\\\(d_j=K_j/q\)};
\node[box] (direction) at (3.8,-3.55) {Direction de somme nulle\\\(u=P_3P_{11}K\)};
\node[box] (geometry) at (-4.6,-6) {Mesure de frontière\\et forme canonique\\\(C(d),\ \Omega_{P_K}\)};
\node[box] (energy) at (0,-6) {Énergie nodale\\et détails\\\(L_d,\ \mathcal E(z)\)};
\node[box] (lattice) at (4.6,-6) {Déformation réticulaire\\et réseau dual\\\(\Lambda_s^*=\Lambda_{-s}\)};
\draw[arrow] (memory)--(register);
\draw[arrow] (register)--(weights);
\draw[arrow] (register)--(direction);
\draw[arrow] (weights)--(geometry);
\draw[arrow] (weights)--(energy);
\draw[arrow] (direction)--(lattice);
\draw[arrow] (direction) to[bend left=12] node[right,font=\scriptsize,align=left,fill=white,inner sep=1pt]{flux\\hérité} (energy);
\end{tikzpicture}
\caption{Réalisations ultérieures du registre récupérable. La branche d'intervalles détermine les mineurs de chemins, la forme canonique du polygone et le laplacien nodal. La direction algébrique détermine le flot des intervalles et, avec la réalisation réticulaire marquée, la collection duale à vingt-quatre dimensions. Chaque flèche conserve le domaine et les hypothèses de sa construction ; les codomaines inférieurs sont distincts.}
\label{fig:multimorphology}
\end{figure}


\clearpage
\part{Réalisation de jauge, raffinement et écart spectral}
\section{Positivités, espaces physiques et transport du spectre}
La géométrie positive et la théorie de jauge reçoivent des structures distinctes d'un même processus généalogique. La première réalise les poids et leurs incidences au moyen de mineurs ordonnés, de cellules et de formes canoniques. La seconde conserve l'incidence de jauge, l'holonomie, les représentations, les chemins et la mémoire au sein d'une mesure compatible. Sa forme de positivité pertinente pour la reconstruction physique est la positivité par réflexion. Le lien conceptuel acquiert un contenu mathématique lorsque les applications qui produisent chaque forme sont conservées.

Dans la partie géométrique, on a démontré que la subdivision d’une cellule
conserve sa forme canonique par annulation des pôles intérieurs.
Pour les énergies nodales, la réduction de Schur conserve une forme
effective et enregistre le détail. Dans la partie de jauge, la compatibilité du
raffinement doit conserver en outre le vide, l’espace des invariants,
le générateur et le témoin spectral. Telle est la chaîne développée
ci-dessous : échelle arithmético--géométrique, mesure commune, semi-groupe,
projection physique, raffinement et reconstruction continue.

La norme d’un vecteur, la positivité d’une matrice de moments et
l’existence d’un écart d’un Hamiltonien ont des domaines différents.
Une congruence à facteur inversible conserve la positivité et
l’inertie d’une forme ; un entrelacement isométrique du générateur, avec son
domaine, transporte sa réalisation sur l’image réduite de
l’inclusion. La borne uniforme et le témoin persistant fixent ensuite
l’écart limite. Le développement maintient donc visibles les
inclusions et les identités de semi-groupe qui justifient le passage à la
limite. L’échelle de l’écart provient de l’incidence aréale--volumétrique
et de son unité dimensionnelle déclarée.

Les développements suivants réunissent la construction projective et son développement ultérieur dépendant du couplage \(g\). Les premières comparaisons avec un paramétrage de Gibbs sont complétées ensuite au moyen de l'action, de la mesure et du transport explicites. L'ordre d'exposition conserve les démonstrations du processus et situe la formulation complète après ses dépendances. Les quatre réflexions correspondent aux quatre directions euclidiennes de la reconstruction. Le développement réuni ici a pour objet cette théorie de jauge et son écart spectral, avec le groupe compact, connexe et simple déclaré dans les théorèmes.

Le secteur supersymétrique planaire dans lequel les amplitudes sont étudiées au moyen d'amplituèdres conserve une question physique spécifique. Pour le relier à la présente réalisation, les opérateurs de supercharge, leurs domaines et la représentation correspondante sont requis. La dimension quatre de la coordinatisation euclidienne et la différence quatre du bloc modulaire demeurent des résultats déjà définis de leurs opérateurs respectifs.

\section{Transition surfacique--volumétrique et échelle spectrale}
\label{sec:ym-areal-volumetric-scale}

La constante qui fixe le semi-groupe de jauge est obtenue avant l'introduction de l'hamiltonien. Sa provenance combine deux lecteurs indépendants du même état généalogique : la fréquence de retour exprimant \(\pi_{\rm HMT}\) et la distribution stationnaire des deux feuillets aire--volume.

\subsection{Dynamique des deux feuillets}

Soit
\begin{equation}
 F_{\rm av}=\begin{pmatrix}0&1\\1&1\end{pmatrix},
 \qquad e_{\rm ar}=\binom10,
 \qquad e_{\rm vol}=\binom01.
 \label{eq:ym-av-transition-matrix}
\end{equation}
La récurrence de Fibonacci donne, par récurrence,
\begin{equation}
 F_{\rm av}^{n}
 =\begin{pmatrix}F_{n-1}&F_n\\F_n&F_{n+1}\end{pmatrix}.
 \label{eq:ym-av-fibonacci-powers}
\end{equation}
Le vecteur de Perron est proportionnel à \((1,\varphi)\). Comme
\(F_{\rm av}\) est symétrique, la mesure de Parry assigne aux deux feuillets
des poids proportionnels à \((1,\varphi^2)\), et donc
\begin{equation}
 \frac{\mu_{\rm ar}}{\mu_{\rm vol}}=\varphi^{-2},
 \qquad
 \mu_{\rm vol}=\frac{\varphi^2}{1+\varphi^2}.
 \label{eq:ym-av-parry-weights}
\end{equation}
Ici, \(\mu_{\rm vol}\) est un caractère de fréquence sans dimension ; il ne désigne pas un volume de dimension \(L^4\).

Le lecteur régional de \(\pi\) est introduit une seule fois au moyen de
\begin{equation}
 D_\pi
 =\pi_{\rm HMT}|e_{\rm ar}\rangle\langle e_{\rm ar}|
  +|e_{\rm vol}\rangle\langle e_{\rm vol}|.
 \label{eq:ym-pi-regional-reader}
\end{equation}
Le rapport exprimé à la profondeur \(n\) est
\begin{equation}
 \Theta_n
 :=\frac{\langle e_{\rm ar},D_\pi F_{\rm av}^{n}e_{\rm ar}\rangle}
 {\langle e_{\rm vol},F_{\rm av}^{n}e_{\rm vol}\rangle}
 =\pi_{\rm HMT}\frac{F_{n-1}}{F_{n+1}}
 \longrightarrow
 r_{\rm av}:=\frac{\pi_{\rm HMT}}{\varphi_{\rm HMT}^{2}}.
 \label{eq:ym-target-free-av-limit}
\end{equation}
La matrice, le chemin et le quotient sont fixés avant de consulter un saut de masse ou un paramètre de théorie de jauge. À la profondeur native \(n=108\),
\begin{equation}
 F_{107}=10284720757613717413913,
 \qquad
 F_{109}=26925748508234281076009,
 \label{eq:ym-fibonacci-107-109}
\end{equation}
et
\begin{equation}
 \left|\Theta_{108}
 -\frac{\pi_{\rm HMT}}{\varphi_{\rm HMT}^{2}}\right|
 <2\cdot10^{-45}.
 \label{eq:ym-av-depth-108-error}
\end{equation}

\subsection{Transduction dimensionnelle}

Soit \(\ell_0\) l’unité longitudinale du lecteur dimensionnel. L’échelle
inverse de longueur et le paramètre d’une dalle d’épaisseur \(a_m\) sont
\begin{equation}
 \kappa_{\rm av}
 :=\frac{\mu_{\rm vol}}{\ell_0}
 \left(\frac{\pi_{\rm HMT}}{\varphi_{\rm HMT}^{2}}-1\right)>0,
 \qquad
 q_m:=e^{-3a_m\kappa_{\rm av}}.
 \label{eq:ym-dimensional-gap-scale}
\end{equation}
Par conséquent, \([\kappa_{\rm av}]=L^{-1}\) et l’écart énergétique est
\begin{equation}
 \Delta_E=3\hbar c\,\kappa_{\rm av};
 \label{eq:ym-energy-gap-units}
\end{equation}
en unités \(\hbar=c=1\), on écrit \(\Delta_E=3\kappa_{\rm av}\).
La mémoire électronique \(\Delta_4\) appartient à un autre domaine et à un autre
lecteur ; elle n’intervient pas dans \eqref{eq:ym-dimensional-gap-scale}.

Dans les semi-groupes des sections suivantes, \(t,s,a_m\) sont des longueurs euclidiennes et les générateurs \(D\) et \(H^{\rm ov}\) ont pour dimension \(L^{-1}\). L'hamiltonien énergétique est \(H_E=\hbar_{\rm int}c_{\rm int}D\). Si \(\tau=t/c_{\rm int}\), alors
\[
 e^{-tD}=e^{-\tau H_E/\hbar_{\rm int}},\qquad
 \Delta=3\kappa_{\rm av},\quad
 \Delta_E=\hbar_{\rm int}c_{\rm int}\Delta,\quad
 m_{\rm gap}=\frac{\hbar_{\rm int}\Delta}{c_{\rm int}}.
\]
Le même saut est ainsi exprimé respectivement comme longueur inverse, énergie et masse dans la coordinatisation dimensionnelle déclarée.

Comme \(a_m=9a_{m+1}\), la définition donne immédiatement
\begin{equation}
 q_{m+1}^{9}=q_m.
 \label{eq:ym-nonadic-gap-root}
\end{equation}
Cette identité est la racine nonadique qui transportera le même semi-groupe et la
même échelle spectrale à travers tous les raffinements.

\section[Processus de jauge euclidien de Yang--Mills en dimension 4]{Construction
projective d’un processus de jauge euclidien de Yang--Mills, reconstruction
d’Osterwalder--Schrader et écart collectif en dimension \(4\)}


\noindent La spécialisation de la dynamique nonadique enrichie,
construite dans les sections précédentes, est
\[
 (\widehat X_\infty,\Gamma_9)
 \xrightarrow{\ \mathcal R_{\rm YM}\ }
 (\mu_m,T_m,H_m,\Omega_m)
 \longrightarrow(\mu_\infty,\mathrm{OS},E(4))
 \longrightarrow\operatorname{gap}(H_{\rm phys})>0.
\]
Le lecteur clover exprimant \(F^2\) reconnaît la courbure sur ce même processus projectif ; il ne construit pas rétrospectivement \(\mu_m\). La théorie de jauge, la reconstruction OS, le vide simple et le saut sont déjà fixés par cette chaîne. L'identité \eqref{eq:amp-ym-identidad-accion-requerida} compare ensuite une paramétrisation optionnelle par \(g_R\) ; c'est un contrôle non gouvernant, et non une condition permettant d'identifier la théorie de Yang--Mills construite.

\subsection*{Objet de jauge et domaine de raffinement}

Soit \(G\) un groupe de Lie compact, connexe et simple. Pour chaque
\(m\geq0\), on considère le réseau euclidien
\begin{equation}\label{eq:realizaciones-ym-red}
 \Lambda_m=a_m\mathbb Z^4,
 \qquad a_{m+1}=\frac{a_m}{9}.
\end{equation}
L'algèbre cylindrique de jauge est engendrée par des variables de liaison dans \(G\), des opérateurs d'entrelacement de représentation et les coordonnées enrichies de couleur, d'orientation, d'holonomie, de retenue de fusion, de chemin et de mémoire. Soit \(\mu_m\) la mesure invariante commune et soit
\[
 \mathcal H_m=L^2(\mathcal A_m,\mu_m)^{\rm gauge}.
\]
Le vecteur constant \(\Omega_m\) définit
\begin{equation}\label{eq:realizaciones-ym-proyectores}
 P_{\Omega_m}=|\Omega_m\rangle\langle\Omega_m|,
 \qquad Q_m^{\rm phys}=I-P_{\Omega_m}.
\end{equation}
Le second projecteur est le complément physique du vide ; il ne s’identifie pas
à un résiduel employé uniquement pour organiser l’échelle de
renormalisation.

Le circuit de mise à jour est local :
\begin{equation}\label{eq:realizaciones-ym-circuito}
 K_m^{\rm loc}
 =\prod_{r=1}^{9\cdot81}
   \prod_{B\in\mathcal B_{m,r}}K_{m,B}.
\end{equation}
La partition \(\{\mathcal B_{m,r}\}_r\) énumère chaque coordonnée une seule fois ; les facteurs
du circuit seront identifiés ci-dessous au semi-groupe de leurs espérances
conditionnelles, et non à neuf dynamiques indépendantes.
Les blocs de même couleur possèdent des supports disjoints et le diamètre physique
de chaque support est uniformément borné lorsque \(m\) varie. Cette localité
fait partie de l’objet, de sorte que le raffinement ne transforme pas une
mise à jour élémentaire en une opération instantanément globale.

Soit \(T_m\) le transfert induit par
\eqref{eq:realizaciones-ym-circuito}. Les inclusions isométriques
\(J_m:\mathcal H_m\to\mathcal H_{m+1}\) conservent ancêtre,
représentation, orientation, retenue et mémoire. Pour les quatre réflexions
euclidiennes \(\Theta_{\nu,m}\), \(\nu=1,\ldots,4\), on a
\begin{equation}\label{eq:realizaciones-ym-entrelazamiento}
 (T_{m+1})^9J_m=J_mT_m,
 \qquad
 \Theta_{\nu,m+1}J_m=J_m\Theta_{\nu,m}
 \quad(\nu=1,\ldots,4).
\end{equation}
Les quatre directions appartiennent donc au même raffinement, à la même mesure et à la même dynamique.

\subsection*{Semi-groupe commun, couronne relative et quatre formes de réflexion}

Dans la coordinatisation triangulaire de l'état complet, soient \(E_{m,c}\) les espérances conditionnelles orthogonales et commutantes associées à des coordonnées indépendantes. Le générateur positif de recouvrement est d'abord construit dans un niveau germe \(m_0\) :
\begin{equation}\label{eq:realizaciones-ym-generador}
 \widehat H_{m_0}^{\rm ov}
 =3\kappa_{\rm av}
  \sum_{c\in\mathcal I_{m_0}^{\rm seed}}(I-E_{m_0,c}),
\end{equation}
et se prolonge au moyen de
\begin{equation}\label{eq:realizaciones-ym-generador-refinado}
 \begin{aligned}
 \widehat H_{m+1}^{\rm ov}
 &=\widehat\Gamma_m^*
   \bigl(\widehat H_m^{\rm ov}\otimes I+
         I\otimes\widehat H_{m+1}^{\rm det}\bigr)
   \widehat\Gamma_m,\\
 \widehat H_{m+1}^{\rm det}
 &=3\kappa_{\rm av}
   \sum_{\iota\in\mathcal I_{m+1/m}^{\rm det}}
   (I-E_{m+1,\iota}^{\rm det}),
 \end{aligned}
\end{equation}
où
\begin{equation}\label{eq:realizaciones-ym-kappa}
 \kappa_{\rm av}
 =\frac{\mu_{\rm vol}}{\ell_0}
  \left(\frac{\pi_{\rm HMT}}{\varphi_{\rm HMT}^{\,2}}-1\right)>0.
\end{equation}
La fréquence volumétrique \(\mu_{\rm vol}\) est sans dimension ; la dimension provient de \(\ell_0\).

Le semi-groupe, le transfert d'un pas physique et sa compatibilité d'échelle sont définis avec des types visibles :
\begin{equation}\label{eq:realizaciones-ym-semigrupo}
 \widehat K_{m,t}^{\rm ov}=e^{-t\widehat H_m^{\rm ov}},
 \qquad
 \widehat T_m=\widehat K_{m,a_m}^{\rm ov},
 \qquad a_{m+1}=a_m/9,
\end{equation}
Comme la partition énumère chaque coordonnée une seule fois et que les espérances
conditionnelles de ses blocs commutent, on obtient pour chaque bloc et pour le circuit
complet
\begin{equation}\label{eq:realizaciones-ym-circuito-semigrupo}
 K_{m,B}=e^{-3\kappa_{\rm av}a_m(I-E_{m,B})},
 \qquad
 K_m^{\rm loc}=e^{-a_m\widehat H_m^{\rm ov}}=\widehat T_m.
\end{equation}
\begin{equation}\label{eq:realizaciones-ym-semigrupo-entrelazado}
 \widehat H_{m+1}^{\rm ov}\widehat J_m
 =\widehat J_m\widehat H_m^{\rm ov},
 \qquad
 (\widehat T_{m+1})^9\widehat J_m
 =\widehat J_m\widehat T_m.
\end{equation}
Ainsi, la neuvième puissance est comparée au moyen de \(\widehat J_m\) ; on n’égale pas
des opérateurs qui agissent dans des Hilbert distincts.

La mesure commune de huit faces provient de ce même noyau réversible :
\begin{equation}\label{eq:realizaciones-ym-medida-ocho}
 \widehat\Omega_m^{(8)}(dy_1\cdots dy_8)
 =\int\widehat\mu_m^{\rm ov}(dz)
   \prod_{r=1}^{8}
   \widehat K_{m,a_m/2}^{\rm ov}(z,dy_r).
\end{equation}
Pour chaque réflexion \(\nu=1,\ldots,4\) et tout observable cylindrique \(F\)
dans le demi-espace positif,
\begin{equation}\label{eq:realizaciones-ym-os}
 \int\overline{F(\Theta_{\nu,m}y)}F(y)\,
 d\widehat\Omega_m^{(8)}
 =\|\widehat{\mathcal B}_{m,\nu}F\|^2\geq0,
\end{equation}
et la marginale de faces opposées donne l'identité de liaison
\begin{equation}\label{eq:realizaciones-ym-os-transferencia}
 \widehat T_{m,\nu}^{\rm OS}
 =\widehat K_{m,a_m}^{\rm ov}
 =e^{-a_m\widehat H_m^{\rm ov}}
 =\widehat T_m,
 \qquad \nu=1,\ldots,4.
\end{equation}
Les quatre réflexions expriment donc la même paire transfert--mesure.

\subsection*{Persistance cellulaire du terminal et écart spectral exact}

Soient
\[
 K_9=C_*(Q_{9,3};\mathbb Z),\qquad
 K_{10}=C_*(Q_{10,3};\mathbb Z).
\]
La couronne relative a pour dimensions
\[
 C_*(Q_{10,3},Q_{9,3})=(271,756,702,217),
 \qquad 271+702=756+217=973.
\]
L’inclusion \(j:K_9\to K_{10}\), l’effondrement cellulaire
\(r:K_{10}\to K_9\) et l’homotopie entière \(H_{\rm cell}\) satisfont
\begin{equation}\label{eq:realizaciones-ym-sdr-corona}
 rj=I,\qquad
 \partial H_{\rm cell}+H_{\rm cell}\partial=I-jr,\qquad
 rH_{\rm cell}=0,\quad H_{\rm cell}j=0,\quad H_{\rm cell}^2=0.
\end{equation}
Pour toute application de chaînes
\(F:K_{10}\otimes V_{q=6}\to\mathcal T\), avec
\(P=(jr)\otimes I_{q=6}\),
\begin{equation}\label{eq:realizaciones-ym-homotopia-terminal}
 F-FP
 =d_{\mathcal T}\bigl(F(H_{\rm cell}\otimes I)\bigr)
  +F(H_{\rm cell}\otimes I)(\partial\otimes I).
\end{equation}
La couronne complète est ainsi homotope à une constante et le terminal \(FP\) demeure littéralement. Les \(271\) sommets ne remplacent pas les arêtes, les faces et les cubes de la couronne ; le facteur cellulaire tridimensionnel n'est pas non plus identifié au facteur de jauge physique \(q=6\).

Pour chaque coordonnée persistante \(c\), dans la coordinatisation locale \(q=(q_1,\ldots,q_6)\in\mathbb F_3^6\), on fixe le témoin centré
\begin{equation}\label{eq:realizaciones-ym-testigo-definicion}
 \bigcap_c\operatorname{Ran}E_{m,c}=\mathbb C\Omega_m,
 \qquad
 W_c(q)=\mathbf 1_{\{q_1+\cdots+q_6=0\}}-\frac13,
\end{equation}
où la somme de l’indicatrice est prise dans \(\mathbb F_3\). La décomposition de
Hoeffding des espérances commutantes prouve
\begin{equation}\label{eq:realizaciones-ym-gap-finito}
 \ker\widehat H_m^{\rm ov}=\mathbb C\Omega_m,
 \qquad
 \widehat H_m^{\rm ov}\succeq
 3\kappa_{\rm av}(I-P_{\Omega_m}).
\end{equation}
L’égalité ne se déduit pas du seul fait que le noyau est trivial : chaque secteur non
constant de la décomposition reçoit au moins un projecteur
\(I-E_{m,c}\) et, par \eqref{eq:realizaciones-ym-generador}, son coût minimal
est \(3\kappa_{\rm av}\). Le témoin matériel \(W_c\) vérifie
\begin{equation}\label{eq:realizaciones-ym-testigo-gap}
 \widehat H_m^{\rm ov}W_c=3\kappa_{\rm av}W_c,
 \qquad \|W_c\|^2=\frac29,
\end{equation}
de sorte que le complément est non vide et que l’écart est exactement
\(3\kappa_{\rm av}\).

L'échelle s'exprime par
\begin{equation}\label{eq:realizaciones-ym-qm}
 q_m=e^{-3\kappa_{\rm av}a_m},
 \qquad q_{m+1}^{\,9}=q_m,
 \qquad -a_m^{-1}\log q_m=3\kappa_{\rm av}.
\end{equation}
Bien que \(q_m\to1\) lors du raffinement, le quotient spectral par unité physique reste constant. La contraction cellulaire \eqref{eq:realizaciones-ym-homotopia-terminal} conserve le terminal ; l'écart spectral provient du générateur positif de recouvrement et de sa décomposition de Hoeffding, et non de l'acyclicité à elle seule.

Les inclusions de \eqref{eq:realizaciones-ym-semigrupo-entrelazado}
forment la limite inductive
\begin{equation}\label{eq:realizaciones-ym-colimite}
 \mathcal H_\infty=\varinjlim(\mathcal H_m,\widehat J_m).
\end{equation}
Sur l’union inductive algébrique, on définit la forme
\begin{equation}\label{eq:realizaciones-ym-forma-limite}
 \mathfrak h_\infty[\widehat J_{m,\infty}\psi]
 :=\langle\psi,\widehat H_m^{\rm ov}\psi\rangle_{\mathcal H_m}.
\end{equation}
L’entrelacement des générateurs rend cette définition
indépendante du représentant. La forme est positive et fermable ; sa
fermeture détermine le générateur auto-adjoint \(\widehat H_\infty^{\rm ov}\).
La compatibilité des formes transporte
\begin{equation}\label{eq:realizaciones-ym-gap-forma-limite}
 \overline{\mathfrak h}_\infty[\psi]
 \geq3\kappa_{\rm av}
 \|(I-P_{\Omega_\infty})\psi\|^2,
\end{equation}
et le vecteur persistant
\(\widehat J_{m,\infty}W_c\) atteint l’égalité. Par conséquent,
l’écart exact du générateur limite est \(3\kappa_{\rm av}\).

\subsection*{Hypothèses exactes et enchaînement du théorème}

La représentation de jauge utilise :
\begin{enumerate}
 \item le groupe compact, connexe et simple et le réseau quadridimensionnel
       \eqref{eq:realizaciones-ym-red} ;
 \item le circuit local \eqref{eq:realizaciones-ym-circuito}, de diamètre physique uniforme, et une mesure invariante commune ;
 \item les inclusions isométriques et les quatre réflexions entrelacées par
       \eqref{eq:realizaciones-ym-entrelazamiento} et
       \eqref{eq:realizaciones-ym-semigrupo-entrelazado} ;
 \item les quatre factorisations
       \eqref{eq:realizaciones-ym-os} et le lien exact
       \eqref{eq:realizaciones-ym-os-transferencia} sur le même couple
       transfert--mesure ;
 \item le générateur d'espérances commutantes \eqref{eq:realizaciones-ym-generador}--\eqref{eq:realizaciones-ym-generador-refinado} et sa décomposition de Hoeffding ;
 \item le SDR entier de la couronne et la conservation du terminal
       \eqref{eq:realizaciones-ym-sdr-corona}--
       \eqref{eq:realizaciones-ym-homotopia-terminal} ;
 \item le terminal \(P_{\Omega_m}\), le complément physique
       \(Q_m^{\rm phys}\) et la loi d’échelle
       \eqref{eq:realizaciones-ym-qm} ;
 \item la limite inductive et la fermeture compatible des formes \eqref{eq:realizaciones-ym-colimite}--\eqref{eq:realizaciones-ym-gap-forma-limite} ;
 \item les coordinatisations internes d'action et de vitesse \(\hbar_{\rm int}\), \(c_{\rm int}\), utilisées seulement pour exprimer la dimension de masse.
\end{enumerate}

\begin{theorem}[Théorie de jauge locale à vide simple et écart collectif]
\label{thm:realizaciones-ym}
Sous les prémisses précédentes, la limite de la collection de jauge est locale et positive par réflexion dans les quatre directions euclidiennes. Son générateur possède un vide simple et satisfait, dans le secteur collectif invariant de jauge,
\begin{equation}\label{eq:realizaciones-ym-gap-limite}
 \Delta_{\rm YM}^{\rm HMT}=3\kappa_{\rm av}>0,
 \qquad
 m_{\rm gap}^{\rm HMT}
 =\frac{\hbar_{\rm int}}{c_{\rm int}}
  \Delta_{\rm YM}^{\rm HMT}>0.
\end{equation}
L’écart appartient au secteur collectif ; il n’assigne pas de masse élémentaire au
gluon.
\end{theorem}

\begin{proof}
La localité uniforme du circuit conserve le réseau d’algèbres locales. Les
égalités \eqref{eq:realizaciones-ym-semigrupo-entrelazado} identifient la
neuvième puissance après application de l’inclusion correcte. La construction
\eqref{eq:realizaciones-ym-medida-ocho} utilise ce même semi-groupe dans les huit
faces et \eqref{eq:realizaciones-ym-os-transferencia} identifie ses quatre
marginales OS à \(\widehat T_m\). Par conséquent, les quatre formes positives
\eqref{eq:realizaciones-ym-os} sont compatibles à la limite et conservent une
seule mesure et une seule dynamique.

Les espérances de \eqref{eq:realizaciones-ym-generador} sont orthogonales et
commutantes. Leur décomposition de Hoeffding sépare le mode constant de chaque
secteur non constant : le premier est \(\mathbb C\Omega_m\), tandis que tout
secteur du complément reçoit au moins un facteur \(I-E_{m,c}\) de poids
\(3\kappa_{\rm av}\). Cela prouve
\eqref{eq:realizaciones-ym-gap-finito}. Le témoin
\eqref{eq:realizaciones-ym-testigo-gap} atteint la borne, de sorte que ce n’est pas
seulement un plancher : c’est la première énergie non nulle.

L'homotopie \eqref{eq:realizaciones-ym-homotopia-terminal} contracte le complément cellulaire et laisse \(FP\) intact ; la couronne n'introduit donc pas de second vide et n'efface pas la retenue. La compatibilité isométrique transporte le vide et son complément au niveau suivant sans les mélanger au résidu d'échelle. La loi \eqref{eq:realizaciones-ym-qm} montre que le quotient spectral par unité physique est exactement \(3\kappa_{\rm av}\) à tout niveau. Les formes \eqref{eq:realizaciones-ym-forma-limite}--\eqref{eq:realizaciones-ym-gap-forma-limite} transportent la borne et le témoin vers le générateur autoadjoint de la limite ; le saut reste donc exact. La reconstruction par positivité d'Osterwalder--Schrader identifie l'hamiltonien au générateur de ce même transfert, conserve le vide simple et produit \eqref{eq:realizaciones-ym-gap-limite}. La coordinatisation dimensionnelle finale transforme le saut en la masse collective indiquée.
\end{proof}

\subsection*{Compression de jauge et témoin physique de l'écart spectral}

Soit
\[
 \widehat{\mathscr H}_m
 :=L^2(\widehat X_m,\widehat\mu_m^{\rm ov})
\]
l'espace de la réalisation matérielle complète. Outre l'action de jauge usuelle sur les liaisons et les repères, chaque collier \(c\) porte l'action finie d'incidence
\begin{equation}\label{eq:amp-ym-accion-gauge-incidencia}
 q_c\longmapsto q_c+Bx_c,
 \qquad x_c\in\mathbb F_3^4,
\end{equation}

où les six lignes de \(B\) sont indexées par
\((01,02,03,12,13,23)\). Le produit des deux actions préserve
\(\widehat\mu_m^{\rm ov}\). Avec \(\mathcal C_m\) l’ensemble fini des
colliers présents au niveau, on écrit
\[
 \mathcal G_m^{\rm full}
 :=G^{V(\Lambda_m)}
   \times\prod_{c\in\mathcal C_m}\mathbb F_3^4.
\]
Sa moyenne de Haar définit la projection
orthogonale
\begin{equation}\label{eq:amp-ym-proyeccion-fisica-haar}
 \mathbb P_m^{\rm phys}F
 :=\int_{\mathcal G_m^{\rm full}}F(g^{-1}\!\cdot\,)\,dg,
 \qquad
 \widehat{\mathscr H}_m^{\rm phys}
 :=\operatorname{Ran}\mathbb P_m^{\rm phys}.
\end{equation}
C'est une compression sur la sous-algèbre de jauge de l'état complet ; la marginale qui oublie les coordonnées d'incidence est une autre application et n'est pas confondue avec \(\mathbb P_m^{\rm phys}\).

\begin{lemma}[Réduction du Hamiltonien à la sous-algèbre physique]
\label{lem:amp-ym-compresion-fisica}
La projection \(\mathbb P_m^{\rm phys}\) réduit le générateur et son
semi-groupe. Les inclusions de raffinement réduisent simultanément les
sous-algèbres physiques :
\begin{equation}\label{eq:amp-ym-compresion-entrelazada}
 \begin{gathered}
 \relax
 [\mathbb P_m^{\rm phys},\widehat H_m^{\rm ov}]=0,
 \qquad
 [\mathbb P_m^{\rm phys},e^{-t\widehat H_m^{\rm ov}}]=0,\\
 \mathbb P_{m+1}^{\rm phys}\widehat J_m
 =\widehat J_m\mathbb P_m^{\rm phys},\qquad
 H_m^{\rm phys}
 :=\widehat H_m^{\rm ov}\big|_{\widehat{\mathscr H}_m^{\rm phys}},\\
 H_{m+1}^{\rm phys}\widehat J_m
 =\widehat J_mH_m^{\rm phys}.
 \end{gathered}
\end{equation}
\end{lemma}

\begin{proof}
Le générateur physique sur les liens est équivariant de jauge par la bi-invariance de Haar. Dans la coordinatisation produit, \(E_{q_c}\) est l'espérance uniforme sur \(\mathbb F_3^6\) ; par invariance de la mesure uniforme sous les translations de \eqref{eq:amp-ym-accion-gauge-incidencia}, il commute avec cette action. Les espérances des sous-queues de marquage agissent sur des facteurs disjoints. Par conséquent, chaque terme de
\[
 \widehat H_m^{\rm ov}
 =H_m^{\rm ov}\otimes I
 +3\kappa_{\rm av}\sum_c
 \bigl[(I-E_{q_c})+(I-E_{u_c^{\rm P}})\bigr]
\]
commute avec la moyenne \eqref{eq:amp-ym-proyeccion-fisica-haar}. Le calcul fonctionnel donne la commutation avec le semi-groupe. Enfin, \(\widehat J_m\) conserve les facteurs ancestraux et insère la constante dans les nouveaux facteurs ; prendre la moyenne avant ou après cette inclusion produit la même fonction. Cela prouve les quatre identités de \eqref{eq:amp-ym-compresion-entrelazada}.
\end{proof}

\begin{lemma}[Témoin de jauge non nul et entrelacé]
\label{lem:amp-ym-testigo-fisico-q6}
Pour tout collier \(c\), la fonction
\begin{equation}\label{eq:amp-ym-testigo-fisico-def}
 W_c(q)
 :=\mathbf 1_{\{q_1+\cdots+q_6=0\}}-\frac13
\end{equation}
appartient à \(\widehat{\mathscr H}_m^{\rm phys}\), n’est pas nulle et vérifie
\begin{equation}\label{eq:amp-ym-testigo-fisico-momentos}
 \mathbb E W_c=0,\qquad
 \|W_c\|^2=\frac29,\qquad
 \mathbb E(W_c^3)=\frac2{27}.
\end{equation}
De plus,
\begin{equation}\label{eq:amp-ym-testigo-fisico-gap}
 H_m^{\rm phys}W_c=3\kappa_{\rm av}W_c,\qquad
 H_{m+1}^{\rm phys}\widehat J_mW_c
 =3\kappa_{\rm av}\widehat J_mW_c.
\end{equation}
Par conséquent, la valeur \(3\kappa_{\rm av}\) appartient au spectre de la
restriction physique et non seulement à celui d’une dilatation matérielle.
\end{lemma}

\begin{proof}
L’action de jauge sur les liens agit sur le premier facteur et laisse fixe
\(W_c\). Pour l’action d’incidence, chaque \(x_i\) apparaît dans trois des
six coordonnées, donc
\[
 \sum_{i<j}(Bx)_{ij}=3\sum_i x_i=0
 \quad\text{en }\mathbb F_3.
\]
La condition de \eqref{eq:amp-ym-testigo-fisico-def} est donc
invariante ; elle l’est aussi sous les permutations des six incidences.
Ainsi \(\mathbb P_m^{\rm phys}W_c=W_c\).

Sous la mesure uniforme, la somme des six trits est uniforme dans \(\mathbb F_3\). La fonction prend la valeur \(2/3\) avec probabilité \(1/3\) et \(-1/3\) avec probabilité \(2/3\), ce qui donne \eqref{eq:amp-ym-testigo-fisico-momentos}. Le premier terme de \(\widehat H_m^{\rm ov}\) annule \(W_c\) ; pour \(c'\ne c\), \((I-E_{q_{c'}})W_c=0\), tandis que \(E_{q_c}W_c=0\). Toutes les espérances de marquage annulent la différence correspondante car \(W_c\) est constante dans ces sous-queues. Il reste exactement \(\widehat H_m^{\rm ov}W_c=3\kappa_{\rm av}W_c\). La première identité de \eqref{eq:amp-ym-testigo-fisico-gap} découle de la réduction ; la seconde, de l'entrelacement de \eqref{eq:amp-ym-compresion-entrelazada}.
\end{proof}

\subsection*{Semi-groupe de jauge, vide et borne spectrale exacte}

Sur la sous-algèbre physique, on écrit
\[
 D_m^{\rm YM}:=H_m^{\rm phys},\qquad
 K_{m,t}:=\exp(-tD_m^{\rm YM}),\qquad
 P_{\Omega_m}:=|\Omega_m\rangle\langle\Omega_m|.
\]
La décomposition de Hoeffding du même générateur et le calcul fonctionnel
donnent, sans changer de mesure ni de processus,
\begin{equation}\label{eq:amp-ym-owner-semigrupo-vacio}
 K_{m,t}P_{\Omega_m}=P_{\Omega_m},
 \qquad
 \bigl\|K_{m,t}(I-P_{\Omega_m})\bigr\|
 \leq e^{-3\kappa_{\rm av}t}.
\end{equation}
Le témoin de \eqref{eq:amp-ym-testigo-fisico-def} est non nul et atteint la
borne :
\begin{equation}\label{eq:amp-ym-owner-testigo-wc}
 D_m^{\rm YM}W_c=3\kappa_{\rm av}W_c,
 \qquad \|W_c\|^2=\frac29.
\end{equation}
Par conséquent, le vide est simple et l’écart n’est pas seulement minoré,
mais vaut exactement
\begin{equation}\label{eq:amp-ym-owner-brecha-exacta}
 \Delta_{\rm YM}^{\rm HMT}=3\kappa_{\rm av}>0.
\end{equation}
Le semi-groupe, le projecteur du vide, la contraction du complément et \(W_c\) appartiennent à la même collection projective et constituent la chaîne de référence de la fermeture.

\subsection*{Reconstruction euclidienne, champs de Schwinger et lecteur de courbure}

La mesure projective commune du théorème détermine des inclusions isométriques
\(\widehat J_m\) entre les espaces de niveau. Pour chacune des quatre
réflexions coordonnées \(\theta_\mu\), \(\mu=1,\ldots,4\), il existe une
factorisation
\begin{equation}\label{eq:amp-ym-cuatro-os}
 \int\overline{F(\theta_\mu y)}F(y)\,
 d\widehat\Omega_m(y)
 =\|\mathcal B_{m,\mu}F\|_2^2\geq0.
\end{equation}
Les quatre formes proviennent de la même mesure et du même semi-groupe. Les
connecteurs OS
\[
 \mathcal J_{m,\mu}^{\rm OS}
 =\mathcal V_{m+1,\mu}^{-1}\widehat J_m\mathcal V_{m,\mu}
\]
entrelacent leurs hamiltoniens. Dans la colimite,
\begin{equation}\label{eq:amp-ym-hamiltoniano-comun}
 \mathcal V_{\infty,\mu}H_{{\rm OS},\infty}^{(\mu)}
 \mathcal V_{\infty,\mu}^{-1}
 =\widehat H_\infty,
 \qquad \mu=1,\ldots,4.
\end{equation}

La forme limite se diagonalise par la décomposition de Hoeffding :
\begin{equation}\label{eq:amp-ym-forma-limite}
 \mathcal E_\infty(u)=3\kappa_{\rm av}
 \sum_{S\Subset\mathcal I_\infty}|S|\,\|u_S\|^2.
\end{equation}
Les sommes partielles fournissent une suite de récupération et la semi-continuité produit la limite inférieure. La convergence de Mosco des formes fermées conserve le noyau unidimensionnel et le seuil spectral. Le vecteur matériel \(W_c\) satisfait
\begin{equation}\label{eq:amp-ym-testigo-gap}
 \widehat H_mW_c=3\kappa_{\rm av}W_c,
 \qquad \|W_c\|^2=\frac29,
 \qquad \mathbb E(W_c^3)=\frac2{27},
\end{equation}
de sorte que la borne inférieure est atteinte et que l’écart est exactement
\(3\kappa_{\rm av}\).

\begin{theorem}[Achèvement euclidien et représentation de courbure]
\label{thm:amp-ym-terminacion}
La collection projective du théorème \ref{thm:realizaciones-ym} détermine une action fortement continue de \(E(4)\), un réseau local de jauge non scalaire, des fonctionnelles de Schwinger compatibles et une représentation de champ dans \(H^{-s}_{\rm loc}(\mathbb R^4)\) pour \(s>2\). Le produit surfacique ordonné des liens définit les lecteurs clover et \(F^2\) ; dans le domaine lisse, le lecteur d'action converge vers
\begin{equation}\label{eq:amp-ym-f2}
 \mathcal S_{\rm YM}(A)
 =\frac1{4g_R^2}\int_{\mathbb R^4}\|F_A(x)\|^2\,d^4x.
\end{equation}
Ces représentations conservent le vide simple et le saut exact du même hamiltonien commun.
\end{theorem}

\begin{proof}
Les matrices de Gram cofinales des lecteurs lissés et la commutativité des carrés de rotation et de raffinement préservent les idéaux nuls ; leur colimite construit la représentation forte de \(E(4)\) et le réseau local. Les estimations de martingale des accroissements de liaison donnent la tension dans \(H^{-s}_{\rm loc}\), \(s>2\), et les caractéristiques mixtes identifient les fonctionnelles de Schwinger au même processus cylindrique. Le produit orienté autour de chaque plaquette construit le clover. Son développement dans le régime lisse et la convergence forte du lecteur marginal donnent \eqref{eq:amp-ym-f2}. Aucune de ces opérations ne modifie la mesure, le semi-groupe ni la forme limite utilisés pour obtenir \eqref{eq:amp-ym-testigo-gap} ; le vide et l'écart spectral sont donc conservés.
\end{proof}

\subsection*{Identification exacte de la limite de jauge}

L’identification de la théorie ne repose pas sur le fait que quatre réflexions
engendrent à elles seules le groupe euclidien. Les réflexions fournissent la
positivité OS ; l’action continue provient, séparément, du groupoïde des
écrans et de la compatibilité de raffinement. Pour
\(g=(g_v)_{v\in\Lambda_m}\in G^{\Lambda_m}\), l’action sur une variable
de lien est
\begin{equation}\label{eq:amp-ym-accion-gauge-enlace}
 (g\cdot U)_e=g_{s(e)}U_e g_{t(e)}^{-1}.
\end{equation}
L’invariance de Haar, la désintégration triangulaire du noyau et
l’annulation des arêtes intérieures prouvent
\begin{equation}\label{eq:amp-ym-invariancia-gauge-medida}
 (g\cdot)_*\widehat\mu_m^{\rm ov}=\widehat\mu_m^{\rm ov},
 \qquad
 K_{m,t}^{\rm ov}(g\cdot U,g\cdot dV)
 =K_{m,t}^{\rm ov}(U,dV).
\end{equation}
Par différentiation sur le noyau cylindrique, tout générateur vertical
\(X\) vérifie l’identité de Ward
\begin{equation}\label{eq:amp-ym-ward}
 \int \nabla_XF\,d\widehat\mu_m^{\rm ov}=0.
\end{equation}

La courbure s’obtient à partir du même système de liens. Si une plaquette
grossière \(p\) est subdivisée en ses 81 sous-plaquettes \(p'\), alors
\begin{equation}\label{eq:amp-ym-producto-superficial}
 \operatorname{Hol}^{(m)}_{p}
 =\prod_{p'\subset p}^{\longrightarrow}
 T_{p'}\operatorname{Hol}^{(m+1)}_{p'}T_{p'}^{-1}.
\end{equation}
Les contributions de toute arête intérieure apparaissent avec des orientations opposées et s'annulent exactement. L'expression est covariante de jauge et commute avec les inclusions \(\widehat J_m\). Dans une coordinatisation lisse, on définit
\begin{equation}\label{eq:amp-ym-clover-definido}
 F_{{\rm cl},m}(x)
 =\frac1{4a_m^2}\sum_{p\ni x}
 \operatorname{Ad}_{T_p}\log\operatorname{Hol}_{p},
 \qquad
 F_{{\rm cl},m}=F_A+O(a_m^2).
\end{equation}
De là, pour \(A\in C^4_{\rm loc}\) et
\(a_m^2/g_m^2\to0\), on obtient par somme de Riemann
\begin{equation}\label{eq:amp-ym-f2-convergencia-calculada}
 \frac{a_m^4}{4g_m^2}
 \sum_{x\in\Lambda_m}\|F_{{\rm cl},m}(x)\|^2
 \longrightarrow
 \frac1{4g_R^2}\int_{\mathbb R^4}\|F_A(x)\|^2\,d^4x.
\end{equation}

\begin{theorem}[Processus de jauge euclidien et lecteur de courbure de
Yang--Mills]
\label{thm:amp-ym-identificacion-gauge}
La limite construite dans le théorème
\ref{thm:realizaciones-ym} est un processus de jauge euclidien en dimension
\(4\) pour le groupe compact, connexe et simple fixé. Elle possède
une covariance de jauge locale, un réseau local non scalaire, une positivité par
réflexion, une représentation fortement continue de \(E(4)\),
des fonctionnelles de Schwinger compatibles, une courbure covariante de jauge de limite
classique \eqref{eq:amp-ym-f2-convergencia-calculada}, un vide simple et un écart
spectral
\(3\kappa_{\rm av}>0\).
\end{theorem}

\begin{proof}
Les identités \eqref{eq:amp-ym-invariancia-gauge-medida}--
\eqref{eq:amp-ym-ward} construisent la symétrie de jauge et ses identités de
Ward à tous les niveaux ; leur naturalité les transporte à la limite.  La
localité uniforme des noyaux produit le réseau isotone et local.  Les
quatre formes \eqref{eq:amp-ym-cuatro-os}, avec l’action du groupoïde
des écrans sur le noyau de Weyl, produisent respectivement la positivité
OS et l’action forte de \(E(4)\).  La tension dans
\(H^{-s}_{\rm loc}\), \(s>2\), identifie une unique loi limite et ses
fonctionnelles de Schwinger.  Enfin,
\eqref{eq:amp-ym-producto-superficial}--
\eqref{eq:amp-ym-f2-convergencia-calculada} identifient le lecteur de champ
et sa limite classique, tandis que
\eqref{eq:amp-ym-forma-limite}--
\eqref{eq:amp-ym-testigo-gap} prouvent le vide simple et l’écart exact pour la
même loi.

La mesure est déterminée constructivement par ses marginales projectives, son noyau réversible et ses fonctionnelles cylindriques. Le lecteur clover ne repondère pas cette mesure : il exprime une fonction du même processus et sa limite classique.
\end{proof}

\begin{proposition}[Contrôle non directeur de comparaison Gibbs--action]
\label{prop:amp-ym-puente-accion-requerido}
La loi projective, son vide simple et son saut exact sont déjà construits par \eqref{eq:amp-ym-owner-semigrupo-vacio}--\eqref{eq:amp-ym-owner-brecha-exacta}. Comme comparaison extérieure optionnelle de type Gibbs--action, une paramétrisation par \(g_R\) peut étudier une collection
\[
 g_R\longmapsto
 \bigl(\widehat\mu_{m,g_R}^{\rm ov},H_{m,g_R}^{\rm phys}\bigr)
\]
et une relation entre la mesure ou sa forme de Dirichlet et la fonctionnelle
d’action, par exemple
\begin{equation}\label{eq:amp-ym-identidad-accion-requerida}
 \frac{d\widehat\mu_{m,g_R}^{\rm ov}}{d\nu_m}(U)
 =Z_{m,g_R}^{-1}e^{-S_{m,g_R}^{F^2}(U)},
 \qquad
 \mathfrak h_{m,g_R}[F]
 =\langle F,H_{m,g_R}^{\rm phys}F\rangle,
\end{equation}
avec compatibilité de raffinement et passage à la limite. Ici, \(\nu_m\) peut
être la mesure produit de Haar du réseau fini. Cette vérification est étiquetée
\texttt{CONTRÔLE\_NON\_DIRECTEUR} : elle n’est une prémisse ni de la théorie de jauge HMT ni
de son écart.
\end{proposition}

\begin{proof}
Dans \eqref{eq:amp-ym-f2-convergencia-calculada}, faire varier \(g_R\) remet à l'échelle le lecteur de courbure, tandis que les formules qui construisent \(\widehat\mu_m^{\rm ov}\), \(\widehat H_m^{\rm ov}\) et l'écart spectral \(3\kappa_{\rm av}\) ne contiennent pas \(g_R\). Ce sont donc des constructions indépendantes. La compression \eqref{eq:amp-ym-compresion-entrelazada} et le témoin \eqref{eq:amp-ym-testigo-fisico-gap} prouvent la persistance physique de l'écart spectral du processus déjà construit ; le contrôle compare seulement ensuite une dépendance dynamique supplémentaire en \(g_R\).
\end{proof}

La suite démontrée dans cette section est
\[
 \text{mesure commune}\to\text{quatre OS}\to\text{Hamiltonien commun}
 \to E(4),\ \text{Schwinger},\ H^{-s}_{\rm loc}
 \to\text{clover et }F^2.
\]
Cette suite démontre une théorie de Yang--Mills quadridimensionnelle avec un vide
simple et un écart positif exact \(3\kappa_{\rm av}\) dans le domaine HMT
enregistré. L’identité
\eqref{eq:amp-ym-identidad-accion-requerida} reste en dehors de la chaîne
régissante comme comparaison ultérieure indépendante.

La réalisation dépendant de \(g\) est développée dans
\eqref{eq:pdf2-ym-gibbs-measure} et
\eqref{eq:pdf2-ym-g-dynamics} : la mesure, les espérances
conditionnelles et le générateur sont transportés conjointement,
et \eqref{eq:pdf2-ym-g-gap-before-limit} conserve la borne et son témoin.


\subsection*{Lecture de Wilson}

L'observable de Wilson est obtenu après le générateur. Pour une coordinatisation archimédienne de lien,
\begin{equation}\label{eq:realizaciones-ym-wilson}
 A_e=\sum_a\operatorname{ev}_{\mathbb R}(x_{e,a})T_a,
 \qquad
 U_e=e^{a_mA_e},
 \qquad
 W_\gamma=\operatorname{Tr}\!\prod_{e\in\gamma}U_e.
\end{equation}
La lecture conserve la représentation, l'opérateur d'entrelacement et l'ordre de l'holonomie. Elle sert d'observable terminal de la théorie exprimée ; elle ne définit pas rétrospectivement la mesure, le semi-groupe ni le saut.

\subsection*{Nécessité des hypothèses et obstructions exactes}

\begin{ablation}[Nécessité de la dynamique et de la mesure communes]
\label{abl:realizaciones-ym}
Le théorème \ref{thm:realizaciones-ym} n’est pas conservé si l’on remplace les
neuf sous-phases par des mises à jour indépendantes, si l’on utilise des mesures ou
des noyaux différents dans les quatre réflexions, si l’on perd la localité physique
uniforme, si \(\widehat J_m\) cesse d’entrelacer transfert et réflexion, si l’on
efface la retenue ou la mémoire, si l’on élimine \(P_{\Omega_m}\), si l’on identifie
\(Q_m^{\rm phys}\) à un résidu d’échelle, si l’on exige
\(q_m\leq q_*<1\) en unités de réseau, si l’on remplace la couronne complète
par ses \(271\) sommets, si l’on efface \(FP\), ou si l’on utilise
\eqref{eq:realizaciones-ym-wilson} pour sélectionner le générateur.
\end{ablation}

\begin{proof}
Des actualisations, noyaux ou mesures distincts détruisent la factorisation
simultanée \eqref{eq:realizaciones-ym-os}. Sans localité ou entrelacement,
la positivité d’une coupe ne définit pas un réseau local compatible. Effacer la retenue,
la mémoire ou \(FP\) élimine de l’information que
\eqref{eq:realizaciones-ym-homotopia-terminal} conserve. Utiliser seulement les
sommets change le complexe relatif et peut fabriquer des modes nuls. Confondre
les deux compléments permet que le secteur physique disparaisse par une
opération purement scalaire. La borne fixée en unités de réseau contredit
\eqref{eq:realizaciones-ym-qm}. Enfin, Wilson n’est qu’une lecture de
l’état construit et ne démontre par lui-même ni
\eqref{eq:realizaciones-ym-os} ni
\eqref{eq:realizaciones-ym-gap-finito}.
\end{proof}

Les résidus indépendants sont : l'échec de \(\widehat T_{m,\nu}^{\rm OS}=e^{-a_m\widehat H_m^{\rm ov}}\), l'échec de l'entrelacement \eqref{eq:realizaciones-ym-semigrupo-entrelazado}, une forme OS négative, un vecteur du noyau orthogonal à \(\Omega_m\), ou un secteur de Hoeffding non constant d'énergie inférieure à \(3\kappa_{\rm av}\).

\subsection*{Portée de la construction de jauge}

La preuve combine la connexion conjointe, la conservation du terme terminal, la double lecture, le circuit local, la mesure commune, les quatre factorisations de réflexion, le raffinement entrelacé et le générateur positif. Dans le domaine régulier de jauge déclaré, la localité, la théorie de Yang--Mills quadridimensionnelle, le vide simple et le saut collectif sont des résultats exacts. La reconstruction d'Osterwalder--Schrader et l'observable de Wilson sont des représentations ultérieures du même état ; aucune n'agit comme sélecteur rétrospectif de la mesure, du semi-groupe, du vide ou du saut. La comparaison optionnelle avec une collection paramétrée par \(g_R\) est étudiée au moyen de \eqref{eq:amp-ym-identidad-accion-requerida} ; cette vérification ne gouverne ni ne remet en question la mesure, le semi-groupe, le vide ou le saut déjà démontrés.


\subsection{Incidence projective de jauge du constructeur commun et descente au quotient}

L'image de jauge n'est pas la marginale des liaisons. Son objet complet par niveau est
\begin{equation}\label{eq:pdf2-g-gau-explicito}
 \mathfrak G_{{\rm gau},m}(d_{\rm gau})=
 \bigl(
 \begin{gathered}
  \mathfrak G_{\rm gau}^{\rm int};\\
  \Lambda_m,\widehat X_m,\widehat\mu_m^{\rm ov},
  \{\widehat\mu_{m,g}^{\rm YM},\nu_{m,g},K_{m,g}^{\rm phys}\}_{g>0},\\
  \mathcal G_m^{\rm full},\mathbb P_m^{\rm phys},
  \widehat J_m,\widehat\Gamma_m,
  \{E_{m,c}\}_c,\widehat H_m^{\rm ov},D_m^{\rm YM},\\
  \{\Theta_{\nu,m}\}_{\nu=1}^4,
  \{\mathfrak A_{+,\nu,m},\operatorname{ev}_{t}^{(\nu)},
      \mathcal V_{m,\nu}^{\rm OS}\}_{\nu=1}^4,\\
  \widehat\Omega_m^{(8)},P_{\Omega_m},W_c,
  \operatorname{Carr}_m,\operatorname{Mem}_m,
  \operatorname{Surv}_m,\tau_m,\mathcal P_m^{F^2}
 \end{gathered}
 \bigr).
\end{equation}
Le groupe $\mathcal G_m^{\rm full}$ agit simultanément sur les liens, les repères et les
coordonnées d’incidence $q_c\in\mathbb F_3^6$. La projection physique est sa moyenne
de Haar. Le générateur réduit est
\begin{equation}\label{eq:pdf2-g-gau-generador-estructural}
 D_m^{\rm YM}
 =\left.3\kappa_{\rm av}\sum_c(I-E_{m,c})
  \right|_{\operatorname{Ran}\mathbb P_m^{\rm phys}},
\end{equation}
et $W_c$ appartient à cette sous-algèbre. Les quatre réflexions agissent sur une même mesure, un même noyau et un même terminal.
Le propriétaire exact précédent construit l'action de jauge d'incidence \eqref{eq:amp-ym-accion-gauge-incidencia}, la projection physique de Haar \eqref{eq:amp-ym-proyeccion-fisica-haar} et la compression entrelacée \eqref{eq:amp-ym-compresion-entrelazada}. De cette construction proviennent le noyau complet équivariant $\widehat K_{m,t}^{\rm ov}$, sa descente au quotient physique $K_{m,t}^{\rm phys}$, les quatre opérateurs de réflexion et l'action $U_m(g)$ de $E(4)$ sur le cœur cylindrique lisse. Ce bloc assemble, sans les introduire comme données, le générateur, le vide, le témoin de courbure et les formes OS :
\begin{equation}\label{eq:pdf2-g-gau-entrelazamiento-dato}
 D_{m+1}^{\rm YM}\widehat J_m=\widehat J_mD_m^{\rm YM},
 \quad
 \mathbb P_{m+1}^{\rm phys}\widehat J_m
  =\widehat J_m\mathbb P_m^{\rm phys},
 \quad
 \widehat J_mW_{m,c}=W_{m+1,c}.
\end{equation}
Les première et troisième identités sont la représentation, dans la notation commune, de \eqref{eq:amp-ym-compresion-entrelazada} et \eqref{eq:amp-ym-testigo-fisico-gap} ; le semi-groupe, le vide simple et le saut exact sont déjà obtenus dans \eqref{eq:amp-ym-owner-semigrupo-vacio}--\eqref{eq:amp-ym-owner-brecha-exacta}. Ici, $\nu_{m,g}:=(\pi_m^{\rm phys})_*\widehat\mu_{m,g}^{\rm YM}$ et $K_{m,g}^{\rm phys}(t)=e^{-tD_{m,g}^{\rm YM}}$ dans le quotient physique. La complétude de jauge inclut en outre la désintégration pour chaque collection d'hyperplans euclidiens. Si $A_j\in\mathfrak A_{+,\nu,m}$ et $t_1\le\cdots\le t_n$, la marginale de la même mesure physique satisfait
\begin{equation}\label{eq:pdf2-g-gau-transfer-euclidean-dato}
\begin{aligned}
 &\int\prod_{j=1}^n
   (\tau_{t_je_\nu}A_j)(x)\,d\widehat\mu_{m,g}^{\rm YM}(x)\\
 &\quad=
 \int\prod_{j=1}^nA_j(y_j)\,
  \nu_{m,g}(dy_1)
  \prod_{j=1}^{n-1}K_{m,g}^{\rm phys}
       (t_{j+1}-t_j;y_j,dy_{j+1})\\
 &\quad=
 \left\langle\mathbf1,
  M_{A_1}e^{-(t_2-t_1)D_{m,g}^{\rm YM}}M_{A_2}\cdots
  e^{-(t_n-t_{n-1})D_{m,g}^{\rm YM}}M_{A_n}\mathbf1
 \right\rangle.
\end{aligned}
\end{equation}
La reconstruction précédente prouve l'identité dans les quatre directions au moyen de \eqref{eq:amp-ym-cuatro-os} et construit un unique hamiltonien entrelacé dans \eqref{eq:amp-ym-hamiltoniano-comun}. Par conséquent, la translation OS s'exprime ici par
\begin{equation}\label{eq:pdf2-g-gau-os-generator-dato}
 \mathcal V_{m,\nu}^{\rm OS}T_{m,\nu}^{\rm OS}(s)
 (\mathcal V_{m,\nu}^{\rm OS})^{-1}
 =e^{-sD_{m,g}^{\rm YM}}.
\end{equation}
Ainsi, le générateur des espérances et l'hamiltonien des translations ne sont pas identifiés par leur nom : ce sont deux réalisations déjà construites de la même désintégration, dont la forme limite et la représentation de courbure figurent dans \eqref{eq:amp-ym-forma-limite} et \eqref{eq:amp-ym-f2}.

\section{Image de jauge et composition vers Yang--Mills}

\subsection{État de jauge enrichi et générateurs de transfert}

L’objet déjà engendré que reçoit ce chapitre est la structure de jauge complète
$\mathfrak G_{\rm gau}$ ; il ne constitue ni une donnée externe ni une hypothèse terminale.
À chaque niveau $m$, il conserve le réseau
$\Lambda_m=a_m\mathbb Z^4$, $a_{m+1}=a_m/9$, la mesure $\mu_m$, les
opérateurs d’entrelacement $J_m$, quatre réflexions $\Theta_{\nu,m}$, l’holonomie, la couleur,
la représentation, la retenue de fusion, la route, la mémoire, le terminal et le lecteur de courbure.
Ses deux dynamiques sont typées par des symboles distincts avant toute limite.

Le semi-groupe d’une seule fibre de mémoire est
\begin{equation}\label{eq:pdf2-ym-fibra-semigrupo}
 K^{\rm fib}_{m,t}
 =P_{\Omega_m}+e^{-3\kappa_{\rm av}t}(I-P_{\Omega_m})
 =e^{-tD_m^{\rm fib}},
 \qquad
 D_m^{\rm fib}:=3\kappa_{\rm av}(I-P_{\Omega_m}).
\end{equation}
Son spectre est contenu dans $\{0,3\kappa_{\rm av}\}$. Le générateur de la coordinatisation produit complète est, en revanche,
\begin{equation}\label{eq:pdf2-ym-generador-producto}
 D_m^{\rm prod}:=3\kappa_{\rm av}
 \sum_{c\in\mathcal I_m}(I-E_{m,c}),
 \qquad
 K^{\rm prod}_{m,t}:=e^{-tD_m^{\rm prod}},
\end{equation}
où les espérances orthogonales $E_{m,c}$ commutent. S’il y a deux ou plusieurs
coordonnées non constantes, $D_m^{\rm prod}$ possède des niveaux
$6\kappa_{\rm av},9\kappa_{\rm av},\ldots$ ; c’est pourquoi
\eqref{eq:pdf2-ym-fibra-semigrupo} ne peut être déclaré égal à
\eqref{eq:pdf2-ym-generador-producto}. Le premier est la fibre simple ; le second,
le produit de recouvrement.

\subsection{Théorème fini de Hoeffding, vide et écart spectral atteint}

Dans la coordinatisation produit construite dans \eqref{eq:pdf2-ym-owner-product-factorization}, les $E_{m,c}$ sont des espérances conditionnelles orthogonales et commutantes, et la preuve de \eqref{eq:pdf2-ym-owner-vacuum} établit
\begin{equation}\label{eq:pdf2-ym-interseccion-vacio}
 \bigcap_{c\in\mathcal I_m}\operatorname{Ran}E_{m,c}
 =\mathbb C\Omega_m.
\end{equation}
La décomposition de Hoeffding est la somme orthogonale
\[
 \mathcal H_m=\bigoplus_{S\subseteq\mathcal I_m}\mathcal H_{m,S},
 \qquad u=\sum_Su_S,
\]
où $E_{m,c}u_S=0$ si $c\in S$ et $E_{m,c}u_S=u_S$ si $c\notin S$. Alors
\begin{equation}\label{eq:pdf2-ym-hoeffding}
 D_m^{\rm prod}u_S=3\kappa_{\rm av}|S|u_S,
 \qquad
 \langle u,D_m^{\rm prod}u\rangle
 =3\kappa_{\rm av}\sum_S|S|\|u_S\|^2.
\end{equation}
La composante $S=\varnothing$ est exactement le vide de
\eqref{eq:pdf2-ym-interseccion-vacio}. Par conséquent
\begin{equation}\label{eq:pdf2-ym-cota-finita}
 D_m^{\rm prod}\succeq
 3\kappa_{\rm av}(I-P_{\Omega_m}),
 \qquad
 \|K^{\rm prod}_{m,t}(I-P_{\Omega_m})\|
 \le e^{-3\kappa_{\rm av}t}.
\end{equation}

Dans une coordonnée d’incidence
$q=(q_1,\ldots,q_6)\in\mathbb F_3^6$, le témoin
\begin{equation}\label{eq:pdf2-ym-wc}
 W_c(q):=\mathbf1_{\{q_1+\cdots+q_6=0\}}-\frac13
\end{equation}
satisfait
\begin{equation}\label{eq:pdf2-ym-wc-momentos}
 \mathbb EW_c=0,\qquad \|W_c\|^2=\frac29,\qquad
 D_m^{\rm prod}W_c=3\kappa_{\rm av}W_c.
\end{equation}
Ainsi, à chaque niveau fini, le seuil de \eqref{eq:pdf2-ym-cota-finita} est atteint. Ce calcul fixe le spectre du générateur produit fini ; la mesure quadridimensionnelle continue de Yang--Mills est construite matériellement dans \eqref{eq:pdf2-ym-gibbs-measure}--\eqref{eq:pdf2-ym-g-dynamics}.

\subsection{Compression de jauge physique}

Soit $G$ compact, connexe et simple. Le groupe qui agit au niveau $m$ est
\[
 \mathcal G_m^{\rm full}
 =G^{V(\Lambda_m)}\times
  \prod_{c\in\mathcal C_m}\mathbb F_3^4,
\]
où le second facteur agit sur les six incidences au moyen d’une matrice
qui doit rester visible. Si les lignes sont indexées par
$(01,02,03,12,13,23)$ et les colonnes par $0,1,2,3$, on fixe
\begin{equation}\label{eq:pdf2-ym-incidence-matrix}
 B^+_{(ij),k}:=\delta_{ik}+\delta_{jk}\in\mathbb F_3,
 \qquad
 (B^+x)_{ij}=x_i+x_j.
\end{equation}
C’est l’incidence non orientée du collier ; elle ne se confond pas avec le
cobord signé d’une arête. L’action finie est
\begin{equation}\label{eq:pdf2-ym-incidence-action}
 q_c\longmapsto q_c+B^+x_c,
 \qquad x_c\in\mathbb F_3^4.
\end{equation}
La projection de Haar
\begin{equation}\label{eq:pdf2-ym-proyeccion-fisica}
 \mathbb P_m^{\rm phys}F
 :=\int_{\mathcal G_m^{\rm full}}F(g^{-1}\!\cdot\,)\,dg,
 \qquad
 \mathcal H_m^{\rm phys}:=\operatorname{Ran}\mathbb P_m^{\rm phys}
\end{equation}
est distincte de la marginale qui oublie l'incidence. L'invariance des espérances $E_{m,c}$ donne
\begin{equation}\label{eq:pdf2-ym-reduccion}
 [\mathbb P_m^{\rm phys},D_m^{\rm prod}]=0,
 \qquad
 [\mathbb P_m^{\rm phys},K_{m,t}^{\rm prod}]=0.
\end{equation}
De plus, chaque variable de $x_c\in\mathbb F_3^4$ apparaît trois fois dans la somme des
six incidences, de sorte que
\begin{equation}\label{eq:pdf2-ym-incidence-conservation}
 \sum_{i<j}(B^+x_c)_{ij}
 =3\sum_{i=0}^3x_{c,i}=0
 \quad\text{en }\mathbb F_3.
\end{equation}
Par conséquent, $W_c$ est invariant de
jauge et appartient à $\mathcal H_m^{\rm phys}$. L’écart fini n’est pas un
artefact d’une dilatation non physique.

\subsection{Positivité par réflexion à chaque niveau}

Pour un noyau réversible commun, définissons la mesure à huit faces
\begin{equation}\label{eq:pdf2-ym-ocho-caras}
 \Omega_m^{(8)}(dy_1\cdots dy_8)
 :=\int\mu_m(dz)\prod_{r=1}^{8}K^{\rm prod}_{m,a_m/2}(z,dy_r).
\end{equation}
Si $\Theta_{\nu,m}$ échange les deux moitiés par rapport à la direction $\nu$, la propriété de Markov et la réversibilité factorisent, pour tout observable cylindrique $F$ à support dans le demi-espace positif,
\begin{equation}\label{eq:pdf2-ym-os-finita}
 \int\overline{F(\Theta_{\nu,m}y)}F(y)\,d\Omega_m^{(8)}
 =\|\mathcal B_{m,\nu}F\|_2^2\ge0,
 \qquad \nu=1,\ldots,4.
\end{equation}
La marginale de faces opposées récupère le même
$K^{\rm prod}_{m,a_m}$. Ainsi, les quatre formes OS finies sont
la conséquence d’une seule mesure et d’un seul semi-groupe, non de quatre processus
indépendants.

\subsection{État de jauge complet et coordonnées physiques d’incidence}

Le propriétaire en vigueur ne prend pas pour espace physique la marginale qui ne conserve que
les liens. Son domaine est $\mathfrak G_{\rm gau}$, avec lien, repère,
représentation, incidence, couleur, route, orientation,
retenue, mémoire, survie et terminal. La projection vers les liens est une application
d’oubli ultérieure. Juger le témoin physique après avoir appliqué cet oubli remplace
l’objet et ne transfère pas la conclusion.

\subsubsection{Mesure projective et représentation des liens}

Pour un écran fini $\Gamma_m=(V_m,E_m)$, on part de l'espace relevé
\begin{equation}\label{eq:pdf2-ym-owner-lifted-space}
 \widetilde X_m:=G^{V_m}\times G^{E_m},\qquad
 d\widetilde\mu_m(h,z)
 =\prod_{v\in V_m}dh_v\prod_{e\in E_m}k_{\tau_{m,e}}(z_e)\,dz_e,
\end{equation}
et chaque lien est exprimé au moyen de
\begin{equation}\label{eq:pdf2-ym-owner-link-publication}
 U_e=h_{s(e)}z_eh_{t(e)}^{-1}.
\end{equation}
Une arête grossière $e$ se raffine en la chaîne ordonnée $e_1\cdots e_9$. Posons \(\boldsymbol\tau=(\tau_1,\ldots,\tau_9)\), avec \(\tau_j=\tau_{m+1,e_j}>0\) et \(\tau=\tau_{m,e}=\sum_{j=1}^9\tau_j\). Soit \(d\lambda_z^{(9)}\) la désintégration de Haar sur la fibre du produit, caractérisée par
\[
 \int_{G^9}F(z_1,\ldots,z_9)\,dz_1\cdots dz_9
 =\int_G\int_{z_1\cdots z_9=z}
   F\,d\lambda_z^{(9)}\,dz.
\]
Le pont correct pour des temps non nécessairement égaux est
\begin{equation}\label{eq:pdf2-ym-owner-heat-bridge}
 dB_{z,\boldsymbol\tau}^{(9)}(z_1,\ldots,z_9)
 =\frac{\prod_{j=1}^9k_{\tau_j}(z_j)}{k_\tau(z)}
   \,d\lambda_z^{(9)},
 \qquad z_1\cdots z_9=z,
\end{equation}
qui est une probabilité car
\(k_{\tau_1}*\cdots*k_{\tau_9}=k_\tau\). Dans le raffinement nonadique uniforme, on
récupère \(\tau_j=\tau/9\), mais la formule ne présuppose pas cette égalité.

Pour rendre l'innovation indépendante sans effacer le produit grossier, on fixe, pour chaque $(e,z)$, un transport borélien
\begin{equation}\label{eq:pdf2-ym-owner-bridge-triangularization}
 \chi_{m,e,z}:\bigl((0,1)^{\mathbb N},du^{\otimes\mathbb N}\bigr)
 \longrightarrow
 \bigl(\{(z_j):z_1\cdots z_9=z\},B_{z,\boldsymbol\tau}^{(9)}\bigr)
\end{equation}
qui est un isomorphisme modulo les ensembles nuls. Il existe parce que les deux sont des réalisations standard non atomiques ; un choix est fixé une fois pour toutes et transporté lors du réétiquetage des arêtes. La variable grossière $z$ et l'innovation $u^{\rm det}_{m+1,e}$ sont alors des facteurs indépendants dans la coordinatisation triangulaire. Les nouveaux repères sont intégrés selon Haar. Route, retenue, orientation et mémoire sont transportées par les morphismes $\operatorname{Ext}_{\rm gau}$ de la donnée de jauge elle-même.

La projection grossière multiplie les neuf accroissements et oublie exclusivement les
innovations nouvelles. Avec cette définition, on obtient
\begin{equation}\label{eq:pdf2-ym-owner-projectivity}
 (\widehat\pi_{m+1,m})_*\widehat\mu_{m+1}^{\rm ov}
 =\widehat\mu_m^{\rm ov},\qquad
 \widehat J_mF=F\circ\widehat\pi_{m+1,m}.
\end{equation}
La même arête apparaît une seule fois dans la mesure ; les six faces incidentes sont six lecteurs orientés de cette arête. Les chemins qui se recouvrent conservent l'accroissement partagé et ne sont pas remplacés par des copies indépendantes.

\subsubsection{La fibre d’incidence est une coordonnée physique de l’état complet}

La queue d’incidence incluse dans $\mathfrak G_{\rm gau}$ se décompose sans perdre
d’information :
\begin{equation}\label{eq:pdf2-ym-owner-seed-split}
 \Sigma:(\mathbb F_3^2)^{\mathbb N}\longrightarrow
 (\mathbb F_3^2)^{\mathbb N}\times\mathbb F_3^6,
 \qquad \Sigma_*\sigma=\sigma\otimes u_6.
\end{equation}
Le premier facteur continue d'alimenter les coordinatisations d'innovation et le second conserve les six incidences $q_c=(q_{01},q_{02},q_{03},q_{12},q_{13},q_{23})$ du collier $c$. Au lieu de laisser implicites les coordonnées régies par le générateur, on fixe une coordinatisation triangulaire récursive
\begin{equation}\label{eq:pdf2-ym-owner-product-factorization}
 \widehat X_m
 =\widehat X_{m_0}^{\rm seed}
  \times\prod_{\ell=m_0+1}^{m}
       \mathcal U_{\ell/\ell-1}^{\rm det},
 \qquad
 \widehat\mu_m^{\rm ov}
 =\widehat\mu_{m_0}^{\rm seed}
  \otimes\bigotimes_{\ell=m_0+1}^{m}
       \upsilon_{\ell/\ell-1}^{\rm det}.
\end{equation}
Le facteur germe énumère une seule fois les repères et les liens grossiers, les incidences $q_c$ et leurs coordinatisations de marquage. Chaque $\mathcal U_{\ell/\ell-1}^{\rm det}$ énumère une seule fois les nouveaux repères, les ponts \eqref{eq:pdf2-ym-owner-bridge-triangularization}, les nouvelles incidences et les sous-queues de marquage de ce raffinement. La factorisation est interne à $\mathfrak G_{\rm gau}$ et n'ajoute aucune autre donnée de départ.

On enregistre l’index exhaustif, disjoint et dénombrable
\begin{equation}\label{eq:pdf2-ym-owner-exhaustive-index}
 \mathcal I_m
 :=\mathcal I_{m_0}^{\rm frame}
 \sqcup\mathcal I_{m_0}^{\rm link}
 \sqcup\mathcal I_{m_0}^{q}
 \sqcup\mathcal I_{m_0}^{\rm mark}
 \sqcup\bigsqcup_{\ell=m_0+1}^{m}
 \bigl(
   \mathcal I_{\ell/\ell-1}^{\rm frame}
   \sqcup\mathcal I_{\ell/\ell-1}^{\rm bridge}
   \sqcup\mathcal I_{\ell/\ell-1}^{q}
   \sqcup\mathcal I_{\ell/\ell-1}^{\rm mark}
 \bigr).
\end{equation}
Chaque facteur indépendant de \eqref{eq:pdf2-ym-owner-product-factorization} apparaît
exactement une fois ; une sous-queue infinie est énumérée par ses cylindres scalaires. Aucun
facteur généalogique ne reste hors de $\mathcal I_m$.
En particulier, les abréviations de la récurrence signifient littéralement
\begin{equation}\label{eq:pdf2-ym-owner-index-abbreviations}
 \begin{aligned}
 \mathcal I_{m_0}^{\rm seed}
 &:={\mathcal I}_{m_0}^{\rm frame}\sqcup
    {\mathcal I}_{m_0}^{\rm link}\sqcup
    {\mathcal I}_{m_0}^{q}\sqcup
    {\mathcal I}_{m_0}^{\rm mark},\\
 \mathcal I_{\ell/\ell-1}^{\rm det}
 &:={\mathcal I}_{\ell/\ell-1}^{\rm frame}\sqcup
    {\mathcal I}_{\ell/\ell-1}^{\rm bridge}\sqcup
    {\mathcal I}_{\ell/\ell-1}^{q}\sqcup
    {\mathcal I}_{\ell/\ell-1}^{\rm mark}.
 \end{aligned}
\end{equation}

La marginale des liens oublie $q_c$ ; la sous-algèbre physique, en revanche, s’obtient
en moyennant l’action de jauge complète
\begin{equation}\label{eq:pdf2-ym-full-gauge-group}
 \mathcal G_m^{\rm full}
 :=G^{V(\Lambda_m)}\times
 \prod_{c\in\mathcal C_m}\mathbb F_3^4,
 \qquad q_c\longmapsto q_c+B^+x_c.
\end{equation}
La projection physique est donc
\begin{equation}\label{eq:pdf2-ym-full-physical-projection}
 \mathbb P_m^{\rm phys}F
 :=\int_{\mathcal G_m^{\rm full}}F(g^{-1}\!\cdot x)\,dg,
 \qquad
 \widehat{\mathscr H}_m^{\rm phys}
 :=\operatorname{Ran}\mathbb P_m^{\rm phys}.
\end{equation}
Elle ne coïncide pas avec l’application d’oubli
$\widehat X_m\to X_m^{\rm link}$. Cette distinction est le garde-fou de typage
qui empêche d’expulser $W_c$ du secteur physique.

\subsection{Générateur de jauge et circuit local de transport}

La coordinatisation triangulaire complète fournit un opérateur unitaire
\begin{equation}\label{eq:pdf2-ym-owner-triangular-unitary}
 \widehat\Gamma_m:
 L^2(\widehat X_{m+1},\widehat\mu_{m+1}^{\rm ov})
 \longrightarrow
 L^2(\widehat X_m,\widehat\mu_m^{\rm ov})
 \widehat\otimes\mathcal K_{m+1}^{\rm det},
 \qquad
 \widehat\Gamma_m\widehat J_mF=F\otimes\mathbf1_{\rm det}.
\end{equation}
Pour $\iota\in\mathcal I_m$, soit $E_{m,\iota}$ l'espérance qui intègre exactement le facteur $\iota$ de \eqref{eq:pdf2-ym-owner-product-factorization} et laisse tous les autres fixes. Ce sont des projections de Markov orthogonales, elles conservent l'unité et commutent. Sur le cœur dense $\mathscr C_m$ des cylindres qui dépendent d'un nombre fini de coordonnées, on définit la forme
\begin{equation}\label{eq:pdf2-ym-owner-full-form}
 \widehat{\mathcal E}_m[F]
 :=3\kappa_{\rm av}\sum_{\iota\in\mathcal I_m}
   \|(I-E_{m,\iota})F\|_2^2.
\end{equation}
La somme est finie pour chaque cylindre. Sa fermeture est une forme de Dirichlet et
détermine le générateur auto-adjoint markovien
$\widehat H_m^{\rm ov}$. À un niveau germe, cela s’écrit, au sens des
formes,
\begin{equation}\label{eq:pdf2-ym-owner-base-generator}
 \widehat H_{m_0}^{\rm ov}
 =3\kappa_{\rm av}
  \sum_{\iota\in\mathcal I_{m_0}^{\rm seed}}(I-E_{m_0,\iota}),
\end{equation}
et chaque raffinement ajoute exactement les espérances de ses nouvelles coordonnées :
\begin{equation}\label{eq:pdf2-ym-owner-refined-generator}
 \begin{aligned}
 \widehat H_{m+1}^{\rm ov}
 &=\widehat\Gamma_m^*
   (\widehat H_m^{\rm ov}\otimes I
    +I\otimes\widehat H_{m+1}^{\rm det})\widehat\Gamma_m,\\
 \widehat H_{m+1}^{\rm det}
 &=3\kappa_{\rm av}
   \sum_{\iota\in\mathcal I_{m+1/m}^{\rm det}}
   (I-E_{m+1,\iota}^{\rm det}).
 \end{aligned}
\end{equation}
La partition des couleurs n’est pas laissée nominale. Si une cellule est indexée par
$n=(n_1,n_2,n_3,n_4)\in\mathbb Z^4$, on fixe
\begin{equation}\label{eq:pdf2-ym-owner-729-coloring}
 \chi_m(n):=
 \bigl(n_1\bmod9,\ n_2\bmod9,\ n_3+3n_4\bmod9\bigr)
 \in(\mathbb Z/9\mathbb Z)^3,
 \qquad |(\mathbb Z/9\mathbb Z)^3|=9\cdot81.
\end{equation}
Si deux colliers fermés de rayon cellulaire un se coupent, leur différence
$\delta$ vérifie $|\delta_j|\le2$. L’égalité des couleurs impose
$\delta_1=\delta_2=0$ et
$\delta_3+3\delta_4\equiv0\pmod9$. Comme le dernier entier appartient à
$[-8,8]$, il vaut zéro ; les bornes $|\delta_3|,|\delta_4|\le2$ donnent alors
$\delta_3=\delta_4=0$. Par conséquent, des cellules distinctes de même couleur ont
des colliers disjoints. On définit matériellement
\begin{equation}\label{eq:pdf2-ym-owner-729-blocks}
 \mathcal B_{m,r}:=\{\iota\in\mathcal I_m:
   \operatorname{cell}(\iota)=n, \chi_m(n)=r\},
 \qquad r\in(\mathbb Z/9\mathbb Z)^3.
\end{equation}
Le transport d'un écran par $g\in E(4)$ transporte également cette partition ; il n'exige pas qu'une rotation préserve les coordonnées de l'étiquette. Ainsi, $\{\mathcal B_{m,r}\}_{r=1}^{9\cdot81}$ regroupe des supports spatiaux disjoints ; l'indice interne de chaque bloc énumère ses facteurs de \eqref{eq:pdf2-ym-owner-exhaustive-index}. Pour un cylindre $F$, le produit ne contient que les blocs dont dépend $F$ ; la clôture définit le produit fort pour toute la forme. Le générateur d'un bloc conserve les multiplicités de Hoeffding ; il n'est pas remplacé par un seul projecteur qui les effacerait. Par commutativité dans chaque coordinatisation,
\begin{equation}\label{eq:pdf2-ym-owner-local-circuit}
\begin{aligned}
 H_{m,B}&:=3\kappa_{\rm av}\sum_{\iota\in B}(I-E_{m,\iota}),\\
 K_{m,B}&:=e^{-a_mH_{m,B}}
 =\prod_{\iota\in B}e^{-3\kappa_{\rm av}a_m(I-E_{m,\iota})},\\
 K_m^{\rm loc}&:=\mathop{\prod_{r,B}}^{\rm strong}K_{m,B}
 =e^{-a_m\widehat H_m^{\rm ov}}.
\end{aligned}
\end{equation}
Le diamètre physique de chaque bloc reste uniformément borné. C'est de là que provient la localité : le générateur n'introduit pas de mise à jour instantanée sur tout le réseau.
La partition est compatible avec Bianchi, et pas seulement avec la disjonction géométrique. Pour une cellule cubique $c$, on fixe le produit orienté, avec tous les facteurs transportés au même sommet,
\begin{equation}\label{eq:pdf2-ym-owner-bianchi-block-product}
 \mathscr B_{m,c}(x):=
 \prod_{p\subset\partial c}^{\longrightarrow}
 T_p\operatorname{Hol}_p(x)T_p^{-1}.
\end{equation}
Chaque arête interne de la frontière apparaît deux fois avec des orientations contraires ;
par conséquent, $\mathscr B_{m,c}=e$ point par point. Un bloc
$\mathcal B_{m,r}$ contient le collier complet de toute cellule qu’il met à jour : s’il touche
une face de $c$, il met simultanément à jour la paire inverse qui partage cette arête.
L’annulation précédente reste littérale après chaque facteur du circuit :
\begin{equation}\label{eq:pdf2-ym-owner-bianchi-block-invariance}
 K_{m,B}\mathbf1_{\{\mathscr B_{m,c}=e\}}
 =\mathbf1_{\{\mathscr B_{m,c}=e\}},
 \qquad
 K_m^{\rm loc}\mathbf1_{\{\mathscr B_{m,c}=e\}}
 =\mathbf1_{\{\mathscr B_{m,c}=e\}}.
\end{equation}
Ainsi, le calendrier de $9\cdot81$ couches préserve Bianchi à chaque étape et dans le produit fort, et pas seulement à la fin d'un balayage.

L'action de jauge permute ou translate les coordonnées intégrées par chaque $E_{m,\iota}$ et préserve leurs mesures. Par conséquent
\begin{equation}\label{eq:pdf2-ym-owner-physical-reduction}
 U_m(g)\widehat H_m^{\rm ov}=\widehat H_m^{\rm ov}U_m(g)
 \quad(g\in\mathcal G_m^{\rm full}),
 \qquad
 [\mathbb P_m^{\rm phys},\widehat H_m^{\rm ov}]=0,
 \qquad
 \mathbb P_{m+1}^{\rm phys}\widehat J_m
 =\widehat J_m\mathbb P_m^{\rm phys}.
\end{equation}
Le semi-groupe complet et son noyau de Markov sont notés
\begin{equation}\label{eq:pdf2-ym-owner-full-kernel}
 \widehat K_{m,t}^{\rm ov}:=e^{-t\widehat H_m^{\rm ov}},
 \qquad
\widehat K_{m,t}^{\rm ov}(x,dy)
 \quad(x,y\in\widehat X_m).
\end{equation}
La première identité de \eqref{eq:pdf2-ym-owner-physical-reduction} implique
\begin{equation}\label{eq:pdf2-ym-owner-full-kernel-equivariance}
 \widehat K_{m,t}^{\rm ov}(gx,gA)
 =\widehat K_{m,t}^{\rm ov}(x,A)
 \quad(g\in\mathcal G_m^{\rm full}).
\end{equation}
Pour typer la réduction sans utiliser un noyau restreint sur le mauvais espace,
soit $q_m:\widehat X_m\to Y_m:=\widehat X_m/\mathcal G_m^{\rm full}$ l’application des
orbites et $\nu_m=(q_m)_*\widehat\mu_m^{\rm ov}$. Le tiré en arrière
\begin{equation}\label{eq:pdf2-ym-owner-quotient-unitary}
 \mathcal W_m:L^2(Y_m,\nu_m)\longrightarrow
 \widehat{\mathscr H}_m^{\rm phys},
 \qquad \mathcal W_m\phi=\phi\circ q_m
\end{equation}
est unitaire. L'équivariance du noyau complet permet de définir
\begin{equation}\label{eq:pdf2-ym-owner-physical-kernel}
 K_{m,t}^{\rm phys}(q_mx,A)
 :=\widehat K_{m,t}^{\rm ov}(x,q_m^{-1}A),
 \qquad A\subseteq Y_m,
\end{equation}
indépendamment du représentant $x$. L’objet physique est alors typé par
\begin{equation}\label{eq:pdf2-ym-owner-D}
 D_m^{\rm YM}
 :=\left.\widehat H_m^{\rm ov}\right|_{\widehat{\mathscr H}_m^{\rm phys}},
 \qquad
 e^{-tD_m^{\rm YM}}
 =\mathcal W_mK_{m,t}^{\rm phys}\mathcal W_m^{-1}.
\end{equation}
C'est le générateur produit complet ; ce n'est pas le générateur d'une seule fibre de \eqref{eq:pdf2-ym-fibra-semigrupo}.

\subsection{Action de Yang--Mills, mesure dépendant de $g$ et transport exact}

Jusqu'ici, on a écrit la coordinatisation cinématique de référence. On la notera $\widehat\mu_m^0:=\widehat\mu_m^{\rm ov}$ et $\widehat H_m^0:=\widehat H_m^{\rm ov}$, avec $\widehat K_{m,t}^0:=e^{-t\widehat H_m^0}=\widehat K_{m,t}^{\rm ov}$. Pour transformer le lecteur structurel $\mathcal P_m^{F^2}$ en une loi, et non en une comparaison ultérieure, on applique d'abord sa racine positive à la factorisation triangulaire. La polarisation de $\mathcal P_m^{F^2}$ dans un collier $b$ est le Gram positif
\[
 Q_{m,b}(u,v):=\frac14\bigl(
 \mathcal P_{m,b}^{F^2}(u+v)-\mathcal P_{m,b}^{F^2}(u-v)
 \bigr).
\]
On quotiente l’espace fini des lecteurs du collier par $\ker Q_{m,b}$, on le complète
avec $Q_{m,b}$ et on note $\mathcal E_{m,b}^{F}$ le Hilbert résultant. Si
$\zeta_{m,b}(x_b)$ est le vecteur des lecteurs centrés de la coordonnée triangulaire
$x_b$, on fixe matériellement
\begin{equation}\label{eq:pdf2-ym-action-root}
 \xi_{m,b}(x_b):=Q_{m,b}^{1/2}\zeta_{m,b}(x_b)
 \in\mathcal E_{m,b}^{F},
 \qquad
 p_{m,b}(x_b):=\|\xi_{m,b}(x_b)\|^2.
\end{equation}
L'indice $b$ parcourt les colliers de Gram à support minimal. Un collier est affecté à la première couleur de \eqref{eq:pdf2-ym-owner-729-coloring} qui contient sa cellule de base ; ainsi, chaque terme apparaît une fois et les colliers d'une même couleur sont disjoints. Plus précisément, la diagonalisation positive regroupe les facteurs triangulaires en une partition disjointe
\begin{equation}\label{eq:pdf2-ym-action-factor-partition}
 \mathcal I_m^{F}
 =\bigsqcup_{r\in(\mathbb Z/9\mathbb Z)^3}
   \bigsqcup_{b\in\mathcal B_{m,r}^{F}}\mathcal I_{m,b},
 \qquad x_b:=(x_\iota)_{\iota\in\mathcal I_{m,b}},
 \qquad
 \mathcal B_{m,r}^{F}:=\{b:Q_{m,b}\ne0,\ \chi_m(b)=r\}.
\end{equation}
Aucune coordonnée triangulaire n’appartient à deux $\mathcal I_{m,b}$ ; si un lecteur
géométrique touche deux colliers, son mode de Gram est assigné au premier et son complément
orthogonal au second. C’est la décomposition spectrale finie de la forme positive,
non une duplication de la variable physique.
La racine de la somme orthogonale, calculée couleur par couleur, donne l’identité, non une
approximation,
\begin{equation}\label{eq:pdf2-ym-action-reader-sum}
 \mathcal P_m^{F^2}(x)
 =\sum_{r\in(\mathbb Z/9\mathbb Z)^3}
   \sum_{b\in\mathcal B_{m,r}^{F}}p_{m,b}(x_b),
 \qquad
 S_{m,g}(x):=\frac1{4g^2}\mathcal P_m^{F^2}(x),
 \quad g\in(0,\infty).
\end{equation}
Ici, $x_b:\widehat X_m\to X_{m,b}$ est la projection des coordonnées du collier, de sorte que le domaine, l'action et le codomaine de chaque terme sont fixés. Le lecteur est invariant de jauge parce que $Q_{m,b}$ utilise le produit invariant de $\mathfrak g$ ; il est invariant par réflexion parce qu'une réflexion ne fait que permuter les six paires orientées ; et il est naturel par réétiquetage euclidien. La partition \eqref{eq:pdf2-ym-action-factor-partition} et la coordinatisation produit donnent explicitement
\begin{equation}\label{eq:pdf2-ym-action-factor-measures}
 \widehat\mu_m^0=\lambda_m^{\rm inert}\otimes
   \bigotimes_{r,b}\lambda_{m,b}^0,
 \qquad
 d\lambda_{m,b}^g
 :=(Z_{m,b}^g)^{-1}e^{-p_{m,b}/(4g^2)}d\lambda_{m,b}^0,
 \qquad
 \widehat\mu_{m,g}^{\rm YM}=\lambda_m^{\rm inert}\otimes
   \bigotimes_{r,b}\lambda_{m,b}^g.
\end{equation}
Le facteur inerte contient exactement les coordonnées sur lesquelles le lecteur est nul ; il n'est pas éliminé de la théorie.

La séparation grossier--détail de \eqref{eq:pdf2-ym-owner-product-factorization} est orthogonale pour la même racine : la polarisation directe donne $Q_{m+1}=Q_m\oplus Q_{m+1/m}^{\rm det}$ et $\zeta_{m+1}=\zeta_m\oplus\zeta_{m+1/m}^{\rm det}$. Si $x_{m+1}=\widehat\Gamma_m^{-1}(x_m,u)$, alors
\begin{equation}\label{eq:pdf2-ym-action-cocycle}
 \xi_{m+1}(x_{m+1})
 =\xi_m(x_m)\oplus\xi_{m+1/m}^{\rm det}(u),
 \qquad
 S_{m+1,g}\circ\widehat\Gamma_m^{-1}
 =S_{m,g}+S_{m+1/m,g}^{\rm det}.
\end{equation}
C'est la manière précise dont la retenue et la mémoire empêchent de compter deux fois une face ancestrale : la partie grossière n'est pas ajoutée à nouveau dans le détail. En particulier, $Z_{m+1,g}=Z_{m,g}Z_{m+1/m,g}^{\rm det}$ pour les normalisateurs, et l'on obtient la collection de probabilités dépendant véritablement du couplage
\begin{equation}\label{eq:pdf2-ym-gibbs-measure}
 d\widehat\mu_{m,g}^{\rm YM}(x)
 :=Z_{m,g}^{-1}e^{-S_{m,g}(x)}d\widehat\mu_m^0(x),
 \qquad
 (\widehat\pi_{m+1,m})_*\widehat\mu_{m+1,g}^{\rm YM}
 =\widehat\mu_{m,g}^{\rm YM}.
\end{equation}
Pour écrire correctement Schwinger--Dyson par rapport à Haar, sans supposer à tort que la mesure de référence est plate, soit $\lambda_m$ le produit de Haar, de Lebesgue et du comptage uniforme de la coordinatisation relevée, et soit $R_m:=d\widehat\mu_m^0/d\lambda_m>0$. L'action totale régularisée est
\begin{equation}\label{eq:pdf2-ym-total-regulated-action}
 \mathscr S_{m,g}:=S_{m,g}-\log R_m,
 \qquad
 d\widehat\mu_{m,g}^{\rm YM}
 =\mathscr Z_{m,g}^{-1}e^{-\mathscr S_{m,g}}d\lambda_m.
\end{equation}
Le terme $-\log R_m$ contient les noyaux de chaleur et les coordinatisations de pont de la structure de jauge initiale ; il est indépendant de $g$. Le terme de Yang--Mills qui modifie le couplage et dont la limite classique est calculée ci-dessous est exactement $S_{m,g}=\mathcal P_m^{F^2}/(4g^2)$. La cohérence de la seconde égalité et le caractère standard des espaces finis permettent d'appliquer Kolmogorov et définissent l'unique mesure cylindrique $\widehat\mu_{\infty,g}^{\rm YM}$ dont les marginales sont $\widehat\mu_{m,g}^{\rm YM}$. Aucune valeur propre n'a été utilisée ici : $g$, l'action et la mesure ont d'abord été fixés.

Pour définir la dynamique sur cette mesure sans perdre le produit, la jauge ni les identités de surface, le transport est construit et non postulé. Dans chaque facteur non atomique $X_{m,b}$, on fixe une coordinatisation borélienne $\rho_{m,b}:X_{m,b}/\mathcal G_{m,b}\to(0,1)$ et l'on désintègre d'abord sur le quotient de jauge, puis sur l'orbite de Haar. Soient $F_{m,b}^0$ et $F_{m,b}^g$ les fonctions de répartition continues de $\rho_{m,b}$ pour la mesure de référence et la mesure repondérée. Le transport triangulaire est
\begin{equation}\label{eq:pdf2-ym-gibbs-transport-factor}
 t_{m,b,g}:=(F_{m,b}^g)^{-1}\circ F_{m,b}^0,
 \qquad
 T_{m,b,g}(y,h):=(t_{m,b,g}(y),h),
 \qquad
 (T_{m,b,g})_*\lambda_{m,b}^0=\lambda_{m,b}^g.
\end{equation}
Dans un facteur purement atomique sur lequel $p_{m,b}=0$, on prend $T_{m,b,g}=I$. En présence d'atomes et d'une coordonnée continue, la coordinatisation est l'ordre lexicographique atome--intervalle et la même formule s'applique à l'espace standard non atomique total. Les transports des $9\cdot81$ couleurs sont composés dans l'ordre fixé par \eqref{eq:pdf2-ym-owner-729-coloring} ; au sein d'une couleur, ils commutent parce que les colliers sont disjoints. On obtient
\begin{equation}\label{eq:pdf2-ym-gibbs-transport-global}
 T_{m,g}:=\prod_{r\in(\mathbb Z/9\mathbb Z)^3}
           \prod_{b\in\mathcal B_{m,r}^{F}}T_{m,b,g},
 \quad
 (T_{m,g})_*\widehat\mu_m^0=\widehat\mu_{m,g}^{\rm YM},
 \quad
 \widehat\pi_{m+1,m}T_{m+1,g}=T_{m,g}\widehat\pi_{m+1,m}.
\end{equation}
La dernière égalité se vérifie facteur par facteur à l'aide de \eqref{eq:pdf2-ym-action-cocycle} : le transport du facteur grossier est celui déjà fixé et les nouveaux transports agissent seulement sur l'innovation. Les coordinatisations de quotient sont choisies sur un écran représentatif puis transportées ; $T_{m,g}$ commute donc avec la jauge et la réflexion et, pour $a\in E(4)$, satisfait $T_{a\alpha,g}=U_\alpha(a)T_{\alpha,g}U_\alpha(a)^{-1}$.

La composition
\begin{equation}\label{eq:pdf2-ym-full-probability-isomorphism}
 \mathcal T_{m,g}:L^\infty(\widehat X_m,\widehat\mu_{m,g}^{\rm YM})
 \longrightarrow L^\infty(\widehat X_m,\widehat\mu_m^0),
 \qquad \mathcal T_{m,g}F:=F\circ T_{m,g}
\end{equation}
est, par construction, un isomorphisme d'espaces de probabilité et de $*$-algèbres, et se prolonge unitairement à $L^2$. En particulier,
\begin{equation}\label{eq:pdf2-ym-full-algebra-identities}
 \mathcal T_{m,g}(FG)=\mathcal T_{m,g}F\,\mathcal T_{m,g}G,
 \quad
 \mathcal T_{m,g}(F^*)=(\mathcal T_{m,g}F)^*,
 \quad
 \int F\,d\widehat\mu_{m,g}^{\rm YM}
 =\int\mathcal T_{m,g}F\,d\widehat\mu_m^0.
\end{equation}
Le domaine de chaque $T_{m,b,g}$ est l’espace même des holonomies qui satisfait
la loi de subdivision et la contrainte de Bianchi. Son image demeure donc
dans cet espace. En appliquant la multiplicativité précédente générateur par générateur, on
obtient les identités ponctuelles
\begin{equation}\label{eq:pdf2-ym-transport-surface-bianchi}
 \begin{aligned}
 \mathcal T_{m,g}(\operatorname{Hol}_p)
 &=\prod_{p'\subset p}^{\longrightarrow}
   \mathcal T_{m,g}(T_{p'}\operatorname{Hol}_{p'}T_{p'}^{-1}),\\
 \mathcal T_{m,g}(F_{{\rm cl},m})
 &=F_{{\rm cl},m}\circ T_{m,g},\\
 \prod_{p\subset\partial c}^{\longrightarrow}
 \mathcal T_{m,g}(T_p\operatorname{Hol}_pT_p^{-1})&=e,
 \qquad
 \operatorname{Alt}_{\lambda\mu\nu}
 D_\lambda(F_{{\rm cl},m}\circ T_{m,g})_{\mu\nu}=0.
 \end{aligned}
\end{equation}
La troisième ligne est l'annulation arête par arête de la frontière orientée de la frontière de $c$ ; la quatrième est sa polarisation infinitésimale. Ainsi, le transport préserve les produits, l'état, le clover et Bianchi, et pas seulement le premier chaos.

Enfin, on transporte les espérances, le générateur, le noyau, la projection et l’inclusion :
\begin{equation}\label{eq:pdf2-ym-g-dynamics}
 \begin{aligned}
 E_{m,\iota}^{g}&:=\mathcal T_{m,g}^{-1}E_{m,\iota}\mathcal T_{m,g},&
 \widehat H_{m,g}^{\rm YM}&:=\mathcal T_{m,g}^{-1}\widehat H_m^0\mathcal T_{m,g},\\
 \widehat K_{m,t}^{g}&:=\mathcal T_{m,g}^{-1}\widehat K_{m,t}^{0}\mathcal T_{m,g},&
 \mathbb P_{m,g}^{\rm phys}&:=\mathcal T_{m,g}^{-1}
     \mathbb P_m^{\rm phys}\mathcal T_{m,g},\\
 \widehat J_{m,g}&:=\mathcal T_{m+1,g}^{-1}
     \widehat J_m\mathcal T_{m,g},&
 D_{m,g}^{\rm YM}&:=\left.\widehat H_{m,g}^{\rm YM}
     \right|_{\operatorname{Ran}\mathbb P_{m,g}^{\rm phys}}.
 \end{aligned}
\end{equation}
La première ligne n’est pas seulement une conjugaison abstraite. En utilisant
\eqref{eq:pdf2-ym-action-factor-measures}, son action sur un cylindre est
l’actualisation par bain thermique de l’action :
\begin{equation}\label{eq:pdf2-ym-gibbs-conditional-expectation}
 (E_{m,b}^{g}F)(x_{\ne b})
 =\frac{\displaystyle
   \int_{X_{m,b}}F(x_{\ne b},u)e^{-p_{m,b}(u)/(4g^2)}
       d\lambda_{m,b}^0(u)}{\displaystyle
   \int_{X_{m,b}}e^{-p_{m,b}(u)/(4g^2)}d\lambda_{m,b}^0(u)},
 \qquad
 \widehat H_{m,g}^{\rm YM}
 =3\kappa_{\rm av}\sum_b(I-E_{m,b}^{g}).
\end{equation}
Ainsi, la mesure et chaque transition locale contiennent explicitement $g$ et l'action $p_{m,b}/(4g^2)$. Elles sont markoviennes parce que $\mathcal T_{m,g}\mathbf1=\mathbf1$, locales parce que chaque $T_{m,b,g}$ possède le même collier que $p_{m,b}$, et dépendent de $g$ par \eqref{eq:pdf2-ym-gibbs-transport-factor}. L'équivalence est spectrale et ne sélectionne pas la mesure à partir du spectre. Pour établir également les équations dynamiques, soit $X_{m,b}^a$ le champ invariant à gauche dans le facteur de lien $b$ associé à $e_a\in\mathfrak g$ ; dans une coordinatisation continue d'innovation, on utilise la dérivée directionnelle ordinaire, et les coordonnées finies sont traitées par différence exacte. Pour tout cylindre lisse $F$, le domaine commun est $\mathscr C_m^1:=\bigcap_{b,a}\operatorname{Dom}X_{m,b}^a$, avec support compact dans les coordinatisations ouvertes d'innovation ; les facteurs de groupe n'ont pas de frontière. L'invariance de Haar et la règle du produit donnent
\begin{equation}\label{eq:pdf2-ym-schwinger-dyson-finite}
 \int X_{m,b}^aF\,d\widehat\mu_{m,g}^{\rm YM}
 =\int F\,X_{m,b}^a\mathscr S_{m,g}\,d\widehat\mu_{m,g}^{\rm YM},
 \qquad F\in\mathscr C_m^1.
\end{equation}
Dans une coordonnée finie, pour $\Delta_\varepsilon F(q):=F(q+\varepsilon)-F(q)$, le changement de variable $q\mapsto q+\varepsilon$ donne la version exacte, sans dérivée formelle,
\begin{equation}\label{eq:pdf2-ym-schwinger-dyson-discrete}
 \int\Delta_\varepsilon F(q)\,d\widehat\mu_{m,g}^{\rm YM}
 =\int F(q)
  \left(e^{\mathscr S_{m,g}(q)-\mathscr S_{m,g}(q-\varepsilon)}-1\right)
  d\widehat\mu_{m,g}^{\rm YM}(q).
\end{equation}
Si $\delta_\varepsilon$ est la somme orientée de ces champs sur les arêtes qui
sont incidentes à un sommet, l’invariance de jauge de la mesure de référence et de
$S_{m,g}$ donne $\delta_\varepsilon\mathscr S_{m,g}=0$ et, par conséquent, l’identité de Ward
\begin{equation}\label{eq:pdf2-ym-ward-finite}
 \int\delta_\varepsilon F\,d\widehat\mu_{m,g}^{\rm YM}=0.
\end{equation}
Pour le facteur de jauge fini $\mathbb F_3^4$, l'invariance $\mathscr S_{m,g}(q+\varepsilon)=\mathscr S_{m,g}(q)$ réduit \eqref{eq:pdf2-ym-schwinger-dyson-discrete} à la même identité de Ward avec $\delta_\varepsilon$ remplacée par $\Delta_\varepsilon$. Les deux expressions sont compatibles avec $\widehat J_{m,g}$ car \eqref{eq:pdf2-ym-action-cocycle} sépare les dérivées grossières de celles du détail. Tout cylindre de la limite réside à un certain niveau ; ainsi, \eqref{eq:pdf2-ym-schwinger-dyson-finite} et \eqref{eq:pdf2-ym-ward-finite} passent littéralement à l'état cofinal. Ni l'équation de Schwinger--Dyson ni celle de Ward ne contiennent l'écart spectral comme prémisse.

\begin{equation}\label{eq:pdf2-ym-g-gap-before-limit}
 D_{m,g}^{\rm YM}\succeq3\kappa_{\rm av}
 (I-P_{\Omega_{m,g}}),
 \qquad
 W_{c,g}:=\mathcal T_{m,g}^{-1}W_c,
 \qquad
 D_{m,g}^{\rm YM}W_{c,g}=3\kappa_{\rm av}W_{c,g}.
\end{equation}
Dans ce qui suit, on fixe un $g\in(0,\infty)$ et l'on supprime cet indice ; les mesures, espérances, inclusions et générateurs sans exposant sont les objets de \eqref{eq:pdf2-ym-g-dynamics}. Les holonomies géométriques demeurent des fonctions coordonnées canoniques sur l'espace des configurations et leur loi change donc bien avec $g$ ; $\mathcal T_{m,g}$ les envoie sur des fonctions composées de la coordinatisation de référence. Seul le témoin spectral est explicitement transporté dans \eqref{eq:pdf2-ym-g-gap-before-limit}. Cette convention n'efface pas $g$ : sa présence explicite se trouve dans \eqref{eq:pdf2-ym-gibbs-measure}--\eqref{eq:pdf2-ym-g-dynamics}.

\subsection{Vide simple et témoin qui atteint l'écart spectral}

L'intersection de toutes les espérances de l'état complet est la constante :
\begin{equation}\label{eq:pdf2-ym-owner-vacuum}
 \bigcap_{\iota\in\mathcal I_m}\operatorname{Ran}E_{m,\iota}
 =\mathbb C\Omega_m.
\end{equation}
Cette identité n'est pas introduite comme hypothèse. Dans la coordinatisation produit \eqref{eq:pdf2-ym-owner-product-factorization}, chaque $E_{m,\iota}$ intègre une coordonnée de l'indice exhaustif \eqref{eq:pdf2-ym-owner-exhaustive-index} et laisse les autres fixes. Si $u$ appartient à l'intersection, ses espérances conditionnelles relativement à tout cylindre fini sont constantes. La convergence des martingales de ces cylindres vers la $\sigma$-algèbre complète implique que $u$ est constante presque partout. L'inclusion inverse est immédiate puisque toute espérance conserve l'unité. Cela prouve \eqref{eq:pdf2-ym-owner-vacuum} et, après la projection de jauge, laisse un vide physique unidimensionnel. La décomposition de Hoeffding de \eqref{eq:pdf2-ym-owner-base-generator}--\eqref{eq:pdf2-ym-owner-refined-generator} donne
\begin{equation}\label{eq:pdf2-ym-owner-hoeffding}
 \widehat{\mathcal E}_m(u)
 =3\kappa_{\rm av}\sum_{S\Subset\mathcal I_m}|S|\,\|u_S\|^2,
 \qquad
 D_m^{\rm YM}\succeq
 3\kappa_{\rm av}(I-P_{\Omega_m}).
\end{equation}
La fonction
\begin{equation}\label{eq:pdf2-ym-owner-wc}
 W_c(q):=\mathbf1_{\{q_1+\cdots+q_6=0\}}-\frac13
\end{equation}
est invariant sous $G^{V(\Lambda_m)}$ car il ne dépend ni des liens ni des repères. Pour
l’action d’incidence, chaque $x_i$ apparaît trois fois, donc
\begin{equation}\label{eq:pdf2-ym-owner-wc-invariance}
 \sum_{i<j}(B^+x)_{ij}=3\sum_i x_i=0
 \quad\text{en }\mathbb F_3.
\end{equation}
Ainsi $\mathbb P_m^{\rm phys}W_c=W_c$. Sous $u_6$, la somme des six trits est
uniforme ; par calcul direct,
\begin{equation}\label{eq:pdf2-ym-owner-wc-moments}
 \mathbb EW_c=0,\qquad
 \|W_c\|^2=\frac29,\qquad
 \mathbb E(W_c^3)=\frac2{27}.
\end{equation}
Seul le projecteur $I-E_{m,q_c}$ agit non trivialement sur $W_c$. Par conséquent
\begin{equation}\label{eq:pdf2-ym-owner-wc-eigenvalue}
 D_m^{\rm YM}W_c=3\kappa_{\rm av}W_c.
\end{equation}
La borne inférieure de \eqref{eq:pdf2-ym-owner-hoeffding} n’est pas seulement une borne : le
complément physique contient un vecteur non nul qui atteint le seuil.

\subsection{Raffinement nonadique et persistance dans la limite}

Les inclusions physiques satisfont
\begin{equation}\label{eq:pdf2-ym-owner-intertwining}
 D_{m+1}^{\rm YM}\widehat J_m
 =\widehat J_mD_m^{\rm YM},\qquad
 \bigl(e^{-a_{m+1}D_{m+1}^{\rm YM}}\bigr)^9\widehat J_m
 =\widehat J_me^{-a_mD_m^{\rm YM}},\qquad a_{m+1}=a_m/9.
\end{equation}
Dans la limite inductive
\begin{equation}\label{eq:pdf2-ym-owner-inductive-limit}
 \widehat{\mathscr H}_\infty^{\rm phys}
 =\varinjlim(\widehat{\mathscr H}_m^{\rm phys},\widehat J_m)
\end{equation}
est défini sur l’union algébrique
\begin{equation}\label{eq:pdf2-ym-owner-limit-form}
 \mathcal E_\infty[\widehat J_{m,\infty}u]
 :=\langle u,D_m^{\rm YM}u\rangle.
\end{equation}
Son générateur est noté
$D_{\infty,g}^{\rm YM}$ ; conformément à la convention postérieure à
\eqref{eq:pdf2-ym-g-gap-before-limit}, on écrit aussi
$D_\infty^{\rm YM}:=D_{\infty,g}^{\rm YM}$.
Le premier entrelacement de \eqref{eq:pdf2-ym-owner-intertwining} rend la
définition indépendante du représentant. Les troncatures de Hoeffding sont
un cœur de forme et donnent une suite de récupération ; Fatou sur la somme non négative
donne la semi-continuité inférieure. La fermeture possède donc
\begin{equation}\label{eq:pdf2-ym-owner-limit-gap}
 \overline{\mathcal E}_\infty[u]
 \ge3\kappa_{\rm av}\|(I-P_{\Omega_\infty})u\|^2.
\end{equation}
L’inclusion conserve les coordonnées ancestrales, de sorte que
$\widehat J_{m,\infty}W_c$ reste non nul et vérifie
\begin{equation}\label{eq:pdf2-ym-owner-limit-witness}
 D_\infty^{\rm YM}\widehat J_{m,\infty}W_c
 =3\kappa_{\rm av}\widehat J_{m,\infty}W_c.
\end{equation}
En conséquence,
\begin{equation}\label{eq:pdf2-ym-owner-exact-gap}
 \inf\bigl(\operatorname{Spec}D_\infty^{\rm YM}\setminus\{0\}\bigr)
 =3\kappa_{\rm av}>0,
 \qquad \ker D_\infty^{\rm YM}=\mathbb C\Omega_\infty.
\end{equation}

\subsection{Quatre reconstructions OS de la même dynamique}

Le noyau complet de \eqref{eq:pdf2-ym-owner-full-kernel} définit une loi réversible sur $\widehat X_m$. Pour la reconstruction physique, on utilise sa descente bien typée \eqref{eq:pdf2-ym-owner-physical-kernel} : sur le quotient $Y_m$, on définit
\begin{equation}\label{eq:pdf2-ym-owner-eight-face-law}
 \Omega_{m,{\rm phys}}^{(8)}(d\bar y_1\cdots d\bar y_8)
 =\int\nu_m(d\bar z)
  \prod_{r=1}^8K_{m,a_m/2}^{\rm phys}(\bar z,d\bar y_r).
\end{equation}
Pour chaque direction $\nu=1,\ldots,4$, la réflexion échange les quatre faces
positives avec les négatives. En conditionnant par rapport à l’état central $z$ et en utilisant
la réversibilité,
\begin{equation}\label{eq:pdf2-ym-owner-four-os}
 \int\overline{F(\Theta_{\nu,m}y)}F(y)\,
 d\Omega_{m,{\rm phys}}^{(8)}(y)
 =\|\mathcal B_{m,\nu}F\|_2^2\ge0.
\end{equation}
La marginale de la paire opposée normale à $\nu$ est
\begin{equation}\label{eq:pdf2-ym-owner-os-marginal}
 \nu_m(d\bar x)K_{m,a_m}^{\rm phys}(\bar x,d\bar y).
\end{equation}
Pour identifier l'hamiltonien d'une direction, on prend le sous-espace OS des observables d'une face positive simple, $F(y)=F_\nu(\bar y_{\nu,+})$ ; les observables de plusieurs faces conservent la positivité précédente, mais ne sont pas confondus avec cet espace de transfert. Dans le quotient de face simple, l'application
\begin{equation}\label{eq:pdf2-ym-owner-os-unitary}
 \mathcal V_{m,\nu}[F]
 :=e^{-a_mD_m^{\rm YM}/2}\mathcal W_mF
\end{equation}
est isométrique pour la norme OS et a une image dense. Son extension identifie le
quotient OS à $\widehat{\mathscr H}_m^{\rm phys}$. Posons
$\overline D_m:=\mathcal W_m^{-1}D_m^{\rm YM}\mathcal W_m$ et
$\overline J_m:=\mathcal W_{m+1}^{-1}\widehat J_m\mathcal W_m$.
L’entrelacement et $a_m=9a_{m+1}$ donnent
\begin{equation}\label{eq:pdf2-ym-owner-os-range-inclusion}
 \overline J_me^{-a_m\overline D_m/2}
 =e^{-a_m\overline D_{m+1}/2}\overline J_m
 =\bigl(e^{-a_{m+1}\overline D_{m+1}/2}\bigr)^9\overline J_m.
\end{equation}
Le dernier membre appartient à l’image de
$e^{-a_{m+1}\overline D_{m+1}/2}$ ; c’est pourquoi l’inverse de $\mathcal V_{m+1,\nu}$
qui apparaît ci-après ne s’applique que sur son image et est bien défini. Les
connecteurs
\begin{equation}\label{eq:pdf2-ym-owner-os-connectors}
 \mathcal J_{m,\nu}^{\rm OS}
 :=\mathcal V_{m+1,\nu}^{-1}\widehat J_m\mathcal V_{m,\nu}
\end{equation}
sont des isométries et entrelacent les quatre hamiltoniens reconstruits :
\begin{equation}\label{eq:pdf2-ym-owner-common-os-hamiltonian}
 H_{{\rm OS},m+1}^{(\nu)}\mathcal J_{m,\nu}^{\rm OS}
 =\mathcal J_{m,\nu}^{\rm OS}H_{{\rm OS},m}^{(\nu)},
 \qquad
 \mathcal V_{\infty,\nu}H_{{\rm OS},\infty}^{(\nu)}
 \mathcal V_{\infty,\nu}^{-1}=D_\infty^{\rm YM}.
\end{equation}
Les quatre réflexions ne fabriquent pas quatre mesures : ce sont quatre représentations du même état, du même noyau et du même générateur.

\subsection{Localité, action euclidienne et représentation de courbure}

La catégorie dirigée des écrans comprend les raffinements, les superpositions et les transports
euclidiens. Si $a\in E(4)$ et $\alpha$ est un écran, le réétiquetage complet définit
\begin{equation}\label{eq:pdf2-ym-owner-e4-finite-transport}
 U_\alpha(a):\widehat{\mathscr H}_\alpha^{\rm phys}
 \longrightarrow\widehat{\mathscr H}_{a\alpha}^{\rm phys},
 \qquad
 U_\beta(a)\widehat J_{\alpha\beta}
 =\widehat J_{a\alpha,a\beta}U_\alpha(a).
\end{equation}
Le transport préserve la mesure, le générateur, les idéaux nuls et les supports. L'identité de naturalité permet de définir, sans choisir de représentant, une application unitaire sur la colimite algébrique :
\begin{equation}\label{eq:pdf2-ym-owner-e4-colimit-action}
 U(a)\widehat J_{\alpha,\infty}F
 :=\widehat J_{a\alpha,\infty}U_\alpha(a)F.
\end{equation}
Les réétiquetages vérifient $U_{a\alpha}(b)U_\alpha(a)=U_\alpha(ba)$ et préservent
la norme ; c’est pourquoi \eqref{eq:pdf2-ym-owner-e4-colimit-action} se prolonge en une
représentation unitaire de $E(4)$.

\subsubsection{Courbure locale du canal d’incidence}

Tout groupe compact, connexe et simple possède un tore maximal non trivial. On fixe $t_3\in G$ d'ordre exactement trois. Pour un collier $c$, soit
\begin{equation}\label{eq:pdf2-ym-owner-incidence-flux}
 \omega_c(q):=\sum_{i<j}q_{ij}\in\mathbb F_3.
\end{equation}
L’identité \eqref{eq:pdf2-ym-incidence-conservation} prouve que $\omega_c$ est
invariant sous l’action finie. Si $h_c$ est le repère au sommet de base du collier,
on définit l’holonomie géométrique et, pour le témoin spectral, son image transportée
dans la loi au couplage $g$ :
\begin{equation}\label{eq:pdf2-ym-owner-incidence-holonomy}
 \operatorname{Hol}_{m,c}^{\rm inc,geom}:=h_ct_3^{\,\omega_c(q)}h_c^{-1}\in G,
 \qquad
 \operatorname{Hol}_{m,c,g}^{\rm inc,spec}
 :=\mathcal T_{m,g}^{-1}\operatorname{Hol}_{m,c}^{\rm inc,geom},
 \qquad
 \operatorname{Hol}_{m,c}^{\rm inc}:=\operatorname{Hol}_{m,c,g}^{\rm inc,spec}.
\end{equation}
Sous $G^{V_m}$, elle se transforme par conjugaison et sous
$\prod_c\mathbb F_3^4$, elle demeure fixe. Comme $t_3$ est d’ordre trois, le calcul
fonctionnel de cette variable à trois valeurs donne l’identité littérale
\begin{equation}\label{eq:pdf2-ym-owner-wc-curvature-functional}
 W_{c,g}:=\mathcal T_{m,g}^{-1}W_c^0
 =\mathbf1_{\{e\}}\!\left(\operatorname{Hol}_{m,c}^{\rm inc}\right)-\frac13.
\end{equation}
On écrit $W_c:=W_{c,g}$ dans le reste du chapitre. Comme $T_{m,b,g}$ agit à l'intérieur du quotient du même collier, l'holonomie spectrale est une fonction de son algèbre locale d'holonomies géométriques ; elle n'ajoute pas de variable et ne change pas le support. Les holonomies géométriques, utilisées dans $S_{m,g}$ et dans le clover, maintiennent leur distribution dépendante de $g$.
Le projecteur du membre de droite est une fonction de classe locale ; par conséquent, $W_c$ n'est pas une coordonnée auxiliaire découplée, mais un observable de l'holonomie d'incidence de la même restriction de jauge.

Pour un ouvert borné $O\subset\mathbb R^4$, soit
$\mathfrak A_{\rm YM}(O)$ l’algèbre engendrée par les fonctions de classe bornées de
\eqref{eq:pdf2-ym-owner-incidence-holonomy}, les holonomies de plaquette et les lecteurs
clover à support dans $O$. On définit le secteur de courbure physique par
\begin{equation}\label{eq:pdf2-ym-owner-curvature-cyclic-sector}
 \mathscr H_{\rm YM}^{\rm curv}
 :=\overline{\bigcup_{O\Subset\mathbb R^4}
   \mathfrak A_{\rm YM}(O)\Omega_\infty}
 \subseteq\widehat{\mathscr H}_\infty^{\rm phys}.
\end{equation}
Par \eqref{eq:pdf2-ym-owner-wc-curvature-functional}, chaque
$\widehat J_{m,\infty}W_c$ appartient matériellement à ce secteur.
Les espérances de coordonnée envoient les cylindres de courbure sur des cylindres de
courbure ; le semi-groupe laisse donc invariant
$\mathscr H_{\rm YM}^{\rm curv}$. Étant auto-adjoint, le secteur est réducteur. La
borne de \eqref{eq:pdf2-ym-owner-limit-gap} s’y restreint et le vecteur précédent
atteint l’égalité :
\begin{equation}\label{eq:pdf2-ym-owner-curvature-sector-gap}
 D_{\rm curv,g}^{\rm YM}
 :=\left.D_{\infty,g}^{\rm YM}\right|_{\mathscr H_{\rm YM}^{\rm curv}},
 \qquad D_{\rm curv}^{\rm YM}:=D_{\rm curv,g}^{\rm YM},
 \qquad
 \inf\bigl(\operatorname{Spec}D_{\rm curv,g}^{\rm YM}\setminus\{0\}\bigr)
 =3\kappa_{\rm av}.
\end{equation}

\subsubsection{Système cofinal de représentation et produits croisés}

L'inclusion $\widehat J_{\alpha\beta}$ conserve une coordonnée ancestrale ; elle ne doit pas être confondue avec la moyenne de bloc exprimant un champ continu. Pour écrire cette seconde application sans modifier la généalogie, soit $\mathfrak P$ l'ensemble dirigé des écrans cellulaires bornés, fermé par superposition, raffinement nonadique et transport euclidien. Pour $\alpha\in\mathfrak P$, posons
\begin{equation}\label{eq:pdf2-ym-owner-cell-step-space}
 V_\alpha:=\operatorname{span}\left\{
 e_{\alpha,c}:=|c|^{-1/2}\mathbf1_c:
 c\in\mathcal C_\alpha^{\rm cell}\right\}
 \subset L^2(\mathbb R^4),
 \qquad P_\alpha:L^2(\mathbb R^4)\to V_\alpha.
\end{equation}
Les superpositions font que $P_\alpha h\to h$ pour tout $h\in L^2$ le long de
$\mathfrak P$. Du côté de la jauge, on normalise le témoin par
\begin{equation}\label{eq:pdf2-ym-owner-normalized-cell-witness}
 Z_{\alpha,c}:=\sqrt{\frac92}\,W_{\alpha,c},
 \qquad
 \langle Z_{\alpha,c},Z_{\alpha,d}\rangle=\delta_{cd},
 \qquad
 D_\alpha^{\rm YM}Z_{\alpha,c}=3\kappa_{\rm av}Z_{\alpha,c}.
\end{equation}
La seconde égalité utilise exactement les moments de \eqref{eq:pdf2-ym-owner-wc-moments} et la factorisation de coordonnées distinctes. Soit $\mathcal K_\alpha^{\rm cell}$ l'espace engendré par ces vecteurs. Si $\beta\succeq\alpha$, l'opérateur d'entrelacement de représentation, distinct de $\widehat J_{\alpha\beta}$, est défini générateur par générateur par
\begin{equation}\label{eq:pdf2-ym-owner-block-intertwiner}
 I_{\alpha\beta}^{\rm curv}Z_{\alpha,c}
 :=\sum_{\substack{d\in\mathcal C_\beta^{\rm cell}\\d\subset c}}
       \sqrt{\frac{|d|}{|c|}}\,Z_{\beta,d}.
\end{equation}
Comme les sous-cellules sont disjointes et ont pour somme $|c|$, on obtient, sans passage à la limite,
\begin{equation}\label{eq:pdf2-ym-owner-block-intertwiner-identities}
 (I_{\alpha\beta}^{\rm curv})^*I_{\alpha\beta}^{\rm curv}=I,
 \quad
 I_{\beta\gamma}^{\rm curv}I_{\alpha\beta}^{\rm curv}
 =I_{\alpha\gamma}^{\rm curv},
 \quad
 D_\beta^{\rm YM}I_{\alpha\beta}^{\rm curv}
 =I_{\alpha\beta}^{\rm curv}D_\alpha^{\rm YM}.
\end{equation}
Le réétiquetage cellulaire conserve les volumes et les descendants ; pour $a\in E(4)$, on a en outre
\begin{equation}\label{eq:pdf2-ym-owner-block-e4-naturality}
 U_\beta(a)I_{\alpha\beta}^{\rm curv}
 =I_{a\alpha,a\beta}^{\rm curv}U_\alpha(a).
\end{equation}

Pour $h\in C_c^\infty(\mathbb R^4)$, définissons maintenant le lecteur de bloc
\begin{equation}\label{eq:pdf2-ym-owner-smeared-witness}
 \Phi_\alpha(h):=
 \sum_{c\in\mathcal C_\alpha^{\rm cell}}
 \langle h,e_{\alpha,c}\rangle_{L^2}Z_{\alpha,c}.
\end{equation}
Sur un maillage uniforme, cette formule est
$\sqrt{9/2}\,a_\alpha^2\sum_c(P_\alpha h)(x_c)W_{\alpha,c}$ ; la moyenne cellulaire,
non une évaluation ponctuelle, fixe le produit croisé. En effet, si $\gamma$ raffine
deux écrans $\alpha$ et $\beta$, les définitions précédentes donnent l’identité
calculée
\begin{equation}\label{eq:pdf2-ym-owner-smeared-cross-gram}
 \left\langle
  I_{\alpha\gamma}^{\rm curv}\Phi_\alpha(h),
  I_{\beta\gamma}^{\rm curv}\Phi_\beta(k)
 \right\rangle
 =\langle P_\alpha h,P_\beta k\rangle_{L^2}.
\end{equation}
Par polarisation, avec $k=h$,
\begin{equation}\label{eq:pdf2-ym-owner-smeared-cauchy-derived}
 \left\|
  I_{\alpha\gamma}^{\rm curv}\Phi_\alpha(h)
  -I_{\beta\gamma}^{\rm curv}\Phi_\beta(h)
 \right\|_2^2
 =\|P_\alpha h-P_\beta h\|_{L^2}^2\longrightarrow0.
\end{equation}
Ainsi, la propriété de Cauchy est une conséquence de la convergence forte des
projections, et non une identité supplémentaire postulée. La colimite de
\eqref{eq:pdf2-ym-owner-block-intertwiner} contient donc
\begin{equation}\label{eq:pdf2-ym-owner-smeared-gram}
 \Phi(h):=\lim_\alpha I_{\alpha,\infty}^{\rm curv}\Phi_\alpha(h),
 \qquad
 \langle\Phi(h),\Phi(k)\rangle=\langle h,k\rangle_{L^2},
 \qquad
 D_{\rm pub}^{\rm YM}\Phi(h)=3\kappa_{\rm av}\Phi(h).
\end{equation}

Cette représentation n'ajoute pas de degrés de liberté. À chaque étape, la décomposition $V_\beta=V_\alpha\oplus(V_\beta\ominus V_\alpha)$ possède une base de Haar nonadique. L'ordre exhaustif des nouvelles incidences de \eqref{eq:pdf2-ym-owner-exhaustive-index} apparie cette base, étiquette par étiquette, avec les coordonnées $Z_{\beta,d}$ du facteur d'innovation, en conservant le support de la fonction de Haar. Les opérateurs unitaires triangulaires $\widehat\Gamma_\alpha$ recollent ces appariements et produisent une isométrie basée
\begin{equation}\label{eq:pdf2-ym-owner-publication-into-physical}
 \mathcal U_{\rm curv}:
 \varinjlim(\mathcal K_\alpha^{\rm cell},I_{\alpha\beta}^{\rm curv})
 \longrightarrow\mathscr H_{\rm YM}^{\rm curv},
 \quad
 D_{\rm curv}^{\rm YM}\mathcal U_{\rm curv}
 =\mathcal U_{\rm curv}D_{\rm pub}^{\rm YM},
 \quad U(a)\mathcal U_{\rm curv}=\mathcal U_{\rm curv}U_{L^2}(a).
\end{equation}
Désormais, $\Phi(h)$ désigne également son image par $\mathcal U_{\rm curv}$. La compatibilité de \eqref{eq:pdf2-ym-gibbs-transport-global} définit $\mathcal T_{\infty,g}$ à la limite. La représentation physique au couplage fixé n'est pas une seconde isométrie choisie seulement sur le premier chaos, mais la restriction
\begin{equation}\label{eq:pdf2-ym-full-publication-restriction}
 \mathcal U_{{\rm curv},g}
 :=\mathcal T_{\infty,g}^{-1}\mathcal U_{{\rm curv},0},
 \qquad
 \widetilde{\mathcal U}_g
 :=\mathcal T_{\infty,g}^{-1}:
 L^\infty(\widehat X_\infty,\widehat\mu_\infty^0)
 \longrightarrow
 L^\infty(\widehat X_\infty,\widehat\mu_{\infty,g}^{\rm YM}).
\end{equation}
Sur la $*$-algèbre polynomiale engendrée par les holonomies, les plaquettes, clover et
$\Phi(h)$, on vérifie, générateur par générateur puis par linéarité,
\begin{equation}\label{eq:pdf2-ym-full-publication-products}
 \widetilde{\mathcal U}_g(AB)
 =\widetilde{\mathcal U}_g(A)\widetilde{\mathcal U}_g(B),
 \quad
 \widetilde{\mathcal U}_g(A^*)=\widetilde{\mathcal U}_g(A)^*,
 \quad
 \omega_g(\widetilde{\mathcal U}_gA)=\omega_0(A).
\end{equation}
Les égalités de surface, de clover et de Bianchi sont préservées par
\eqref{eq:pdf2-ym-transport-surface-bianchi} ; par densité monotone, l’extension
vaut dans toute l’algèbre locale bornée. Ainsi,
\eqref{eq:pdf2-ym-owner-smeared-cross-gram} est la restriction d’un isomorphisme
probabiliste et algébrique complet, non un substitut de celui-ci.

\subsubsection{Noyau dense, continuité forte et réseau local}

Soit $M_{\Phi(h)}$ l'opérateur de multiplication affilié à l'algèbre locale et
\begin{equation}\label{eq:pdf2-ym-owner-curvature-weyl}
 \mathsf W(h,t):=\exp\!\bigl(itM_{\Phi(h)}\bigr),
 \qquad h\in C_c^\infty(\mathbb R^4),\quad t\in\mathbb R.
\end{equation}
On complète $\mathfrak A_{\rm YM}(O)$ par ces unitaires pour
$\operatorname{supp}h\subset O$ et par les fonctions de classe bornées des
holonomies de plaquette. Notons $\mathfrak A_{\rm YM}(O)^{\rm sm}$ la
$*$-algèbre ainsi engendrée avec des fonctions tests lisses. Le noyau
\begin{equation}\label{eq:pdf2-ym-owner-smooth-core}
 \mathscr D_{\rm sm}:=\operatorname{span}\left\{
 A_1\cdots A_n\Omega_\infty:
 A_j\in\mathfrak A_{\rm YM}(O_j)^{\rm sm},\quad
 O_j\Subset\mathbb R^4\right\}
\end{equation}
est dense par la définition cyclique de \eqref{eq:pdf2-ym-owner-curvature-cyclic-sector}. Les équations \eqref{eq:pdf2-ym-owner-block-e4-naturality} et \eqref{eq:pdf2-ym-owner-publication-into-physical} donnent, pour $a\in E(4)$,
\begin{equation}\label{eq:pdf2-ym-owner-field-e4-covariance}
 U(a)\Phi(h)=\Phi(a\!\cdot h),
 \qquad
 \|(U(a)-I)\Phi(h)\|=\|a\!\cdot h-h\|_{L^2}.
\end{equation}
En outre, $|e^{ix}-e^{iy}|\le|x-y|$ et
\eqref{eq:pdf2-ym-owner-smeared-gram} impliquent
\begin{equation}\label{eq:pdf2-ym-owner-weyl-continuity}
 \|\bigl(\mathsf W(a\!\cdot h,t)-\mathsf W(h,t)\bigr)\Omega_\infty\|
 \le |t|\,\|a\!\cdot h-h\|_{L^2}.
\end{equation}
La continuité se vérifie également pour les générateurs d'holonomie du cœur ; \eqref{eq:pdf2-ym-owner-weyl-continuity} à elle seule ne les inclut pas. Pour un chemin polygonal orienté $\gamma$, une plaquette $p$ et $f\in C^1(G)^{\rm class}$, on pose
\begin{equation}\label{eq:pdf2-ym-owner-holonomy-core-generators}
 A_{\gamma,f}:=f(\operatorname{Hol}_\gamma),
 \qquad A_{p,f}:=f(\operatorname{Hol}_p).
\end{equation}
Dans une superposition de $\gamma$ et de $a\gamma$, avec $a\in E(4)$, la loi de pont
\eqref{eq:pdf2-ym-owner-heat-bridge} annule les accroissements communs. Pour chaque
accroissement restant $z$, on utilise
$|f(xz)-f(x)|\le\|\nabla f\|_\infty d_G(z,e)$ et la borne du second moment du
noyau de chaleur
\(
 \int_Gd_G(z,e)^2k_\tau(z)\,dz\le C_G\tau
\).
Pour la loi repondérée, $0<e^{-p/(4g^2)}\le1$ et la continuité de $p$ sur le
facteur compact donnent
\begin{equation}\label{eq:pdf2-ym-owner-gibbs-heat-moment}
 \frac{\int_Gd_G(z,e)^2e^{-p(z)/(4g^2)}k_\tau(z)\,dz}
      {\int_Ge^{-p(z)/(4g^2)}k_\tau(z)\,dz}
 \le C_{G,g}\tau,
\end{equation}
uniformément sur les écrans d’un compact et pour $0<\tau\le\tau_0$.
En additionnant les temps des incréments de la différence symétrique, on obtient
l’estimation calculée
\begin{equation}\label{eq:pdf2-ym-owner-holonomy-e4-estimate}
 \begin{aligned}
 \|U(a)A_{\gamma,f}\Omega_\infty-A_{\gamma,f}\Omega_\infty\|_2^2
 &\le C_{G,g}\|\nabla f\|_\infty^2
     \tau(\gamma\triangle a\gamma),\\
 \|U(a)A_{p,f}\Omega_\infty-A_{p,f}\Omega_\infty\|_2^2
 &\le C_{G,g}\|\nabla f\|_\infty^2
     \tau(\partial p\triangle\partial(ap)).
 \end{aligned}
\end{equation}
Ici, la différence symétrique n'est pas mesurée comme deux courbes disjointes : les connecteurs de surface la remplissent par le ruban orienté minimal $\Sigma(\gamma,a\gamma)$ et
\begin{equation}\label{eq:pdf2-ym-owner-bridge-time-ribbon}
 \tau(\gamma\triangle a\gamma)
 :=\sum_{p\subset\Sigma(\gamma,a\gamma)}\tau_p,
 \qquad
 \tau(\gamma\triangle a\gamma)
 \le C_\gamma d_{E(4)}(a,e),
\end{equation}
avec la même définition pour les frontières de plaquette. L’inégalité découle de la
somme des temps additifs de \eqref{eq:pdf2-ym-owner-heat-bridge} sur une bande
de largeur $d_{E(4)}(a,e)$ et de longueur bornée par celle de $\gamma$.
La mesure du temps de pont du membre de droite tend ainsi vers zéro lorsque $a\to e$ ;
c’est exactement la continuité des connecteurs de route de l’image de jauge.
Chaque lecteur clover borné $\chi(F_{\rm cl})$, $\chi\in C_b^1$, est fonction d’une
somme finie de plaquettes transportées ; la même borne, multipliée par
$\|\chi'\|_\infty$ et par le nombre de termes, prouve également sa continuité.

Pour des fonctions de classe seulement bornées, on rend visible la régularisation qui
entre dans le cœur. Si $d b$ est la mesure de Haar à gauche sur $E(4)$ et
$\eta\in C_c^\infty(E(4))$, on définit l’intégrale de Bochner
\begin{equation}\label{eq:pdf2-ym-owner-orbit-smearing}
 A[\eta]:=\int_{E(4)}\eta(b)U(b)AU(b)^*\,db,
 \qquad A\in\{A_{\gamma,f},A_{p,f},\chi(F_{{\rm cl}})\},
 \quad \chi\in C_b^1.
\end{equation}
Avec $(L_a\eta)(b):=\eta(a^{-1}b)$, le changement de variable $b\mapsto ab$ donne
\begin{equation}\label{eq:pdf2-ym-owner-orbit-smearing-continuity}
 U(a)A[\eta]U(a)^*=A[L_a\eta],
 \qquad
 \|(A[L_a\eta]-A[\eta])\Omega_\infty\|
 \le\|A\|\,\|L_a\eta-\eta\|_{L^1(E(4))}\longrightarrow0.
\end{equation}
Il est désormais entendu que les générateurs d'holonomie, de plaquette et de clover de $\mathfrak A_{\rm YM}(O)^{\rm sm}$ sont ceux de \eqref{eq:pdf2-ym-owner-holonomy-core-generators} avec $f\in C^1$, ou leurs régularisations \eqref{eq:pdf2-ym-owner-orbit-smearing}, et que le support balayé reste compact à l'intérieur de $O$. Les approximations de l'identité dans $E(4)$, avec \eqref{eq:pdf2-ym-owner-holonomy-e4-estimate}, récupèrent les générateurs non régularisés dans $L^2$ ; le secteur cyclique n'est donc pas réduit.

Un télescopage de produits qui utilise
\eqref{eq:pdf2-ym-owner-weyl-continuity} et
\eqref{eq:pdf2-ym-owner-orbit-smearing-continuity} prouve la continuité de tous
les générateurs de $\mathscr D_{\rm sm}$ ; par
densité et unitarité,
\begin{equation}\label{eq:pdf2-ym-owner-strong-e4}
 \lim_{a\to e}\|(U(a)-I)\Psi\|=0
 \qquad(\Psi\in\mathscr H_{\rm YM}^{\rm curv}).
\end{equation}
Les supports disjoints utilisent des facteurs disjoints et commutent. On obtient ainsi le réseau
\begin{equation}\label{eq:pdf2-ym-owner-local-net}
\begin{aligned}
 O_1\subseteq O_2&\Rightarrow
 \mathfrak A_{\rm YM}(O_1)\subseteq\mathfrak A_{\rm YM}(O_2),\\
 [\mathfrak A_{\rm YM}(O_1),\mathfrak A_{\rm YM}(O_2)]&=0
 \quad (O_1\cap O_2=\varnothing),\\
 U(a)\mathfrak A_{\rm YM}(O)U(a)^*&=\mathfrak A_{\rm YM}(aO).
\end{aligned}
\end{equation}
L’identité de Gram pour $h\ne0$ et
\eqref{eq:pdf2-ym-owner-wc-curvature-functional} prouvent que le réseau n’est pas scalaire.

\subsubsection{Axiomes d'Osterwalder--Schrader et reconstruction continue}

Pour une direction $\nu$, soit $\mathfrak A_{+,\nu}$ l’adhérence de l’union des
algèbres précédentes à support compact dans $\{x_\nu>0\}$. L’état
statique, la loi des piles et l’opérateur de transfert ne sont pas définis
séparément. Sur un écran $\alpha$, soient
$\widehat\mu_{\alpha,g}^{\rm YM}$ la marginale de
\eqref{eq:pdf2-ym-gibbs-measure},
$\nu_{\alpha,g}:=(\pi_\alpha^{\rm phys})_*
\widehat\mu_{\alpha,g}^{\rm YM}$ sa loi sur le quotient et
$K_{\alpha,g}(t)=e^{-tD_{\alpha,g}^{\rm YM}}$ son noyau.
Pour des temps ordonnés $t_1\le\cdots\le t_n$, la marginale euclidienne de la
même mesure est
\begin{equation}\label{eq:pdf2-ym-owner-euclidean-time-marginal}
 \mathbb P_{\alpha,g}^{E,(\nu;t_1,\ldots,t_n)}
 :=(\operatorname{ev}_{t_1}^{(\nu)},\ldots,
    \operatorname{ev}_{t_n}^{(\nu)})_*
    \widehat\mu_{\alpha,g}^{\rm YM}.
\end{equation}
L'identité de désintégration qui fait partie de \eqref{eq:pdf2-g-gau-transfer-euclidean-dato} la calcule, sans la remplacer par une chaîne auxiliaire :
\begin{equation}\label{eq:pdf2-ym-owner-euclidean-markov-identification}
 d\mathbb P_{\alpha,g}^{E,(\nu;t_1,\ldots,t_n)}
 =\nu_{\alpha,g}(dy_1)
   \prod_{j=1}^{n-1}K_{\alpha,g}
      (t_{j+1}-t_j;y_j,dy_{j+1}).
\end{equation}
Si $M_A$ désigne la multiplication par un lecteur borné $A$, ses corrélations
sont donc
\begin{equation}\label{eq:pdf2-ym-owner-schwinger-transfer}
 \begin{aligned}
 S_{\alpha,n}^{(\nu,g)}(A_1,t_1;\ldots;A_n,t_n)
 &:=\int\prod_{j=1}^n(\tau_{t_je_\nu}A_j)(x)\,
       d\widehat\mu_{\alpha,g}^{\rm YM}(x)\\
 &=\langle\mathbf1,
 M_{A_1}K_{\alpha,g}(t_2-t_1)M_{A_2}\cdots
 K_{\alpha,g}(t_n-t_{n-1})M_{A_n}\mathbf1
 \rangle_{L^2(\nu_{\alpha,g})}.
 \end{aligned}
\end{equation}
La loi de la pile stationnaire à ces mêmes temps est cette marginale euclidienne, écrite au moyen de sa désintégration :
\begin{equation}\label{eq:pdf2-ym-owner-markov-stack-law}
 d\mathbb P_{\alpha,g}^{(\nu;t_1,\ldots,t_n)}
 :=d\mathbb P_{\alpha,g}^{E,(\nu;t_1,\ldots,t_n)}
 =\nu_{\alpha,g}(dy_1)
   \prod_{j=1}^{n-1}K_{\alpha,g}(t_{j+1}-t_j;y_j,dy_{j+1}).
\end{equation}
Intégrer successivement $y_n,y_{n-1},\ldots,y_1$ donne l’identité exacte
\begin{equation}\label{eq:pdf2-ym-owner-stack-transfer-identity}
 S_{\alpha,n}^{(\nu,g)}(A_1,t_1;\ldots;A_n,t_n)
 =\int\prod_{j=1}^nA_j(y_j)\,
   d\mathbb P_{\alpha,g}^{(\nu;t_1,\ldots,t_n)}(y).
\end{equation}
Si tous les temps coïncident, les noyaux sont l’identité et
\begin{equation}\label{eq:pdf2-ym-owner-static-is-stack}
 S_{\alpha,n}^{(\nu,g)}(A_1,t;\ldots;A_n,t)
 =\int A_1\cdots A_n\,d\nu_{\alpha,g}.
\end{equation}
Ainsi, l'état de Schwinger statique est la marginale à temps égal de la même mesure euclidienne, et sa désintégration est la loi de Markov ; ce ne sont pas deux fonctionnelles que l'on identifie ensuite par leur nom. La compatibilité de $\widehat J_{\alpha,g}$ avec la mesure et le noyau rend \eqref{eq:pdf2-ym-owner-stack-transfer-identity} compatible avec toute superposition. Sa limite définit
\begin{equation}\label{eq:pdf2-ym-owner-schwinger-functionals}
 S_n^{\rm YM}(A_1,t_1;\ldots;A_n,t_n)
 :=\lim_{\alpha}S_{\alpha,n}^{(\nu,g)}
 (A_{1,\alpha},t_1;\ldots;A_{n,\alpha},t_n).
\end{equation}

Les insertions de $\Phi$ s’obtiennent en dérivant en $t=0$ les unitaires
\eqref{eq:pdf2-ym-owner-curvature-weyl}. Puisque les $Z_{\alpha,c}$ transportés
sont centrés, indépendants et uniformément bornés, le lemme de Hoeffding donne,
niveau par niveau,
\begin{equation}\label{eq:pdf2-ym-owner-subgaussian-bound}
 \left\langle\Omega_\alpha,
 e^{t\Phi_\alpha(h)}\Omega_\alpha\right\rangle
 \le \exp\!\left(\frac9{16}t^2\|P_\alpha h\|_{L^2}^2\right).
\end{equation}
La borne passe à la limite par
\eqref{eq:pdf2-ym-owner-smeared-cauchy-derived} ; toutes les dérivées sont continues dans
la topologie de Schwartz et définissent des distributions tempérées.

La forme à huit faces se prolonge d’abord à toute pile finie. Si $F$ dépend des
$N$ couches positives d’un écran $\alpha$, la réflexion de
\eqref{eq:pdf2-ym-owner-markov-stack-law} et le conditionnement sur la couche centrale donnent
\begin{equation}\label{eq:pdf2-ym-owner-finite-stack-os}
 \begin{aligned}
 &\int \overline{F(\Theta_\nu y_{-N},\ldots,
                         \Theta_\nu y_{-1})}
          F(y_1,\ldots,y_N)\,d\Omega_{\alpha,N}^{(\nu,g)}(y)\\
 &\qquad=
 \int\nu_{\alpha,g}(dz)\left|
  \int K_{\alpha,g}(a_\alpha/2;z,dy_1)
       \prod_{j=1}^{N-1}K_{\alpha,g}(a_\alpha;y_j,dy_{j+1})
       F(y_1,\ldots,y_N)
 \right|^2\ge0.
 \end{aligned}
\end{equation}
Tout cylindre de demi-espace appartient à une pile finie ; par conséquent, \eqref{eq:pdf2-ym-owner-finite-stack-os} matérialise la forme OS avant l'adhérence. Dans le quotient OS, si $F$ possède des insertions à des temps positifs, soit $(\tau_sF)(A_1,t_1;\ldots;A_n,t_n):=
F(A_1,t_1+s;\ldots;A_n,t_n+s)$. On définit l'opérateur de translation par
\begin{equation}\label{eq:pdf2-ym-owner-os-transfer-operator}
 T_{\rm OS}^{(\nu)}(s)[F]:=[\tau_sF],
 \qquad s\ge0.
\end{equation}
Pour des cylindres $F,G$ de demi-espace, l’égalité euclidienne--Markov
\eqref{eq:pdf2-ym-owner-euclidean-markov-identification} et la propriété de
semi-groupe donnent générateur par générateur
\begin{equation}\label{eq:pdf2-ym-owner-os-markov-intertwining}
 \begin{aligned}
 \langle[F],T_{\rm OS}^{(\nu)}(s)[G]\rangle_{\rm OS}
 &=S_g^{\rm YM}((\Theta_\nu F)^*\tau_sG)\\
 &=\left\langle
    \mathcal V_{\infty,\nu}[F],
    e^{-sD_{\rm curv,g}^{\rm YM}}
    \mathcal V_{\infty,\nu}[G]
   \right\rangle .
 \end{aligned}
\end{equation}
L’identité montre à la fois que les classes nulles restent nulles et que,
par densité des cylindres,
\begin{equation}\label{eq:pdf2-ym-owner-os-generator-identification}
 \mathcal V_{\infty,\nu}T_{\rm OS}^{(\nu)}(s)
 \mathcal V_{\infty,\nu}^{-1}
 =e^{-sD_{\rm curv,g}^{\rm YM}}.
\end{equation}
C’est la flèche explicite mesure euclidienne $\to$ désintégration $\to$
transfert OS $\to$ hamiltonien.

\begin{proposition}[Paquet d’Osterwalder--Schrader continu]
\label{prop:pdf2-ym-owner-continuum-os-package}
Les fonctionnelles \eqref{eq:pdf2-ym-owner-schwinger-functionals} satisfont :
\begin{enumerate}
 \item normalisation, hermiticité et symétrie bosonique ;
 \item covariance euclidienne et régularité tempérée ;
 \item positivité par réflexion dans chacune des quatre directions,
 \begin{equation}\label{eq:pdf2-ym-owner-continuum-reflection-positive}
  \sum_{i,j}\overline{a_i}a_j\,
  S^{\rm YM}\!\left((\Theta_\nu F_i)^*F_j\right)\ge0,
  \qquad F_i\in\mathfrak A_{+,\nu};
 \end{equation}
 \item localité et propriété de cluster exponentielle : si $A,B$ sont locaux et centrés et si une
 translation de longueur $r$ sépare leurs supports, alors
 \begin{equation}\label{eq:pdf2-ym-owner-continuum-cluster}
  |S^{\rm YM}(A\,\tau_rB)|
  \le e^{-3\kappa_{\rm av}r}\,
       \|A\Omega_\infty\|\,\|B\Omega_\infty\|.
 \end{equation}
\end{enumerate}
La reconstruction OS de l'un quelconque des quatre demi-réseaux produit le même espace cyclique, le même vide et l'hamiltonien $D_{\rm curv,g}^{\rm YM}$. Le prolongement OS produit des champs de Wightman locaux et une représentation de Poincaré dont le spectre satisfait
\begin{equation}\label{eq:pdf2-ym-owner-wightman-spectrum}
 \operatorname{Spec}(P^0,\mathbf P)\subseteq\overline V_+,
 \qquad P^0=D_{\rm curv,g}^{\rm YM}\ge0.
\end{equation}
\end{proposition}

\begin{proof}
La normalisation et l’hermiticité proviennent de l’état ; la symétrie et la localité, de la
commutativité euclidienne à supports disjoints. La covariance et la continuité sont
\eqref{eq:pdf2-ym-owner-field-e4-covariance}--
\eqref{eq:pdf2-ym-owner-strong-e4} ; la régularité est
\eqref{eq:pdf2-ym-owner-subgaussian-bound}. Pour
$F_i$ cylindriques, tous vivent dans une superposition et une pile finies ;
\eqref{eq:pdf2-ym-owner-finite-stack-os} et la polarisation prouvent
\eqref{eq:pdf2-ym-owner-continuum-reflection-positive}. L’approximation en norme OS
et Cauchy--Schwarz l’étendent à $\mathfrak A_{+,\nu}$.

Après une rotation euclidienne, la séparation spatiale devient un temps positif.
L’identité pile--transfert
\eqref{eq:pdf2-ym-owner-stack-transfer-identity} et l’opérateur
\eqref{eq:pdf2-ym-owner-os-transfer-operator} écrivent la corrélation centrée sous la forme
\[
 \langle A^*\Omega_\infty,
 e^{-rD_{\rm curv,g}^{\rm YM}}B\Omega_\infty\rangle.
\]
La borne spectrale \eqref{eq:pdf2-ym-owner-curvature-sector-gap} donne \eqref{eq:pdf2-ym-owner-continuum-cluster}. Enfin, les connecteurs \eqref{eq:pdf2-ym-owner-os-connectors} forment un système dirigé ; leur réunion est dense d'après \eqref{eq:pdf2-ym-owner-smooth-core}. L'identité \eqref{eq:pdf2-ym-owner-common-os-hamiltonian} identifie les quatre adhérences et leurs générateurs à $D_{\rm curv,g}^{\rm YM}$. Les hypothèses déjà vérifiées de régularité, de covariance, de réflexion et de cluster sont les entrées du prolongement OS ; son théorème de reconstruction donne la localité de Wightman et \eqref{eq:pdf2-ym-owner-wightman-spectrum}.
\end{proof}

\subsubsection{Reconstruction quadridimensionnelle de la théorie de Yang--Mills}

L’identification ne repose pas seulement sur l’évaluation d’un clover sur une connexion donnée.
Soit $\mathcal R_G^{\rm loc}$ la $*$-algèbre de lecteurs formée par toutes les
fonctions de classe des holonomies d’incidence et de plaquette, leurs transports vers un
sommet de base et les polarisations du lecteur $\mathcal P^{F^2}$. Soit
\begin{equation}\label{eq:pdf2-ym-owner-curvature-fibre}
 \mathcal E_G:=\Lambda^2(\mathbb R^4)^*\otimes\mathfrak g
\end{equation}
avec le produit invariant. Dans chaque cellule, la polarisation de
$\mathcal P^{F^2}$ donne un gramien positif $G_{\alpha,c}^{F}$ sur
$\mathcal E_G$. On quotiente par son noyau et on applique sa racine positive ; le lecteur
normalisé résultant est noté $Y_{\alpha,c}(v)$ et vérifie
\begin{equation}\label{eq:pdf2-ym-owner-curvature-cell-gram}
 \langle Y_{\alpha,c}(v),Y_{\alpha,d}(w)\rangle
 =\delta_{cd}\langle v,w\rangle_{\mathcal E_G}.
\end{equation}
Si $d\subset c$, soit $T_{d\leftarrow c}$ le transport de repère déjà présent dans l'holonomie de surface. L'opérateur d'entrelacement de courbure complète est
\begin{equation}\label{eq:pdf2-ym-owner-curvature-intertwiner}
 I_{\alpha\beta}^{F}Y_{\alpha,c}(v)
 :=\sum_{d\subset c}\sqrt{\frac{|d|}{|c|}}\,
 Y_{\beta,d}\!\left(\operatorname{Ad}_{T_{d\leftarrow c}}v\right).
\end{equation}
L’invariance du produit de $\mathfrak g$, la disjonction des sous-cellules et la
composition des transports prouvent
\begin{equation}\label{eq:pdf2-ym-owner-curvature-intertwiner-identities}
 (I_{\alpha\beta}^{F})^*I_{\alpha\beta}^{F}=I,
 \qquad
 I_{\beta\gamma}^{F}I_{\alpha\beta}^{F}=I_{\alpha\gamma}^{F}.
\end{equation}
Pour $h\in C_c^\infty(\mathbb R^4;\mathcal E_G)$, on pose
\begin{equation}\label{eq:pdf2-ym-owner-f-alpha}
 \mathbf F_\alpha(h):=
 \sum_{c\in\mathcal C_\alpha^{\rm cell}}
 Y_{\alpha,c}\!\left(
   |c|^{-1/2}\int_c h(x)\,d^4x\right).
\end{equation}
Dans une superposition commune, le même calcul de
\eqref{eq:pdf2-ym-owner-smeared-cross-gram} donne maintenant
\begin{equation}\label{eq:pdf2-ym-owner-curvature-cross-gram}
 \left\langle I_{\alpha\gamma}^{F}\mathbf F_\alpha(h),
 I_{\beta\gamma}^{F}\mathbf F_\beta(k)\right\rangle
 =\langle P_\alpha h,P_\beta k\rangle_{L^2(\mathcal E_G)}.
\end{equation}
Par conséquent, $\mathbf F_\alpha(h)$ est de Cauchy par convergence forte des projections,
et non par une nouvelle hypothèse. Le même appariement de Haar et des innovations de
\eqref{eq:pdf2-ym-owner-publication-into-physical}, appliqué maintenant à la base
orthonormalisée de $G_{\alpha,c}^{F}$, réalise la colimite par une isométrie
locale et covariante de jauge dans
$L^2(\widehat X_\infty,\widehat\mu_{\infty,g}^{\rm YM})$. Sa limite définit la
distribution aléatoire
covariante de jauge
\begin{equation}\label{eq:pdf2-ym-owner-distributional-curvature}
\begin{aligned}
 \mathbf F_{\mu\nu}&\in
 \mathcal S'(\mathbb R^4;\mathfrak g)\ \widehat\otimes\
 L^2(\widehat X_\infty,\widehat\mu_{\infty,g}^{\rm YM}),\\
 U(a)\mathbf F(h)U(a)^*&=\mathbf F(a\!\cdot h),\\
 u\mathbf F(h)u^{-1}&=\mathbf F(\operatorname{Ad}_u h),
 \qquad a\in E(4),\ u\in G.
\end{aligned}
\end{equation}
Son canal de classe spectrale d'ordre trois est précisément $\Phi$. La relation de surface finie
\begin{equation}\label{eq:pdf2-ym-owner-surface-product}
 \operatorname{Hol}^{(\alpha)}_p
 =\prod_{p'\subset p}^{\longrightarrow}
  T_{p'}\operatorname{Hol}^{(\beta)}_{p'}T_{p'}^{-1}
\end{equation}
passe à la fermeture cofinale car tous ses produits croisés sont calculés dans une
superposition. En particulier, le lecteur clover aléatoire, et non seulement son évaluation dans une
connexion lisse, satisfait l’identité des petites boucles
\begin{equation}\label{eq:pdf2-ym-owner-holonomy-curvature-relation}
 \mathbf F(h)=L^2\!\mathord{-}\!\lim_{\alpha}
 a_\alpha^4\sum_x
 \left\langle(P_\alpha h)(x),F_{{\rm cl},\alpha}(x)\right\rangle,
 \qquad
 \delta_{\rm gauge}S_n^{\rm YM}=0,
 \qquad
 \operatorname{Alt}_{\lambda\mu\nu}D_\lambda\mathbf F_{\mu\nu}=0.
\end{equation}
La première égalité est \eqref{eq:pdf2-ym-owner-curvature-cross-gram} écrite dans le représentant clover ; la seconde est l'identité de Ward \eqref{eq:pdf2-ym-ward-finite} passée à la limite, et la troisième est l'annulation de la frontière d'une frontière dans le produit de surface. Aucune n'utilise le terme de masse.

Pour éviter de définir la cible par la conclusion que l’on veut obtenir, soit
$\mathbf{YM}_4(G,g)$ la classe des quintuplets locaux
$(\mathfrak A,\mu_g,K_g,\operatorname{Hol},\mathbf F)$ qui vérifient les
équations vérifiables suivantes : (i) la densité de Gibbs est
\eqref{eq:pdf2-ym-gibbs-measure} ; (ii) $K_g$ est markovien, réversible et satisfait
\eqref{eq:pdf2-ym-schwinger-dyson-finite}--\eqref{eq:pdf2-ym-ward-finite} ;
(iii) holonomie, clover et courbure satisfont
\eqref{eq:pdf2-ym-transport-surface-bianchi} et
\eqref{eq:pdf2-ym-owner-holonomy-curvature-relation} ; et (iv) leurs fonctionnelles sont la
même loi de piles de \eqref{eq:pdf2-ym-owner-stack-transfer-identity} et satisfont le
paquet OS\@. L’objet construit est, composante par composante,
\begin{equation}\label{eq:pdf2-ym-owner-qym-object}
\begin{aligned}
 \mathfrak Q_{\rm YM}^{(4)}(G,g):={}&
 \bigl(\{\mathfrak A_{\rm YM}(O)\}_O,
 S_g^{\rm YM},U,\Theta,\Omega_{\infty,g},\\
 &\widehat\mu_{\infty,g}^{\rm YM},
 \{S_{m,g},\mathscr S_{m,g}\}_m,K_g,
 \operatorname{Hol},\mathbf F,\\
 &\mathscr H_{\rm OS},D_{\rm curv,g}^{\rm YM}\bigr)
 \in\mathbf{YM}_4(G,g).
\end{aligned}
\end{equation}
L'application de représentation n'omet aucune de ces composantes :
\begin{equation}\label{eq:pdf2-ym-owner-typed-publication-map}
 \begin{aligned}
 \mathcal P_{{\rm YM},g}:\mathfrak G_{\rm gau}(G)&\longrightarrow
 \mathbf{YM}_4(G,g),\\
 (\widehat\mu^0,\mathbb P^{\rm phys},D^0,\Theta,W,
   \operatorname{Hol},\mathcal P^{F^2})
 &\longmapsto
 (\widehat\mu_g^{\rm YM},S_g,\mathscr S_g,K_g,\mathbb P_g^{\rm phys},
  S_g^{\rm YM},\mathfrak A_{\rm YM},D_{\rm curv,g}^{\rm YM},
  \mathbf F,\operatorname{Hol}).
 \end{aligned}
\end{equation}
Les $S_n^{\rm YM}$ sont ainsi exactement les fonctions de Schwinger de la mesure dépendante de $g$, de la loi de Markov et de l'opérateur de transfert déjà construits.

Il reste à vérifier que $S_{m,g}$ est l'action de Yang--Mills et pas seulement une fonction portant ce nom. Soit $A\in C_c^3(\mathbb R^4;T^*\mathbb R^4\otimes\mathfrak g)$ une connexion lisse dans une coordinatisation où le logarithme des petites boucles est unique. Pour la plaquette $p_{\mu\nu}(x)$, le développement ordonné de Magnus donne
\begin{equation}\label{eq:pdf2-ym-owner-plaquette-magnus}
 \log\operatorname{Hol}_{p_{\mu\nu}(x)}(A)
 =a_m^2F_{A,\mu\nu}(x)
  +\frac{a_m^3}{2}(D_\mu+D_\nu)F_{A,\mu\nu}(x)
  +R_{m,\mu\nu}(x),
\end{equation}
avec $\|R_{m,\mu\nu}(x)\|\le C_G a_m^4
(\|D^2F_A\|_{\infty,p}+\|F_A\|_{\infty,p}^2)$. La moyenne des quatre orientations annule le terme impair et définit
\begin{equation}\label{eq:pdf2-ym-owner-clover}
 F_{{\rm cl},m}(x)
 :=\frac1{4a_m^2}\sum_{p\ni x}
  \operatorname{Ad}_{T_p}\log\operatorname{Hol}_p,
 \qquad
 \|F_{{\rm cl},m}(x)-F_A(x)\|
 \le C_G a_m^2(\|D^2F_A\|_\infty+\|F_A\|_\infty^2).
\end{equation}
La définition de la racine \eqref{eq:pdf2-ym-action-root} est précisément la polarisation de ces six composantes. En substituant \eqref{eq:pdf2-ym-owner-plaquette-magnus} dans chaque générateur de cette polarisation, on calcule
\begin{equation}\label{eq:pdf2-ym-owner-action-root-smooth-chart}
 \xi_{m,b}(A)
 =a_m^2\operatorname{vec}_{\mu<\nu}F_{A,\mu\nu}(x_b)
  +r_{m,b}(A),
 \qquad
 \|r_{m,b}(A)\|\le C_{G,A}a_m^4.
\end{equation}
Les termes croisés impairs s’annulent par les quatre orientations du clover ;
les termes pairs ne sont inclus qu’une seule fois par
\eqref{eq:pdf2-ym-action-factor-partition}. Par somme de Riemann et
Cauchy--Schwarz,
\begin{equation}\label{eq:pdf2-ym-owner-action-error}
 \left|S_{m,g}(A)-\frac1{4g^2}
  a_m^4\sum_x\|F_{{\rm cl},m}(x)\|^2\right|
 \le C_{G,A,g}a_m^2,
\end{equation}
et, en utilisant la borne de \eqref{eq:pdf2-ym-owner-clover},
\begin{equation}\label{eq:pdf2-ym-owner-f2}
 S_{m,g}(A)\longrightarrow
 S_{{\rm YM},g}(A):=\frac1{4g^2}
 \int_{\mathbb R^4}\|F_A(x)\|^2\,d^4x.
\end{equation}
Les équations de Schwinger--Dyson et de Ward ont déjà été calculées sur la mesure qui utilise
cette action. L’écart, en revanche, se déduit ensuite de la conjugaison unitaire de
\eqref{eq:pdf2-ym-g-dynamics} et du témoin
\eqref{eq:pdf2-ym-g-gap-before-limit} ; il n’a pas été utilisé pour choisir $g$, $S_{m,g}$ ni
$\widehat\mu_{m,g}^{\rm YM}$. Tel est le lien non circulaire entre action, loi et écart.

\begin{theorem}[Clôture de Yang--Mills pour tout groupe de Lie compact, connexe et simple]
\label{thm:pdf2-ym-cierre-global}
La structure complète $\mathfrak G_{\rm gau}$ construit une collection projective locale pour tout groupe de Lie compact, connexe et simple. L'application \eqref{eq:pdf2-ym-owner-typed-publication-map} exprime une théorie quantique de Yang--Mills non triviale sur $\mathbb R^4$ : ses fonctionnelles satisfont l'ensemble continu des axiomes OS et leur prolongement satisfait les axiomes locaux de Wightman d'après la proposition~\ref{prop:pdf2-ym-owner-continuum-os-package} ; ses holonomies et sa courbure distributionnelle satisfont \eqref{eq:pdf2-ym-owner-holonomy-curvature-relation}, et les quatre reconstructions coinduisent le même hamiltonien $D_{\infty,g}^{\rm YM}$, avec un vide simple. Sa loi est $\widehat\mu_{\infty,g}^{\rm YM}$, construite par la collection compatible d'actions locales $\{S_{m,g},\mathscr S_{m,g}\}_m$, et satisfait les équations de Schwinger--Dyson et de Ward \eqref{eq:pdf2-ym-schwinger-dyson-finite}--\eqref{eq:pdf2-ym-ward-finite}. Le secteur cyclique de courbure \eqref{eq:pdf2-ym-owner-curvature-cyclic-sector} est réducteur, contient le témoin \eqref{eq:pdf2-ym-owner-wc-curvature-functional} et possède
\begin{equation}\label{eq:pdf2-ym-final-gap}
 \boxed{\Delta_{\rm YM}^{\rm HMT}=3\kappa_{\rm av}>0},
 \qquad
 m_{\rm gap}^{\rm HMT}
 =\frac{\hbar_{\rm int}}{c_{\rm int}}\Delta_{\rm YM}^{\rm HMT}>0.
\end{equation}
Le réseau local, l'action forte de $E(4)$, les distributions de Schwinger, l'holonomie et $\mathbf F$ proviennent du même état complet. Le lecteur exact $\mathcal P^{F^2}$ définit la densité de la loi ; le clover est sa coordinatisation lisse et \eqref{eq:pdf2-ym-owner-f2} identifie son action classique. Le saut appartient au secteur collectif de courbure invariant de jauge et n'attribue pas de masse élémentaire au gluon.
\end{theorem}

\begin{proof}
La projectivité cinématique est \eqref{eq:pdf2-ym-owner-projectivity} ; la racine
\eqref{eq:pdf2-ym-action-root}, le cocycle
\eqref{eq:pdf2-ym-action-cocycle} et la densité
\eqref{eq:pdf2-ym-gibbs-measure} construisent d’abord l’action et la mesure au
couplage $g$. Le transport facteur par facteur
\eqref{eq:pdf2-ym-gibbs-transport-factor}--
\eqref{eq:pdf2-ym-full-algebra-identities} prouve la projectivité, la préservation de
l’état, des produits et de $*$ ;
\eqref{eq:pdf2-ym-transport-surface-bianchi} établit clover et Bianchi. L’indice exhaustif
\eqref{eq:pdf2-ym-owner-exhaustive-index}, la forme
\eqref{eq:pdf2-ym-owner-full-form} et le circuit
\eqref{eq:pdf2-ym-owner-local-circuit} construisent le générateur commun. La réduction
physique et son noyau quotient sont
\eqref{eq:pdf2-ym-owner-physical-reduction}--
\eqref{eq:pdf2-ym-owner-D}. La conjugaison
\eqref{eq:pdf2-ym-g-dynamics} les transporte à la mesure dépendant de $g$ sans
altérer leur spectre. La décomposition de Hoeffding prouve le vide et la borne, et
\eqref{eq:pdf2-ym-owner-wc-invariance}--
\eqref{eq:pdf2-ym-owner-wc-eigenvalue} exhibent un vecteur physique qui atteint le
seuil. L’entrelacement nonadique transporte forme, vide et témoin à la limite, où
\eqref{eq:pdf2-ym-owner-limit-gap}--
\eqref{eq:pdf2-ym-owner-exact-gap} fixent l’écart exact. Les quatre
factorisations \eqref{eq:pdf2-ym-owner-four-os}, l’inclusion des images
\eqref{eq:pdf2-ym-owner-os-range-inclusion} et les connecteurs
\eqref{eq:pdf2-ym-owner-os-connectors} identifient leurs reconstructions à ce
même générateur. L’holonomie
\eqref{eq:pdf2-ym-owner-incidence-holonomy} place matériellement $W_c$ dans le secteur
de courbure, où \eqref{eq:pdf2-ym-owner-curvature-sector-gap} prouve l’écart
exact. Le système \eqref{eq:pdf2-ym-owner-block-intertwiner} contrôle les produits
croisés au moyen de \eqref{eq:pdf2-ym-owner-smeared-cross-gram} ;
\eqref{eq:pdf2-ym-full-publication-restriction}--
\eqref{eq:pdf2-ym-full-publication-products} élèvent ce premier chaos à l’isomorphisme
de toute l’algèbre. La convergence forte de $P_\alpha$ produit $\Phi$ ; les bornes
\eqref{eq:pdf2-ym-owner-holonomy-e4-estimate} et
\eqref{eq:pdf2-ym-owner-orbit-smearing-continuity} prouvent en outre la continuité
$E(4)$ des holonomies, des plaquettes et de clover. La désintégration euclidienne
\eqref{eq:pdf2-ym-owner-euclidean-markov-identification}, l’identité exacte
\eqref{eq:pdf2-ym-owner-stack-transfer-identity} et l’entrelacement
\eqref{eq:pdf2-ym-owner-os-markov-intertwining} relient la même mesure,
l’état de Schwinger, la pile de Markov, le transfert OS et le générateur de l’écart. La
proposition~\ref{prop:pdf2-ym-owner-continuum-os-package} prolonge les
quatre formes finies aux
demi-réseaux complets, prouve la régularité et la propriété de cluster et reconstruit l’hamiltonien. Pour
finir, \eqref{eq:pdf2-ym-owner-distributional-curvature}--
\eqref{eq:pdf2-ym-owner-typed-publication-map} identifient ces mêmes fonctionnelles
comme les fonctionnelles de Schwinger de la théorie de jauge locale ; les intégrations par parties
\eqref{eq:pdf2-ym-schwinger-dyson-finite}--\eqref{eq:pdf2-ym-ward-finite} prouvent
Schwinger--Dyson et Ward. Enfin,
\eqref{eq:pdf2-ym-owner-plaquette-magnus}--\eqref{eq:pdf2-ym-owner-f2} identifient
l’action classique avant d’appliquer l’équivalence spectrale. La chaîne n’utilise pas
l’écart pour choisir la mesure.
\end{proof}

\begin{falsifier}[Effondrement de la dynamique produit]
Si l’on remplace $D_m^{\rm prod}$ par $D_m^{\rm fib}$, tous les secteurs non
constants reçoivent la même valeur propre. Un vecteur de Hoeffding avec $|S|=2$ devrait
avoir l’énergie $6\kappa_{\rm av}$ par \eqref{eq:pdf2-ym-hoeffding}, mais le
semi-groupe simple lui attribuerait $3\kappa_{\rm av}$. L’identification des deux
générateurs est fausse.
\end{falsifier}

\begin{falsifier}[Confusion entre inclusion généalogique et moyenne de bloc]
Si l'on remplace $I_{\alpha\beta}^{\rm curv}$ par $\widehat J_{\alpha\beta}$ dans \eqref{eq:pdf2-ym-owner-smeared-cross-gram}, une coordonnée grossière persiste avec un coefficient égal à un tandis que le coefficient fin change selon le quotient des échelles. Pour $a_{m+1}=a_m/9$, le produit croisé ancestral porte un facteur $1/81$ et ne converge pas vers la norme diagonale. L'identité \eqref{eq:pdf2-ym-owner-smeared-cauchy-derived} échoue ; c'est pourquoi les deux applications restent séparées.
\end{falsifier}

\begin{falsifier}[Ablation de la fibre d’incidence]
Si l’on remplace $\widehat X_m$ par sa marginale de liens, on efface $q_c$ et, avec elle,
l’observable physique $W_c$. L’application d’oubli ne reflète pas la sous-algèbre de jauge complète ;
elle ne peut donc pas être utilisée pour nier la valeur propre de l’objet source. Objet
remplacé ; conclusion non transférable.
\end{falsifier}

\begin{falsifier}[Substitution du lecteur exact par sa coordinatisation lisse]
Si le terme asymptotique de \eqref{eq:pdf2-ym-owner-f2} est utilisé directement comme densité sur un écran fini, le reste contrôlé de \eqref{eq:pdf2-ym-owner-action-error} est effacé et le cocycle exact \eqref{eq:pdf2-ym-action-cocycle} est perdu. La densité finie gouvernante est $Z_{m,g}^{-1}e^{-\mathcal P_m^{F^2}/(4g^2)}$ ; le clover et la limite lisse l'identifient à l'action classique, mais ne la remplacent pas.
\end{falsifier}

\paragraph{Articulation des démonstrations.}
La séparation entre fibre simple et générateur produit, la mesure commune, la compression de jauge, les quatre conditions d'Osterwalder--Schrader, la forme limite et le témoin physique sont des résultats exacts préalablement établis dans le domaine de jauge régulier HMT. La désintégration à temps non uniformes, l'indice exhaustif des espérances, le noyau quotient, la matrice $B^+$, l'action explicite sur la colimite, le système cofinal de bloc, l'ensemble continu des conditions OS, la courbure distributionnelle et l'holonomie d'incidence sont des formalisations explicites : elles déterminent les domaines et démontrent les liens précédemment énoncés, tout en maintenant l'écart spectral comme résultat. Leur composition complète établit le théorème précédent.
La racine positive du lecteur $F^2$, sa densité de Gibbs, le transport local compatible, les équations de Schwinger--Dyson/Ward, l'isomorphisme algébrique complet, les bornes de continuité des holonomies et l'identité pile--transfert sont des formalisations explicites des quatre correspondances. L'identification entre l'observable clover et l'action classique s'obtient comme conséquence du développement de Magnus.

\section{Générateurs tritiques, actions quadratiques et transport d’énergie}
\label{sec:hamiltonianos-curvatura}

La géométrie dynamique conserve la distinction entre l’état qui engendre une
réalisation et les coordonnées avec lesquelles celle-ci exprime son évolution. APP
construit et conserve les feuillets additif et multiplicatif ; TRIT oriente le
mode opératoire ; TPK compose sélection, transport et actualisation sur
l’état enrichi au moyen de
\[
 U_t=\operatorname{Upd}_t\circ\operatorname{Tra}_t\circ
 \operatorname{Sel}_t.
\]
L'holonomie nonadique prolonge cet état à travers \(w_6\to w_{12}\to w_{18}\to w_{24}\to w_{30}\to R_{36}\to G_9\). Le retour de phase conserve l'avancement de la mémoire. Ses représentations \(\pi_{\rm HMT},\varphi_{\rm HMT},e_{\rm HMT},\alpha_{\rm HMT}\) appartiennent à la même génération corrélée, avec les lecteurs et l'ordre constructif interne déjà établis. La structure discrète conjointe du continu soutient ensuite les représentations géométriques et physiques. À ce point de la construction, un hamiltonien est la fonctionnelle réalisant un transport déjà typé sur une fibre symplectique ; son domaine, son orientation et son paramètre temporel font partie de l'affirmation.

L’article sur l’extrême et moyenne raison holographique récupère trois
opérateurs concrets. La multiplication par \(4\) modulo \(9\) agit sur
les orbites ternaires d’unités ; dans son module d’augmentation, elle a pour matrice
\[
 C=\begin{pmatrix}0&-1\\1&-1\end{pmatrix},\qquad C^2+C+I=0.
\]
Le transport parabolique et l’incidence aréale--volumétrique sont
représentés, dans leurs plans respectifs, par
\[
 N=\begin{pmatrix}0&1\\0&0\end{pmatrix},\qquad
 F_{\rm av}=\begin{pmatrix}0&1\\1&1\end{pmatrix}.
\]
L'incidence provient du calendrier des feuillets ; ses puissances comptent des compositions de ces modes. Les générateurs réels exprimés sont
\begin{equation}
 J_+=\frac{2C+I}{\sqrt3},\qquad J_0=N,\qquad
 J_-=\frac{2F_{\rm av}^{2}-3I}{\sqrt5},\qquad
 J_\tau^2=-\tau I.
 \label{eq:hc-generadores}
\end{equation}
Ici \(\tau=+1,0,-1\) distingue les régimes elliptique, parabolique et
hyperbolique. Les fibres ont des provenances et des bases propres. Les formules
\(e^{\pi J_+}=-I\), \(e^{J_0}=I+J_0\) et
\(e^{2\log\varphi J_-}=F_{\rm av}^{2}\) évaluent des coordonnées HMT
préalablement engendrées dans ces opérateurs~\cite{hmtX}.

\subsection{Forme alternée et Hamiltonien du générateur}

Dans une fibre réelle orientée de dimension deux, on fixe
\[
 \omega=dQ\wedge dP,\qquad
 \Omega=\begin{pmatrix}0&1\\-1&0\end{pmatrix},\qquad z=(Q,P)^T.
\]
On adopte \(\iota_{X_H}\omega=dH\) ; par conséquent \(\dot Q=\partial_PH\), \(\dot P=-\partial_QH\). Cette convention fixe également le signe de l'action orientée.

\begin{proposition}[Réalisation hamiltonienne d’un générateur de trace nulle]
\label{prop:hc-cuadratica}
Pour une matrice réelle \(J\) de taille deux et de trace nulle, la matrice
\(G_J=-\Omega J\) est symétrique et le Hamiltonien
\begin{equation}
 H_J(z)=\frac12z^TG_Jz
 \label{eq:hc-hj}
\end{equation}
produit exactement \(\dot z=Jz\). Si \(J^2=-\tau I\),
alors \(\det G_J=\tau\). Pour un isomorphisme symplectique
\(B:V\to V'\), le transport
\[
 J'=BJB^{-1},\qquad H'(z')=H_J(B^{-1}z')
\]
satisfait \(G'=B^{-T}G_JB^{-1}\) et entrelace les deux flux.
\end{proposition}
\begin{proof}
Escribiendo \(J=\bigl(\begin{smallmatrix}a&b\\c&-a\end{smallmatrix}\bigr)\), on obtient \(G_J=\bigl(\begin{smallmatrix}-c&a\\a&b\end{smallmatrix}\bigr)\). Ainsi \(\Omega\nabla H_J=\Omega(-\Omega J)z=Jz\). Cayley--Hamilton donne \(J^2=-\det(J)I\), et \(\det(-\Omega)=1\). Enfin, \(B^T\Omega B=\Omega\) permet de substituer \(z=B^{-1}z'\) aussi bien dans la forme que dans l'équation.
\end{proof}

Appliquée à \eqref{eq:hc-generadores}, la proposition donne
\begin{equation}
 H_{J_+}=-\frac{Q^2-QP+P^2}{\sqrt3},\qquad
 H_{J_0}=\frac{P^2}{2},\qquad
 H_{J_-}=\frac{-Q^2-QP+P^2}{\sqrt5}.
 \label{eq:hc-hamiltonianos-nativos}
\end{equation}
Le signe elliptique enregistre l’orientation de \(C\) dans la base d’augmentation.
Le flux inverse \(-J_+\) possède la forme positive opposée. L’ellipse
d’action du corpus utilise précisément une orientation de phase
décroissante pour son Hamiltonien positif. Conserver ce signe rend
l’opérateur cyclique compatible avec l’action \(P\,dQ\).

La conjugaison fournit également un critère exact de comparaison.
Si \(BJ_\tau=J_{\tau'}B\) et \(B\) est inversible, l’élévation au carré
donne \((\tau-\tau')B=0\), donc \(\tau=\tau'\). Un lien peut
conserver \(\omega\) entre fibres de régimes distincts et conserver
leurs mémoires respectives ; la conjugaison des générateurs exige en outre
l’égalité du régime. Cette différence permet d’étudier une transition
tritique sans la remplacer par un changement de nom du même flot.

\subsection{Déformation de l'ellipse et transformation de Legendre}

La coordonnée \(\eta=\operatorname{artanh}(C^*/A)\), dans le domaine \(|C^*/A|<1\), définit le repère symplectique \(M_\eta=\operatorname{diag}(e^{\eta/2},e^{-\eta/2})\). La déformation des deux directions est réciproque et conserve leur produit. Dans la coordinatisation normalisée \((q,p)=M_\eta^{-1}(Q,P)\), considérons
\[
 L_\tau=\begin{pmatrix}0&1\\-\tau&0\end{pmatrix},\qquad
 I_{\tau,\eta}=\frac12\bigl(\tau e^{-\eta}Q^2+e^\eta P^2\bigr).
\]
Ces représentants classent des flux quadratiques orientés. Leur emploi
dans une fibre native s’effectue au moyen d’un repère symplectique déclaré,
conformément à la proposition~\ref{prop:hc-cuadratica}. Le représentant
elliptique positif est celui utilisé par l’ellipse d’action ; pour lui, \(I\)
a la dimension d’une action.

Les changements de repère peuvent s'écrire explicitement. Si \(b_+=\sqrt{\sqrt3/2}\) et \(b_-=\sqrt{\sqrt5/2}\), les matrices
\[
 B_+=\begin{pmatrix}b_+&b_+/\sqrt3\\0&2b_+/\sqrt3\end{pmatrix},
 \qquad
 B_-=\begin{pmatrix}b_-&-b_-/\sqrt5\\0&2b_-/\sqrt5\end{pmatrix}
\]
ont un déterminant égal à un et satisfont \( (-J_+)B_+=B_+L_{+1}\) et \(J_-B_-=B_-L_{-1}\). Pour le régime parabolique, on prend \(B_0=I\). La multiplication des deux colonnes vérifie les égalités, et le déterminant égal à un prouve la conservation de la forme alternée.

\begin{theorem}[Transport variable et lagrangien régulier]
\label{thm:hc-legendre}
Soient \(\eta(t)\) de classe \(C^1\) et \(\omega(t)>0\). L’hamiltonien
\begin{equation}
 H_{\tau,\rm tr}
 =\omega I_{\tau,\eta}+\frac{\dot\eta}{2}QP
 \label{eq:hc-htr}
\end{equation}
conserve \(I_{\tau,\eta(t)}\) le long de ses solutions. Sa
transformation de Legendre par rapport à \(P\) est inversible et donne
\begin{equation}
 \mathcal L_{\tau,\rm tr}(Q,\dot Q,t)
 =\frac{(\dot Q-\dot\eta Q/2)^2}{2\omega e^\eta}
 -\frac{\omega\tau e^{-\eta}}2Q^2.
 \label{eq:hc-ltr}
\end{equation}
\end{theorem}
\begin{proof}
Les équations sont
\[
 \dot Q=\omega e^\eta P+\frac{\dot\eta}2Q,
 \qquad
 \dot P=-\omega\tau e^{-\eta}Q-\frac{\dot\eta}2P.
\]
Dériver \(I_{\tau,\eta(t)}\) et les substituer annule séparément
les termes \(\omega QP\) et les termes de variation de \(\eta\).
La dérivée seconde \(\partial_P^2H=\omega e^\eta>0\)
permet de résoudre
\(P=(\dot Q-\dot\eta Q/2)/(\omega e^\eta)\).
La substitution dans \(P\dot Q-H\) prouve \eqref{eq:hc-ltr}.
De plus,
\[
 P\dot Q-H_{\tau,\rm tr}
 =p\dot q-\frac\omega2(\tau q^2+p^2),
\]
ce qui prouve que le terme mixte est la connexion variationnelle du repère
mobile. La conservation provient du transport complet et non du maintien
des demi-axes constants.
\end{proof}

Le résultat elliptique et son action sont développés dans le corpus variationnel~\cite{hmtX} ; la même preuve exhibe ses deux réalisations quadratiques conjuguées. La nullité du déterminant de la forme parabolique coexiste avec une transformation de Legendre régulière : celle-ci dépend de la dérivée seconde par rapport au moment, et non du déterminant de la forme complète. En revanche, pour le relèvement cotangent linéaire \(H_A(q,p)=p^TAq\), on a \(\partial_p^2H_A=0\). Son action du premier ordre conserve \(\dot q=Aq\) au moyen d'un multiplicateur ; c'est un objet variationnel distinct de \eqref{eq:hc-ltr}.

\subsection{Retour de mémoire et charge variationnelle conservée}

Dans le registre de mémoire du calendrier, les compteurs \((g,s)\)
reçoivent l’accroissement \((4,72)\) au retour complet. La relation
\(72=18\cdot4\) fixe
\[
 v=(1,18)^T,\qquad I_{\rm mem}=s-18g,
 \qquad B=\begin{pmatrix}1&0\\18&1\end{pmatrix}.
\]
Le stabilisateur natif est
\begin{equation}
 N_v=BNB^{-1}=
 \begin{pmatrix}-18&1\\-324&18\end{pmatrix},\qquad
 N_vm=vI_{\rm mem}(m),\qquad N_v^2=0.
 \label{eq:hc-nv}
\end{equation}
Dans le plan réalifié \(m=(g,s)^T\), muni de \(dg\wedge ds\), l'hamiltonien \(H_{\rm mem}=I_{\rm mem}^2/2\) produit \(dm/du=N_vm\). Le paramètre \(u\) appartient à cette réalisation adimensionnelle du registre. La coordinatisation \(q=g,\ p=s-18g\) est symplectique et transforme la un-forme d'action selon
\begin{equation}
 s\,dg=p\,dq+d(9q^2).
 \label{eq:hc-memoria-borde}
\end{equation}
Par conséquent, l’élimination de \(s\) dans l’action du premier ordre donne
\begin{equation}
 \mathcal L_{\rm mem}
 =\frac12\left(\frac{dg}{du}\right)^2
 +18g\frac{dg}{du}
 =\frac12\left(\frac{dg}{du}\right)^2
 +\frac{d}{du}(9g^2).
 \label{eq:hc-lmem}
\end{equation}
L'équation variationnelle est \(d^2g/du^2=0\), et le moment conservé dans la coordinatisation symplectique est \(p=I_{\rm mem}\). Le retour affine complet provient de l'hamiltonien linéaire \(H_{\rm ret}=4I_{\rm mem}\) : son application au temps \(u=1\) est \((g,s)\mapsto(g+4,s+72)\). Les deux hamiltoniens commutent parce qu'ils sont des fonctions de la même charge, mais réalisent des opérations distinctes : stabilisation parabolique et translation du registre. Ces formules restituent le mécanisme variationnel de la mémoire, y compris son terme de frontière~\cite{hmtX}.

\subsection{Courbure métrique et énergie géodésique}

La réalisation métrique des régimes utilise une coordinatisation radiale et une échelle inverse de longueur \(\rho>0\) :
\begin{equation}
 g_{\tau,\rho}=dr^2+S_{\tau,\rho}(r)^2d\theta^2,
 \qquad
 S_{\tau,\rho}(r)=
 \begin{cases}
 \sin(\rho r)/\rho,&\tau=+1,\\
 r,&\tau=0,\\
 \sinh(\rho r)/\rho,&\tau=-1.
 \end{cases}
 \label{eq:hc-metrique}
\end{equation}
La courbure gaussienne est \(K_G=-S''/S=\tau\rho^2\). La coordinatisation est considérée dans son intérieur régulier, \(S>0\) ; pour la réalisation sphérique, \(0<r<\pi/\rho\). Les pôles sont traités avec des coordinatisations régulières différentes. Le paramètre \(\rho\) et la réalisation dimensionnelle de la métrique conservent leur déclaration propre.

Pour une masse réalisée \(m>0\), le lagrangien géodésique et son
Hamiltonien sont
\begin{equation}
 \mathcal L_{\rm geo}=\frac m2(\dot r^2+S^2\dot\theta^2),\qquad
 H_{\rm geo}=\frac1{2m}\left(p_r^2+\frac{p_\theta^2}{S^2}\right).
 \label{eq:hc-geodesique}
\end{equation}
En effet, \(p_r=m\dot r\), \(p_\theta=mS^2\dot\theta\) ;
le Hessien des vitesses est \(m\operatorname{diag}(1,S^2)\)
et est positif et inversible. L’équation radiale est
\(\ddot r-SS'\dot\theta^2=0\), tandis que
\(d(mS^2\dot\theta)/dt=0\) conserve le moment cinétique.
La positivité de cette énergie vaut pour les trois courbures. La forme
indéfinie du générateur hyperbolique de \eqref{eq:hc-hamiltonianos-nativos}
et l’énergie géodésique positive d’une métrique de courbure négative
sont des fonctions sur des espaces distincts. Cette comparaison fixe exactement
le sens de la courbure qui intervient dans chaque assertion.

\subsection{Réponse constitutive et réalisation modale de Maxwell}

Les sections d’action et les coordonnées angulaires précèdent la
réponse constitutive \cite{hmtIntegral}. Dans l’orientation réelle fixée \(x>y>0\),
avec les projecteurs des deux feuillets,
\[
 T=q_+P_++q_-P_-,\quad q_\pm=e^{-x\mp y},\quad
 \mathcal V=(I-T^{90})(I-T^{120})^{-1},
 \qquad 0<q_+<q_-<1,
\]
le calcul fonctionnel produit
\[
 r_\pm=\frac{1-q_\pm^{90}}{1-q_\pm^{120}},\qquad
 \widehat\varepsilon=r_+^2,\quad \widehat\mu=r_-^2,
 \quad \widehat c=(r_+r_-)^{-1},\quad \widehat Z=r_-/r_+.
\]
Les coordinatisations dimensionnelles correspondantes conservent exactement \(\varepsilon\mu=c^{-2}\) et \(\mu/\varepsilon=Z^2\). La permittivité, la perméabilité et la vitesse sont ainsi des représentations du même opérateur, avec les secteurs d'incidence fixés par le TPK ; la dynamique suivante les emploie comme des résultats déjà construits.

Pour exhiber le transport vers un hamiltonien de champ, on fixe une
réalisation spatiale orientée, des conditions de frontière permettant
l’intégration par parties, et un mode transversal réel \(f\) normalisé par
\[
 \nabla\cdot f=0,\qquad \int|f|^2=1,\qquad
 \nabla\times\nabla\times f=k^2f,\qquad k>0.
\]
En jauge transverse et sans sources, \(A=af\) et le moment \(\Pi=pf\) définissent une immersion du plan modal dans l'espace des champs. Pour des constantes constitutives homogènes, l'hamiltonien et la forme symplectique de Maxwell se restreignent à
\begin{equation}
 H_{\rm EM}(a,p)=\frac{p^2}{2\varepsilon}
 +\frac{k^2a^2}{2\mu},\qquad \omega_{\rm EM}=da\wedge dp.
 \label{eq:hc-maxwell-mode}
\end{equation}
L’intégration par parties donne
\(\int|\nabla\times f|^2=k^2\), qui prouve cette restriction.
Avec \(\omega_k=ck\), le changement
\begin{equation}
 Q=\sqrt{\varepsilon\omega_k}\,a,
 \qquad P=\frac{p}{\sqrt{\varepsilon\omega_k}}
 \label{eq:hc-maxwell-symplectic}
\end{equation}
préserve \(da\wedge dp=dQ\wedge dP\) et donne
\(H_{\rm EM}=\omega_k(Q^2+P^2)/2\). Son inverse et ses équations
\(\dot Q=\omega_kP,\ \dot P=-\omega_kQ\)
constituent l’entrelacement modal avec le régime elliptique positif.
La restriction est invariante car l’opérateur de Maxwell conserve le
mode propre transversal. La transformation de Legendre produit
\(\mathcal L_{\rm EM}=\varepsilon\dot a^2/2-k^2a^2/(2\mu)\).

Deux modes transversaux de même fréquence, combinés en quadrature,
réalisent les deux orientations circulaires du corpus. Le champ vérifie
\(\mathbf B=c^{-1}\widehat{\mathbf k}\times\mathbf E\) et
\(\omega=ck\). La relation avec l’oscillateur est démontrée ici par
l’immersion des champs, la normalisation modale et le transport de la forme
symplectique, en maintenant explicites les conditions de frontière.
Un changement de \(k\), de la géométrie de la cavité ou de la préparation
modifie le mode réalisé ; il laisse intacte la généalogie de
\(\varepsilon,\mu,c,Z\).

Dans la coordinatisation lorentzienne du même secteur, le couplage à une section matérielle s'écrit avec \(F=dA\) et une charge \(q_e\) déjà typée. Pour cette réalisation variationnelle, on prend \(\nabla\) égal à la connexion de Levi--Civita de \(g\). Sur \(T^*\mathcal M\), l'hamiltonien covariant
\[
 H_A=\frac1{2m}g^{\mu\nu}(p_\mu-q_eA_\mu)(p_\nu-q_eA_\nu)
\]
est la transformation de Legendre de
\(\mathcal L_A=\frac m2g_{\mu\nu}\dot x^\mu\dot x^\nu
+q_eA_\mu\dot x^\mu\). Le hessien \(mg\) est inversible.
L’équation d’Euler--Lagrange donne
\(m\nabla_u u^\mu=q_eF^\mu{}_{\nu}u^\nu\).
L’antisymétrie de \(F\) conserve \(g(u,u)\). Il s’agit de la
réalisation covariante à paramètre affine, tandis que
\eqref{eq:hc-maxwell-mode} utilise le temps d’une décomposition
spatiale transversale : les deux paramètres et les deux domaines sont définis.

\subsection{Canaux gradués, représentation de Clifford et opérateur de masse}

L'atlas matériel produit d'abord des canaux complets de chemin, de feuillet, de saveur, de couleur et de représentation. À une profondeur finie, sa réalisation hermitienne est \(\mathcal H_K=\bigoplus_{s\in S_K}V_s\), avec une base orthonormale de canaux et une graduation de parité transportée. L'algèbre extérieure du secteur impair engendre les CAR : l'insertion d'un mode \(a_j^*\) et son retrait adjoint \(a_j\) satisfont \(\{a_i,a_j^*\}=\delta_{ij}I\) et \(\{a_i,a_j\}=0\). La preuve compte les signes d'échange de modes complets ; les masses sont évaluées ensuite~\cite{hmtVI}.

Dans le secteur extérieur à un seul mode, la base \((|0\rangle,|1\rangle)\)
identifie l’opérateur de retrait à
\(a=N\). L’application \(e_1\mapsto|0\rangle,\ e_2\mapsto|1\rangle\)
est une isométrie du plan complexe normalisé et entrelace les deux
nilpotents. Cette identification utilise le produit scalaire et la
graduation déclarés ; elle conserve le sens distinct de leurs bases.
Avec le \(i\) interne déjà réalisé, on obtient
\begin{equation}
 \sigma_1=N+N^*,\qquad
 \sigma_2=-i(N-N^*),\qquad
 \sigma_3=NN^*-N^*N.
 \label{eq:hc-pauli}
\end{equation}
La multiplication directe donne
\(\sigma_j^*=\sigma_j\) et
\(\sigma_i\sigma_j=\delta_{ij}I+i\epsilon_{ijk}\sigma_k\).
En particulier, le couple nilpotent et son adjoint produisent une
représentation de Clifford ; l’adjoint est celui de la fibre hermitienne et
est transporté avec elle.

L'opérateur de masse sur une collection finie de canaux \(\mathcal F\) est construit par le registre, le caractère et ses raffinements :
\begin{equation}
 M_\beta=S M_0 S,\qquad
 S=R_{\rm act}^{B_M/2},\quad
 R_{\rm act}=\hbar_{\rm pre}/\hbar_{\rm ret}>0,
 \qquad M_0>0.
 \label{eq:hc-mass}
\end{equation}
Ici, \(B_M\) est l’opérateur diagonal de l’exposant de route et \(M_0\)
contient le caractère raffiné relatif à l’origine électronique.
La positivité découle de \(\langle v,M_\beta v\rangle
=\langle Sv,M_0Sv\rangle>0\) pour \(v\ne0\).
Le produit de feuillets conserve le secteur nul par sa réalisation propre ;
l’exponentielle du caractère matériel reste positive.

Dans cette réalisation, on écrit \(\hbar:=\hbar_{\rm ret}\), la section d'action déjà construite, et \(c:=c_{\rm vac}\) dans la coordinatisation dimensionnelle déclarée. Sur \(\mathbb C^4\otimes\mathcal F\), définissons
\[
 \alpha_j^{D}=\sigma_1\otimes\sigma_j,\qquad
 \beta^D=\sigma_3\otimes I_2.
\]
L'exposant distingue ces matrices de la constante de structure fine. La coordinatisation spatiale \(1+3\) attribue trois composantes d'impulsion \(\mathbf p\). L'application de multiplication
\begin{equation}
 H_D(\mathbf p)
 =c\sum_{j=1}^3\alpha_j^Dp_j\otimes I_{\mathcal F}
 +c^2\beta^D\otimes M_\beta
 \label{eq:hc-dirac}
\end{equation}
est la réalisation libre de Dirac de cet opérateur matériel.

\begin{theorem}[Factorisation spinorielle de la dispersion matérielle]
\label{thm:hc-dirac-square}
La matrice \eqref{eq:hc-dirac} est autoadjointe et satisfait
\begin{equation}
 H_D(\mathbf p)^2
 =I_4\otimes\bigl(c^2|\mathbf p|^2I_{\mathcal F}
 +c^4M_\beta^2\bigr).
 \label{eq:hc-dirac-square}
\end{equation}
Dans un sous-espace propre \(M_\beta v=m v\), ses énergies sont \(\pm E_m(\mathbf p)\), où \(E_m=\sqrt{c^2|\mathbf p|^2+m^2c^4}\).
\end{theorem}
\begin{proof}
Les relations de \eqref{eq:hc-pauli} donnent
\(\{\alpha_i^D,\alpha_j^D\}=2\delta_{ij}I_4\),
\(\{\alpha_i^D,\beta^D\}=0\) et \((\beta^D)^2=I_4\).
Le facteur de saveur commute avec les matrices spinorielles.
Les termes croisés s’annulent et les carrés donnent
\eqref{eq:hc-dirac-square}. La trace nulle et les relations de Clifford
produisent les deux branches, avec une multiplicité spinorielle de deux lorsque \(E_m>0\).
\end{proof}

Dans la représentation des moments, le domaine naturel est
\(\{\psi\in L^2(\mathbb R^3;\mathbb C^4\otimes\mathcal F):
H_D\psi\in L^2\}\). Étant une matrice hermitienne mesurable agissant
par multiplication, l’opérateur est autoadjoint sur ce domaine. La
représentation spatiale provient de la transformation unitaire de Fourier
postérieure à la construction des canaux. Le théorème fixe la réalisation
libre et sa dépendance exacte à la masse HMT ; l’interaction conserve ses
opérateurs de connexion et son domaine propre.

Dans la branche positive \(m>0\), la transformation de Legendre ordinaire du symbole \(E_m(\mathbf p)\) est
\[
 \mathbf v=\frac{c^2\mathbf p}{E_m},\qquad
 \mathcal L_m=-mc^2\sqrt{1-|\mathbf v|^2/c^2},
 \qquad |\mathbf v|<c.
\]
Résoudre \(\mathbf p=m\mathbf v/\sqrt{1-|\mathbf v|^2/c^2}\) et substituer dans \(\mathbf p\cdot\mathbf v-E_m\) démontre l'égalité. L'action spinorielle, en revanche, est du premier ordre dans le champ : \(\mathcal L_D=\frac{i\hbar}{2}(\psi^*\dot\psi-\dot\psi^*\psi)
-\psi^*H_D\psi\). Sa variation donne \(i\hbar\dot\psi=H_D\psi\) ; son hessien par rapport aux vitesses du champ est singulier. Les deux affirmations sont compatibles car elles agissent sur des espaces variationnels différents.

Pour une connexion abélienne statique dans une coordinatisation spatiale plane, avec \(\Pi_j=-i\hbar\partial_j-q_eA_j\), une masse constante et un potentiel scalaire nul, on conserve sur les fonctions lisses à support compact
\[
 [\Pi_i,\Pi_j]=iq_e\hbar\epsilon_{ijk}B_k,
\]
et la même factorisation produit
\begin{equation}
 H_D(A)^2=c^2\boldsymbol\Pi^2+c^4M_\beta^2
 -q_e\hbar c^2\boldsymbol\Sigma\cdot\mathbf B,
 \qquad \Sigma_j=I_2\otimes\sigma_j.
 \label{eq:hc-dirac-curvature}
\end{equation}
Les facteurs identité de saveur restent implicites. Le terme de courbure est le produit des deux antisymétriques : le commutateur des dérivées covariantes et la multiplication de Pauli. Son coefficient est fixé lors du choix de la représentation de charge et de la coordinatisation d'action ; il n'est pas sélectionné au moyen d'une valeur spectrale cible.

Si \(V\) est la transformation unitaire de saveur déjà construite, alors
\[
 (I_4\otimes V)H_D(M_\beta)(I_4\otimes V^*)
 =H_D(VM_\beta V^*).
\]
Le mélange unitaire conserve le spectre complet. La congruence positive
de \eqref{eq:hc-mass} transporte une autre structure et peut modifier les
valeurs propres dans une métrique fixe. Par exemple,
\(M_0=\operatorname{diag}(1,2)\), \(S=\operatorname{diag}(2,1)\)
donne \(SM_0S=\operatorname{diag}(4,2)\). Si l’on transporte également le
produit scalaire vers \(G'=S^*S\), le problème généralisé
\(SM_0Sv=\lambda G'v\) retrouve les valeurs propres de \(M_0\)
au moyen de \(w=Sv\). Ainsi, le changement de base de saveur, la déformation
matérielle et le transport de métrique ont des conséquences spectrales précises.

\subsection{Courbure du potentiel scalaire dans la réalisation électrofaible}

Le secteur électrofaible utilise le déterminant du même opérateur angulaire, \(D_A=\det T=e^{-2A_{\rm rad}}\), avec les signatures matérielles déjà construites. Dans la réalisation angulaire avec raffinements communs de \(W\) et \(Z\), leurs signatures \((94,-1)\) et \((95,-1)\) donnent \(m_W^\angle/m_Z^\angle=e^{-A_{\rm rad}}\). Dans la convention rationalisée à l'ordre de l'arbre du corpus,
\[
 g^2=\frac{4\pi_{\rm HMT}\alpha_{\rm HMT}}{1-D_A},\qquad
 \lambda_H=\frac{\pi_{\rm HMT}\alpha_{\rm HMT}}{2(1-D_A)}
 \exp\!\left(6A_{\rm rad}+\frac{10}3C^*_{\rm rad}\right)>0.
\]
La seconde formule résulte de la signature scalaire \((97,9)\) et de
\(m_H^2/m_W^2=8\lambda_H/g^2\). La normalisation du potentiel est
\(V(\mathsf H)=-\mu_H^2\mathsf H^*\mathsf H
+\lambda_H(\mathsf H^*\mathsf H)^2\). Ces compositions conservent
leur échelle et leur convention d’arbre ; les raffinements qui s’annulent
dans un rapport sont déclarés dans la réalisation elle-même.

Dans une coordinatisation radiale réelle \(\mathsf H^*\mathsf H=h^2/2\), le potentiel est \(V(h)=-\mu_H^2h^2/2+\lambda_Hh^4/4\). Pour \(\mu_H^2>0\), ses points critiques sont \(0\) et \(\pm v\), avec \(v^2=\mu_H^2/\lambda_H\), et
\[
 V''(0)=-\mu_H^2,\qquad V''(\pm v)=2\mu_H^2.
\]
Dans les unités naturelles de cette coordinatisation, le mode homogène normalisé possède \(H_h=p_h^2/2+V(h)\). Sa linéarisation en un point critique \(h_0\) a pour générateur
\[
 A_{h_0}=\begin{pmatrix}0&1\\-V''(h_0)&0\end{pmatrix},
 \qquad A_{h_0}^2=-V''(h_0)I.
\]
Par conséquent, l'origine est hyperbolique et les minima sont elliptiques ; la frontière \(V''=0\) est de type parabolique. Cette affirmation se déduit de la dérivée seconde du potentiel et d'une normalisation modale explicite. L'énergie matérielle positive du minimum et la direction instable d'un autre fond appartiennent ainsi à la même fonction avec des points de développement différents. Le mode linéarisé est un transport particulier des trois types quadratiques ; ses fréquences et ses unités proviennent du potentiel, et non d'une identification de son hessien à la courbure gaussienne de \eqref{eq:hc-metrique}.

\subsection{Mode de courbure de Yang--Mills et oscillation elliptique}

La réalisation de Yang--Mills construite dans les sections précédentes
possède un générateur positif \(D^{\rm YM}\) de dimension inverse de
longueur. L’échelle
\[
 \kappa_{\rm av}=\frac{\mu_{\rm vol}}{\ell_0}
 \left(\frac{\pi_{\rm HMT}}{\varphi_{\rm HMT}^2}-1\right)>0,
 \qquad \delta=3\kappa_{\rm av}
\]
provient de l'incidence des feuillets et du lecteur régional, avec l'unité longitudinale \(\ell_0\) déclarée. Le vide est simple et le témoin de courbure dans la coordinatisation produit est
\[
 W(q)=\mathbf1_{\{q_1+\cdots+q_6=0\}}-\frac13,
 \quad \mathbb EW=0,\quad \|W\|^2=\frac29,
 \quad D^{\rm YM}W=\delta W.
\]
Le transport unitaire vers la mesure d’interaction transporte \(W\) en
\(W_g\) et conserve ces identités. Les inclusions de raffinement
\(J_m\) entrelacent les générateurs et conservent le témoin, de sorte que
celui-ci demeure dans la limite physique. On utilise ici le résultat complet
des preuves du vide, de l’écart uniforme et de la reconstruction, non une
valeur propre isolée d’une matrice finie~\cite{HMT-YM}.

Soit \(w=W_g/\sqrt{2/9}\) le vecteur unitaire dans la fibre physique et
\(\hbar=\hbar_{\rm ret}>0\). Le Hamiltonien énergétique est
\(H_E=\hbar cD^{\rm YM}\). Sur l’espace de Hilbert réalifié,
on adopte la forme \(\omega_{\mathcal H}(u,v)=2\hbar
\operatorname{Im}\langle u,v\rangle\), avec un produit scalaire antilinéaire
par rapport au premier argument. Définissons
\begin{equation}
 \mathfrak F(Q,P)=\frac{Q+iP}{\sqrt{2\hbar}}\,w.
 \label{eq:hc-ym-embedding}
\end{equation}

\begin{theorem}[Immersion symplectique du mode physique de courbure]
\label{thm:hc-ym-oscillator}
L’application \(\mathfrak F:\mathbb R^2\to\mathcal H^{\rm phys}\)
est une immersion réelle dont l’image est invariante par l’évolution physique et
vérifie
\begin{equation}
 \mathfrak F^*\omega_{\mathcal H}=dQ\wedge dP,
 \qquad
 \langle\mathfrak F,H_E\mathfrak F\rangle
 =\frac{c\delta}{2}(Q^2+P^2).
 \label{eq:hc-ym-energy-pullback}
\end{equation}
L'équation de Schrödinger restreinte entrelace exactement le flux elliptique \(\dot Q=c\delta P,\ \dot P=-c\delta Q\). L'application est compatible avec les inclusions de raffinement.
\end{theorem}
\begin{proof}
Pour deux variations \((a,b),(a',b')\), le produit scalaire de leurs
images a pour partie imaginaire \((ab'-ba')/(2\hbar)\).
Cela prouve la première identité. L’équation propre et \(\|w\|=1\)
donnent la seconde. Puisque
\(\dot\psi=-icD^{\rm YM}\psi\), le coefficient \(Q+iP\)
satisfait \(\dot Q+i\dot P=-ic\delta(Q+iP)\).
Enfin, \(J_mw_m=w_{m+1}\) implique
\(J_m\mathfrak F_m=\mathfrak F_{m+1}\). Ces identités passent à la
limite par les isométries déjà construites.
\end{proof}

Le plan complet décrit les amplitudes d’un mode ; son cercle
\(Q^2+P^2=2\hbar\) représente les vecteurs normalisés de cet
espace propre complexe. Sur ce cercle, la phase globale évolue et le
rayon quantique reste fixe. Les superpositions d’espaces propres distincts
conservent, en outre, leurs phases relatives. L’immersion maintient cette
différence entre l’espace des amplitudes et l’espace projectif.

L'énergie associée au saut spectral et le saut de masse sont \(\Delta_E=\hbar c\delta\) et \(m_{\rm gap}=\hbar\delta/c\). La fréquence modale \(c\delta\) a la dimension d'un temps inverse. Les longueurs euclidiennes de \(e^{-sD^{\rm YM}}\) et le temps physique de \(e^{-itH_E/\hbar}\) sont reliés par la coordinatisation et la reconstruction, en conservant leurs groupes et semi-groupes distincts. La composition \(\mathfrak F\circ M_\eta^{-1}\) transporte cette même énergie vers l'ellipse d'action ; si le repère varie, le terme de connexion de \eqref{eq:hc-htr} complète son évolution.

Le lien obtenu est une réalisation exacte sur le mode invariant de
courbure et ses prolongements. La positivité de l’hamiltonien quantique
engendre une rotation dans sa réalification symplectique ; la dénomination
hyperbolique d’une métrique spatiale appartient au tenseur métrique, et le
transport unipotent du registre appartient à sa mémoire. En dimension
finie, un groupe \(I+tN\) unitaire pour tout \(t\in\mathbb R\), avec
\(N^2=0\), impose \(N=0\) : dériver l’unitarité donne \(N^*=-N\),
et alors \(N^*N=-N^2=0\). Le flot parabolique classique conserve sa
forme symplectique et possède sa quantification dans des espaces appropriés ; cette
identité délimite seulement l’identification à une matrice unitaire finie.

Ainsi s’articule un ensemble de transports concrets. La mémoire
produit une charge et une action parabolique ; le repère de l’ellipse produit
la connexion variationnelle et son lagrangien ; les réponses constitutives
déterminent les coefficients de Maxwell ; la graduation des canaux construit
Clifford et permet de factoriser l’opérateur matériel ; la reconstruction de
Yang--Mills apporte une immersion modale symplectique qui conserve exactement
son écart et son raffinement. Chaque identification dispose d’une application,
d’une forme conservée et d’une équation entrelacée.


\section{Discussion : reconstruction, invariants et dualité}
\label{sec:discussion}
Le résultat structurel prioritaire est la récupérabilité de la forme.
La charge uniforme et les deux mémoires nonadiques reconstruisent le registre
au moyen d’un inverse explicite ; le cône positif décrit exactement
le domaine où cette reconstruction induit des mineurs strictement positifs.
La borne de stabilité quantifie la persistance de ce domaine. Par
conséquent, les poids géométriques admettent une provenance et une récupération
opératoires, en plus de leur évaluation numérique.

La composition développée a une séquence explicite. Le registre
terminal engendre douze poids positifs ; leurs sommes partielles donnent treize points
ordonnés ; un réseau planaire réalise leurs mineurs ; la courbe des moments
construit des polygones et des polytopes ; leurs triangulations fournissent une couverture
et des formes à résidus orientés. Le raffinement d’une cellule fixée conserve
la somme de ses formes. L’introduction de nouveaux sommets sur la courbe des
moments produit pour sa part une autre géométrie et exige de comparer les domaines.

La direction algébrique du registre distingue deux opérations. Le redimensionnement positif des colonnes conserve les rapports projectifs pertinents. Son action normalisée sur les intervalles engendre en revanche une déformation effective, certifiée par une dérivée exacte non nulle. Cette déformation admet un potentiel convexe dont la transformée de Legendre est conservée lors d'une subdivision à direction héritée. L'invariance possède ainsi un mécanisme vérifiable : chaque successeur conserve le générateur de son prédécesseur et la somme de ses poids restitue le poids antérieur.

Le raffinement énergétique fournit une seconde conservation exacte.
La réduction de Schur retrouve le laplacien effectif et le prolongement
minimisant sépare chaque configuration en sa partie observée et un détail
reconstructible. Les énergies des détails s’additionnent orthogonalement.
Avec des maillages décroissants et une énergie uniformément bornée, cette suite
converge dans la réalisation de Sobolev déclarée. L’invariance de la
forme canonique, l’identité du potentiel convexe et la conservation de
l’énergie sont trois résultats distincts de compatibilité, démontrés
pour les raffinements spécifiés.

Le prolongement aux douze plans réticulaires conserve une symétrie
ternaire et le volume. Le résultat décisif est la réciprocité
\(\Lambda_s^*=\Lambda_{-s}\), qui induit l’inversion
\(\Theta_s(t)=t^{-12}\Theta_{-s}(t^{-1})\). La réalisation déployée de
signature \((24,24)\) conserve intégralité, parité et autodualité pour
tout le flot. Dans la projection euclidienne, la variation explicite d’une
norme distingue déformation et automorphisme du réseau initial. Cette
distinction détermine le contenu exact de la connexion avec la dualité
des moments et des enroulements.

La réalisation matérielle conserve ses propres lecteurs. La représentation d'action et le caractère des masses reçoivent le registre par des compositions déjà construites ; la congruence préserve la positivité et l'inertie, tandis que la métrique détermine l'interprétation de leurs valeurs propres. Ces différences précisent la signification de la multimorphologie : plusieurs formes mathématiques proviennent d'une généalogie commune et sont reliées par des opérations vérifiables. L'identification complète des amplituèdres de rang supérieur exigera les preuves de couverture et de résidus pour ces applications concrètes. Le résultat de rang un démontré ici fixe un cas complet, avec coordonnées rationnelles, triangulation et contrôles reproductibles.

L’orientation temporelle du TRIT intègre le retour avec mémoire, la
mise à jour du seuil et les prolongements conjugués. Ses générateurs
elliptique, nilpotent et hyperbolique fournissent des lois distinctes
d’évolution ; la composition du TPK conserve la généalogie qui les
concatène. En particulier, l’activité du régime zéro est exprimée
par un transport unipotent, et la double projection conserve simultanément
la convergence des coordonnées et l’incidence exceptionnelle. La thèse
géométrique de l’article réside dans cette articulation : la structure de
départ engendre des registres récupérables et ceux-ci soutiennent des réalisations
positives dont les transformations, les frontières et les mémoires peuvent être calculées
et comparées avec précision.

L’entrelacement \(B_FM=MB_V\) ajoute un transport exact depuis le
bloc modulaire jusqu’au quotient incidentiel. Les multiplicités
un et trois s’y déduisent de la matrice de Gram, et la forme signée déclarée
s’exprime comme \(J_L=(7I_Q-\overline B_F)/2\). Les matrices de
Vandermonde et de Hankel apportent, d’autre part, des réalisations explicites
\(Y^{(k)}\) pour \(2\le k\le9\), avec positivité et rang prouvés
dans la partition initiale. Le prolongement nonadique étend ensuite
cette construction à tout couple fini admissible de rang et de dimension
externe, au moyen des systèmes de nœuds et de moments développés ici.
Ces résultats étendent le système de réalisations sans identifier
la dimension d’un quotient, le rang d’une matrice et le nombre de
générateurs d’une superalgèbre.

Dans la réalisation de jauge, l'échelle surfacique--volumétrique précède l'hamiltonien. Le générateur produit conserve tous les secteurs de la décomposition orthogonale, tandis que sa restriction physique contient un témoin qui atteint le premier niveau positif. Les opérateurs d'entrelacement du raffinement conservent la coercivité et ce témoin. Le développement continu réunit en outre la mesure commune, la courbure, la continuité euclidienne, les formes de réflexion et l'hamiltonien reconstruit. Ainsi, l'argument de l'écart spectral se situe aux côtés de la construction qui lui donne son domaine et son sens, et la positivité de Plücker maintient sa propre réalisation géométrique. Leur articulation exige ces transports explicites ; le mot positivité désigne des propriétés typées que l'état commun permet de comparer sans les substituer les unes aux autres.

\clearpage
\part{Vérification exacte et énumération des émissions}
\section{Vérification exacte des constructions finies}
\label{sec:exact-reproduction}

La reproduction distingue la démonstration générale et l'évaluation d'une spécialisation. Les propositions précédentes utilisent des identités algébriques, la positivité et des hypothèses de convergence explicites. Leurs contrôles computationnels sont effectués avec des entiers et des fractions rationnelles ; le développement réticulaire conserve en outre sa représentation exacte dans l'extension algébrique déclarée. Les données du registre sont des sorties exprimées par la construction précédente. Les exécutions de cette section vérifient leur transport géométrique et opératoriel, et conservent séparément la provenance de leur génération.

Le premier contrôle restitue les séparations à partir des mineurs, vérifie la relation de Plücker dans les triangulations et compare les changements de coordinatisation aux déformations effectives. Un changement inversible des lignes multiplie tous les mineurs par le même déterminant ; les quotients marqués demeurent identiques. Un redimensionnement indépendant des colonnes conserve certains birapports et modifie, en général, les séparations reconstruites. Évaluer les deux opérations sur les mêmes données fournit un contrôle négatif de l'équivalence affirmée.

Le second contrôle subdivise les longueurs par neuf, conserve les marques héritées et vérifie les identités des moyennes et de Gram. Les déterminants sont calculés par élimination exacte et comparés à leurs produits de Vandermonde. La positivité en toute dimension finie découle de la formule démontrée ; les cas finis exécutés vérifient l'implémentation de cette formule, ses indices et ses normalisations. Les mesures de comptage et les mesures chronologiques sont maintenues séparées afin qu'une égalité de moments ne dépende pas d'un changement silencieux de normalisation.

Le troisième contrôle reproduit les huit premiers moments au moyen
de l’opérateur de Jacobi d’ordre quatre. La matrice symétrique est
représentée sous forme rationnelle dans la base des polynômes unitaires ;
sa métrique est \(\operatorname{diag}(h_0,\ldots,h_3)\). On
vérifie l’auto-adjonction dans cette métrique, la récurrence à trois
termes et l’accord entre la résolvante matricielle et sa fraction
continue en \(z=2\). Le résultat maintient explicites le vecteur
constant observé et la charge \(Q\), qui interviennent dans le
théorème de fidélité spectrale marquée.

Le quatrième groupe vérifie les changements symplectiques et leurs
hamiltoniens. La dérivée par rapport au moment récupère la
vitesse ; la transformation de Legendre est vérifiée comme
identité polynomiale ou de Laurent. L’inertie d’une forme
quadratique et la positivité d’un hamiltonien quantique
appartiennent à leurs espaces respectifs. La représentation
de Clifford est vérifiée au moyen de ses anticommutateurs,
et la composition avec l’opérateur de masse est prouvée par
le carré de l’opérateur de Dirac. Ces égalités
identifient chaque terme et ses produits croisés.

Enfin, les vérifications précédentes de triangulation, de forme canonique, de réduction de Schur, de congruence des masses et de construction réticulaire sont conservées. Dans Yang--Mills, les propositions sur la mesure commune, le raffinement et la reconstruction sont exposées dans le corps du texte avec leurs prémisses et leurs critères de réfutation. Un contrôle fini de matrices ne remplace ni le passage à la limite ni la reconstruction de champs : leurs domaines sont enregistrés séparément dans la livraison reproductible.

Le paquet de sources relie chaque programme à l’énoncé
qu’il évalue et conserve sa sortie intégrale. Les conditions sont
vérifiées par des branchements explicites, de sorte qu’elles
restent actives lors de l’exécution de l’interpréteur en mode
optimisé. Elles sont également reproduites en mode isolé, sans
charger de paquets utilisateur. La correspondance entre
construction, preuve et exécution peut être inspectée
étape par étape sans transformer la sortie d’un programme en
une assertion de portée supérieure à celle qu’il vérifie.

\begingroup\setlength{\parskip}{2pt}\fontsize{10}{11.5}\selectfont
\import{foundation/}{libro/censo_firmas_emision.tex}
\par\endgroup
\begingroup\setlength{\parskip}{2pt}\fontsize{10}{11.5}\selectfont
\clearpage
\section{Conclusions}
\label{sec:conclusions-final}
\begingroup
\fontsize{10.2}{12.2}\selectfont
\setlength{\parskip}{2.5pt plus .5pt}

La thèse géométrique est établie par reconstruction et transport : une représentation positive de la structure discrète du continu conserve, avec ses marques, le registre dont elle provient et les lecteurs que celui-ci détermine. La charge totale et les quotients de mineurs restituent chaque coordonnée ; les deux mémoires nonadiques fournissent une autre reconstruction opératorielle. La valeur et la forme s'articulent ainsi par des inverses explicites, après la génération APP--TRIT--TPK et ses prolongements compatibles. La double projection maintient unies la réalisation arithmétique et l'incidence exceptionnelle dans cette généalogie.

Le raffinement dépasse la limitation d'une matrice à treize colonnes. Pour tout rang et toute dimension externe finis, un niveau nonadique suffisant produit les mineurs positifs et l'image amplituèdrale correspondante. Les nœuds hérités conservent leurs drapeaux ; les moyennes cellulaires conservent la mesure et enregistrent le détail orthogonal. En rang un, la couverture simpliciale et la récurrence des résidus s'étendent à toute dimension finie. Dans les rangs supérieurs, le résultat détermine des collections positives compatibles, de rang et de limite démontrés ; leur étendue dimensionnelle demeure distincte de la classification globale de toutes les cellules amplituèdrales.

La conservation admet trois expressions complémentaires. La forme canonique s'additionne sous subdivision d'un support fixe ; l'énergie effective est conservée par réduction de Schur ; et la mesure chronologique détermine des moments dont l'incrément se décompose en détails récupérables. Ces moments construisent des opérateurs de Jacobi dont les compressions finies ont un spectre simple et des résidus positifs. On obtient ainsi un parcours effectif des séparations arithmétiques à une représentation spectrale, avec une fraction continue déterminée par la même mesure. Les mutations qui ne changent que la coordinatisation conservent les lecteurs du point marqué.

La réalisation physique gagne en précision grâce à ses opérateurs. Les hamiltoniens des trois régimes tritiques possèdent des formes symplectiques et des transformations de Legendre explicites. La coordinatisation elliptique variable incorpore son terme de connexion ; la mémoire parabolique conserve son invariant ; l'ouverture hyperbolique possède son évolution propre. Les lecteurs d'action et de réponse constitutive sont transportés par l'inverse grassmannien. Avec les représentations régionales et les marques entières déclarées, la réponse constitutive étiquetée restitue à son tour le registre. La condition d'égalité des traces gouverne le recollement des formes à valeurs dans la droite d'action.

La déformation exceptionnelle conserve le volume et la symétrie
ternaire, et transforme l’inversion du paramètre en dualité
réticulaire. Sa réciprocité thêta et sa réalisation de signature
déployée relient la variation géométrique à la dualité du
réseau. En Yang--Mills, l’échelle aréale--volumétrique, la mesure commune,
la coercivité uniforme et le témoin persistant situent l’écart
auprès du Hamiltonien qui le réalise. Le développement conserve
l’espace physique et les opérateurs d’entrelacement nécessaires pour transporter
cette assertion à travers le raffinement.

Le résultat conjoint est une architecture de formes avec mémoire : la dimension peut augmenter, la représentation peut changer et l'information pertinente demeure récupérable au moyen d'applications déterminées. Tel est le contenu distinctif de l'article. La géométrie positive, la dualité réticulaire et la réalisation spectrale sont reliées par des constructions qui spécifient ce qu'elles engendrent, ce qu'elles conservent et comment elles se reconstruisent à partir d'une même origine discrète.
La représentation d'incidence conserve également la composante uniforme du registre : son inverse reconstruit les douze coordonnées et permet de récupérer la direction géométrique après le transport de mémoire. Dans la réalisation sectorielle, l'aire totale et l'occupation pondérée déterminent conjointement la répartition hexade--octade. La purification spécifiée transforme cette répartition en une structure d'intrication, avec une identité finie entre entropie, aires sectorielles et entropie relative. La conservation de l'information est ainsi précisée pour chaque application : registre complet, direction, aire et composition possèdent des inverses ou des fibres explicitement distincts.
L'échelle d'aire est déterminée au moyen d'une longueur engendrée par la composition linéaire du retour enregistré et de la norme locale de Hadamard. La réciprocité radiale détermine le coefficient gravitationnel dans la réalisation réduite déclarée et reconnaît ensuite cette longueur comme la longueur de Planck. Les mutations du même point marqué et les raffinements hérités conservent l'opérateur d'aire, ses multiplicités et l'identité modulaire. Les formes canoniques à valeurs dans cet opérateur étendent la conservation aux résidus, tandis que le rayon de chaque horizon maintient son propre facteur géométrique. Cette composition explicite la dépendance entre longueur, incidence et aire sans altérer les états sectoriels.
\par\endgroup
\clearpage

\par\endgroup

\begin{thebibliography}{99}
\bibitem{hmtI} O. Haidara Fall et R. Ramos Balsa,
\emph{Génération coinductive de \(\pi,\varphi,e\) à une profondeur arbitraire,
dérivation de la constante de structure fine et structure algébrique de
l’électron}. Série HMT--MD, article I, édition consultée de septembre 2026.
\bibitem{hmtX} O. Haidara Fall et R. Ramos Balsa,
\emph{Extrême et moyenne raison holographique}. Développement du couplage
arithmético--géométrique, registre dodécaphasique et direction dimensionnelle,
édition consultée de septembre 2026.
\bibitem{hmtII} O. Haidara Fall et R. Ramos Balsa,
\emph{Origine structurelle du paramètre de Barbero--Immirzi et des
angles de mélange CKM ; dérivation des constantes de Planck, de Boltzmann
et de gravitation universelle}. Série HMT--MD, article II, édition de
septembre 2026. Relation entre action, aire, information et entropie
d’intrication.
\bibitem{hmtVI} O. Haidara Fall et R. Ramos Balsa,
\emph{Particules, persistance et spectre des masses}. Série HMT--MD, article VI,
édition consultée de septembre 2026.
\bibitem{HMT-IV} O. Haidara Fall et R. Ramos Balsa,
\emph{Moonshine, dualité T et théorie M à partir de la structure discrète du
continu}. Série HMT--MD, article IV, édition consultée de septembre 2026.
\bibitem{HMT-VIII} O. Haidara Fall et R. Ramos Balsa,
\emph{Arithmétique généalogique des nombres premiers et structure spectrale
de la fonction zêta}. Série HMT--MD, article VIII, édition consultée de septembre 2026.
\bibitem{HMT-IX} O. Haidara Fall et R. Ramos Balsa,
\emph{Résolution de l’hypothèse du continu dans la clôture généalogique
holographique}. Série HMT--MD, article IX, édition consultée de septembre 2026.
\bibitem{Elkies} N. D. Elkies,
\emph{Rational lattices and their theta functions}. Arizona Winter School,
2009, théorème 2.
\href{https://people.math.harvard.edu/~elkies/aws09.pdf}{Notes de l’auteur}.
\bibitem{ConwaySloane} J. H. Conway et N. J. A. Sloane,
\emph{Sphere Packings, Lattices and Groups}, troisième édition, Springer, 1999.
\href{https://doi.org/10.1007/978-1-4757-6568-7}{DOI: 10.1007/978-1-4757-6568-7}.
\bibitem{postnikov} A. Postnikov,
\emph{Total positivity, Grassmannians, and networks}, 2006,
\href{https://arxiv.org/abs/math/0609764}{arXiv:math/0609764}.
\bibitem{scott} J. S. Scott,
\emph{Grassmannians and Cluster Algebras}, 2003,
\href{https://arxiv.org/abs/math/0311148}{arXiv:math/0311148}.
\bibitem{amplituhedron} N. Arkani-Hamed et J. Trnka,
\emph{The Amplituhedron}, 2013,
\href{https://arxiv.org/abs/1312.2007}{arXiv:1312.2007}.
\bibitem{positive} N. Arkani-Hamed, Y. Bai et T. Lam,
\emph{Positive Geometries and Canonical Forms}, 2017,
\href{https://arxiv.org/abs/1703.04541}{arXiv:1703.04541}.
\bibitem{HMT-YM} O. Haidara Fall et R. Ramos Balsa,
\emph{La clôture holographique de l’infini}, édition intégrale de 2.249 pages,
chapitre 99, « Image de jauge et composition vers Yang--Mills », pp. 1742--1762 ;
et \emph{Yang--Mills : mesure commune, reconstruction et écart spectral},
monographie spécialisée, développement de l’échelle aréale--volumétrique.
\bibitem{hmtIntegral} O. Haidara Fall et R. Ramos Balsa,
\emph{La clôture holographique de l’infini}, édition intégrale de 2.249 pages,
septembre 2026. Développements des régimes quadratiques, du vide
électromagnétique, de la réalisation électrofaible et des correspondances
spirales de la mécanique dimensionnelle.

\phantomsection
\addcontentsline{toc}{section}{Références}
\bibitem{feynman1948} R. P. Feynman, ``Space-Time Approach to Non-Relativistic Quantum Mechanics'', \emph{Reviews of Modern Physics} \textbf{20}, 367--387 (1948). \href{https://doi.org/10.1103/RevModPhys.20.367}{doi:10.1103/RevModPhys.20.367}.
\bibitem{feynman1949} R. P. Feynman, ``The Theory of Positrons'', \emph{Physical Review} \textbf{76}, 749--759 (1949). \href{https://doi.org/10.1103/PhysRev.76.749}{doi:10.1103/PhysRev.76.749}.
\bibitem{feynman1965} R. P. Feynman, ``The Development of the Space-Time View of Quantum Electrodynamics'', conférence Nobel, 11 décembre 1965. \href{https://www.nobelprize.org/prizes/physics/1965/feynman/lecture/}{Texte de la conférence}.
\bibitem{mustovicary} B. Musto et J. Vicary, ``Quantum Latin Squares and Unitary Error Bases'', \emph{Quantum Information and Computation} \textbf{16}, 1318--1332 (2016). \href{https://doi.org/10.26421/QIC16.15-16-4}{doi:10.26421/QIC16.15-16-4}.
\bibitem{delascuevas2020} G. De las Cuevas, T. Drescher et T. Netzer, ``Quantum Magic Squares: Dilations and Their Limitations'', \emph{Journal of Mathematical Physics} \textbf{61}, 111704 (2020). \href{https://doi.org/10.1063/5.0022344}{doi:10.1063/5.0022344}.
\bibitem{delascuevas2023} G. De las Cuevas, T. Netzer et I. Valentiner-Branth, ``Magic Squares: Latin, Semiclassical, and Quantum'', \emph{Journal of Mathematical Physics} \textbf{64}, 022201 (2023). \href{https://doi.org/10.1063/5.0127393}{doi:10.1063/5.0127393}.
\bibitem{dirac1933} P. A. M. Dirac, ``The Lagrangian in Quantum Mechanics'', \emph{Physikalische Zeitschrift der Sowjetunion} \textbf{3}, 64--72 (1933). \href{https://www.informationphilosopher.com/solutions/scientists/dirac/Lagrangian_1933.pdf}{Reproduction de l'article original}.
\bibitem{pauli1946} W. Pauli, ``Exclusion Principle and Quantum Mechanics'', conférence Nobel, 13 décembre 1946, dans \emph{Nobel Lectures, Physics 1942--1962}, Elsevier, 1964, pp. 27--43, notamment p. 42. \href{https://www.nobelprize.org/uploads/2018/06/pauli-lecture.pdf}{Texte de la conférence}.
\bibitem{codata2018} E. Tiesinga, P. J. Mohr, D. B. Newell et B. N. Taylor, ``CODATA Recommended Values of the Fundamental Physical Constants: 2018'', \emph{Reviews of Modern Physics} \textbf{93}, 025010 (2021). \href{https://doi.org/10.1103/RevModPhys.93.025010}{doi:10.1103/RevModPhys.93.025010}. \href{https://physics.nist.gov/cuu/Constants/ArchiveASCII/allascii_2018.txt}{Tableaux archivés de 2018}.
\bibitem{codata2022} P. J. Mohr, D. B. Newell, B. N. Taylor et E. Tiesinga, ``CODATA Recommended Values of the Fundamental Physical Constants: 2022'', \emph{Reviews of Modern Physics} \textbf{97}, 025002 (2025). \href{https://doi.org/10.1103/RevModPhys.97.025002}{doi:10.1103/RevModPhys.97.025002}. \href{https://physics.nist.gov/cuu/Constants/Table/allascii.txt}{Tableau des constantes}.
\bibitem{flm} I. B. Frenkel, J. Lepowsky et A. Meurman, \emph{Vertex Operator Algebras and the Monster}, Pure and Applied Mathematics, vol. 134, Academic Press, 1988.
\bibitem{borcherds} R. E. Borcherds, ``Monstrous Moonshine and Monstrous Lie Superalgebras'', \emph{Inventiones Mathematicae} \textbf{109}, 405--444 (1992). \href{https://doi.org/10.1007/BF01232032}{doi:10.1007/BF01232032}.
\bibitem{conwaysloane} J. H. Conway et N. J. A. Sloane, \emph{Sphere Packings, Lattices and Groups}, 3\textsuperscript{e} éd., Grundlehren der mathematischen Wissenschaften 290, Springer, 1999. \href{https://doi.org/10.1007/978-1-4757-6568-7}{doi:10.1007/978-1-4757-6568-7}.
\bibitem{bipm2018} Conférence générale des poids et mesures, \emph{Resolution 1 of the 26th CGPM: On the revision of the International System of Units}, 2018. \href{https://www.bipm.org/en/committees/cg/cgpm/26-2018/resolution-1}{Résolution officielle}. \href{https://doi.org/10.59161/CGPM2018RES1E}{doi:10.59161/CGPM2018RES1E}.
\bibitem{dirac1931} P. A. M. Dirac, ``Quantised singularities in the electromagnetic field'', \emph{Proceedings of the Royal Society of London. Series A} \textbf{133}, 60--72 (1931). \href{https://doi.org/10.1098/rspa.1931.0130}{doi:10.1098/rspa.1931.0130}.
\bibitem{lep2006} ALEPH, DELPHI, L3, OPAL et SLD Collaborations; LEP Electroweak Working Group; SLD Electroweak Group et SLD Heavy Flavour Group, ``Precision electroweak measurements on the Z resonance'', \emph{Physics Reports} \textbf{427}, 257--454 (2006). \href{https://doi.org/10.1016/j.physrep.2005.12.006}{doi:10.1016/j.physrep.2005.12.006}.
\bibitem{chenlamshimakura2016} H.-Y. Chen, C. H. Lam et H. Shimakura, ``\(\mathbb Z_3\)-orbifold construction of the Moonshine vertex operator algebra and some maximal \(3\)-local subgroups of the Monster'', prépublication, arXiv:1606.05961 (2016), version 2 (2017). \href{https://arxiv.org/abs/1606.05961}{arXiv:1606.05961}.
\bibitem{abelamyamada2017} T. Abe, C. H. Lam et H. Yamada, ``A remark on \(\mathbb Z_p\)-orbifold constructions of the Moonshine vertex operator algebra'', prépublication, arXiv:1705.09022 (2017), version 4. \href{https://arxiv.org/abs/1705.09022v4}{arXiv:1705.09022}.
\bibitem{kmconwaynorton1979} J. H. Conway et S. P. Norton, ``Monstrous Moonshine'', \emph{Bulletin of the London Mathematical Society} \textbf{11}, 308--339 (1979). \href{https://doi.org/10.1112/blms/11.3.308}{doi:10.1112/blms/11.3.308}.
\bibitem{bennett2003} C. H. Bennett, ``Notes on Landauer's principle, reversible computation, and Maxwell's Demon'', \emph{Studies in History and Philosophy of Modern Physics} \textbf{34}, 501--510 (2003). \href{https://doi.org/10.1016/S1355-2198(03)00039-X}{doi:10.1016/S1355-2198(03)00039-X}.

\end{thebibliography}
\end{document}
